% =====================================================================
%  R. Nielsen, OM HINDRINGER OG BETINGELSER FOR DET AANDELIGE LIV I
%  NUTIDEN. Sexten Forelæsninger, holdte ved Universitetet i Christiania
%  September--Oktober 1867. Kjøbenhavn: Gyldendal (F. Hegel), 1868.
%  / On Obstacles to and Conditions for Spiritual Life in the Present Age.
%    Sixteen Lectures Delivered at the University of Christiania,
%    September--October 1867.
%
%  TRANSLATION (English)
%  Translated from transcription.tex (Danish), which is COMPLETE and
%  collated at the image (see its header §4). Mirrors it 1:1: every
%  printed page carries % --- p. N --- and \opage{N} at the corresponding
%  point in the English. Method: ../../../TRANSLATION-PLAYBOOK.md.
%
%  ---- TERMINOLOGY ----
%  Aligned with the 1868 Brøchner--Nielsen exchange translations
%  (../brochners-kritik, ../../brochner/problemet-tro-viden):
%    Aand = spirit · aandelig = spiritual · Erkjendelse = cognition
%    Viden = knowledge · Videnskab = science · Tro = faith
%    Villie = will · Bevidsthed = consciousness · Forestilling = representation
%    Personlighed = personality · Begreb = concept · Tilværelse = existence
%    Ueensartethed = heterogeneity · Samvittighed = conscience
%    Evangelietro(en) = the faith of the Gospel · Nutiden = the present age
%    Forelæsning = lecture · Tilhørere og Tilhørerinder = listeners, men and women
%
%  ---- CONVENTIONS ----
%  „…“ -> ``…''; \emph{} (letterspacing) carried over one for one; the
%  print's true italic (\textit{}) likewise. Latin, German, French and
%  Greek quotations stay in their language, formatted as in the Danish.
%  ɔ: -> i.e.; Danish work titles stay Danish. Printer's errors: the
%  Danish carries them as printed; the English translates the SENSE and
%  logs the divergence in a % note at the site. The four Rettelser
%  (pp. 50, 70, 109, 174) are followed in the English, with a % note.
%  Heads, toc entries and running heads as in transcription.tex (the
%  latter two editorial).
% =====================================================================
\documentclass[11pt]{book}
\usepackage[utf8]{inputenc}
\usepackage[T1]{fontenc}
\usepackage{libertinus}
\usepackage{libertinust1math}
\usepackage{textalpha}
\usepackage{graphicx}
\DeclareUnicodeCharacter{0254}{\reflectbox{c}}
\usepackage{amsmath}
\usepackage[english]{babel}
\usepackage[protrusion=true,expansion=true]{microtype}
\usepackage{enumitem}
\usepackage{fancyhdr}
\setlength{\headheight}{28pt}
\addtolength{\topmargin}{-13.5pt}
\pagestyle{fancy}
\fancyhf{}
\fancyhead[LE,RO]{\thepage}
\fancyhead[RE]{\textit{On Obstacles to and Conditions for Spiritual Life}}
\fancyhead[LO]{\textit{\leftmark}}
\renewcommand{\thefootnote}{\ifcase\value{footnote}\or *)\or **)\or ***)\else
  \arabic{footnote})\fi}
\usepackage{setspace}
\setstretch{1.3}
\usepackage{marginnote}
\newcommand{\opage}[1]{\marginnote{\footnotesize\textit{[#1]}}}
\newcommand{\apage}[1]{\marginnote{\footnotesize\textit{[#1]}}}
\usepackage{hyperref}
\hypersetup{pdftitle={On Obstacles to and Conditions for Spiritual Life in the Present Age},
  pdfauthor={R. Nielsen}, hidelinks}

\begin{document}

\frontmatter
% --- front matter: title page (PDF 6), Preface (PDF 8), dedication (PDF 10),
% in the order of the leaves in this copy. Unpaginated: no page markers.
% --- PDF 6: title page ---
\begin{titlepage}
\centering
On\par
{\Large Obstacles to and Conditions for}\par
{\LARGE Spiritual Life in the Present Age.}\par
\bigskip
{\large Sixteen Lectures,}\par
delivered at the University of Christiania\par
\emph{September---October 1867}\par
\medskip
by\par
\emph{R. Nielsen.}\par
\vfill
\vfill
Copenhagen.\par
Published by the Gyldendal Bookshop (F. Hegel).\par
{\small Printed by J. H. Schultz.}\par
1868.\par
\bigskip
{\small Translated from the Danish}\par
\end{titlepage}

% --- PDF 8: Forord ---
\chapter*{Preface.}
\addcontentsline{toc}{chapter}{Preface}

Although these lectures in their present form cannot be said to be in
each and every respect a direct reproduction of the spoken word, I have
nevertheless, by making use partly of my own notes and partly of more
detailed reports in Norwegian papers, sought to come as close as possible
to the lectures as originally and freely delivered. Here and there,
however, the analyses of concepts have been carried somewhat further, and
certain parts, notably in the last two lectures but one, have been
inserted; I have likewise found myself called upon to meet several
misgivings and epilogical misconstructions in a separate Postscript.

\medskip
\noindent\hspace*{2em}Copenhagen, 17 January 1868.

\begin{flushright}
R. Nielsen.
\end{flushright}

% --- PDF 10: dedication ---
\cleardoublepage
\thispagestyle{empty}
\vspace*{\fill}
\begin{center}
To my Norwegian\par
\bigskip
{\Large Listeners, Men and Women,}\par
\medskip
{\large these Lectures}\par
\medskip
{\small are dedicated}\par
\bigskip
\emph{with deep gratitude}\par
\bigskip
by\par
\bigskip
the Author.\par
\end{center}
\vspace*{\fill}
\cleardoublepage

\tableofcontents

\mainmatter

\chapter*{I.}
\addcontentsline{toc}{chapter}{Lecture I}
\markboth{Lecture I}{Lecture I}
\label{ch:i}
% --- p. 1 ---
\opage{1}\setcounter{footnote}{0}%
A thought, a great and fruitful thought, has gathered here a numerous
circle of honor. It is teachers and students at this university, old
Norway's young university, who in fair union have set the thought in
motion, the surely momentous thought for the future, that a closer
scientific intercourse among the Nordic universities must contribute in
particular to the mutual elevation of our spiritual life, the awakening
of spiritual powers. So the invitation went out, to me as well; and
since I, gripped by this fair thought, could not but prize the rare
honor that on so significant an occasion it was granted me to come
forward with my contributions and to take the word in an assembly such
as this: I knew of no other or more fitting task to choose than an
attempt at a general presentation of \emph{Hindrances and Conditions
for the Spiritual Life in the Present Age}.

Yes, in the present age the spiritual life certainly has peculiar,
favorable conditions, but for that very reason it also meets great and
peculiar hindrances. If we look away from the hindrances in order to fix
our gaze exclusively on
% --- p. 2 ---
\opage{2}\setcounter{footnote}{0}%
the conditions, then bright prospects open up; of the bright prospects
is born a glittering, fleeting hope, but a fleeting hope can beguile. If
we forget the conditions and lose ourselves in contemplating the manifold
hindrances, then the prospects grow dark; of the dark prospects comes
despondency; but despondency is like a mist before the eye. We will not,
then, give ourselves up to a one-sided consideration either of the
conditions alone, so that we thereby forget the hindrances, or of the
hindrances alone, so that we overlook the conditions; no, we shall in
all sobriety seek to weigh hindrances against conditions, conditions
against hindrances; so there comes equilibrium.

It is an old saying that he who wills the end must will the means; but
the saying has a double meaning. It can be understood thus: that nothing
can come forth in the end which does not lie in the means, that the
realized end is the result of the means, so that everything first and
last depends on the means: that is the understanding's view, that is
the counsel of prudence, and wisdom itself cannot give up prudence.
``For''---says the Holy Word---``which of you, intending to build a
tower, sitteth not down first, and counteth the cost, whether he have
sufficient to finish it; lest haply, after he hath laid the foundation,
and is not able to finish it, all that behold it begin to mock him,
saying, This man began to build, and was not able to finish.'' But it
can also be understood thus: that the means are the subordinate, the
end the superordinate; the true, self-valid goal has, in an eternal
sense, the means in its power: first the end, the first decisive thing,
the One Thing Needful, and the conditions, the means, the finite
necessities, come as of themselves! This is the demand of spirit. The
kingdom of spirit is in the highest sense the kingdom of righteousness,
the kingdom of God. In the kingdom of God, of
% --- p. 3 ---
\opage{3}\setcounter{footnote}{0}%
righteousness, prudence must serve wisdom, care for the temporal must
vanish in striving toward the eternal; for---says the Holy
Word---``seek ye first the kingdom of God, and his righteousness; and
all these (temporal) things shall be added unto you.''

The reciprocal determination of means and end just stated can now serve
to clarify for us the relation between hindrances and conditions for
the spiritual life. In active union with the overarching, all-determining
end, the finite means are nothing but favorable conditions; the sensuous
serves the spiritual by being permeated, inspirited and moved by spirit.
But if the conditions are taken in a finite sense as a sum of means
which of themselves, added together, are to make up the end, then they
stand against spirit and become for the life of spirit nothing but
snares, nothing but barriers, nothing but hindrances. From this
standpoint it cannot then be difficult to characterize the present age.

The present age is in a peculiar sense to be called the age of
prudence, an age that looks to the means and works with the conditions.
Hence this age has at once its bright and its shadow side. Our age
rightly calls itself an age of the sciences, of enlightenment and of
culture; the eye of culture is an eye for conditions; the practical
consciousness of culture understands thoroughly how to conquer reality,
increase the means and bring them into use. What has the modern age not
already achieved in this direction? The progress of the natural
sciences, these ingenious inventions, these magnificent, truly
astonishing discoveries---what fertility in means, what wealth of
conditions! That ships go against wind and current; that the steam
horse, like a snorting Leviathan, rushes through the lands from place to
place and makes the enormous distances vanish; that the express
messenger of thoughts is sent out
% --- p. 4 ---
\opage{4}\setcounter{footnote}{0}%
and reaches in the instant of thought to another world: this and more
bears sufficient witness to the power of ingenuity, to the fruit of
research; yes, the searching eye is an eye for conditions.

The lively sense for conditions and the conditioned contains naturally a
temptation to overlook the absolute, the unconditioned, and thus a
temptation to misjudge spirit. To misjudge spirit is to draw it down
under nature-bound terms in such a way that the laws of freedom simply
coincide with the laws of natural necessity. What holds of the outward
natural, the consciousness of culture will only too gladly transfer to
the inward and spiritual. By a sharpened eye for conditions we discover
the laws of nature; by the laws of nature we cognize the forces of
nature; by insight into the forces of nature we become more and more
able to master matter and show ourselves lords in the realm of nature.
It then lies near at hand to ask whether something similar could not
happen with the spiritual too? ``Should it not at last be possible for
our searching eye to discover the laws of spirit too, by the laws of
spirit to determine and measure the powers of spirit, by the powers of
spirit to master spirit itself and raise ourselves to autocrats in the
realm of spirit---on conditions, of course?'' Here is the presumption.
It is impossible for man to set himself up as lord of spirit; it is, on
the contrary, his highest destiny to be its servant.

That the present age misjudges the right of spirit is expressed, in all
indefiniteness, by saying that the age is too materialistic. In what
does its materialism consist? It surely cannot consist in its regarding
matter as something real, for in that it is of course right; nor can it
be placed in its laying weight on material goods, working to increase
the material conditions; for in that it certainly does right. No, the
materialism of the age consists
% --- p. 5 ---
\opage{5}\setcounter{footnote}{0}%
precisely in this, that it measures the eternal with the measure of
temporality, makes the absolute relative and draws the unconditioned
down under finite conditions. For sheer means it is unable to see the
goal. What one-sided weight the present age lays on the means is seen,
among other things, from the immense zeal with which, in small things
and great, by all roads and in all ways, it strives to acquire the
means of means, not the self-valid but the silver-valid end (nervus
rerum gerendarum).
% Danish pun: "ikke det selvgyldige, men det sølvgyldige Øiemed"
The power of the present age is in a particular sense the power of
capital. And yet the character of the present age is not properly
avarice. The mark of avarice is that it makes gold the goal, everything
else a means; but the present age makes gold the highest goal only
insofar as it makes it the highest means. The age of means is the age
of interests; everything has its interest insofar as it has its relative
value; everything has its relative value insofar as it can pay. Now, can
freedom pay? What does war cost? How many percent does justice yield
when it is put out at interest? If really everything is to be able to
pay, then---it lies in the nature of the case---everything must also be
able to be paid for. How much, then, does one pay for a man's honor, a
woman's virtue? By considering how consistently these questions follow
from the whole fundamental view, and where they lead, one will at least
find it explicable that poets and deeper natures regard the golden age of
interests as the iron age of spirit, despair of getting the
unconditioned along in a world of sheer conditions, and complain
despondently of the spiritual slackness of the age. It is like a sigh of
the soul in the poet's lament that the race is growing old, that youth
is already old, characters pale, that the busy individuals, interested
in freedom and independence, would much rather go with the wind, leaving
it to the ships to go against the wind. There
% --- p. 6 ---
\opage{6}\setcounter{footnote}{0}%
is a sigh of the soul in the poet's lament that the spirit of the
present age is faint and cold, that the coldness of spirit in the autumn
of life is the harbinger of winter. There is a sigh of the soul in the
poet's lament that the days of creation are shortening, that the light
gives way and the shadows grow, that ``the sun sinks lower and lower,
that the nights grow longer and longer, and the cold dark powers step
more and more into life.''

Yet what good is it to complain; selfish dejection, overwhelming
despondency, betray an offense against spirit too. If the hindrances are
great, well then, the conditions are all the richer. For hindrances and
conditions relate to one another in quite another way in the spiritual
than in the sensuous world. The more sensuously the opposition is
conceived, the more strongly will the opposed terms stand against one
another and exclude one another; the more spiritual, on the other hand,
the conception is, the more clearly will it be seen that the opposed
terms relate either as consequence to ground or as power to impotence.
In the sensuous meaning it holds, indeed, that hindrance is not
condition, condition not hindrance: the stronger the hindrances, the
weaker the conditions, and conversely; but spiritually understood it
holds that the hindrance is in possibility condition, the condition
hindrance: the greater the hindrances, the more essential the
conditions, and conversely.

We can illustrate this by some examples. Only set adversity against
prosperity, suffering against acting, fear against hope, and the
difference between sensuous and spiritual understanding will be
conspicuous. Sensuously seen, adversity stands against prosperity as
poverty against wealth, sickness against health, sorrow against joy. How
sharply exclusive the opposition is, shines out immediately from the
fact that an increase of the one term necessarily presupposes a decrease
of the other: the more adversity, the less prosperity, the greater the
sorrow,
% --- p. 7 ---
\opage{7}\setcounter{footnote}{0}%
the less the joy. In the world of spirit it is otherwise. Here
prosperity is not simply a decrease of adversity, wealth of poverty,
sickness of health, sorrow of joy; no, the relation is in all
originality quite the contrary: that intensified adversity becomes the
ground and presupposition of a correspondingly heightened prosperity;
the power of spirit reveals itself precisely in this, that it transforms
adversity into prosperity, founds a spiritual prosperity on sensuous
adversity, a spiritual wealth on sensuous poverty, a spiritual health on
sensuous sickness, a spiritual joy on sensuous sorrow. Thus speaks the
Holy Word: ``Blessed are ye that mourn; ye shall be comforted. Blessed
are ye that hunger; ye shall be filled. Blessed are ye that weep; for ye
shall laugh.'' Sorrow as sorrow is certainly not itself joy, for the joy
is in the comforting; but the joy of comforting is conditioned by the
sorrow. Here, then, there is talk not merely of a present adversity and
a coming prosperity, a present sorrow and a coming joy; but, when we
look more closely, also of a present prosperity conditioned by
adversity, a present joy conditioned by the sorrow itself. So far is it
from being the case that the prosperity thus present should have been
diminished because it coincides with adversity, that the corresponding
joy is on the contrary raised to its highest, to blessedness. Spirit
cannot, then, prove its power more clearly by transforming hindrances
into conditions than when it so wondrously refashions adversity into
prosperity that it even becomes blessed to mourn, blessed to hunger,
blessed to weep, blessed to suffer.

For spirit---we mean the absolutely divine spirit, because a spiritual
life without it is as little possible in the present age as it has ever
been in the past---for spirit itself nothing in the natural
% --- p. 8 ---
\opage{8}\setcounter{footnote}{0}%
world is simply either hindrance or condition; but everything depends
on a transformation. Therefore the divine can in a certain sense use
everything and everyone. In the kingdom of spirit no one is too weak
when spirit transforms his weakness into strength for him, and no one
too strong when spirit transforms his strength into weakness. From this
standpoint there can be no talk here of any essential difference between
good and evil, convenient and inconvenient times; every time is at
bottom equally good for the end that produces its own means, always
convenient for the unconditioned that itself creates its conditions.
Spirit in its working is by no means indifferent to conditions, not
condition-less, nor does it create the conditions simply out of nothing;
no, it forms them out of natural presuppositions, it forms them by a
transformation. Whatever in the natural states with their natural terms
lets itself be transformed becomes, however much it may look outwardly
like a hindrance, nevertheless a condition; whatever remains in its
naturalness and does not let itself be transformed becomes, though to a
sensuous view it looks like the most favorable condition of all,
nevertheless essentially a hindrance.

But, it will perhaps be objected, if there really is something in the
natural world that the divine spirit cannot transform, something that
can offer it resistance, set it a barrier and place itself as a
hindrance, then spirit is surely not almighty, and so not truly divine
either. If spirit is truly divine, and so almighty, able everywhere to
lift every resistance, break through every barrier and create conditions
out of sheer natural hindrances: why then does it not come with equal
strength to every generation in the history of mankind, why is it not
equally mighty in every age? If it needed servants who
% --- p. 9 ---
\opage{9}\setcounter{footnote}{0}%
could go before and prepare its way; if it now and then had to lie still
and wait for favorable conditions, just as a skipper lies still and
waits for a favorable wind, then we could understand it. But when it
must itself inspirit its servants and itself bring the conditions along,
is it not then incomprehensible that it will not always and everywhere
work with equally powerful presence? Here, then, there must surely in
reality exist something that can set spirit barriers, something that can
prepare serious hindrances for it and bring it about that it, though
almighty, seems to be impotent. We answer: there really is such a
something! But this something does not exist outside spirit; no, it is
for spirit itself, in its own kingdom: it lies in freedom, and the
essence of spirit is freedom.

If we grasp freedom in brief concept as the will's unity of wisdom and
power, then we grasp at once that the eternally free spirit must be
raised above dark natural necessity and blind arbitrariness. Free in
itself, spirit will work on the terms of freedom in free organs; such
free organs are the human wills. If it depended solely on using power
and ruling as a despot, then there would certainly be nothing either in
heaven or on earth able to set the almighty spirit barriers; but by
ruling despotically over free wills, the spirit of freedom would
contradict itself. Here, then, appears the singular sight that what in
the realm of creation must seem the weakest of all, the human will,
really can do what the hard rocks, the mighty mountains, the high stars
cannot: prepare hindrances for the divine will, set a barrier to
omnipotence itself. To a fleeting, superficial consideration this must
seem to us quite absurd. For what, after all, is the human will when we
look away from drive and inclination? The pure willing, the so-called
pure
% --- p. 10 ---
\opage{10}\setcounter{footnote}{0}%
self-determining, is in itself an empty shadow-being, an ideal phantom:
what can well be more shadowy than the willing that wills itself, more
empty of content than the determining that determines itself? And yet it
is literally true that pure self-determining is the only thing in the
whole world able to set a hindrance for spirit. Or if anyone thinks that
the hindrance does not originally proceed from the will in man, let him
tell us, then, whence it proceeds. We know well that our will in reality
can determine itself to something only insofar as it is determined by
something; but we also know that the will could not possibly become free
if it could not originally determine itself to let itself be determined.
The matter is not so involved as it might seem at first glance; it
clears up at once when we consider what difference it makes whether the
will is determined by drive or by spirit. The natural will is in the
beginning one with the natural drive; sensuously moved by sensuous
drive, the will reaches for the sensuous. If the drive is now further
determined as inclination, and inclination rises to passion, then the
sensuous will becomes strong by the strong passion; but the will ruled
by passion is, with all its sensuous strength, unfree. Why unfree?
Because in its naturalness it has not itself decided for the passion,
but is simply possessed, inflamed by the passion. It is an easy matter
for inclinations and passions to get the natural will into their power,
because they can, if we may say so, come upon it from behind, swoop down
on it, seize it and move it despotically. The strong passions are the
despots of the sensuous will. Already from this one sees why spirit
cannot, like drive, master the will. The spirit of freedom will set the
will free; but the will can be set free only by its itself inwardly
willing freedom. Over against drive and
% --- p. 11 ---
\opage{11}\setcounter{footnote}{0}%
inclination the will can do nothing when it is left to itself; what it
can do, it can do solely and only by the aid of spirit. If the aid of
spirit were immediate like the promptings that proceed from drive, then
it could simply come to a struggle, independent of the will, between
the free spirit and the nature-bound drive; if drive then got the upper
hand and swept the will along, it was not the will's but spirit's fault:
the divine spirit would then have been too weak. If a man has leave to
see the matter from this standpoint, he can easily excuse himself: ``How
can I help that I am born with sensuous drives and impure desires? I
cannot possibly change my own nature. If spirit wants the desires purged,
it must purge them itself; if it would set my will free, it must itself
seize and move it. I am quite willing to let myself be set free; my will
offers absolutely no resistance; spirit, for that matter, may take it,
turn it, twist it and do with it what it will, if only it will do
something. But when the divine spirit itself either cannot or will not
immediately do something, what then am I supposed to be able to do?''
Here is an enormous misunderstanding. If spirit really were to move the
will quite immediately, in the manner of drive, then it would have to
begin the will's liberation by making it unfree---what a contradiction!
The will's first act of freedom, the fruit of its very first
inspiriting, is precisely the heightened self-determination with which
it inwardly longs for and, in struggle against drive, gives itself up to
the influences of spirit. As these influences are then received and
appropriated by the will that, in appropriating, frees itself, the
relation is by no means this, that the will, in self-separation from
spirit, should have any power for any independent working whatever; on
the contrary! it can move only
% --- p. 12 ---
\opage{12}\setcounter{footnote}{0}%
after spirit, as it is moved with freedom by spirit. The moving power of
freedom is solely and only the power of spirit; but it rests with the
will itself whether it will, in inwardness, in true originality,
appropriate this power. If it will not, then it ``grieves the spirit'';
then its impotence is a negative power which, sadly enough, sets a
barrier to omnipotence itself.

And what, then, is the practical result of the considerations here set
forth? What does spirit demand of the man who works with the conditions
of reality; what sort of fundamental demand is it that it makes in
particular of this present age? Is our age worldly beyond all previous
ages? Is it renunciation of the world in great outward demonstrations,
medieval states of affairs with monkery, praying, fasting, church
processions and monastic good works, that spirit demands in order to
transform the natural hindrances into favorable conditions? By no means!
The demand of spirit is not, and by its principle cannot be, any finite
demand for anything finite and outward whatever. As to its outward
side, the present age may well be said to be, if not better, then at
least as good as any preceding one. No, the fundamental demand that
spirit makes of this as of all other times is first and last a demand
of inwardness; it demands faith, faith in its eternally free, liberating
power; it demands will, will to make oneself free by letting oneself be
set free; in short and definitely: it demands earnestness.

Every finitely determinate demand the present age, with its restless
reflections and its almost feverishly overstrained criticism, can easily
rewrite, halve and, especially when the demand is put forward singly,
easily protest. But there is one demand it cannot rewrite, cannot halve,
cannot protest, and that is the demand: earnestness!

% --- p. 13 ---
\opage{13}\setcounter{footnote}{0}%
This demand is strict enough, but it is not too strict. There may
perhaps be a race that boasts of having given up eternal life in order
to live the temporal all the more vigorously; but there is surely no
race that would in earnest boast of having given up all earnestness.
There may perhaps be born a generation that makes it an honor to lack
a sense for God and his holy Word; but could there well be born a
generation that would make it its honor to lack earnestness? When,
then, spirit in this present demands only essential and true
earnestness, it demands at bottom nothing else of the age than what the
age must necessarily demand of itself: here, then, an understanding must
be possible. To be in earnest about something is, as one says, to take
something to heart, to lay it on one's mind; by earnestness there is
posited a unity of matter and mind. But from this it follows in turn
that the mind, which takes the matter into itself and makes itself one
with it, must itself in turn become of the same essence as the matter.
If the matter is unessential, in itself insignificant, then the mind is
paltry; only a paltry mind can make much of the insignificant. If the
matter is not simply unessential but still of finite nature, then the
mind gives itself up as a prey to the finite, makes its essence finite
by taking the matter to heart and giving itself wholly up to it. The
finite is by its nature changeable; but earnestness demands, after all,
that the mind in relation to the matter be like itself, know itself
raised above the changes in the midst of the changeable. The world of
interests and means is inconstant. The earnest man must be prepared for
vicissitudes; his fear and hope are determined by his always being able
to expect the opposite: after adversity prosperity, after gain loss, and
so forth. Here the wise words fit: Sperat infestis, metuit secundis
alteram sortem
% --- p. 14 ---
\opage{14}\setcounter{footnote}{0}%
bene præparatum pectus. See, that is the earnestness of prudence.

But the earnestness of prudence cannot possibly be true earnestness; it
suffers from an inborn frailty, it drags along a heavy contradiction,
the contradiction that the mind is at once both to give itself up and
yet not to give itself up to the finite essence. Certainly the finite
has its interest as condition, as means. But if one becomes so earnest
about the means that the goal is excluded, then the self-deception is
manifest. The condition is seized; but it has significance only for
something that was to be conditioned by it. Now what sort of something
is this? Within finitude itself, a new condition! And the new condition?
Has significance only for a something that was to be conditioned by it!
The real value of the condition falls infinitely outside the condition
itself, just as the value of the means falls outside the means. He who
on these terms will live on his means and sate his soul with sheer
conditions can indeed be said to be an earnest man---earnest, yes, like
Tantalus when he would pluck the fruits and drink from the stream. The
earnestness that makes a Tantalus old and gray may well make a whole
race old before its time. If it is really true that mankind is growing
old, then it cannot be nature's fault, for of nature holds what the
poet says: ``every age and every year has its own, its blossoming
spring.'' Still less can it be spirit's fault, for spirit as spirit has
springs of eternity, and the eternal cannot grow old. The fault must
then fall on the false earnestness, the finite earnestness of interests,
of means, of conditions.

If false earnestness has made a race old, then it is to be hoped that
true earnestness can make it young again. True earnestness is a gift of
spirit, but on the terms of freedom. The condition is that the
inspirited will will receive, preserve
% --- p. 15 ---
\opage{15}\setcounter{footnote}{0}%
and use the gift. True earnestness lays decisive weight on the end, but
in such a way that it at the same time lays weight on the means. To
rave for empty ideals that fly over reality is not earnestness. The
earnestness of ideals is that they will be realized. Realization
requires exertion, endurance, constancy. But endurance and constancy
under great exertions correspond precisely to essential earnestness.
Earnestness is a power of the soul that holds out in trials, a
fundamental mood of the mind that is secured by firm resolution;
earnestness is the renewal of the self, the consecration of character,
the baptism of the will, that it is baptized with spirit. Therefore true
earnestness is not merely something with which the spiritual life is to
be completed and concluded, but also something with which it is to be
begun. Earnestness is not for the old alone, but also for the young. A
gifted youth has essential earnestness because it has essential spirit.
There is indeed something cool, hidden and severe in the depths of
earnestness, just as there is something cool in the depths of the sea;
but just as well as the deep sea can be irradiated by the sun, so can
earnestness itself be surrounded by the radiance of joy. If a sense for
the earnestness of life did not originally belong to the rich endowment
of youth, how would things then look for the future? What good, then,
would all our striving with upbringing, instruction and education do, if
that which in due time was to be applied and used in earnest were not
grasped and understood in earnest? What good would---to keep to what now
lies so near at hand---an endeavor to bring richer variety into the
academic lectures and to initiate closer intercourse among the different
universities do, if the student youth were not itself livingly
interested in all that furthers and heightens the spiritual life.

It is, then, with lively joy and a good hope that we dare
% --- p. 16 ---
\opage{16}\setcounter{footnote}{0}%
try to carry out our undertaking, because we recognize that the younger
have as much---let us just say it, fully as much---share in its carrying
out as the elder. It is the young who have willed that now it should be
earnest, and what is thus willed in earnest will, with the aid of
Providence, also be carried out in earnest. A youth that listens when
ideal strings are touched, and itself works along so that spiritual aims
are furthered, an academic youth with clear eyes, wakeful sense and
spring-fresh thoughts, will soon understand what is at stake when it is
a matter of the choice of means and goal, of the transformation of
hindrances by the power of spirit, of the victory of freedom in the
strife of life: it will understand it in earnest.

\begin{center}\rule{3cm}{0.4pt}\end{center}

\chapter*{II.}
\addcontentsline{toc}{chapter}{Lecture II}
\markboth{Lecture II}{Lecture II}
\label{ch:ii}
% --- p. 17 ---
\opage{17}\setcounter{footnote}{0}%
It is told of Alexander, son of Philip of Macedon, that as a youth he
wept when his father had performed a great deed. They were tears of
admiration the youth wept, for he looked up to his father with filial
reverence; he saw in him the renowned hero. But they were also tears of
envy; ``for,'' said he, ``my father performs all the great things; what
is left over for me?''

This so characteristic relation between the youth Alexander and his
father can in a certain sense---from one side at least---serve to
illuminate the successive ages of history. \textit{Le présent est gros
de l'avenir}, says the philosopher Leibnitz, and rightly; the time
immediately following is, if I may so express myself, always a son of
the preceding one. But if the present age has the preceding one for
father, and so, to stay in the image, the one before that for
grandfather, then, when we look back over the long series of
generations, over great-grandfather, great-great-grandfather,
great-great-great\ldots{} and so on, there results a mighty family tree.
The generation now living is a noble race with many ancestors and great
fathers; when, then, the present age turns and contemplates the past,
its deeds and heroes, it must involuntarily exclaim: ``My
% --- p. 18 ---
\opage{18}\setcounter{footnote}{0}%
fathers have performed the great things!'' and could then easily be
tempted to ask with Alexander: ``What is left over for me?''

It is not, earnestly seen, youthful ambition, not envy, still less empty
vanity, that breathes in such a question; it is on the contrary a quiet,
deep concern; it is spiritual anxiety for one's livelihood. It does not
help that someone will say to us: ``Where are you going with your
asking about the great? If, like Alexander, you have the great powers,
well then, like Alexander you will yourselves know how to perform great
deeds; if not, then you do best to keep to the ground, learn to crawl
until you can walk, and make do with the lesser.'' This rebuff is
unwarranted. For our own part we will gladly content ourselves with the
lesser; but in order to have the lesser with spirit and truth, we must
necessarily be helped by the great. If the great so vanishes from a
given contemporary age that it no longer has the power of presence,
then it becomes a spiritual torment to live in the lesser. The past must
indeed furnish presuppositions for the present; but the present wastes
away when it must feed exclusively on a pastness; every age must have
sources of nourishment in itself: the present lives on a present and the
lesser on the greater. It goes with the movements in world history as
in the world ocean. If wind and storm and ebb and flow and strong
currents did not ceaselessly move the waters in the great sea, how would
things then look with bays and fjords?

It is a true saying that no one can live unless he has something to
live on and something to live for. Sensuously the two demands can be
separated, but not spiritually. In material respects the first necessity
is undeniably this, that he who is to live must have something
% --- p. 19 ---
\opage{19}\setcounter{footnote}{0}%
to live on. Yes, the urgency of this demand is even so conspicuous that
many take it for the only valid one and think to themselves: if only we
had enough to live on, we should easily find something to live for. But
in the spiritual sense the two demands are inseparable: what a man
spiritually lives for, he also lives on. An artist, for example, who
lives exclusively for his art also lives---certainly not always
materially, but still ideally---on his artistic production, has in this
production his spiritual food, his life, his joy. He who unconditionally
lives for faith lives spiritually on faith; he who lives for knowledge,
on knowledge, and so on. In the ideal as well as in the material it is a
matter of working capital.

In the commercial world, and generally in the practical world of
workers and of civic interests, watch is kept that the great capitals
are not scattered and divided into mere small change. It is as much in
the interest of those without means as of those with means that the
capitals be preserved. So also in history. Capitals of spirit are
needed for every age to have something to live for and to live on; at
critical turning points great stakes must be laid down. The great stakes
mark the epochs of history; the employment and distribution of the
capital in greater and lesser, until it has given its yield and is used
up, is what characterizes a period. The greater the epoch, the greater
is the life-capital staked, and the longer will the corresponding period
be too. The absolute epoch of spirit is, in the language of religion,
the turning point called the fullness of time, the turning point when
not merely one of the great thoughts of eternity but the Eternal
himself, he who says: ``before Abraham was, I am,'' revealed in flesh
and blood, stepped into time. Seen from this summit, the Word of Life is
the spiritual fundamental capital on which the generations
% --- p. 20 ---
\opage{20}\setcounter{footnote}{0}%
are to live; for this epoch all following times are to be regarded as a
single, continuous period.

It is otherwise, on the other hand, with the relative epochs; to these
belong naturally only relative periods, longer or shorter in proportion
to the greater or lesser capital staked. Such relative epochs we have in
the ideal breakthroughs on which the division between antiquity, the
Middle Ages and the modern age originally rests; as regards the life of
spirit they may on the whole be said to represent: the folk-spirit, the
church-spirit, the world-spirit. It is necessary in this connection to
bethink oneself of the distinction between the relative forms under
which spirit works in human life, and the divine spirit itself. In the
history of culture spirit does not work immediately on man's thoughts
and powers; here there is an inner intermediary between spirit absolute
and human self-consciousness, an intermediate being in which spirit
each time deposits a peculiar content, determined by a peculiar form;
this intermediate being we call the idea. Only as man's consciousness,
in mind and thoughts, in willing and working, appropriates the idea can
the human spirit receive its peculiar stamp.

By bringing out the difference between idea and spirit we can at once
characterize the opposition between revelation and mythology. In
revelation the idea is so perfectly transparent that in many cases it
looks as though it had quite vanished. God speaks to man, teaches man,
places himself over against man as will against will, spirit against
spirit. At the stage of mythology the matter stands the other way
round; here the idea as idea has become opaque to the human spirit; the
proof of this is that the ideas themselves are set up in figurative
form as gods: the mythic gods
% --- p. 21 ---
\opage{21}\setcounter{footnote}{0}%
are personified ideas. The opposition between Jews and Greeks is in
this respect characteristic; the Jewish people is the people of spirit,
the Greeks the people of ideas. Among the Jews everything aims at
satisfying spirit; therefore the humane ideas of art and science could
not come to development; not even the state as state got leave to
develop out of its own idea: the theocratic state is immediately
subject to the power of spirit. Otherwise among the Greeks. Here it is
so far from being the case that spirit overwhelms the idea, that the
content of life in all directions
on the contrary develops out of the idea, is formed and individualized
into a peculiar idea-content. Herein lie the strong impulses to the
development of art and science, politics and rational ethics; but from
this, then, is explained also the fundamental defect, that the
supernatural everywhere melts together with the natural, the will with
drive, freedom with necessity, in such a way that not a single demand
of spirit comes to hold in absolute unconditionedness.

As the most gifted cultured people of antiquity, the Greeks furnish a
striking example of what a folk-spirit signifies in its originality,
what inner limits it has, how it struggles with hindrances and
conditions. In every gifted people there stirs from the beginning a
circle of peculiar ideas with the force of nature. How characteristic
such ideas are of the national peculiarity leaps to the eye when we
compare Greek mythology with Nordic. Zeus and Odin, Frigga and Here,
Helheim and Hades, Olympus and Valhalla---what a difference! The
difference is twofold: partly outward, insofar as the gods' form,
individuality, deeds and fate are depicted otherwise by Greek fantasy
than by Nordic; partly inward, insofar as each of the different worlds
of gods
% --- p. 22 ---
\opage{22}\setcounter{footnote}{0}%
presupposes its own degree of self-consciousness, its own view of the
world and of life in the people concerned. The Greek gods are gods of
beauty, the Nordic gods of battle. This must of course not be
understood as though the Greek gods could not fight and the Nordic
could not be beautiful; but the point is that the Greek have their most
perfect manifestations in beauty, the Nordic in battle. How far-reaching
this opposition is in determining the peculiar spiritual direction of
the two peoples and their role in world history can be illustrated by a
couple of traits.

That a deity's manifestation is completed in beauty means, in other
words, that its peculiarity, its ideal independence, lies essentially in
aesthetic individuality. But when individuality in its immediate
phenomenon is the highest expression of the independence of the gods,
then it must also be so for that of men; for the divine is the pattern
for the human. Homer will, he says (Iliad. I, 1---5), have the goddess
sing the wrath of Achilles, which sent the souls of many heroes to
Hades, but made them themselves a prey for the dogs
(αὐτοὺς δὲ ἑλώρια τεῦχε κύνεσσιν). Certainly we must at this ``them
themselves'' involuntarily think of the corpses; but what is
characteristic in the expression is still this, that what designates
the heroes themselves is not their souls but their bodies. The bodily
form is, namely, the aesthetic individuality of the hero; to have
individuality is common to gods and heroes. But now comes the
difference: the gods are immortal, men mortal; by the division between
immortal and mortal races fate has set a tragic barrier between gods and
heroes. Forsaken by bodily individuality, the bloodless shades in Hades
drag out a pitiful existence. It is said of Hector's death (Iliad. XXII,
362): ``Out of his limbs flew his
% --- p. 23 ---
\opage{23}\setcounter{footnote}{0}%
soul and hastened to Hades, wailing over the lot of parting from valor
and youth.'' Otherwise with the heroes in the old North; there another
lot awaited them as Einheriar in Valhalla, and why? Because Odin was so
different from Zeus, the gods of Valhalla different in mind and thoughts
from those of Olympus. Odin and Thor did not place the personal essence
in beautiful individuality but in the strength of will that lives on
battle and grows in battle; it is such a strength of will, which the
heroes have in common with the gods, that ennobles the deed and
overcomes death. That death triumphs, that the light is quenched, and
the souls tremble in the realm of shades, is not the fault of fate; the
fault is in cowardice and the weak will. Therefore Thor is ``so wroth
with the fearful men of Helheim,'' and addresses them harshly: ``You
wretched fools, who feared wound and death! Now Hel shall wound you
forever with torment and bitter need.''

From the difference between the aesthetic individuality and the
battling character in the mythic patterns there flows yet another
consequence, which is truly characteristic of the role each of the
peoples named was to play in history. The ``blessed gods'' of Olympus
could certainly not---it lies in the nature of individuality---break
through with the clear, completed form and get the better of the old,
half-formless elemental gods sprung from chaos, except after great
upheavals and hard battles; but while the Æsir, the strong-willed
battle-gods of Valhalla, must wage a constant war with the giants, a war
that begins with the slaying of Ymer and ends with Ragnarok, the Olympian
individualities are finished in a comparatively short time with the
whole war of the gods and so get occasion, in plastic beauty and
classical repose, to enjoy all the pleasure of life with the blessed joy
of ideality:
% Danish prints "Ro et nyde" (printer's defect, et for at); translated as "to enjoy".
``Here the heart lacked not an ample meal, nor the lovely lyre struck by
Phoebus Apollo, nor the alternating songs
% --- p. 24 ---
\opage{24}\setcounter{footnote}{0}%
of the Muses' beautiful voice.'' While the Northmen with their divine
leaders are so caught up in the whirl of endless battle that they could
only begin to come to rest in the advanced Middle Ages and develop any
civilization on the foundation of Christianity, the Greeks reach,
through comparatively brief crises and grandiose catastrophes, the
turning point of the time of rest, when the ideas are clarified and the
material formed, so that not only the laws of the state but also those
of thought and beauty are fulfilled. It is now easy to see that the
folk-spirit determines the actual life of the people, that its
historical rise, culmination and decline are inwardly conditioned by the
ideas at work in the spirit itself. Greek history is in a peculiar sense
to be called catastrophic. The catastrophic is suited precisely to an
all-round development of aesthetic individuality, and gives us
conspicuous examples of how a folk-spirit first proves its power by
transforming hindrances into conditions, and then again betrays its
weakness when the conditions turn about and are transformed into
hindrances. Greek individuality in its heroic immediacy is combative;
already in the struggle for Troy the strife between Achilles and
Agamemnon is of great catastrophic effect. Had the combative
individualities now got leave to tumble about freely among one another,
then the state, this universal being, could not possibly be formed and
come to be at one's disposal; but just as Plato's God at a glance upon
chaos originally saw that order is better than disorder, so the Greek
eye saw that disorder was in equal degree contrary as much to the
beautiful as to the good. But as the individuals bowed under the state
as under a deity, they still did not bow so slavishly that the stamp of
individuality was effaced; on the contrary! the individualities were
taken up into the state precisely in such a way that the essence of the
state thereby itself became individualized. Instead of one single
% --- p. 25 ---
\opage{25}\setcounter{footnote}{0}%
general Greek state there forms itself from the beginning a plurality
of individual states, mutually at strife in accordance with their
immediate, and so likewise aesthetic, individuality.

From the ground of the rich individualities the springs of ideas well
up; here, then, everywhere in the outward as in the inward, in the
colonies as in the mother states, impulses, powers, elements are in
motion enough. If, however, something whole and great was to come out of
this motley manifoldness, something on which the spirit of beauty could
set its stamp, and so something that was in truth of one casting: then
above all an enormous pressure had to be brought about; a condition was
needed that carried weight. And the condition came; but at the moment it
came, it certainly by no means looked like a condition; no, it looked on
the contrary like a crushing hindrance. It is Xerxes who with his
hundreds of thousands got the task, the fair historical task, of
procuring the condition for the rich, wondrously strong development of
the Greek folk-spirit. And Xerxes must indeed arm himself and hasten; for
the condition is pressing. In the bay at Salamis strife and discord have
already prepared so many hindrances for Themistocles that he knows how
to do only one thing: transform these hindrances themselves into one
single great, decisive condition. It is natural, after all, that
individuality and freedom must correspond to one another. The
aesthetically strong individuality is inflamed by enthusiasm for freedom
when the danger outwardly impends; in the moment of enthusiasm it has an
enormous elasticity; but it is also a matter of using the moment; and
the Greek heroes used the moment: Marathon, Thermopylae, Salamis and
Plataea speak of it to all generations. The Greek struggle for freedom
was an enormous catastrophe; the blow that was to have crushed
individuality brought about
% --- p. 26 ---
\opage{26}\setcounter{footnote}{0}%
an uprising whereby it came to unfold powers that have astonished the
world and surprised itself. It is in this struggle that the capital is
staked, the life-capital that was to yield great spiritual, incalculable
interest.

Greek political life, Greek poetry and above all Greek art would surely
never have reached the summit at which a pure, strict and completed form
is in so originally simple a manner enlivened and filled by the content
of the idea that the latest time will look upon it with wonder, if the
catastrophic struggle for freedom with its terrors and upheavals had not
brought hindrances with it that spirit could transform into conditions.
It was, it is said, a Greek saying that he could not be blessed who had
not seen Zeus at Olympia. What, then, was there to see at Olympia? The
transfigured individuality in manly youthful strength, in might and
majesty, in repose and loftiness, the ``father of gods and men,'' whose
lightning had struck the Persians as it once struck the Titans. We read
in Homer of the holy sign with which Zeus confirms his promise to Thetis
(Il. I, 528): ``He spoke, and Zeus nodded with his dark-blue brows
therewith, and from the immortal head the lord's ambrosial locks fell
waving down, and he made high Olympus tremble.'' In reading this we get
an inkling of the vision of the god that must have hovered before
Phidias in the moment of enthusiasm; but for us---who do indeed hope to
be blessed without such a vision---it is only as an inkling; for the
Greek artist, on the other hand, it was real and living presence. It was
the Greek folk-spirit that revealed itself in the vision, and that with
a clarity and immediacy which bore witness that it had now really come
to full consciousness of its divine power. If we might not presuppose
this, where
% --- p. 27 ---
\opage{27}\setcounter{footnote}{0}%
should we then seek the condition for what history calls ``the great
art''?
% Danish ends this sentence with a full stop although it is a question.

Upon the catastrophic development with its rapid flowering there
followed, under a likewise catastrophic discord (the Peloponnesian War),
a correspondingly rapid decline, withering and wasting away. When
Demosthenes thundered against Philip, the lightning in the hand of Zeus
was already faint; the powers of spirit were exhausted, for while states
and tribes weakened one another in mutual strife, the idea bound to the
classical forms exhausted its original productivity. It was a tragic
conflict between the ethical and the political in which Greek freedom
perished: ethically Demosthenes was indeed right against Philip; but
that Philip, politically seen, was right against Demosthenes, his son
Alexander proved.

In the Middle Ages the folk-spirit is not so much permeated and
essentially reborn as rather immediately overwhelmed and systematically
ruled by the church-spirit. What, now, is the church-spirit? By no means
one with the original spirit of the congregation, but a derived and yet
mighty being, an ascetic ruling spirit disciplined by dogmas and
traditions, a Greco-Roman spirit, proceeding not from the Father and the
Son but from bishops and councils, hermits and monks, cardinals and
popes. In the church-spirit the immediate folk-spirit is reduced to a
moment, the natural cowed and transformed into a foundation. The
opposition in principle between the flesh and the spirit, the profane
and the holy, the earthly and the heavenly, the present and the coming,
grips feeling, kindles fantasy, and breaks out in all directions, at all
points and under all forms, into mighty, visible dualism. Between the
strictest, most unconditional demands of spirit on the one side and
worldliness with its sensuous drives, its raw lusts and strong
% --- p. 28 ---
\opage{28}\setcounter{footnote}{0}%
passions on the other, the ideas surface anew, and an enormous
fermentation sets in. By the fermentation of the ideas the divine becomes
worldly, the worldly divine, and yet the opposed poles, despite mutual
attraction, go on repelling one another. The church becomes a state, and
the state becomes churchly, and yet they go on, the church as the kingdom
of God and the state as the kingdom of the world, counteracting one
another amid ceaseless strife and reconciliation. The clash between
emperor and pope, so characteristic of the dualism of the Middle Ages,
begins from the moment when Leo the Third anointed and crowned Charles
the Great. It is not long before the clash becomes so strong that the
pillars of society are shaken. While the Carolingian political system
was splitting apart, and the many-sided alternation of private lordship
and personal dependence that characterizes hierarchy and feudalism
confused the concepts of right and gave individuals up as a prey to
violence and arbitrariness, the Northmen stormed from the north, the
Magyars from the east, the Saracens from the south, and completed the
distress and misery of the time. The peoples, whose eyes the new faith
had opened to the blessings of peace and righteousness, and who had
just taken their first look into civilization, regarded the situation in
the light of religion, were seized with terror at the vanity of the
world and awaited the Day of Judgment.

With what, now, is all this to end? Not with the end of the world but
with a victory of civilization, an advance of culture: the strong
fermentations are the birth pangs that bring a new time with a new form
to the light. The fermentation is a ceaseless conversion of hindrances
into conditions and conditions into hindrances. The fermentation itself
cannot produce anything actual; but it can, if I may say so, knead and
duly prepare the possibilities. What is waited for is a great
breakthrough,
% --- p. 29 ---
\opage{29}\setcounter{footnote}{0}%
a mother-idea which, by the spirit of the church, conceives and bears
the many daughters. With this idea the stake is then laid down; its
dowry is the life-capital that the Middle Ages have long been gathering,
and on which they are long to live. The turning point is that there
appears a goal of spirit which can take up into itself and embrace all
other aims, that a work is found out which can enclose and include all
other works, that an enthusiasm is awakened so rich and strong that its
effects can extend to everything that draws nourishment from the ideal;
in short: it is the Holy Sepulchre that is to be conquered.

The idea of the Crusades may with good reason be called the fundamental
idea of the Middle Ages; the goal set hereby was at once so earthly that
it could be grasped with hands, and yet so supra-earthly that it stood
immediately as a gate to heaven. However many good works might be
enumerated, there was nevertheless now no one among the Christians of
the West who went wrong when the question was of the first and
greatest. By the Crusades the cause of the church first became in
earnest a cause of the people: by the cross of Christ against the moon
of Mohammed the free men of the West would measure themselves in battle
with the slaves of the East; by the Crusades the emperor and the princes
first came into the right position, in the church's view, toward the
pope; by the Crusades the monk and the knight got a common task to such
a degree that the monk became a knight and the knight a monk; by the
Crusades the social life, otherwise so particularized, divided into
provinces, estates, towns and corporations, was flooded through by one
feeling, moved by one thought. Adventure and devotion, calculating
self-interest and mystical piety, ambitious striving and romantic
longing, outward-striving sensuousness and ascetic inwardness, in short:
promptings of the most varied origin, endeavors sprung from the
% --- p. 30 ---
\opage{30}\setcounter{footnote}{0}%
most heterogeneous sources, flowed together here into one single
enormous stream. There can surely be named no direction of spirit, no
manifestation, either ideal or material, of the life of the Middle Ages,
be it in trade or in science, in industry or in poetry, in religious
ceremony or in knightly gallantry, which did not, mediately or
immediately, stand in connection with the Crusades.

The idea of the Crusades is an idea of dualism, an idea that unfolds the
essence of the Middle Ages. When the time of the Crusades was over, the
pope took a fever; it was over with his sound constitution, his power
was essentially broken. It is not without significance that Dante begins
his descent into hell, and with it his Last Judgment upon the Middle
Ages, in the very same year in which Boniface the Eighth gathered the
devout in crowds about the graves of the apostles and, by instituting the
jubilee year, founded the traffic in indulgences on a grand scale. Just
as crusade and crusade and again crusade was a sign of the church-spirit
upon the vigorously battling Middle Ages, so now, after St. Peter's star
has for a time stood at the zenith, indulgence and indulgence and again
indulgence becomes a clear sign of the sinking of the Middle Ages.

To acquire the forgiveness of sins by fighting as a crusader is still
always to stake life in order to win life; but to barter oneself the
forgiveness of sins for finite values, yes, for sausages and hams and
ready money---that is more than mysticism, that is materialism itself,
the most brazen that has been seen in history.

How hindrances by the power of spirit transform themselves into
conditions is now seen clearly enough in the Reformation. It was not
because the good and uncorrupted in the papal church still had the
preponderance over the bad and corruptible that the thought of the
Reformation at last broke through; no, it was on the contrary because the
eruption of corruption was so
% --- p. 31 ---
\opage{31}\setcounter{footnote}{0}%
strong, so conspicuous, that it could not be hidden by any hierarchical
artifice. What it means that a central idea takes up, fertilizes and
moves a system of relative ideas, that the spiritual life is awakened,
reborn and furthered as spirit, through the end, shapes the means---this
we learn from the course of the Reformation. For it is true, indeed,
that the reformatory movements on the whole signify an uprising of the
nationalities, a reaction against the foreign yoke, a revolt of the
states against clerical despotism and against the extortions from Rome;
but it nevertheless stands firm that this national and
political-economic uprising would not have led to any decisive result,
still less have gained any lasting significance for the life of spirit,
had it not been taken up as a relative moment into the absolutely
religious fundamental thought. It is true, indeed, that the intellectual
weapons with which the Reformers fought were for a great part fetched
from the well-stocked arsenals of humanism and of the reawakened
scientific spirit; but it nevertheless stands firm that neither the
classical studies nor the humanist-critical reflections would by
themselves have been sufficient to steer the way of thinking of the age,
least of all to fill up the hollowness and water the roots of the tree
of life with streams from a rich spring. It is true, indeed, that the
Reformation neither would nor could give anything originally new in the
religious sense; it could and would only purge out the harmful by
exposing misinterpretations and removing arbitrary additions; but it
nevertheless stands firm that the thought of the Reformation was a
thought fertilized by spirit itself and for that very reason fruitful,
so fruitful and strong that it lifted a generation by the power of
enthusiasm and refreshed the blood in all its veins. With new
awakenings of heroic self-sacrifice, with new thoughts in new ideas, the
capital was then staked, the life-capital whose rich means would have
given the life of the people a lasting upswing, had these means not
% --- p. 32 ---
\opage{32}\setcounter{footnote}{0}%
in sorry measure been sacrificed to dogmatic formulas, scholastic
doctrinal edifices and absolutist state churches, until the capital at
last dwindled away and was consumed in the Thirty Years' War.

With the Reformation the thought of faith came into its right insofar as
it got its Gospel cleansed and its account settled with the Middle Ages;
but a clear development of the principle of faith as a peculiar
principle of cognition in its own right was not yet carried through. In
the midst of the endless disputes over the testimony of Scripture and of
the spirit, the scholastic \textit{Credo ut intelligam} still preserved
its whole ambiguity, an ambiguity that has had not only its dogmatic
but has even to this day its ethical, far-ramifying consequences.

As the original idea of the Reformation grew dim, as the opposition in
principle between the kingdom of God and the kingdom of the world---an
opposition about whose muddling Luther himself exclaims: ``der
lebendige Teufel hört ja nie auf, mir diese zwei Reichen in einander zu
kochen und zu brauen''---was more and more blurred in existence itself,
it became evident too that the great Reformers, under the spiritual
terms of their age, had had to leave the rational question of faith and
knowledge, the greatest of all intellectual tasks, standing unsolved.
The consequence of this was that while in orthodox Protestant dogmatics
reason was constantly taken captive under the obedience of faith, in
Protestant politics faith began more and more to be taken captive under
the obedience of reason. For the sake of a human order, and so for the
sake of reason, churchly orthodoxy was made a civic concern. That
princes and governments, confusing their humane respect for the holy
cause of the Gospel with living faith in the Gospel itself, took it upon
themselves to make the government of the state churchly and to provide
% --- p. 33 ---
\opage{33}\setcounter{footnote}{0}%
for their subjects' salvation by churchly laws of compulsion, looks
indeed, at a fleeting glance, as though it were a practical carrying out
of the principle of taking reason captive and letting faith rule; but it
is an illusion of the senses. The world of politics is a world of
finite interests; in the world of finite interests, it follows of
itself, the finite considerations must rule; by being drawn down into a
ruling finitude, faith as faith becomes not goal but means. To use faith
as a means is to take it captive. When, then, in the orthodox state
churches it looks as though faith sat on the high seat and wielded the
scepter of government: it is not the free faith that rules over a
worldly reason; no, it is the captive faith, whose apparent exaltation
is all the more degrading as it does not even get leave to show its
captivity, but must, as a cloak for the reason that in reality rules,
figure as mistress.

The hidden contradiction was already exposed in the further course of
the Thirty Years' War, particularly by the Peace of Westphalia; but as
yet only for the clear-sighted, the properly initiated, the diplomats.
But diplomats must be able to keep silent; the diplomatic secret was
betrayed only gradually, as the more enlightened came to the
cognition that it can only be worldly goals that are immediately
furthered by worldly means. So the way of thinking changed, and, without
one's being able to point to any finite why, the period of the wars of
religion was at an end. As religious concerns were pushed more and more
into the background, the demands of reason, the theoretical as well as
the practical, naturally pushed themselves forward into the foreground.
If faith, in the form of objective, blind authority, had formerly
wanted to tyrannize over reason, the emancipated reason, through
philosophy,
% --- p. 34 ---
\opage{34}\setcounter{footnote}{0}%
criticism and politics, amid ever stronger clashes with the
established, reached at last the turning point when the thought of
freedom broke through: the epoch-making breakthrough was the Revolution.

The spirit of the Revolution stands in decisive strife with the
church-spirit; they are each other's sworn enemies. The revolutionary
spirit of freedom, however great the passions it roused in its day
among the peoples, is by no means a national spirit. What it has in
common with the natural folk-spirits is indeed this, that it identifies
the idea with the deity; but while the folk-spirit in its first
naturalness is grounded in mythic presuppositions and wastes away when
its chosen gods grow pale, the spirit of the Revolution is on the other
hand so reflective that it makes world-reason itself the deity; it is a
cosmopolitan spirit, which knows with itself that it neither can nor
will strive for other aims than such as belong to the finite-infinite
world of conditions. The spirit of the Revolution is the world-spirit
awakened to self-consciousness, a practical spirit whose bright and
shadow sides are chiefly to be sought in politics.

Seen from its bright side, this political spirit is a spirit of the
rights of man, and so a spirit of humanity; its watchword is liberty and
equality; its earnestness is that it will, if only after many struggles,
through many hindrances, by many detours, put the rights of man in force
and make civil liberty an actual truth. It is a characteristic trait
that when the Constituent Assembly in Paris (1789) was discussing the
rights, and Mirabeau, impatient that the discussions were taking too
philosophical a turn and dragging on, exclaimed: \textit{N'employez pas
le mot de droits, mais dites: Dans l'intérêt de tous, il a été déclaré
\dots}, the assembly held fast with the sureness of instinct
% --- p. 35 ---
\opage{35}\setcounter{footnote}{0}%
to this, that it was precisely rights, and so expressly demands of
principle, that were to be carried through.

But precisely in the demands of principle one sees an enormous strife
and yet a remarkable agreement between the spirit of the Revolution and
that of the Reformation. Can the goal be reached, and the rights of man
be put in force, if the religious principle of 1530 is rejected? We
answer unconditionally: No! Can the goal be reached, and the rights of
man be put in force, if the political principle of 1789 is rejected? We
answer with full conviction: No! When, now, the two principles are
nevertheless so mutually heterogeneous that an immediate and simple
agreement between them has hitherto proved to be an impossibility, then
there is certainly set before the nearest future so significant a
problem, a task so decisive for the temporal as for the spiritual life,
that an anxious present age, which sees with admiration that the fathers
have accomplished the great things, shall not waver irresolutely and
ask: What is left over for me?

\begin{center}\rule{3cm}{0.4pt}\end{center}

\chapter*{III.}
\addcontentsline{toc}{chapter}{Lecture III}
\markboth{Lecture III}{Lecture III}
\label{ch:iii}
% --- p. 36 ---
\opage{36}\setcounter{footnote}{0}%
When rumors of war are heard, the Exchange grows restless; when the
northern lights are kindled, the magnet grows restless; when the elements
of spirit are stirred, and thoughts are roused to strife, the mind grows
restless. Just as one can speak of a storm of war, of a magnetic storm,
so one can speak, and with good reason, of a storm of the spirit. If we
look at the sky when thunder is gathering, we see the clouds moving,
some with the wind, others, as it seems, against the wind: whence come
these opposite movements? They come from the attractions of opposite
poles, from opposite currents in lower and higher layers of the air. The
outbreak itself surprises us all the more, the more hidden and quiet the
preparation has been. So too with the spiritual storm! The preparation
takes place in this way, that opposite ideas, heterogeneous thoughts, in
mutual repulsion and attraction, move with rising tension in opposite
currents through lower and higher regions. In the beginning the
movements may be so hidden, the whole stir smoldering so obscurely, that
conscience dozes, and immediate attention notices nothing. But under the
increasing tension a restlessness is felt; what formerly expressed
itself as fleeting assertions, flashing fancies, abstract representations
indifferent to the practical
% --- p. 37 ---
\opage{37}\setcounter{footnote}{0}%
world, gradually becomes importunate, spreads ever further and gains
more weight; and still the whole, seen from without, is like a riddle.
At last the critical turning point is reached; the contending ideas meet
in strife: there is a collision, a spark is kindled, a lightning flash
gleams, and a trembling passes through souls.

With the religious breakthrough of the Reformation and the political
breakthrough of the Revolution as its presupposition, the present era is
continually moved in two diametrically opposite directions, subject to
two spiritual powers, of which the one
% Danish prints "aandige" (sense "aandelige"); translated "spiritual"
unceasingly works against and combats the other. The clashes are not, to
be sure, equally palpable at all times; on the contrary! there are
intervals in which the spiritual strife almost wholly dies away, and a
tacit agreement seems to herald a peace, if not exactly eternal, yet
long-lasting. But how little one can rely on that sort of peace, the
history of past ages, with its epochs and periods, must surely have
taught us sufficiently. It goes with the resting conditions in history as
it goes in volcanic regions. The volcano is at rest; not, indeed,
without all activity, for vapors and columns of smoke do remind us that
there is something stirring in there. But the gorge is closed, and the
crater itself, with its heaps of lava and ash, a hollow one can walk
upon. Then the whole surroundings look peaceful, inviting and reassuring:
the lava hardens, the vine grows from the warm ground, and men build at
the foot of the mountain and tell of old days, when the earth split open
and the crater spewed fire. But that is long ago, the times of eruption
are past; now there is no danger: thus a peaceful present, even if it is
pregnant with the terrible, always has for the immediate sense something
bewitching, indeed something stupefying. For it smolders in there; it
burns
% --- p. 38 ---
\opage{38}\setcounter{footnote}{0}%
in the depths: suddenly, before anyone suspects it, a rumbling is heard,
and the earth quakes, and the flame breaks out, and the catastrophe is
there again with all its terrors.

It goes with the calm conditions in the world of spirit as in the world
of nature; what the calm and the equilibrium signify, and how long it
will last, depends on how, under what presuppositions and on what
conditions it has been brought about. The equilibrium that at the
present time apparently obtains between the principles of 1789 and the
articles of faith of 1530, and in general between the religious and the
humane endeavors in the civilized world, may rightly be called an
unstable equilibrium. It has by no means proceeded from any true mutual
understanding between the contending powers. Far from it! When decisive
questions are raised, the principles still stand, in this century, just
as sharply and irreconcilably against each other as at the end of the
last. How, then, has the apparent equilibrium been brought about? By a
changed mood, by a finite agreement guided by practical considerations:
the age of conditions, of means, of interests, understands admirably
well how to deal in practical considerations and finite agreements. The
Revolution went too far; its bloody consequences soon showed themselves
to be revolting inconsistencies. It overthrew thrones because it tore
down altars; that was abominable: then the sense for the despised
religion was awakened again; the need for the faith of the fathers
coincided with the need for civil order: then the prince of armies
himself had to concede that a people cannot possibly live without
positive religion. With this the power of the Revolution had for its
part acknowledged religion, though, of course, not as an unconditional
and self-valid end, but as a higher political means. On the other side,
the champions of religion, in persistent, though by no means victorious,
% --- p. 39 ---
\opage{39}\setcounter{footnote}{0}%
strife against ``sound reason'' and negative knowledge, had at length
learned that there must be a civil right and standing in the state
independent of ecclesiastical confession of faith, and now found
themselves referred chiefly to working on people's hearts by means of
feeling, fantasy and romanticism. After some violent scenes, during
which a fanatical reaction, in order to quench the last revolutionary
flames, had given ecclesiastical fury free rein, prudence at last gained
so much of the upper hand that the contending parties found it advisable
to tolerate one another.

What in the meantime has already been done to kindle the light of faith
again and awaken the congregation that had fallen asleep, we shall
appreciate with full acknowledgment; but we must expressly add that what
has been achieved in this respect has been achieved, not by the
scientific way, but solely by religious means. The theologians have, to
be sure, in their own fashion endeavored partly to ward off the attacks
of science on religion, partly to mediate between faith and knowledge;
but that these endeavors, far from benefiting the cause of religion,
have done it irreparable harm, of this the literature of the last
century has furnished a proof that cannot be contradicted.

It is not the relation between church and state or between religion and
politics, no, it is the relation between religion and culture, more
precisely, between religion and science, between the believing and the
knowing consciousness, in which the knot was originally tied. As a
contribution toward illuminating the cardinal point around which the
spiritual life of the present age, with its hindrances and conditions,
seen from the intellectual point of view, must be said to move, we shall
in a few strokes indicate the principled opposition between the
believing and the knowing consciousness, between the consciousness that
will appropriate the spirit of the Gospel
% --- p. 40 ---
\opage{40}\setcounter{footnote}{0}%
and the consciousness that is filled and moved by the spirit of the
world.

If we ask ourselves wherein man's advantage over the animal must
originally be sought, the answer lies close at hand: it must be sought
in consciousness. The animal's self-feeling loses itself in the sensuous,
the momentary; its arbitrariness is led by instinct and drive; its
understanding and memory are so taken up with sense-images and shifting
representations that it cannot possibly get hold of itself, make its
own being the object of a special intuition, or, in other words, come to
engage in any self-contemplation whatever. If a horse, for example, or
the dog with the clever eyes, could only a single time think to itself
or say to itself in its own language: \textit{I am an animal}, then such
an animal would have an I, and the boundary between animal and man would
be effaced. Man's consciousness is, in contrast to the animal's,
expressly self-consciousness; herein lies this, that man has what the
animal does not have, a free, personal self, that man is what the animal
is not, a spiritual being, an I. It is now easy to understand that no
human being can live with an entirely empty consciousness: consciousness
must be filled and moved by a content; as the content is, so is the
consciousness. The human being whose consciousness is filled up with
sensuous representations, sensuous desires and temporal interests is a
crude, sensuous and spiritless human being. Only when consciousness takes
up and appropriates a content of an essential, universally valid nature,
a distinctive content of ideas, so that the self can be said to live in
it, to live from it and to live for it, does man become spirit. The idea
by itself is not yet spirit, any more than the self by itself, without
an idea, is spirit: spirit as spirit is the unity of the idea with the
self. From this it follows that spirit and consciousness must determine
each other mutually: as the spirit is,
% --- p. 41 ---
\opage{41}\setcounter{footnote}{0}%
so is the consciousness; as the consciousness is, so is the spirit.

Relying on these presuppositions, we shall now see whether we can get
the opposition between the ideal powers of the age, the spirit of the
world and the spirit of the Gospel, sufficiently cleared up. Let us
begin with the spirit of the world.

When consciousness is filled by the idea, be it the idea of the state,
of science, of art, or any contentful idea whatever, the self is freed
from its natural limitation, its private egoism, is saturated with
substance, is enlarged and uplifted. That is why all noble natures have
an unmistakable need to give themselves up to the promptings of ideas
and to live for something ideal. But when the universal satisfaction of
the idea is, without further ado, to count as the absolute satisfaction
of the self, then comes the contradiction. The self is temporal, the
idea eternal; the unity of the temporal self and the eternal idea is an
ambiguous unity, and the spirit of the world, as a consequence, an
ambiguous spirit. At first it looks as if the personal self, by
breathing in the idea, were to become eternal; but since the
organic-physical conditions to which its individual existence is
essentially bound are by no means secured by the movement of the idea,
but on the contrary are exposed to danger: so it ends with this, that
the individuality which has ventured the utmost in the service of the
idea is, as it were, abandoned and betrayed by the idea, even though the
sacrifice, which brings about the individual's downfall in the tragic
moment, has something uplifting and infinitely satisfying for the
individual itself. The contradiction is that the highest satisfaction of
the self ends in the selfless; and the ambiguity is that this
self-dissolution satisfies the self. Homer already saw what is typical
here. That Hector, who fights so valiantly to save his fatherland by
saving the city, is fighting for an idea, is clear enough;
% --- p. 42 ---
\opage{42}\setcounter{footnote}{0}%
just as it is also clear that in the midst of the terrible battle he
feels a higher satisfaction: in the temporal strife he enjoys the
eternity of the idea. Then comes the last great moment; but in the great
moment, when in battle with Achilles he must venture the utmost, then
he stands, seen from without, forsaken by the gods and beguiled by
Athena, in other words: abandoned and betrayed by the idea. ``Woe is me!
now truly the gods of heaven call me to death! Fully I believed that
brave Deiphobus stood at my side; but he is yonder in the city, and
Athena has beguiled me'' (Iliad XXII). It can almost make one indignant
to witness the indifference and partiality of the gods toward the noble
Hector; Athena's conduct in particular must seem to us unworthy. But
such, now, are the ideal powers, the mightily ruling world-ideas;
already before Homer's time they ended the game by giving up the self,
and they do the same still to this very day. And yet, let us not wrong
the high powers; there is in the gods of the world something divine
that does not deceive. Hector understands it; he understands it
heroically, just as it is to be understood in the great tragic moment.
The gods withdraw their aid from him in outward things, for his goal is
set, and his hour has come. Hector understands it: ``Now is my rescue
impossible; in former days Zeus and Phoebus Apollo, his son, saved me,
who both of old graciously shielded me in danger, but now fate overtakes
me.'' Yes, fate is the last, all-powerful world-power; the hero feels it;
but does he then despair over fate? No, for he bows with resignation
under fate. Does he perhaps mock the gods? No, for he prays to the gods:
``Let me not then without deed, without glory perish here! Let me
perform a feat that generations to come will remember!'' And the gods
hear his prayer. The divine
% --- p. 43 ---
\opage{43}\setcounter{footnote}{0}%
is infinitely near to him; for it is in himself, it strengthens him in
the strife! ``Speaking thus he drew from its sheath his sharpened sword,
which was girt over his hip, a strong and mighty blade, then stormed
forward like an eagle that, flying high through the air, swoops down
from darkening clouds upon the plain and seizes either a new-born lamb or
else a timid hare; so Hector stormed forward and swung his sharpened
sword.'' Already in Homer's time, then, the high point had been
reached, the heroic high point to which the spirit of the world is able
to lift the self that sacrifices itself for the idea, that dissolves
itself in the selfless essence: it was mythical in Homer's time, and it
is mythical still to this very day.

If the self is to be saved by consciousness being moved, filled and
streamed through by the idea of life, then the idea of life must not
lose itself in the selfless universal, but must in the spirit be one
with the eternal, almighty self, the self that is the Word of life, the
Word that from the beginning is one with God. ``In the beginning was the
Word, and the Word was with God, and the Word was God. The same was in
the beginning with God. All things were made by him, and without him was
not any thing made that was made. In him was life, and the life was the
light of man'' (John 1:1).
% Danish prints "Menneskens" (the Danish Bible has "Menneskenes", "of men"); rendered literally "of man"
So speaks the spirit of the Gospel, not in order to promote a
speculative knowledge, but in order to determine the principle of faith
and to make clear its opposition to the spirit of the world.

The spirit of the Gospel does not deny the actual world with its actual
laws; on the contrary! It is precisely it that makes the actual world
with its strict laws into an educational institution for the human
spirit, a school of self-development for the self. He, the Word of life,
of whom it is said that he came unto his own, and his own received him
not, did not come with pious fantasies and foggy representations of a
world without laws; no,
% --- p. 44 ---
\opage{44}\setcounter{footnote}{0}%
he came to the world of laws, so that he, who is eternally one with the
author of the laws, might at last, in the natural as in the spiritual,
get the laws of selfhood fulfilled. ``Whosoever will save his life shall
lose it; and whosoever loseth his life for my sake shall save it.'' By
these words the fundamental law of the Gospel is set in sharp opposition
to the law of the world. According to the law of the world it holds,
to be sure, that he who in cowardly fashion would save his life shall
lose it; for only by venturing life for an idea can man win a life that
is worth living. So the brave man stakes his life, and the soldier sings
in \textit{Wallensteins Lager}: ``Und setzet Ihr nicht das Leben ein, nie
wird Euch das Leben gewonnen sein''. But the promise attached to the law,
``whosoever loseth his life for my sake shall save it'', the spirit of
the world cannot utter; the moment it tries to, it becomes ambiguous and
mythical. And what, indeed, has it to offer the hero who ventures the
decisive thing? Only the old dilemma: ``either Caesar or nothing''! These
alternatives are well worth noting. The idea can, on given terms, afford
the self the highest selfish satisfaction; at the heroic high point of
self-satisfaction the hero is then a Caesar: from this side it looks as
if the law of the world really were an original law of selfhood. But it
is a semblance; the glorification of the heroic self bears the stamp of
contingency, is dependent on fortune, on Caesar's fortune. Against
fortune stands misfortune; and the law of the world, indifferent toward
the single self, is as much bound up with misfortune as with fortune.
The law of the world is the law of changes, of outwardness, of
circumstances, of chance encounters, of shifting conditions, in a word:
the law of otherness; and when it turns this side forward, then fate
presses, then the hero turns pale, then even a Caesar becomes a nothing.
In its
% --- p. 45 ---
\opage{45}\setcounter{footnote}{0}%
fulfillment, continually hovering between selfhood and otherness, between
the self and the selfless, the law of the world is as ambiguous as the
spirit of the world.

But in the law of the Gospel there is no ambiguity; it is from first to
last an absolutely decisive law of selfhood, and as such a law of
freedom. The law of the Gospel exempts no one from the burdens of
actuality, but it transforms the outward hindrances into conditions for
the development of the self, and makes the burdens light; it exempts no
one from pain and death, but it gives patience in pain and victory over
death: that lies in its essence as an unconditional law of selfhood.
Through its own negation, its inner resistance, its hostile other, the
self gets itself fully into its own power; through corruptibility it
enters into incorruptibility; therefore it behooved him who fulfills the
law, the perfect law of freedom, to enter into glory through suffering
and death. Where the self conquers under the law of the world, there it
conquers egoistically, for if it is sacrificed, then it vanishes in the
selfless and thus does not conquer: it cannot possibly both be
sacrificed and yet conquer. When, on the other hand, the self conquers
under the law of the Gospel, then it conquers precisely by
self-sacrifice. Self-sacrifice is so far from annihilating the free self
that on the contrary it strengthens, purifies and glorifies it: so
divine is the law of selfhood, now that the exemplar, the eternal self,
who alone has been able to fulfill it, has put it into force.

``The tragic hero,'' says a great author, ``needs tears, and he demands
tears.'' Right! for he needs compassion and demands compassion: his
suffering expresses the universal, that the dissolution of the self is
its transfiguration. He demands tears. Right! for he awakens melancholy.
The dissolution of the self is its vanishing: it vanishes for those who
live on after it, and it vanishes for itself; the final transfiguration
is the melancholy result,
% --- p. 46 ---
\opage{46}\setcounter{footnote}{0}%
the contradiction softened in mood, that the one transfigured is no
longer himself, but another, a spiritual afterglow of himself, a memory.
It is otherwise with the exemplar of faith; he rejects all tragic tears,
for he rejects compassion. His suffering does not threaten him with
self-dissolution, it only confirms the superiority of the self; just
because the way of sufferings is his way to glory, he can say to the
weeping women: ``Daughters of Jerusalem, weep not for me, but weep for
yourselves, and for your children.''

That two principles as opposed as the Gospel and the spirit of the world
must at all times move the race in opposite directions is a matter of
course. It is not only the Middle Ages that have been dualistic; the
modern age too is, and every following age will continue to be,
dualistic. How indeed should the dualism
% Danish prints "Horledes" (for "Hvorledes"); translated per the sense
be able to vanish? It certainly cannot happen through the spirit of the
world, with its ideas, its tasks of actuality and its interests,
vanishing; for that vanishes only with the destruction of the world.
Then it would have to happen through the spirit of the Gospel, with its
divine exemplar, its gifts of grace and its promises, vanishing from the
earth; but then every human self would first have had to learn to defy
the law of freedom and to prefer death to life. No, the point is not
that the dualism is to vanish; but the point is that each age must
conceive and determine the character of the dualism in its own
distinctive way. In the present age the opposition is veiled, the
dualistic subdued. This comes from the fact that the consciousness of
the world, with its interests in the present, has stepped into the
foreground. The truth which the Gospel tacitly presupposes, but which
the Middle Ages could not rightly bring out, that a living for what is
to come is inseparable from a vigorous living in the present, that the
development
% --- p. 47 ---
\opage{47}\setcounter{footnote}{0}%
of humane ideas belongs essentially to the actual kingdom of the spirit,
our age has given such preponderance that the heavenly is close to
being absorbed into the earthly, as a prime number divides into its
product. If the heavenly and the earthly, the spirit of the Gospel and
the spirit of the world, are to keep to their mutual heterogeneity, to
go on forming extremes without yet being separated by a gaping chasm,
then the extremes must have a middle; such a middle is precisely the
spirit of the people.

How the spirit of the people and the life of the people are determined
by each other has already been discussed in general terms; here we shall
further attempt, if only provisionally, to determine the popular a
little more closely. What is the popular? It is the fundamentally human
in its naturally vigorous distinctiveness. It is heard often enough in
recent times that man is a rational being; more rarely is due emphasis
laid on this, that man according to his nature is originally a religious
being; and yet not only must both fundamental determinations be derived
from this, that man is a self-conscious being, a personal self, but the
first must even be capable of being regarded as derived from the last.
As a religious being, man stands through the world in relation to God;
for the personal self the relation to God is the original one. As a
rational being, man stands in relation to the natural order of things,
so that thorough insight into the relations of the world becomes an
essential task for him. The cognition of the relation to God is a
religious cognition and rests on faith; cognition of the idea of the
world and of the nature of things is a cognition of reason and rests on
knowledge. Man, then, by nature as man, and thus as a popular individual,
has an inward, essential need both for faith and for knowledge, as surely
as he has by nature an inward, essential need both for cognition of God
and for cognition of the world.
% --- p. 48 ---
\opage{48}\setcounter{footnote}{0}%
Self-consciousness is, and according to its nature must necessarily be,
a consciousness doubled in itself: consciousness of another, and
consciousness of itself. As its other, self-consciousness has in the
absolute sense God, in the relative sense the world; in accordance with
its fundamental doubling, human consciousness will then, with regard to
this, divide itself into consciousness of God and consciousness of the
world. What thus holds of every popular individual, when we look at the
natural endowments and the corresponding need, must incontestably hold
of the people as a people. Every people, then, begins its spiritual life
with a doubled fundamental consciousness, a doubled endowment, a doubled
need, a need for religious and for rational self-development, a need for
the salvation of souls and a need for culture.

Now if a popular development of spirit either were, or in actual history
could be, normal, then it would surely show how the religious life and
the life of culture, separated according to the fundamental doubling
into two absolutely heterogeneous spheres of consciousness, would, by
virtue of the unity of the personal self-consciousness with itself,
develop through all doublings in harmonious interaction; but since no
history of a people is normal, we must conduct the proof the other way
around. The life of the people is, in its first spiritually vigorous
beginning, predominantly religious; only gradually, as the development
of reason and with it culture advances, does the consciousness of the
world gain preponderance over the consciousness of God. How long
religion will be able to counterbalance rising culture will prove to
depend precisely on the depth in which the relation of opposition was
conceived from the beginning, and thus on the dualistic fundamental
doubling. The difference between the Greeks and the Jews is
characteristic in this respect.

It has been aptly said of the Greek gods that in all their
supernaturalness and superhuman bliss
% --- p. 49 ---
\opage{49}\setcounter{footnote}{0}%
they were all too natural, all too human. Herein lies this, that among
the Greeks the consciousness of God coincided too immediately with the
consciousness of the world. The Greek faith was without depth, without
inwardness and ethical seriousness. The more the divine figures pressed
forward, so to speak, to have their revelations aesthetically perfected,
the more the religious faith was transformed into aesthetic faith, into
an aesthetic mood determined by the intuition of fantasy. We see it in
Homer. The domestic Olympian scene (Iliad I), in which Hera's taunts put
Zeus into such a paternal rage that both men and gods must burst out
laughing, is surely by no means designed to give the gods over---as a
Lucian did much later---to mockery and contempt; no, it is a proof that
the poet, without troubling himself about religious reverence, has made
use of his poetic motif and transformed religion into poetry. While
religion was thus gradually transformed into poetry among the cultured,
the immediate, popular religious faith sank down into crude, fantastic
superstition, and philosophy obtained, in the name of reason, the right
to apply criticism and let knowledge dissolve faith. Although the
condemnation of Socrates must surely have made Plato cautious, he still
cannot keep silent that the Homeric gods ought to be excluded from the
Republic.

While the religious development among the Greeks was, even fairly early,
hindered and pushed back by culture, conversely the development of
culture among the Jews was in several respects hindered and pushed back
by religion. If the Greek consciousness of God was absorbed immediately
into the consciousness of the world, the Jewish consciousness of the
world was in return absorbed immediately into the consciousness of God.
Among the Jews, too, we miss the thoroughgoing opposition between the
kingdom of God and the kingdom of the world to which the original
doubling of consciousness must consistently lead; but
% --- p. 50 ---
\opage{50}\setcounter{footnote}{0}%
admittedly in the way we had to expect it among
% Danish prints "vente at finde den" (Rettelser: "vente det"); translated per the Rettelser
the chosen people, at the first stage of revelation, where the divine
emphasis falls on religion. In the Jewish theocracy the difference
between religious and civil legislation is abolished; the kingdom of
David is a kingdom of God, the morality of the prophets a religious
politics. The faith in Jehovah's requital by temporal punishments and
temporal rewards corresponds to this. The consequence of this was that
the life of the spirit, with its free inwardness, lost itself more and
more in a crystallized outwardness. After the Exile, when the prophetic
voices at last fell silent, the religious content was petrified in
traditions; the people, without inwardness of faith and without culture,
was led by Pharisees and scribes.

If we cast a glance at the Roman people, we get a truly uncanny
representation of what a finite mixture of religiosity and worldliness
signifies. The Roman gods are gods of the world in the lowest, most
prosaic meaning of the word; these gods of the understanding, spiritless
in themselves, pettily attentive to pettily calculated ceremonies and
endless formulas, have no other aim than a boundless expansion of the
power of the Roman Empire; nothing else to think about than how they
might be of use and benefit to the Roman citizens in agriculture, cattle
raising, sheep breeding, etc. That so official a religiosity must at
last become comic to itself, once it became conscious of its
contradiction, was natural. In the end, after all, two augurs, despite
all their Roman dignity, could not meet each other without laughing.
What were the augurs laughing at? It was a doubling of consciousness; it
was only formal, and had never been real.

And what is the result of all this? That the life of the people perishes
when religion is undermined! At the vigorous beginning, when all the
elements of the spirit of the people are in full activity, religion is
a spiritual
% --- p. 51 ---
\opage{51}\setcounter{footnote}{0}%
power; the holy is guarded with reverence: there is earnestness in the
% Danish prints "de er Alvor" (sense perhaps "det er Alvor"); translated per the sense
worship of the gods. But with growing culture a separation takes place
in society; the separation is at once civil and religious. From the
civil side the separation shows itself in this, that the great crowd is
more and more excluded from participation in the government of the
state and in public affairs. The exclusion is not arbitrary; it follows
from the nature of the circumstances. It is difference in station in
life, in wealth, in endowment, and above all, difference in culture,
on which it turns. The development of culture is never founded, and can
never be founded, directly on religion; it is a development, not of the
consciousness of God, but of the consciousness of the world; it is
reflection and criticism, knowledge and science in manifold, partly
theoretical, partly practical ramifications and consequences, that
originally furnish the ideal foundation of the life of culture. With
this foundation the multitude never becomes familiar: the light shines
in its eyes, but it is dazzled by the glare and cannot see. Were the
multitude to keep pace, its eye would have to be cleared by another
light; then a religious enlightenment with a corresponding cultivation
of spirit would have to step in and keep pace with culture. Now when
this does not happen, the great crowd sinks down into crudeness,
ignorance and superstition; then the mass lives in the country in the
condition of a peasantry, in the larger cities in the condition of a
rabble, and the development of the people is brought to a halt. The
state may be organized and enlarged, manners refined, science, art and
literature make admirable progress: it does not help; the people's fate
is decided, the pulse of life cut.

It is a tragic spectacle that each of the great states of antiquity
presents as they approach the goal of the development of culture: in the
middle a small circle of the truly cultured, such as have by upbringing
and self-development
% --- p. 52 ---
\opage{52}\setcounter{footnote}{0}%
truly appropriated culture; around them a multitude of the half-cultured,
who take part in civilization at second and third hand, and around these
again the thoughtless multitude, whose crudeness always forms the
extreme that has over-refinement as its opposite. Think of Rome in the
time of Augustus, the Roman knights, the Roman manners, the Roman luxury,
and then a Roman rabble, hungry for bread and spectacles (panes et
circenses), and one gets an impression of what it means that a spirit of
a people can be extinguished, and a people perish.

The principle of popular culture is essentially religious. Only in the
light of religion can the people as a people advance on the way of
culture and be made a partaker in civilization; if this light is
extinguished, the people sits in darkness. The nature religions vanish,
one after the other; the lights of antiquity are put out, one after the
other: only one light is not extinguished, and that is the light of the
Gospel.

It is characteristic of Christianity that from its very beginning it
announces itself as a common cause of the people. Christ's birth is
proclaimed by the angels as a great joy that shall be to all the people;
the aged Simeon, with the Christ child in his arms, praises God that his
eyes have seen the salvation which is prepared for all peoples, the
light that shall lighten the Gentiles. Christ's appearance is popular in
the most proper sense of all. He speaks in the synagogue, he teaches from
the ship, he preaches on the mount, and everywhere crowds stream to him,
and the people gather round him and hear him gladly, for ``he taught them
as one having authority, and not as the scribes''. The whole public life
of Jesus is sacrificed to awakening and preparing the people. But so
strong, again, is the dualism of the holy and the sinful that the
sacrifice ends with an atoning sacrifice. Only when the atoning
sacrifice has been made, and he who sacrificed himself has entered into
his glory, is the Spirit sent out and the congregation founded.
% --- p. 53 ---
\opage{53}\setcounter{footnote}{0}%
The unity of the religious and the popular has precisely the fundamental
doubling, the most decisive separation between the kingdom of God and
the kingdom of the world, between the consciousness of God and the
consciousness of the world, as its presupposition. Just because
Christianity is absolutely heterogeneous with everything worldly, it can
in a spiritual manner take up and use everything worldly; it does not
ask what the world-ideas are worth as relative ends in themselves; for
to the absolute purpose they are all means. Why can being a fisherman
and being a farmer be united, in popular fashion, with being a
Christian? Because fishing and farming are not merely relatively, but
absolutely heterogeneous with Christianity itself. For the same reason,
then, being an artist and being a cultivator of science can also be
united, in popular fashion, with being a Christian. Should the man of
knowledge have reason to complain that his consciousness is halved
because the theoretical cognitions of science do not coincide with the
religious cognitions of faith, then the farmer could, after all, also
complain that agronomic cognitions cannot be developed out of the
Gospel.

On the terms of the fundamental doubling, then, religious enlightenment
will be able to advance together with the enlightenment of culture, and
both unite in a true enlightenment of the people. That the doubling must
make stricter demands on the life of the spirit than the formal unity
does is natural; that the stricter demands may provoke the consciousness
of culture to resistance, and this resistance bring about a tension that
at last discharges itself in a storm of the spirit, is possible: we
shall see in what immediately follows on what the tension rests, and how
it rises.

\begin{center}\rule{3cm}{0.4pt}\end{center}

\chapter*{IV.}
\addcontentsline{toc}{chapter}{Lecture IV}
\markboth{Lecture IV}{Lecture IV}
\label{ch:iv}
% --- p. 54 ---
\opage{54}\setcounter{footnote}{0}%
It is the spirit of the world, with its suggestions of thought and its
movements of ideas in the cultural life of the present age, that we shall
attempt to characterize in a few, but large and vivid, strokes. We shall
consider this spirit, not from its bright side, but from its shadow side;
for the bright side is already sufficiently acknowledged. That culture
and enlightenment, education and science, are essential goods, every
rational human being sees, and what is fully acknowledged by everyone
needs no praise or commendation of ours. But the dark side! The spirit
of the world, not as the spirit of culture and progress,
% Danish prints a stray full stop after "Aand" (AUDIT doubtful: may be a speck); ignored in translation
but as the spirit of delusion, of strong errors, of fateful necessity,
has this being now also been sufficiently seen through, is it perfectly
transparent to the consciousness of culture, or should there not be
obscurities here in the depths, for the clearing up of which a stronger
light might be needed? Here, of course, there is to be no talk,
individual by individual, of any particular ungodliness in the single
individuals of this time---to inveigh against the excesses of
freethinkers and the blasphemies of deniers of God we must leave to the
preachers of repentance---here, on the whole, there is to be no talk of
the unbelief of individuals, but of the spirit of unbelief in the
individuals, a spirit that poisons the air of life, deposits nodes of
poison, and from the poisonous nodal points
% --- p. 55 ---
\opage{55}\setcounter{footnote}{0}%
spreads out and attacks what is noblest and best in society. It is
chiefly science that this spirit misuses; through natural philosophy,
% Danish prints "denné" (stray accent); no bearing on the sense
history and psychology it seeks, by misleading criticism, to distort the
truth of faith; through knowledge and science it breathes poison into
souls. But then is it perhaps the fault of science? By no means! Science
is pure and noble and good when it is allowed to count for what it
originally is, in its objective scientific character. To cast a shadow
on science little becomes an academic teacher, and would be in conflict
with the truth: no branch of science is either dangerous or corrupting
in and for itself. Only an outlived papism can, in desperate struggle
against the strong light-bringing ideas, hit upon pronouncing a curse on
the scientific progress of our time. No, it is neither the sciences nor
the knowing individualities that we are speaking of; we are speaking of
something moving that perverts the idea and gives it a false direction.
We are speaking of a negative, impersonal-half-personal monstrosity which
brings it about that the semblance of science beguiles the man of
knowledge; we are speaking of a thinking without a thinker, a willing
without a will; it looks somewhat mystical---what, one wonders, can it
be? We need not exactly sink down into the innermost darkness of the
God-denying will, or lose ourselves in broodings over what Scripture
calls the depths of Satan, in order to get hold of our mystical
monstrosity; we keep to what is nearer, less inhuman and more
comprehensible, and then designate the spirit of the world, seen from
its shadow side, as the juggling spirit of the age. In order to get
sight of the spirit of the age we must direct our gaze particularly to
the consciousness of culture and to the fundamental idea ruling in the
era. The age of means, of conditions, of interests, is not forsaken by
ideas. It is, to be sure, tempted in practical matters to forget the
absolute, the unconditioned, but
% --- p. 56 ---
\opage{56}\setcounter{footnote}{0}%
theoretically it does not forget it; the age of culture can never forget
that its clearest, purest light proceeds from knowledge and science; in
science it has its absolute: the idea of knowledge is the ruling
fundamental idea. Now that the idea of knowledge can be said to be both
half-personal and impersonal is easy to see. If there were no knowers,
there could not, after all, be any knowledge either: the idea of
knowledge must be borne and moved by knowing individualities. But the
knowing individualities must then also at the same time find themselves
so filled, penetrated and moved by the idea that it is not their own
thoughts that subjectively determine the essence of the idea, but
conversely the essence of the idea, with its necessary laws, the eternal
truths, that objectively determines the thoughts. But ``the eternal
truths'' and universally valid laws are, after all, impersonal: the
fundamental movement is thus this, that the impersonal determines the
personal, that the selfless gets power over the self. In this hovering
between the self and the selfless, self-reflection at last becomes
uncertain of itself. The man of knowledge thinks ``the eternal truths''
and is by this thinking sunk in the universal, and by the universal moved
to intellectual resignation. But the man of knowledge at the same time
exists in a world of finite, actual interests. Interest is selfish; for
the sake of interests the existing thinker is most emphatically called
upon to assert his personal right and to lay selfish weight on his own
self. At once infinitely resigned and yet finitely interested, the
knowing individuality is thus moved both impersonally and yet personally:
where, now, does the absolute truth lie?

The absolutely personal truth lies outside the boundary of knowledge and
can be grasped in faith alone. On these terms the man of knowledge would
then have absolutely to acknowledge a principle of truth that is higher
than knowledge; but precisely that which
% --- p. 57 ---
\opage{57}\setcounter{footnote}{0}%
in the deepest sense alone is able to satisfy the self's demand, the man
of knowledge now regards as a contradiction in the self's essence, a
desperate self-contradiction: here is the deception. From this deception
the immediately existential temptations, and all the confusion of
thought with which the spirit of the age corresponding to the
fundamental idea dazzles and beguiles the man of culture, must
essentially be derived.

It is by no means the idea as idea that beguiles the knowing
individuality; it is on the contrary the individuality that, by deifying
its knowledge and taking the idea in vain, dazzles and beguiles itself.
The absolutely personal claim of the man of culture science cannot and
shall not satisfy; from the objective standpoint of the deification of
knowledge such a claim must quite rightly be rejected. But an absolutely
personal claim the personality cannot, after all, reject; each time it
is pushed back, it presents itself anew. The absolutely personal is the
ideal of freedom; but the ``eternal truths'' of knowledge rest in
necessity. Captive of necessity, the man of culture must then ceaselessly
yearn for freedom. The spirit of the time that dazzles him is to be
called an unstable, fickle spirit precisely for this reason, that it
continually hovers between two poles, between the bright region of
freedom and the dark region of necessity; in order to hide the deception
and complete its juggling play, it first falls in love with the light of
freedom and turns away from dark necessity, but then freedom is in turn
to be semblance, and necessity itself the source of enlightenment.

We are reminded of Asa-Loke from the old myths. Loke is of giant nature,
but richly gifted, full of wit and cleverness, lovely to look at when he
turns toward the light, so that Odin himself in the morning of time
entered into sworn brotherhood with him. Though a giant, he feels
himself drawn to the Aesir, he lets himself be used in the service of
light;
% --- p. 58 ---
\opage{58}\setcounter{footnote}{0}%
but he never becomes truly at home in the region of light; his inner
being is too dark for that. Seen from the shadow side he is hideous, and
this hideousness testifies precisely to his giant nature. Thus with
Angerbode he has begotten the Fenris wolf, the Midgard serpent and Hel,
and stands constantly in league with the realm of darkness. He lacks
neither liveliness nor criticism; but his liveliness is dissolute, and
his malicious reflections consume his will. By these and similar traits
Loke may well show himself as an image of the spirit of the age, but
yet only of the spirit of the age in general, as the false, ambiguous
spirit of the world. Our task is expressly this, to determine the spirit
of the age in the present by showing how, from a mystified knowledge as
from a shadowy background, it carries on its juggling play.

To begin with sheer lies cannot be done; a pure lie is moreover a
contradiction, for in every lie there are always a few grains of truth:
a lie is truth bungled, a fraction of truth that passes itself off as a
whole number. Every intellectual mystification must begin with something
unambiguous in order to introduce the ambiguous. It is, for example, an
unambiguous proposition that he who has said A must also say B; on this
proposition the emphasis now falls when the spirit of the world, with
the help of knowledge, sets about confusing culture. If one has granted
a presupposition, then one is compelled, whether one will or not---if,
that is, one does not want to stand as a thoughtless person in
self-contradiction---to grant its consequences. It is perfectly true and
unambiguous that he who says A must say B; but if, over this unambiguous
thing, one forgets to examine whether the presuppositions are homogeneous
or heterogeneous with the consequences drawn from them, then the
ambiguity comes. The proposition is valid only on the condition that
presuppositions
% --- p. 59 ---
\opage{59}\setcounter{footnote}{0}%
and consequences belong to the same sphere: from presuppositions in one
sphere one cannot possibly derive consequences that are to hold in
another. Thus there is something called harmony, both in the sphere of
tones and in that of colors; now if someone places the presupposition in
the harmony of tones, and derives from it consequences for the colors,
then he becomes a sophist, and yet one can say: harmony is always
harmony. One can speak of the hardness of stones and of the hardness of
hearts, and one can say in general: hardness is always hardness; but one
cannot, from determinations about the hardness of stones,
derive consequences for the hardness of hearts, or conversely. One can
speak of light in the spirit and of light in nature, and one can say
that the concept light is always light; but one cannot derive from the
natural light consequences for the spiritual. Here, then, is the
subtlety: before one is well aware of it, everything is turned around;
consequences in one sphere are slipped in under presuppositions in
another. If one goes along with such a logic, one lets oneself be
fooled, and once the sophistry has got going, there will be no lack of
reflection and dialectic to complete the deception. So it is when
religion and science are in question.

With religious presuppositions one must bow to religious consequences,
and likewise in science to scientific ones; in this there is no
ambiguity. But if one now turns it around so that one would derive
scientific consequences from religious presuppositions, and conversely,
then we have the sophistry. It is in this sophistry of the spirit of
the age that the corruption lies; it is this that attacks personal life
from the root; it is this that makes the powers of the spirit decline,
characters grow weaker, and individualities weakened by doubt flutter
up toward the light of faith without yet being able to
% --- p. 60 ---
\opage{60}\setcounter{footnote}{0}%
believe. We are, as has been said, poisoned; the poison is a sophistical
mixture of heterogeneous elements that mutually consume one another.
One can, in order to demonstrate this, proceed from natural science,
from history, from psychology. In the present connection we shall dwell
on natural science.

When the spirit of the age dazzles the consciousness of culture through
natural science, the sophistry begins roughly thus: ``It is certain that
we all, without regard to subjective wishes, ought to seek the truth,
for truth is, after all, the highest, the eternal.'' In this there is no
ambiguity. ``Further, we ought to acknowledge that where there is truth,
there must also be actuality. We know, to be sure, from the higher logic
that the concept of actuality does not fully exhaust the concept of
truth, but where there is complete truth there must necessarily be a
corresponding actuality. If the truth conflicts with the facts of
actuality, then something is lacking; the truth is then only half, and
in this halfness to that extent untrue.'' Here, then, the consciousness
of culture, insofar as it turns negatively against the highest, lays a
special stress on the word actuality, and points out---in a certain
sense with full right---that what we call actuality is something for
which we must in practical life have unconditional respect. If we had no
respect for the actual facts, we could not, after all, even walk on our
own or move freely in the world of things; we would then, for example,
not respect the separation between rock and abyss or between fire and
water. Everyone knows that when he puts his hand into the fire he must
burn himself; he need not repeat the experiment many times, no, a single
experiment is enough, even for a Mucius Scaevola. It is only skeptical
perversity and speculative affectation to want to
% --- p. 61 ---
\opage{61}\setcounter{footnote}{0}%
make experience uncertain by letting it rest on a sum of many single
cases, or depend on rules that always admit exceptions. Certainly there
are facts whose distinctive natural constitution can be secured to our
experience only after many-sided observations and careful tests; but
doubt about particular empirical propositions not yet sufficiently
tested is essentially different from doubt about the whole
reason-determined foundation of experience.

Those deep-seeing men who, in the name of a more ideal knowledge, think
they can raise themselves above ``the patchwork of the empirical
sciences'' by a remark such as this, that the ``so-called laws'' of
natural philosophy are determined only by induction, and that induction,
as an imperfect inference, stands far below the categorical syllogism,
would do well to think a little before they go further. If it is an
uncertain inference that the attraction of masses varies inversely as
the squares of their distances, then it is also an uncertain inference
that he who jumps into the water must get wet. If it is an uncertain
inference that oxygen and hydrogen under determinate conditions must
combine into water, then it is also an uncertain inference that he who
is beheaded loses his life. Now should anyone, for the sake of
induction, doubt the universal validity of either the mechanical or the
chemical inference, let such a deeper-seeing man, for the sake of
induction, once try whether, in spite of the merely empirical laws, he
might not succeed in living without a head.

To refute an induction of natural science it is of little use to raise
doubts, in empty generality, about the abstract logical validity of
induction; no, the demand, the true scientific demand, is that one
demonstrate in concreto cases, if only a single one, that manifestly
conflict with the result determined by the
% --- p. 62 ---
\opage{62}\setcounter{footnote}{0}%
determinate induction. And when such a case occurs, what then? Why,
even then it is not the general scientific validity of induction as
induction that is shaken, but it is a particular induction, applied to
particular cases, that is corrected. The scientific doctrine of
experience has a foundation as secure in its kind as any a priori
science of reason whatever. It is not what is uncertain and wavering in
the results of the natural sciences that is to comfort us in our
anxieties for the religious truths; no, it is something else; it is, in
the first place, an unmasking of a peculiar deception that lies in the
spirit of the age.

We shall now see the deception as, in order to characterize the spirit
of the age, we set about deriving religious consequences from
scientific presuppositions and, conversely, scientific consequences from
religious presuppositions; we shall see in some examples how the
sophistry gets into motion.

An example from the Old Testament! Jehovah, it is said (Gen.\ IX),
expressly sets the rainbow in the clouds as by a solemn act he makes a
covenant with Noah. The religious presupposition is now this, that
Jehovah really by a determinate act of will set this sign, thus worked a
miracle, and, what is after all also a miracle, spoke to Noah and
interpreted the sign to him. ``I will set my bow in the cloud, and it
shall be for a token of a covenant between me and the earth. And it
shall come to pass, when I gather the cloud over the earth, that the bow
shall be seen in the heaven, and I will remember my covenant, which is
between me and you and every living soul of all flesh, and the waters
shall no more become a flood to destroy all flesh; but the bow shall be
in the heaven, and I will look upon it, and remember the everlasting
covenant between God and all flesh that is upon the earth.''

% --- p. 63 ---
\opage{63}\setcounter{footnote}{0}%
Now if we are to be allowed to draw a scientific consequence from this
religious presupposition, then we had surely best, the sooner the
better, throw optics into the fire. For optics cannot and will not in
any way see in the rainbow either a miracle or an express token of
remembrance of God's covenant with all flesh. Optics explains the
rainbow simply and solely by the angle of refraction of the rays of
light, thus as a naturally necessary consequence of natural conditions.

Let us then try going the opposite way; let us make quite clear to
ourselves how the rainbow is formed under certain conditions, and how we
can see rainbows not only in the clouds but also in fountains: then we
have here a secure scientific presupposition. Now if from this scientific
presupposition one is to draw religious consequences, one comes into
conflict with religion. For, say the knowing ones, we know that what is
said in the Old Testament about the rainbow is nothing actual, but a
fantasy; the Old Testament does not know the formation, conditions and
laws of the rainbow. It is a childlike spirit, unacquainted with the
actual facts, that speaks in the words put into Jehovah's mouth. But if
this is so, then it cannot be true either that Jehovah himself spoke to
Noah of any covenant; the covenant itself is here just as imagined as
the explanation given of its phenomenon. When the unknown author of the
first book of Moses has Jehovah say: ``And it shall come to pass, when I
gather the cloud over the earth'', or: ``when I bring the cloud forth'',
then he conceives Jehovah the cloud-gatherer just as figuratively,
childishly and mythically as Homer conceives Zeus the cloud-gatherer
(Νεφεληγερέτα Σεύς).
% Greek as printed: Danish has capital sigma for zeta (Ζεύς); copied verbatim
Now if it is on the whole characteristic of the spiritual life of the
era that movements of thought in such opposite currents pass back and
forth through the consciousness of culture, what wonder, then, if the
doubts
% --- p. 64 ---
\opage{64}\setcounter{footnote}{0}%
grow with a quiet growth; what wonder that those who are touched by them
content themselves with palliatives of aesthetically staged worship, with
a little of faith and a little of knowledge, and feel but slight need to
engage with deeper religious problems.

Another example! In the first chapter of the first book of Moses we meet
the highest, the first of all religious presuppositions, in the account
of the creation. ``In the beginning God created the heaven and the
earth.'' Attention fastens at once upon the earth: ``And the earth was
without form, and void,'' and so on. Around the earth as center the
heaven must vault itself as an outstretched firmament in the midst of
the waters. ``And God made the firmament, and divided the waters which
were under the firmament from the waters which were above the firmament:
and it was so.'' It is said here expressly that the heavenly bodies came
into being for the earth's sake. ``And God made two great lights; the
greater light to rule the day, and the lesser light to rule the night:
he made the stars also. And God set them in the firmament of the heaven
to give light upon the earth.''

The earth, seen from the standpoint of religion, is thus not merely the
physical center of cosmic space but, if we may say so, the apple of
God's eye, because there is something on the earth that God prizes
unconditionally, namely man, who bears his own image. Just as the starry
heaven with sun and moon exists for the earth's sake, so the earth in
turn exists for man's sake: man is not merely the lord of the earth, but
the crown of creation. Precisely for this reason it is regarded as a
cosmic catastrophe that man succumbed in temptation and lost his image of
God; and precisely for this reason it is regarded as a new creation that
God himself, in his only-begotten Son, by taking upon himself flesh and
blood, raises up the fallen race and gives it back the lost image.
Precisely for this reason it is said
% --- p. 65 ---
\opage{65}\setcounter{footnote}{0}%
that when Christ's kingdom has been proclaimed to all the heathen, then
follows the last catastrophe, the destruction of the world and the
judgment. When Christ foretells the destruction of the world, he says
that there shall be signs in the sun and moon, that the sun shall be
darkened and the stars fall. This is the fundamental view of revealed
religion, and this fundamental view is not accidental. Everything from
first to last is so coherent that, if we shake one thing, we end by
dissolving the whole: the world that is to be a stage for God's own
becoming man must be a central world, the highest center point in the
universe.

If one is now, without further ado, to draw scientific consequences from
this religious presupposition, then one must break the staff over the
physics and astronomy of the present age. For astronomy does not admit,
cannot and will not admit, that this earth is the center of the world.
On the contrary! if one will hear the voice of science, one gets another
view of the world, a real view of the real.

What astronomy teaches us about the structure of the world is not loose
opinions and empty conjectures; it is real facts and real laws,
something so incontrovertible that one cannot toss it to the winds. It
holds here that the more cultivated a man is, the greater respect he has
for actual nature and its laws. The lower he stands, the more easily he
settles down in the fancy that the learned do not after all really know
the matter, and that one gets the true physics and the right astronomy
only by taking the words of the Bible literally. Yes, on the lazy
cushion of faith in the letter one slumbers sweetly off; but one does
not sleep one's way to the cognition of truth. If the eye is sound and
the understanding awakened to see with what certainty and consistency
science proceeds in
% --- p. 66 ---
\opage{66}\setcounter{footnote}{0}%
determining the relation between earth and heaven, one will recognize
that the earth, far from being the center of the world, is not even the
center of the small part of the world that is called our solar system.
The propositions belonging here are not merely set up; they are also
grounded. Whoever, face to face with science at its present standpoint,
would be able to doubt the incontrovertibility of a proposition such as
this, that it is not the sun that moves around the earth but the earth
that moves around the sun, that it is therefore the sun and not the
earth that is the center of the planetary system,---let him doubt the
incontrovertibility of the proposition that a headless man cannot live.
That the earth in our little world is a comparatively slight magnitude
beside globes such as Jupiter, Saturn, and Uranus is clear to anyone who
has gained an insight into science. Now let one consider those holy
words: ``There shall be''---for the earth's sake---``signs in the sun,
and in the moon, and in the stars: the sun shall be darkened, and the
moon shall not give her light, and the stars shall fall from heaven, and
the powers of heaven shall be shaken.'' And why? Because the goal on
earth has been reached; because the course of mankind's life is at an
end, because the highest has been proclaimed to all peoples and either
appropriated in faith or rejected in unbelief.

How it will go now when one is to draw religious consequences from
scientific presuppositions, anyone can easily tell himself. It stands
written in the Book of Revelation (VI, 12---13): And I beheld when he
(the Lamb)
% Printer's defect: the opening quotation mark before "Og jeg saae" is not
% printed in the Danish (the closing one after "Veir." is); carried as printed.
had opened the sixth seal, and, lo, there was a great earthquake; and
the sun became black as sackcloth of hair, and the moon became as blood;
and the stars of heaven fell unto the earth, even as a fig tree casteth
her untimely figs, when she is shaken of a mighty wind.'' At the root of
the assumption that the stars could fall to the earth there lies
% --- p. 67 ---
\opage{67}\setcounter{footnote}{0}%
evidently a twofold representation: partly that the stars are very
small in comparison with the earth, partly that the distance between
heaven and earth cannot, after all, be so extraordinarily great.

Among the ancients something is told that testifies to how little the
human imagination is able to grasp the enormous distances in the
universe. Homer (Iliad. I, 586) has Hephaistos, who would so gladly
stand by Hera against the angry Zeus, address his mother: ``The Olympian
lord is hard to fight against. Once before, when I would come to help
thee, he seized me by the foot and hurled me from the threshold of the
gods. The whole day long I floated about, but when the sun went down I
fell upon Lemnos, and there was little life left in me.'' One shudders
to imagine the falling Hephaistos floating a whole day between heaven
and earth; in Hesiod this floating is even extended to nine days and
nine nights, but what does that amount to? How far, then, might it be
from the earth to a star outside the solar system, to a fixed star?

If we consider the height of one of the nearest fixed stars, then we
must use the speed of light, 40,000 miles a second, as our measure, and
not the fall of Hephaistos. The distance from the sun to the earth,
which is 20 million miles, is traversed by light in 8 minutes and 13
seconds; but from
the nearest fixed star, Alpha Centauri, light is 3\textonehalf{} years
on its way, from Sirius 14 years, from the Pole Star 42 years, before it
reaches the earth. The Pole Star could thus be extinguished and yet
still shine for us for 42 years after it had gone out. These are
enormous distances, and yet they dwindle to insignificant magnitudes
when we raise our gaze to the stars of the Milky Way, to stars of which
it can be said without exaggeration that they shine for us with
pre-Adamite light.
% --- p. 68 ---
\opage{68}\setcounter{footnote}{0}%

If, to humor the spirit of the age, we permit ourselves the folly of
drawing religious consequences from scientific presuppositions, then in
truth it will be religion that pays the price. If the religious view of
the starry heaven cannot have true religious validity unless it agrees
scientifically with the astronomical, then let us be honest and say to
Christ in all candor: Lord, thou hast been mistaken, thou hast not
studied thy astronomy; and since thou art not especially sure in thy
astronomy, neither dare we believe unconditionally that thou art always
entirely sure in thy theology. Thou speakest of the stars falling from
heaven on the day of judgment. How should they fall from heaven, unless
thy meaning were that they should fall upon the earth? But Sirius and
the Pole Star and the still more distant suns cannot possibly fall upon
the earth, since the earth is a vanishing point in comparison with them.

Let us leave the heavens and dwell on the earth, and see whether it is
not permitted here to draw religious consequences from scientific
presuppositions, or conversely; it sounds so fine that science should be
a little religious, and religion a little scientific. The holy word of
revelation speaks of the creation of the earth; the speech is pithy and
brief. The earth was without form and void, and darkness was upon the
deep, and the Spirit of God moved upon the waters; then the word of the
Almighty is heard, and day by day the work advances to completion. After
the completion of creation, after man, formed of earth, has been created
in God's image, there is then told of Paradise and of the state in which
the spirit of peace and blessing rested over what had been created.
``And God saw every thing that he had made, and, behold, it was very
good.'' This is a true and essential religious presupposition, and
whoever undermines this presupposition, seen as principle, undermines
religion.
% --- p. 69 ---
\opage{69}\setcounter{footnote}{0}%
But if scientific consequences are to be drawn from it, then one must
necessarily break the staff over geology. Geology is a new science,
which as such still has many obscurities; but it advances boldly: it
already has its secure method, its measure, and its correctives. By
breaking the staff over so thoroughly investigating a science, we would
come into decisive conflict with the whole of culture. Conversely, if
one is to draw religious consequences from scientific presuppositions,
that is, to start from the doctrine of the terrestrial body and of the
development of organic life on the earth, then one arrives at the result
that creation cannot have come about as is told in the first book of
Moses. Long, long intervals must then have passed before the one kind of
creatures followed upon the other; the earth itself must have undergone
many transformations before man could appear upon it. And what becomes
of the Garden of Eden: modern science surely has no room at all for a
Paradise? No, the literature of the Middle Ages knows all about
Paradise; we see it in Dante's divina comedia. In Dante's time Jerusalem
with Mount Zion was still the center of the earth on this side; for its
antipode it then had, on the other side, the mountain of Purgatory, on
whose summit Paradise lay, the blessed earthly one, where Dante met his
Beatrice. But if in our time one is to draw religious consequences from
scientific presuppositions, then there has never been any earthly
Paradise. One thing more! According to Scripture sin came into the world
through man, and death through sin; hence death is the wages of sin, and
whoever violates this proposition violates the religious. But if one is
to draw scientific consequences, one must violate this presupposition;
for geology shows that many lower
% --- p. 70 ---
\opage{70}\setcounter{footnote}{0}%
% Danish prints "stridt mod" (Rettelser: "stridt med"); the sense is the
% same in English.
and higher races of animals fought with one another and devoured one
another long before man came into being, that there has thus been death
and pain before man, that corruptibility did not come because man
sinned, but that man sinned because he came into being in
corruptibility.
% No paragraph break here in the Danish ("Dersom" flush left).
If it is now really true that the world-spirit, under present conditions
more precisely determined as the spirit of the age, conjures up its
illusions precisely in this, that scientific and religious rays of
light, in wildly crossing interferences, form juggling phantoms of mist;
if this, we say, is what is bewitching in the nature of the age, then
such a nature must surely also become manifest in corresponding
phenomena, and these phenomena contribute especially to a
characterization of the spiritual life of the present age. We do not
set ourselves up as judges over the present in its immediate actuality;
we consider only the phenomena insofar as they must be assumed, with
inner necessity, to accompany the spirit of the age.

Here, then, two extremes must first present themselves, two spiritually
opposed conditions. For the farther the individuals stand from
scientific enlightenment, the less will the results of science affect
them, the more easily will they be able to hold fast the religious
truths handed down in the Church and determined in the tradition; this
is the condition in which the preponderant majority of a country's
population naturally finds itself. Among the common people and the
half-educated who have still preserved some essential earnestness, the
word of revelation will then on the whole coincide with the word of the
Church and of the Bible and hold with divine authority. The nearer the
individuals stand, on the other hand, to the scientific light of
culture, the more many-sided their culture is, and the more criticism is
awakened, the more strongly will they be repelled by the wonders of Holy
Scripture, by its simple words about the omnipotence of spirit. Fully
assured of the
% --- p. 71 ---
\opage{71}\setcounter{footnote}{0}%
unconditional right of the scientific view of the world against the
biblical, these knowing ones will then leave positive religion to be the
affair of the common people, while they themselves, partly openly,
partly in silence, labor with all their might to undermine the
foundation of faith among the cultivated, to whom they then, when it is
finally demanded, offer as recompense and full compensation a
water-clear-unclear religion of thought, a religion of white on white
and gray on gray, ``the Absolute itself'' in logical shadow-outline and
metaphysical twilight.

Between these extremes surges the great mass, chiefly in the great
cities, the mobile multitude of the more and less cultivated, who make
up the age's knowing public, reflecting upon everything. If such a
public is considered from its shadow side---and it is expressly the
shadow side that we have before us here---then its mobility is to be
explained most immediately by the fact that everything sits loose in
consciousness: knowledge sits loose, and faith sits loose. Certainly
some faith still remains here, a dim impression of the mystical, a
practical respect for what is ecclesiastically established. But, it
goes without saying, the religion of the cultivated must of course be a
``perfectible religion,'' the priest of the cultivated a perfectible
man, cultivated faith revised by knowledge; then the thing is to hit on
the best arrangement and get an ecclesiastical order fixed firmly. But
to fix something firmly where the ground is loose has its difficulties;
so the public is worried, and its leaders worried, the unlearned worried
for knowledge, the learned for faith: yes, here are worries of every
kind; worries in the public and for the public, worries for the edifying
in worries for the scientific, worries enough for ``silken priests'' and
speculative theologians!

But in these worries there is no earnestness. And why not? Because it
lies in the nature of the case that there
% --- p. 72 ---
\opage{72}\setcounter{footnote}{0}%
can be none. When the hindrances to the spiritual life of the present
age are traced originally to an intellectual confusion, an unclear
mixture of religious and scientific considerations, then earnestness
must correspond to this; if it is really to correspond to it, its
fundamental task must be of an intellectual nature. The age must be
willing to fix its gaze on the first hindrances in order to discover in
them the first conditions; it must be willing to take the opposition
between the religious and the scientific sharply into view. However
great the gulf between faith and knowledge, religion and science, may
be, so great that it looks like a yawning abyss: still there is no
danger, if only thinking can be awakened to earnestness. The
consciousness of culture can, when it will, understand that here,
despite the opposition, there is something which the believer and the
knower must have in common, and that is a true conviction. The
consciousness of culture can, when it will, understand that true
conviction cannot possibly form itself, whether on the side of faith or
on the side of knowledge, without clear and thorough thinking, that
faith, therefore, as surely as it is an independent principle of
cognition, must have thinking with it, just as well as knowledge. The
cultivated man must, when he will, be able to see that thinking, even if
in relation to the absolutely heterogeneous principles it moves
differently in the different spheres, will nevertheless everywhere be
able to be understood by thinking, and the thinking consciousness
understand itself.
% No paragraph break here in the Danish ("Her er" flush left).
Here, in the many hindrances, is the possibility of just as many
conditions; here is reflection enough, and the task of spirit requires
reflection; here is culture enough, and the task of spirit requires
culture; here is freedom enough, and the task of spirit requires
freedom; but one thing is still needful: on this one thing everything
will depend, and it is---earnestness.

\begin{center}\rule{3cm}{0.4pt}\end{center}

\chapter*{V.}
\addcontentsline{toc}{chapter}{Lecture V}
\markboth{Lecture V}{Lecture V}
\label{ch:v}
% --- p. 73 ---
\opage{73}\setcounter{footnote}{0}%
``Ye shall know,'' says Christ, ``the truth, and the truth shall make you
free.'' Modern knowledge takes up and appropriates these words in a
rather peculiar way, and in this appropriation the world-spirit plants
its ambiguities. We cognize the truth by testing the difference between
true and untrue; the difference between true and untrue is tested by
thinking: against this there is nothing to object; but now comes the
distortion. Thinking, it is said, is thorough, clear, and pure only
insofar as it is scientific. To cognize the truth, then, means to
cognize in science what is valid and invalid in religion. The religious
that withstands the test in science can be acknowledged; what does not
pass the test must be rejected. According to this conception, then, the
light of knowledge is the highest light. Science sits on Pilate's
judgment seat; religion is summoned, and the judge asks: what is truth?
Yes, now the consciousness of culture understands what it means: ``Ye
shall know the truth, and the truth shall make you free''; it means, of
course: ``Ye shall know science, and science shall make you free.'' That
science, scientifically seen, works, has worked, and will continue to
work in a purifying and liberating way upon mind and thought is not only
to be
% --- p. 74 ---
\opage{74}\setcounter{footnote}{0}%
admitted, but, in order after all to avert foolish misunderstandings so
far as possible, to be impressed again and again. But liberation and
liberation are two things. One thing is the liberation of consciousness
from false opinions and ingrown prejudices; another is the liberation of
the soul from the bonds of corruptibility to an eternal life; one thing
is scientific, another personal truth. The truth of which Christ speaks
is not the universally valid, impersonal, rationally necessary truth of
science, but the personal, free, liberating truth. In speaking of the
truth, he speaks of himself about himself, and says expressly: ``I am
the way, the truth, and the life.'' Here is the knot that no science is
able to untie, the knot that an existing personality is the truth, yes,
that the eternal truth itself is this personality existing in flesh and
blood. Let the science that with a telling expression calls itself the
science of sciences, philosophy, exert itself with all its might; at
the very most it will bring matters so far as to develop the concept of
personality and of truth, and the unity of the concepts in the idea.
Confusing the universal idea of personality with personality itself, one
may then perhaps imagine that one has grasped the Absolute in the pure
concept and furthered the truth by saying in Christ's name: ``Ye shall
know philosophy, and philosophy shall make you free.'' But is that the
untying of the knot? It is surely a great misunderstanding. If
philosophy, in the name of knowledge, wishes to subject the Gospel to a
critique, to dissolve the historical Christ and transform the existence
of the God-man into an idea, it shall be allowed to do so; although, as
regards fruitfulness, one might well be tempted to the remark that long
straw that has been threshed so often will hardly yield any grain. But
what philosophy is in no way allowed to do is to pass off its critical
dissolution of Christ's
% --- p. 75 ---
\opage{75}\setcounter{footnote}{0}%
personal existence as a solution of the riddle of Christianity. That
Christ, according to the evangelical conception---and this must surely
count as the originally religious one---did not come into the world to
help philosophy to a general concept of the idea of personality, but to
give his life, his actual, personal life, for the salvation of the
world, he has plainly proved in word as well as in deed. We hear him say
in the Gospel of John (6th chap.): ``I am the living bread which came
down from heaven: if any man eat of this bread, he shall live for ever:
and the bread that I will give is my flesh, which I will give for the
life of the world.'' This is certainly a hard saying, for the knowing as
well as for the unknowing. And yet the speaker himself lays decisive
weight on his word. For when the Jews strove among themselves as to how
he could give them his flesh to eat, he said to them: ``Verily, verily,
I say unto you, Except ye eat the flesh of the Son of man, and drink his
blood, ye have no life in you.'' It is not the concept of bread but the
actual bread that satisfies the hungry, and it is not the idea of
personality but the actual personality that saves the believers. In
Christ's speech no distinction is made between knowing and not-knowing,
for to both kinds of listeners the speech can be equally offensive; but
a distinction is expressly made between believing and not-believing:
only for those who do not believe is it offensive. ``Doth this offend
you? What and if ye shall see the Son of man ascend up where he was
before? It is the spirit that quickeneth; the words that I speak unto
you, they are spirit, and they are life. But there are some of you that
believe not.''

But, it is said, granted that Christ gives himself, not as the idea of
personality but as actual personality, for the salvation of believers,
still he will surely,
% --- p. 76 ---
\opage{76}\setcounter{footnote}{0}%
as we hear from his speech, not be appropriated in a sensuous manner by
flesh-eating disciples and blood-drinking followers. No, his words,
which are life only insofar as they are spirit, must be appropriated in
spirit and in truth. But if in the end it comes down to the spirit and a
spiritual appropriation of Christ, then the content of the spirit must
surely be clarified by thinking, and by thinking the door of science
opens, and behold, then we are deep inside knowledge! Yes, we are deep
inside the misunderstanding, in full course of confounding the spirit of
revelation with the spirit of science. For the rest, the knowing ones of
the spirit of the age may judge of the spirit of revelation as they
will; but one thing they must be able to see without critical
overexertion, and that is that this spirit, in being and working, in
coming and vanishing, has not the least likeness to the spirit of
science. How did the spirit come upon the apostles at the feast of
Pentecost? ``And when the day of Pentecost was fully come (Acts ch.~2),
they were all with one accord in one place. And suddenly there came a
sound from heaven as of a rushing mighty wind, and it filled all the
house where they were sitting.'' If a scoffer here permitted himself the
remark that a spirit which comes with such a blast cannot be scientific:
upon whom, then, would the scoff fall back? Not upon those who separate
faith and knowledge! ``And there appeared unto them cloven tongues like
as of fire, and it sat upon each of them. And they were all filled with
the Holy Ghost, and began to speak with other tongues, as the Spirit
gave them utterance.'' It is necessary that a spirit's peculiarity be
characterized by its peculiar mode of working. But that the mode of
working of the language-giving spirit is here not merely somewhat
different but infinitely different from the scientific, no one will
venture to deny. For further confirmation we will hear from the apostle
Peter's own
% --- p. 77 ---
\opage{77}\setcounter{footnote}{0}%
mouth how the spirit comes and how it works. How, then, does it come,
and how does it work? ``And it shall come to pass in the last days,
saith God, I will pour out of my Spirit upon all flesh: and your sons
and your daughters shall prophesy, and your young men shall see visions,
and your old men shall dream dreams.'' Is there in such utterances the
least hint that the apostolic spirit will ever be able to identify
itself with the spirit of science? Certainly Christ will be appropriated
in spirit and in truth, but the appropriation is not scientific: on the
contrary, he says in Matthew, chap.~2, the well-known words:
% Danish "Matthæi 2det Kap." as printed; the passage is Matt. 11:25.
``I thank thee, O Father, Lord of heaven and earth, because thou hast
hid these things from the wise and prudent, and hast revealed them unto
babes.''
A stronger protest against the conception of science he could not lodge
than when he expressly says that only babes understand him, and that the
wise and prudent cannot understand him at all. The opposition between
the babes and the wise is not an opposition between the thoughtless and
the thinking; Christ has never said that he could not be understood by
the thinking, but only by the thoughtless; how should he well be able to
say to the thoughtless: ``Ye shall know the truth, and the truth shall
make you free''? No, by the babes he means the simple, who begin with
faith, whose cognition must therefore from first to last become a
cognition of faith; by the wise and prudent he has in mind the knowing;
and then his word holds literally, that the heavenly Father has hidden
the truth of the Gospel from the knowing and revealed it to the
believing.

That Christ will be appropriated not by the knowing but by the
believing is seen also in his apostles. We must surely admit that the
apostles appropriated Christ's personal life in a
% --- p. 78 ---
\opage{78}\setcounter{footnote}{0}%
true, but not in a scientific way. They did not, by thinking about his
personality, clarify for themselves the objective, universally valid
concept of personality. Certainly Paul is a thinker, one of the greatest
in the world, but his thoughts are not scientific. It is a downright
mistreatment of the great thinker to offer his doctrinal conception in
the form of a scientific system, to take, if I may say so, the doctrinal
conception up into the lecture chair while the apostle himself does not
come into the lecture chair. Of a scientific thinker, a philosopher for
example, it holds that when one tests his knowledge, his person is a
matter of indifference. But when one tests an apostle's teaching, one
tests at the same time his life. The apostle's thoughts are thoughts of
personality, his so-called doctrinal conception a personal explanation
of his faith, his life in Christ, his self-sacrifice for the
congregation. How the spirit, where it works in power, ceaselessly
transforms hindrances into conditions, while the world conversely
presents these conditions as sheer hindrances, can, as we have already
seen, be developed from the thought of faith, and the thought of faith
be developed from the principle of faith, without its thereby being
settled that he who thus develops the thought of faith is himself a
believer. It is otherwise with an apostle; he develops the thought of
faith in explaining his life. We choose an example from Paul: ``Now ye
are,'' he says to the ``puffed up'' Corinthians (1 Cor.~4), ``full, now
ye are rich, ye have reigned as kings without us: and I would to God ye
did reign, that we also might reign with you. For I think that God hath
set forth us the apostles last, as it were appointed to death: for we
are made a spectacle unto the world, and to angels, and to men. We are
fools for Christ's sake, but ye are wise in Christ; we are weak,
% --- p. 79 ---
\opage{79}\setcounter{footnote}{0}%
but ye are strong; ye are honorable, but we are despised. \dots\ Being
reviled, we bless; being persecuted, we suffer it; being defamed, we
entreat: we are made as the filth of the world, and are the offscouring
of all things unto this day.'' See, that is a part, a very essential
part, of ``the Pauline doctrinal conception''; but such a doctrinal
conception cannot be filtered into scientific systematics unless one
filters away the Pauline existence. But when one has got that unity of
thinking and acting which characterizes the apostolic existence
scientifically purified out, what then? Why, then the doctrinal
conception is perhaps theological, but it is no longer apostolic.

Another distinguishing mark by which we can easily make the separation
between scientific and religious thinking is the miracle. It is a
characteristic mark of revelation: without miracle no revelation, and
without revelation no faith. If the miracle is taken away or pushed
aside, then the character of revelation is blurred, then faith is
weakened. Let science say of the miracle what it will; it cannot deny
that the miracle is the characteristic mark of religion. To want to
explain and ground religion scientifically without wanting to explain
and ground the miracle scientifically is nonsense: if there is a science
of religion, then there must also be a science of miracles. But the
miracle is the barrier of knowledge; it is the skandalon of
science. Science can, when it is made judge, deny the miracles, and
with them religion; there is sense in that; but it cannot cognize
religion and blur out the miracle. Christ's personality is a miracle. It
is not enough to seek his peculiarity in this, that he has done many
miracles: no, Christ is the miracle itself, the great miracle of
revelation. With this it agrees consistently that the spirit which is to
explain Christ comes
% --- p. 80 ---
\opage{80}\setcounter{footnote}{0}%
by a miracle, that faith, in which the gifts of the spirit are
appropriated, is given by a miracle; from every side knowledge here
strikes against its original barrier, its skandalon, namely the
miracle.

It is true that people have wanted to take the matter from another side,
saying that the miracle of religion is indeed not taken up into worldly
knowledge, but that precisely for that reason a distinction has been
made between worldly and sacred science. As sacred science, theology
does not presume to want to comprehend the miracles, for the miracles
are incomprehensible; no, not by comprehending, but by presupposing the
incomprehensible and bringing out the facts of revelation and of faith
in sharp strokes, theology would furnish a doctrine of faith in
scientific form and reconcile faith with knowledge.

This standpoint has certainly been chosen with the best intentions, in
the guileless presupposition that it would one day succeed, on the
foundation of faith, in raising a new scientific edifice. But it will
never succeed. Already this, that a foundation is laid here consisting
of supernatural facts, hence of unscientific presuppositions, and that
upon this foundation one would erect a doctrinal edifice of something
scientific,---already this must, as has been shown in what preceded,
arouse misgiving.

It has been said that the believing dogmatician relates himself to the
sacred facts in the same way as the natural scientist to the facts of
nature; what an enormous misunderstanding! The natural facts show
themselves, when the question is one of knowledge and science, as
entirely heterogeneous with the sacred. The facts of nature are subject
to the laws of nature; here one has the universally valid, the
rationally necessary, that which makes testing and critique possible.
What such a rationally necessary, universally valid foundation means can
be seen
% --- p. 81 ---
\opage{81}\setcounter{footnote}{0}%
even in material forms. If it is asked, for example, what granite and
flint are, these concepts, with the corresponding facts, can be fixed so
definitely that everyone who has sound understanding, preliminary
training, and a sense for scientific investigations---be he, for the
rest, Jew, Greek, Mohammedan---without regard to national differences
and faith, is compelled to admit the validity of the concepts: ``This
must be granite, this must be flint.'' So also when the question is of
sulfuric acid, carbonic acid, and so on. If, on the other hand, one
looks to such supernatural facts as turning water into wine, walking on
the sea, raising the dead, then it is impossible to present such facts
in such a way that they take on the stamp of something universally valid
or incontrovertible. Greeks, Jews, Roman Catholics, Protestants surely
see the sacred phenomena with different eyes. Those who believe science
cannot make unbelieving; those who do not believe it cannot make
believing: here between faith and unbelief is a personal opposition
which science cannot master. The essential difference between the
natural and the sacred facts is seen, in this connection, from a new
side. The natural facts are present. The objects of nature lie ready for
ever-renewed testing, for the realm of nature is an ever-present realm.
It is otherwise with the sacred facts; they can no longer become the
object of renewed seeing for oneself. He who did the miracles is, after
all, no longer sensuously present; the time of miracles is past:
``There happen,'' people say, ``no miracles in our days.'' The
supernatural facts of revelation belong to the past. They are not to be
compared with the facts of nature, but rather with those of history.
The proposition is therefore altered to this, that the believing
dogmatician
% --- p. 82 ---
\opage{82}\setcounter{footnote}{0}%
relates himself to the sacred facts in the same way as the historian to
the historical facts. Will this work? The cognition of the facts of
history does indeed presuppose a historical faith; but the facts of the
Gospel expressly demand a religious faith, and between religious faith
(fides religiosa) and general historical faith (fides historica) there
is a great gulf fixed. The historical facts science subjects to a
critique; may we now also, asks the consciousness of culture, apply
critique to the sacred facts? Yes or no? If the answer is no, it is of
course instantly clear that the facts of sacred history acquire an
entirely different position for consciousness from the facts of profane
history. If biblical criticism is to be kept out by a dogmatic decree,
then we must necessarily return to the old doctrine that every word,
every tittle in Holy Scripture is inspired by God, that the Testaments,
the Old as well as the New, are in each and every respect infallible
documents of revelation. Scientific reflection must suddenly halt---by
what? By authority---blind authority! And yet it is to be scientific!
Instead of reasons, fully valid, scientific reasons, our dogmatician
begins his ``sacred science'' by invoking blind authority. Let us
consider blind authority a little more closely! The Mohammedans, for
example, also have their document of revelation, their sacred facts,
their authority. Besides the Koran the Sunnites acknowledge a living
tradition, and in this tradition a multitude of divine fairy tales about
the Prophet. He was, it is told, once mocked by the unbelievers because
he could not at once prove his divine mission by a miracle. Long the
mocked one kept meekly silent; but at last his hour struck. The
half-moon
% --- p. 83 ---
\opage{83}\setcounter{footnote}{0}%
stood just then shining bright on the vault of heaven; he called to it,
and it fell with a silver ring to the earth. Then he took it up and
carried it under his arm, as one carries a chapeau-bas. When this was
done, he threw it up into the air; it then flew back to heaven and shone,
as before, on the vault. That is a sacred fact; woe to the Sunnite who
cannot believe in such a miracle. By blind authority critique can
certainly be halted; but, we ask, can such a halting be regarded as the
introduction to a scientific doctrinal conception? Can blind authority be
united with Christ's words: ``Ye shall know the truth, and the truth
shall make you free''? One would then have to interpret them thus: ``Ye
shall know blind authority, and blind authority shall make you free.''

Without scientific critique there can be no talk of science. Critique,
then, is conceded---with moderation, caution, reverence, of course. But
here the proposition holds: if one has said A, one must also say B. If
critique is to be applied, it must be applied consistently, not merely
to natural events, but also to miracles and prophecies. Once the
scientific point of view has been fixed firmly; once one has credulously
entered upon the radically false presupposition that the sacred facts
are to be seen in the same light, have
% Danish prints "Virkeligked" (printer's error); read "Virkelighed".
their actuality in the same sphere, and establish their reality in the
same way as the profane: then personal truth says good night, for now
what is revealed is to be inspected inside out. It is an easy matter for
``modern science,'' after it has, with the assistance of theology, won a
secure foothold in the sphere of religious cognitions, to show how the
miracles make up a web of contradictions, how the prophecies aim at
something entirely other than their alleged fulfillment,
% --- p. 84 ---
\opage{84}\setcounter{footnote}{0}%
and how the biblical accounts conflict not merely with other historical
accounts, but even with one another.

We choose as an example the miracle by which Christ raises Lazarus, a
miracle that is to prove the truth of his words: ``I am the resurrection,
and the life.'' The impression the miracle makes on the many witnesses is
extremely various, some honoring Christ in faith, while others go as
secret informers to the Pharisees. And what then do the Pharisees do?
``Then,'' it says (John 11),
% Chapter numeral blotted in the Danish copy; read 11 (John 11:47--53). Doubtful.
``gathered the chief priests and the Pharisees a council, and said, What
do we? for this man doeth many miracles. If we let him thus alone, all
men will believe on him: and the Romans shall come and take away both
our place and nation \dots\ Then from that day forth they took counsel
together for to put him to death.'' But what kind of psychology is this?
The chief priests and the whole council begin, of course, as believers,
with the presupposition that Jesus really ``doeth many miracles,'' indeed
works so wonderful that if he were allowed to go on thus, then ``all men
would believe on him.'' The raising of Lazarus thus stands before the
consciousness of these scribes not as a deception, a sham miracle
performed by a crafty juggler, but as a fact known to many, a wonder
which it cannot occur to them to deny. But how in all the world can it
then occur to such clever men to want to kill a man who can raise the
dead? If a man is really assumed to have power to raise the dead, then
no religious faith is needed, but only a little sound common sense, to
see how useless, indeed how senseless, it is to seek such a man's life.

The Pharisees and the scribes, one will perhaps
% --- p. 85 ---
\opage{85}\setcounter{footnote}{0}%
say, believed only historically, but could not believe religiously, for
God had hardened their heart. To get even only a glimmer of science out
of all this, one would then have to be able to explain how much of the
hardening belongs to these scribes and how much comes from God. But if
one makes earnest of such a task, surely everyone must be able to see
that the hardening of hearts will not resolve into any scientific
concept. What a swarm of contradictions can scientific critique not
point out here! It says in the Epistle of James (1,~13---14): ``Let no
man say when he is tempted, I am tempted of God: for God cannot be
tempted with evil, neither tempteth he any man.'' But to harden a man's
heart is surely to tempt it to evil. ``Therefore,'' says Christ, ``speak
I to them'' (the Jewish people) ``in parables: because they seeing see
not; and hearing they hear not, neither do they understand. And in them
is fulfilled the prophecy of Esaias, which saith, By hearing ye shall
hear, and shall not understand; and seeing ye shall see, and shall not
perceive.'' Theology would level out the contradictions in such a way
that the explanation should at one and the same time satisfy faith and
content knowledge; in all obligingness it would serve two masters, but
gets ingratitude from both. Ceaselessly confusing the psychology of
faith with a scientific doctrine of consciousness, it succeeds only in
preparing an easy victory for modern science and in driving its negative
critique to the extreme.

By wanting to defend the sacred facts scientifically, theology
contributes to keeping up the deception that these facts can in earnest
be threatened by scientific attacks. But if the sacred facts can be
threatened by scientific attacks, then there lies in this an admission
that they might possibly one day fall before such attacks. Modern
science
% --- p. 86 ---
\opage{86}\setcounter{footnote}{0}%
has turned the admission to its account and answered with triumphant
scorn: ``They have fallen!''

There runs through religious psychology a peculiar series of doublings
of reflection; of these doublings of reflection, with their immediate
outcome in apparent contradictions, science has no concept, and
can---so long as faith's own thinking has not yet been allowed a clear
and many-sided development of its own cognitions, but every time it
would begin on a systematic self-development is instantly twisted round
into the scientific---cannot possibly have any concept. To indicate what
is aimed at here, we compare a couple of passages of the Old Testament.
In the second book of Samuel (XXIV) the following is related: ``And
again the anger of the Lord was kindled against Israel, and he moved
David against them to say, Go, number Israel and Judah!'' And the king
does expressly as Jehovah has commanded him, and bids Joab, who was
captain of the host, to go through all the tribes of Israel from Dan to
Beersheba and number the people. Nevertheless David, just as the census
is ended, gets a notion that he has done something evil. ``But,'' it
says, ``David's heart smote him after that he had numbered the people.
And David said unto the Lord, I have sinned greatly in that I have done:
and now, I beseech thee, O Lord, take away the iniquity of thy servant;
for I have done very foolishly!'' Thus David, by simply acting according
to Jehovah's express command, has committed a sin; thus God has
expressly tempted David to evil. If that is not what is said here in
plain words, what then is said here? Theological interpretation so
gladly proves its scientific character by collecting parallel passages
and letting ``the Bible interpret itself'' by letting the one passage
explain the other. The parallel passage to second Samuel (XXIV) is, as
is well known, first
% --- p. 87 ---
\opage{87}\setcounter{footnote}{0}%
Chronicles (XXI), which begins thus: ``And Satan stood up against
Israel, and provoked David to number Israel. And David said to Joab,''
and so on. Let one turn about hermeneutically and twist oneself
exegetically however one will, it stands expressly in the one passage
that Satan does what in the other passage is said to be done by God
himself. And that God (2~Sam.) does not let Satan tempt David, but does
it himself immediately, shines out plainly enough from the unambiguous
words: Go, number Israel and Judah!

If we cast a glance at the scientifically critical conception of the
prophecies, theology fares not a hair better. The freethinker
criticizes, and the theologian criticizes, and the more lights critique
kindles, the paler religion becomes. The prophecies in the Old Testament
that are fulfilled in the New have been carefully enumerated and
classified, but seen through the spectacles of critique they fall away
altogether; for what the prophets say is something other than what the
evangelists make them say. A single example! In the Gospel of Matthew
(II), where the flight into Egypt is told, it is said that this happened
that it might be fulfilled which was spoken of the Lord by the prophet,
saying: ``Out of Egypt have I called my son.'' It is Hosea (XI, 1) that
is aimed at here. But in Hosea it stands: ``When Israel was a child,
then I loved him, and called my son out of Egypt.'' For critique there
is no prophecy here; for by the son of whom the prophet speaks he means
not Jesus, born of Mary, but the people of Israel. Now, to be sure, one
gets the ingenious explanation that the people of Israel is a type of
Christ, but ingenious circumlocutions we have enough of. What critique
sees at the first glance is that the prophet's thought goes in one
direction, the evangelist's in another. The prophet speaks of something
past, but
% --- p. 88 ---
\opage{88}\setcounter{footnote}{0}%
the evangelist takes him as if he had spoken of something to come. There
is no prophecy in the prophet's speech, and yet the evangelist would see
in it the fulfillment of a prophecy.

What, now, is the consequence of these scientific attempts? One knows
names such as Strausz, Rénan, and others.
% "Rénan": accent as printed in the Danish.
It is common knowledge that Strausz, with scientific thoroughness and
critique, has dissolved the whole evangelical history and made Christ's
life on earth into a myth. It is the world-spirit that through him
completes its destructive consequences; the old theology has long said
A, and now Strausz comes at last and says B. With the help of the
Straussian critique one then becomes free not only of the many
scientific contradictions with which religion has hitherto been
encumbered, but one becomes scientifically free of religion itself; then
critique is true, and truth critical; then we interpret Christ's words
in the spirit of the age: ``Ye shall know critique, and critique shall
make you free.''

If it were now the task, the goal we had set ourselves here, to clear up
the psychology of faith and to develop the religious concepts from
religion's own principle, then we should surely be obliged in the
present connection to indicate how the misapprehensions pointed out
could be removed point by point. But since we have not proposed to
impart either a psychology of faith or a philosophy of religion in
detailed lectures, but rather to speak in general about hindrances to
and conditions for the spiritual life: a pointing to the mistreatments
that religion must suffer from the scientific character introduced by
theology will already be able to make a contribution. It is, from what
we have already seen, by no means an immediate enthusiasm for the ideal,
but at the most an intellectual striving after the cognition of truth,
that can be said to mark
% --- p. 89 ---
\opage{89}\setcounter{footnote}{0}%
the spiritual merit of the age. But an age whose main strength is
intellectual must first and foremost labor to clear itself out of
intellectual confusions. We say it freely and with full conviction that
the first and greatest hindrance to the life of spirit in the present
time is a theological bungling of the sacred facts, a dogmatic mixing
together of religious and scientific cognitions, a critically hollowed
out but officially propped up system of authority, which leans on faith
when knowledge runs short, and defends itself with knowledge where faith
fails it. We say it freely and with full conviction that the first and
greatest of the many intellectual tasks of our intellectual time is, as
surely as it would preserve and further its spiritual life, this task:
to get religion, I mean revealed religion---not by a decree of
orthodoxy, not by ecstatic revivals and mystical visions, but by
thinking, a thorough, clear, illuminating thinking corresponding to
faith---once for all delivered and rescued from the claws of science.
Until this task is thoroughly solved, hindrances and conditions must go
on clashing in dull fermentation; but with the solution of the task
truth will conquer by the power of spirit, not the impersonal but the
personal; not the cold truth that is grounded in necessity, but the
living one that shines in freedom, the very one that is heard in
Christ's words: ``Ye shall know the truth, and the truth shall make you
free.''

\begin{center}\rule{3cm}{0.4pt}\end{center}

\chapter*{VI.}
\addcontentsline{toc}{chapter}{Lecture VI}
\markboth{Lecture VI}{Lecture VI}
\label{ch:vi}
% --- p. 90 ---
\opage{90}\setcounter{footnote}{0}%
The struggle of the consciousness of culture against the spirit of the
Gospel has already been characterized from two standpoints, namely
insofar as the movement of thought passes through natural science and
through history. It now remains to grasp and characterize this struggle
from the standpoint of the doctrine of spirit itself. Here the
battlefield is inward; the struggle reaches its decisive turning point,
in that it ends as a struggle within consciousness itself. If we ask the
man of culture whether, in consequence of the negative world-spirit, he
will deny the Highest, the ideal, the eternal, he answers No; if we ask
him whether he will deny the Godhead itself, he answers No; but if we
add yet the question whether he will then believe in God and his
revealed Word, he also answers No. For the consciousness of culture
grasps its Godhead not as almighty willing personality, but as the
absolute idea of knowledge. In faith God is taken to be a personal
being; but, says the knower, if God is really to be thought of as a
personal being, then he must surely also be thought of as a limited
being. ``You ascribe to God,'' says Fichte to the believers,
``personality and consciousness. What then do you understand by
personality and consciousness? Surely that which you have found in
yourselves and designated by this name? But that you in no
% --- p. 91 ---
\opage{91}\setcounter{footnote}{0}%
way think, or can think, this without limitation and finitude, the
slightest attention to your construction of this concept can teach
you. By ascribing that predicate to this being, you therefore make it a
finite one, a being like yourselves, and you have not, as you wished,
thought God, but only multiplied yourselves in thought.'' The
self-consciousness of the absolute spirit is indeed, according to Hegel,
religion; but if anyone concludes from this that speculative philosophy
really ascribes to God an eternal, self-subsistent self-consciousness,
he will think otherwise when, by penetrating deeper into the system, he
learns to see the meaning of the remarkable words that ``man's
knowledge of God is God's own knowledge of himself.'' From this
philosophical summit the consciousness of culture will then continue its
intellectual struggle to root out faith.

And why will the consciousness of culture root out faith? If its
standpoint is valid, then the endeavor is not merely consistent, but
noble and good. Here there are, outwardly as well as inwardly, spiritual
phenomena enough that could give offense to the enlightened man. A
thinker who was zealous for faith has somewhere raised the question:
``If what we understand by being a Christian really is being a
Christian, what then is God in the heavens?'' And the answer is given:
``He is the most ridiculous being that has ever lived, his Word the most
ridiculous book that has ever come to light: to set (as he indeed does
in his Word) heaven and earth in motion, to threaten with hell, with
eternal punishments, in order to attain what we understand by being
Christians (and we are, of course, true Christians)---no, anything so
ridiculous has never occurred.'' The glaring contradiction between the
ideals of faith and actual life, which characterizes the age of means,
of conditions, of interests,
% --- p. 92 ---
\opage{92}\setcounter{footnote}{0}%
stands as clear to the knowing critic as to the believing thinker. When
the believing thinker, wounded by the sting of the contradiction, bursts
out: ``The most dreadful kind of blasphemy is that of which Christendom
is guilty: to turn the God of spirit into ridiculous twaddle; and the
most spiritless kind of worship, more spiritless than all that is and
was found in paganism, more spiritless than this: to worship as God a
stone, an ox, an insect, more spiritless than all is: to worship under
the name of God---a twaddler!''---when, we say, such outcries sound from
the camp of the believers, then negative criticism answers Yea and Amen.

The fact of the age, that faith has lost its inspiring power, the man of
culture does not grasp as a sign of a spiritual retrogression, but on
the contrary as proof of an essential progress. The reason why the power
of faith is waning is, according to his conviction, that the light of
knowledge is waxing. It is not merely, he says, the objects of faith,
the holy facts, that dissolve into the fabulous and lose themselves in
the mythical; no, faith itself is of a fabulous nature, the believing
consciousness is essentially mythical. We shall now hear the closer
explanation.

How the mythical consciousness comes to grasp the various deities can be
seen in Homer. Already in the first book of the Iliad, where Achilles
and Agamemnon are in their well-known quarrel with one another, there
occurs a turning point characteristic of a revelation of a god.
Achilles' passion is inflamed; he is so furious that he calls Agamemnon
a coward and uses the strongest words of scorn; there is every sign that
he will forget himself. ``Under his shaggy breast his heart beat
doubtfully, whether he should straightway draw his whetted blade from
his thigh, scatter the thronging crowd, and cut
% --- p. 93 ---
\opage{93}\setcounter{footnote}{0}%
him down, the son of Atreus, or master his hot temper and curb his
wrath''; then the revelation takes place. ``Just as he drew the sword
from the sheath, Pallas Athene came down from heaven\dots\ Behind
Achilles she stepped, and pulled at his light-yellow locks, visible to
him alone; of the others none saw her.'' The conversation which now
passes between Achilles and the goddess is highly characteristic.
``Daughter of aegis-bearing Zeus, why hast thou come?'' To which Athene
answers: ``I have come to curb thy wrath, if thou wilt obey.''

How characteristic is every feature of this revelation! What Achilles in
the critical moment so sorely needs is self-possession, and behold,
there stands self-possession itself visibly before him in Pallas Athene.
And he knows her, yes, he knows her at once; her ``dreadful glance
flashes towards him.'' Self-possession is a personal matter: so the
goddess is visible only to him; she has an errand to him alone. And yet
she surprises him and makes him start! Yes, self-possession comes with
divine power, with a calm exalted above the passions, with authority,
with ideal superiority. Athene is no allegory; she is a daughter of
aegis-bearing Zeus: she is a goddess. Here, then, we have the mythical
doubling; in other words: here we have ``faith.'' Achilles believes in
the goddess, and just for that reason she reveals herself to him.
Without himself knowing it, he refers the self-possession awakening in
his own mind to a divine being existing outside him, and then this being
comes ``down from high Olympus'' and brings him what is originally his
own. Here is the doubling of consciousness, the doubling which
characterizes faith. The believing consciousness sets the gods up
outside itself, and does so from inner necessity, because in its obscure
profundity it
% --- p. 94 ---
\opage{94}\setcounter{footnote}{0}%
does not yet know that it has the divine in itself. It is precisely the
obscurity of the mind, with the heaviness and tension corresponding to
the obscurity, that causes the believer to receive the revelation. The
believing consciousness lives in feeling and in the intuition of
fantasy; therefore it has so irresistible an urge to present its divine
content to itself in corresponding images; and when one of these images,
in striking immediacy and as with a magical power, forces itself upon
consciousness, then man imagines that the divine which he sees in his
own vision is a real god's revelation: on this imagining faith rests.
All religions originally give themselves out as revealed; this is from
the beginning no simple priestcraft; far from it! it is a
self-deception of the soul, necessary by the idea; he who has seen the
primitive vision of the gods really believes, fully and firmly, that he
has received a revelation. % AUDIT doubtful in the Danish: "idee-nødvendig" (compound) vs "ideenødvendig"; sense unaffected
Then he reflects upon his vision; he recalls it, and in his recollection
it becomes ever clearer, richer, more wonderful; then he describes his
vision and tells of the revelation, and those who hear the account with
devotion are seized with wonder, perceive within themselves a testimony
that what is told them must be divine truth. These devout ones now
become the believers proper, in that, without themselves having received
any revelation---be it that the originality in their nature is too weak,
or that they live at a later time, after the primitive breakthroughs
have already taken place---they believe in the revelation that has been
granted to another, a worthier one. In the believers proper,
recollection then works further; the legend of the gods (Μῦθος) is
expanded, altered, completed; the idea, working in the form of fantasy,
soon makes the myth a subject for poetry; now we stand again at Homer.

To prove that there is a consciousness of revelation
% --- p. 95 ---
\opage{95}\setcounter{footnote}{0}%
that is not mythical, one appeals to the morally pure, the truly
ethical, to the necessity of grounding the ethical upon the religious.
But how insufficient such a line of proof is can be seen expressly from
the example we have cited. Or can one well conceive of any purer fusion
of the religious and the ethical than precisely in the lines exchanged
here between Achilles and Athene? ``Why hast thou come?''\,\dots\ ``I have
come to curb thy wrath, if thou wilt obey, come from heaven; the
lily-armed Hera sent me, she loves you both alike, and tenderly she
cares for both.'' So exalted is the religious-ethical thought here that
Achilles in his seething anger must be lifted above himself and take
note that there is in heaven a divine love which embraces him and his
enemy alike. And this divine love, how human it is, how humane! The
goddess with ``the dreadfully flashing glance'' does not threaten him,
does not compel him; but gives him counsel, if he will obey. The demand
for obedience is, then, not excessively strict either. He is merely to
refrain from using force, but may otherwise give vent to his wrath and
rebuke Agamemnon with words as much as he pleases. By obeying the
goddess's counsel Achilles here shows himself a model of Greek piety and
true religiousness. ``It is fitting, goddess, to take your word to
heart, however wroth one be in mind, for this is most profitable; whoso
is obedient to the command of the gods, him they gladly hear.''

This, then, was the man of culture's conception of faith and mythology
at the Greek stage of consciousness; how does he now make out the
doctrine of Jehovah, the revelation in the Old Testament? At bottom in
the very same way.

It is an apt saying of a modern author
% --- p. 96 ---
\opage{96}\setcounter{footnote}{0}%
that man plunders himself in order to adorn his God. This holds fully as
much in Judaism as in paganism. The best in the Greek was his beautiful
individuality. In this individuality there was something corruptible;
that he could keep for himself; but the incorruptible essence of
individuality, the divine, that he dared not keep for himself; it
belonged to the blessed gods alone. The best in the Jew is his
personally thinking I, his ethically willing self. In this I there is
something limited, unclear and transitory; that man may keep for
himself; but in the I as I, in spirit as spirit, there is a claim to
holy wisdom and eternal power which man cannot satisfy; infinite wisdom,
holy will, eternal power belong to the God alone. Out of the burning
bush Jehovah speaks (Exodus III) to Moses, and gives himself the
significant name: \emph{I am that I am}. In this almighty Jehovah-I the
people of Israel now intuits the eternal ideal of its spiritual essence,
in a manner corresponding to that in which the Greek intuits the ideal
of individuality in the Olympian gods.

There is indeed a difference here, yes, when we look more closely, even
a very essential difference; but the difference lies, scientifically
viewed, by no means in this, that Jehovah is real while the Greek gods
are unreal: the difference between ``the idols'' and ``the only true
God'' is only a difference of faith, a difference posited by the Jewish
imagination. No, the real difference is that which follows from the
diversity of the ideals. The Greek gods are many; for individuality has
many shapes; the God of the Jews is only one; for the I as I is, despite
all the changes of the world, the same. The Greek gods are, though
invisible to a mortal eye, yet visible according to their essence; they
are image-gods, true gods of art;
% --- p. 97 ---
\opage{97}\setcounter{footnote}{0}%
for in the form of art the individual is transfigured. Jehovah, on the
other hand, even if he now and then assumes a visible shape, is
invisible according to his essence; for thought as thought, will as
will, spirit as spirit is in all eternity invisible. But just as little
as the believing Jew knows that in the imageless Jehovah he intuits the
ideal of his own spiritual self, just so little does the believing Greek
know that in his image-gods he intuits the ideal of his own human
individuality. If the Jewish cognition of God is to be called faith,
then the Greek is faith too; if the Greek faith is to be called mythical,
then the Jewish faith is mythical too. Seen in faith, it is indeed
Jehovah who created man in his image; but seen in knowledge, it is the
Jew himself who created Jehovah in his image. The holy God of
omnipotence, of wisdom, of goodness, of justice is a true man-god, the
Creator of the world a creature of man's spirit.'' % printer's defect: closing quotation mark with no opening mark (the "modern author" citation opens unmarked on p. 96); carried as printed

According to this way of viewing things, it is so far from being the
case that the Jewish faith should be less mythical than the Greek, that
the mythical in the doctrine of Jehovah, even in many striking features,
comes out more strongly; this is because the phenomena of revelation in
Jewish mythology stand in a far more glaring conflict with the concept
of God itself than in the Greek.

God is almighty; the almighty God is the Creator of heaven and earth;
now comes the mythical. For the Almighty does not create as an Almighty;
no, he goes to work piece by piece; he conducts himself like a human
master-builder, rather like Plato's God in the \emph{Timaeus}. First he
talks with himself and considers what he will do; then he does it, and
then afterwards he looks at what he has done, and sees that it is well
done. That such a creation must be an object of faith
% --- p. 98 ---
\opage{98}\setcounter{footnote}{0}%
goes without saying; for as soon as it is made an object of knowledge
and seen in the concept, it betrays its contradiction and vanishes. If
Jehovah is almighty, then his power is absolute; but absolute power can
just as little engage in finite actions as in finite resolutions. It is
said that Jehovah walked in the garden of Paradise in the cool of the
day. How cozy! The almighty Creator of heaven and earth takes a stroll
in the evening cool under the shade of the trees which he himself has
lately created. That a Zeus now and then takes a stroll in the grove is
after a fashion in keeping with his erotic individuality; but that an
Omnipresent One takes a stroll is a monstrous contradiction. How long
can such a thing continue to be an object of faith? Answer: as long as
consciousness can continue to be mythical. That the Jewish people, this
so peculiar, spiritually vigorous people of antiquity, has had a
religious childhood, a stage of consciousness in which such narratives,
attractive by their guilelessness and childlike grace, came forth with
psychological necessity and just for that reason worked with
psychological truth, is understandable. There is in these religious
nursery tales of ``the chosen people'' something so bold and so
grandly simple that it must impress the imagination even to this day.
``And the Lord God made unto Adam and to his wife coats of skins'';
just imagine: he sewed their clothes himself!

What for the childlike intuition of antiquity had its psychological
truth becomes in our time psychological untruth, affectation and
self-deception, when, in defiance of science and sound reason, under the
name of ``childlike faith,'' it is still to be taken straightforwardly
as revelation. We will not dwell on the many surprises and violent
emotions by which Jehovah, forgetting his omniscience and omnipotence,
lets himself be carried away---when
% --- p. 99 ---
\opage{99}\setcounter{footnote}{0}%
he, for example (Genesis VI), is so surprised by the wickedness of men
that ``it repented him that he had made man on the earth, and it grieved
him at his heart''---we will not, as we said, dwell on what is known to
all; only as an example of how far psychological untruth can go in our
days, we will cite one more single feature. It is told (Genesis VII)
that when Noah with all the animals had safely come into the ark, then
``Jehovah himself shut him in.'' To this one of the present age's
childlike believers has now added the following interpretation, as
learned in the Scriptures as it is edifying: ``Jehovah \emph{shut}, for
from outside Noah could not, of course, close it (the ark) any more, so
that the water could not get in; also he must, of course, be secured
against the attempts of the ungodly to be able to break the ark open in
order to get in.'' Yes, the faith of the present age---is childlike
faith!

If we look back on what has been cited here, how then can we doubt that
the Jewish faith is just as mythical as the Greek? Or what should
entitle us to declare Athene a false goddess, when we acknowledge
Jehovah as a true God? That her existence is in conflict with the laws
of nature? But so is Jehovah's existence! That in Homer she sometimes
lets herself be carried away by passion? But so does Jehovah in Moses!
That Jehovah's fundamental character, when we disregard his particular
weaknesses, is ethical? But Athene's fundamental character, when we
disregard her particular weaknesses, is ethical too! Her revelation to
Achilles is then at least as ethical as Jehovah's repentance at having
created men. No, if Jehovah is to be a true God, then let Athene too be
a true goddess. Nor need we then, for faith's sake, fear any logical
contradiction. There is, if only we will believe in
% --- p. 100 ---
\opage{100}\setcounter{footnote}{0}%
revelations and apparitions, absolutely nothing to prevent Athene and
Jehovah from being two different revelations of one and the same
``almighty Creator of heaven and earth.'' We know, after all---we know
it in faith---that Jehovah himself together with two angels ate curds
and milk and veal with Abraham under the tree in the grove of Mamre
(Genesis XXIII); % "XXIII" as printed (the Mamre meal is Gen. XVIII)
if the Almighty can condescend to eat milk and veal when he will reveal
himself to an Abraham, why should not the same Almighty also be able to
condescend to appear as a daughter of aegis-bearing Zeus when he will
reveal himself to an Achilles?---So far the man of culture.

That in the often acute attacks of modern knowledge upon the faith in
revelation there lies an essential hindrance to the development of the
religious life shall not be denied; but if we look more closely, it
might yet turn out that precisely this hindrance, if only the solution
of the fundamental intellectual task of the present age could be taken
in earnest, lets itself be transformed into an essential condition. The
sharper and more penetrating the thinking that attacks faith, the
sharper and more penetrating must also be the thinking by which faith
defends itself against the attack. Sharp against sharp, and then it
will surely come to a decision! Here, then, there is in the attack of
triumphant science, in the first place, an arbitrariness which must be
repelled most emphatically. Here it is presupposed without further ado
that faith in itself is thoughtless, that all clear, consistent and
thorough thinking belongs exclusively to science: to apply the measure
of thinking is to apply the measure of science. This presupposition is
utterly false. Relying on an utterly false presupposition, scientific
criticism begins, quite unembarrassed, by treating the concepts of faith
as
% --- p. 101 ---
\opage{101}\setcounter{footnote}{0}%
immature, botched concepts of knowledge, and can then easily prove that
these concepts contain scientific contradictions. With the false
presupposition that faith is thoughtless, criticism naturally finds it
beneath its dignity to acquaint itself properly with the peculiar
essence and self-development of the believing consciousness; it imposes
upon it, without ceremony, the obligation to be a knowing consciousness,
and can then, since faith as faith is absolutely heterogeneous with
knowledge, easily prove that the believing consciousness, grasped as
knowing, is in contradiction with itself. With the utterly false
presupposition that the believing consciousness is in original
contradiction with itself, it is naturally an easy matter for the
criticism of knowledge to caricature the holy facts and get revelation
lumped together with mythology.

We will now see whether it is not after all possible to disentangle this
confusion. It is, as we said, an utterly false presupposition that faith
as faith is thoughtless; we prove this by presenting the principle of
faith as a principle of thought peculiar in itself. That the principle
of faith is itself a principle of thought does not mean that faith
begins with the pure thinking of pure being; but it means that the
believing consciousness is able to think its own original fundamental
act, by this thinking to distinguish itself from the knowing
consciousness, and by this distinction to become transparent to itself.
The point of departure is in the idea of the Good. The idea of the Good
is the idea of the will: to the absolute Good corresponds the absolute
will, the absolute Self. The absolute Self is the holy God. Over against
the absolute Self there then stands, with the world between them, the
relative self, the self that is not good, but obligated to the Good, the
relative will that is not holy, but sinful. With these presuppositions
the believing consciousness can not only become
% --- p. 102 ---
\opage{102}\setcounter{footnote}{0}%
clear to itself, insofar as it can see their absolute indispensability
for the self, but it can even think its essence through so completely
that each of the presuppositions transforms itself into a moment in the
concept of faith.

Once the principle of faith is grasped and determined in the concept of
faith as a principle of thought, there is no danger from the mythical
interpretation and the mythical. To dissolve the divine images of the
natural religions into personified ideas and to demonstrate the
psychological origin of the unconscious illusions is a worthy task for
knowledge and scientific criticism; through the progress of knowledge
self-cognition in this respect then also makes a real advance: the
Aristotelian concept of God is in truth far preferable to Homeric images
of the gods. But spurred on by this advance, ``modern science,''
especially tempted and seduced by theology, has in lamentable fashion
come to overstep its absolute limit. If Zeus---so one concludes---is a
personification of the individuality of the Greek spirit, then Jehovah
must necessarily also be a personification of the strong, self-willed I
of the Jewish spirit. This conclusion is utterly false. He who says
that Jehovah is only a personification of the human essence of selfhood,
what does he do? He makes the human self into a personification of
subjective thinking, and subjective thinking into an empty reflex of
organic natural forces. By such a procedure things by no means go
upward, but on the contrary downward with the self and its
self-cognition. ``It goes backward, Hakon''; yes, it goes backward! For
surely a reflective age must be able to muster so much reflection that
anyone who earnestly wills it can also very easily---without speculative
overexertion---see that the human self must stand and fall with the
divine, the
% --- p. 103 ---
\opage{103}\setcounter{footnote}{0}%
powerless with the almighty, the relative with the absolute.

It shall willingly be conceded that the miracles of revelation in the
Old and the New Testament bear an unmistakable likeness to the myth
miracles of the natural religions, when they are seen and judged from
the standpoint of knowledge. By making the Bible into a stock-dogmatic
document of revelation, theology has provoked scientific criticism to a
demonstration of this so conspicuous likeness, a demonstration often
wounding, no doubt, but in its main lines victorious and irrefutable.
The cognition of knowledge is, and according to its peculiar character
must necessarily be, a rationally coherent cognition of the world: from
the standpoint of cognition of the world a Mosaic miracle is just as
offensive, improbable, absurd and self-contradictory as a Homeric one.
But the concept of the miracle is a concept of faith, an unconditioned
concept of selfhood, and can never become a concept of knowledge; the
contradictions which ``modern criticism'' so easily discovers in the
miracles can all be traced back to an uncritical confusion of concepts
of faith and concepts of knowledge.

The question of the difference between false and true miracles coincides
originally with the question of the difference between false and true
revelation, between false and true faith. With this the principle of
myth is excluded once and for all, not only from the New, but also from
the Old Testament. There may indeed, especially in the holy legends of
the Old Testament, be talk of analogies to the mythical, but, according
to the principle of faith, by no means of myths proper. To develop the
difference between myths and mythical analogies in detail is a special
task for a philosophical-religious treatment of the sacred history; in
the present context it shall only be emphasized in general that the
difference must
% --- p. 104 ---
\opage{104}\setcounter{footnote}{0}%
be determined in principle by a clear and sharp development of the
relation between the human word and the divine primal Word. The word of
the Bible, as it lies before us, has its human as well as its divine
side; it is not a word of God alongside a word of man, but from first to
last a divine-human word, grasped in faith, presented in faith, handed
down in faith. Religious faith is from the beginning a faith of the
people, the religious word a word of the people. Instead of making the
Bible into a document of revelation dictated word for word, we must on
the contrary remember that the Bible is a people's book just as much as
Homer. But if the Bible is a people's book, what wonder then that it
expresses itself in popular fashion in popular images, images which
came forth when the holy light of revelation refracted its rays in the
natural media of the spirit of the people!

With a changed principle everything is changed. If it may be regarded as
the most essential hindrance to the spiritual life of the present age
that science has taken the concept from faith and made revelation
thoughtless: then everyone can now tell himself what the requirement is
if spirit is to transform such a hindrance into a condition.

\begin{center}\rule{3cm}{0.4pt}\end{center}

\chapter*{VII.}
\addcontentsline{toc}{chapter}{Lecture VII}
\markboth{Lecture VII}{Lecture VII}
\label{ch:vii}
% --- p. 105 ---
\opage{105}\setcounter{footnote}{0}%
If anyone comes forward with the gospel of enlightenment that no one can
believe in the holy, since the holy is precisely that which is not, and
no one can believe in that which is not, he will surprise moral society
in the highest degree. But if only he is allowed to speak his mind, his
view of life will show itself from more than one side to be a product of
a way of thinking that prevails in the present age. ``The holy,'' he
says, ``is a concept which dissolves itself before a clarified
reflection. Why? Because the holy is an obscure power that always gives
way before the light of day, a mystical representation, a riddle of the
fantasy, that awakens dread. Just consider the holy places, the altars
of the gods, the sacred groves! If you see them by the light of day, in
human company, you feel nothing of fear or shuddering, that is to say:
you do not come into contact with the holy. But visit these places in
the darkness of night, alone, by the light of the moon, and you will
shudder; this shuddering, awakened by an obscure representation of an
obscure power, is a sensing of the holy. The holy needs a curtain; its
essence is, like the veiled goddess at Sais, a riddle of the fantasy.
But the cold understanding draws the curtain aside; a bold thought lifts
the veil, and the youth overcomes his
% --- p. 106 ---
\opage{106}\setcounter{footnote}{0}%
terror, and the holy vanishes. We, the true children of the
Enlightenment, the freeborn sons of the Revolution, have learned to
respect the civil laws, the political institutions, the conduct of men
and their pursuits; but---forgive us a bold word---we do not believe in
visions and revelations, we do not bow the knee to the holy. The holy
will, by its mystical nature, always pass for something unconditioned;
but in the world of conditions, of means, of interests, everything is
conditioned. That something is holy to me means only that something is,
under certain conditions, precious to me, perhaps even so precious that
I could sacrifice my life for it. But in the sense of a hidden divine, a
mystical unconditioned, I acknowledge nothing holy, no feeling, no
promise, no oath, no altar, no God;---I say what I mean: the holy is
that which is not, and no one can believe in that which is not.''

Upon such a man the whole of society would now pass the severest
judgment; it would cast him out and see in him a dissolute and brazen
person. But, we must ask, why? Is it because he says such a thing, or
because he thinks it, or for both? Chiefly only for the first. Cultured
society breaks the staff over the man who expresses himself so openly and
inconsiderately, for this inconsiderateness is precisely an injury to
society. If, on the other hand, one looks at the inward and asks whether
one can then regard a man with such a way of thinking as respectable,
then society will doubtless divide into two parties. The one party will
keep to ``the bad heart,'' the other to ``the good head''; yes, it would
perhaps not be impossible that between those who went to extremes a
middle party might arise, which would undertake to prove how difficult
it is to unite a good heart with a good head, that the most interesting
people are most often suspected of having a little
% --- p. 107 ---
\opage{107}\setcounter{footnote}{0}%
flaw in the heart, while the amiable ones are sometimes accused of having
it in the head. But with all this the whole of society will easily
agree that the offensive lies in the utterance; while the disputants
will hardly agree that the essential brazenness lies in the thought. And
yet the difference between him who thinks in silence and him who utters
the thought is only this, that the one who speaks is sincere while he is
brazen, whereas the silent one is brazen without being sincere, and so a
hypocrite, and as a hypocrite perhaps far more dangerous to society.
Religious hypocrisy in the ordinary sense of the word is by no means the
fault of the present age; on the contrary, it is an advantage which
growing enlightenment and greater freedom in all spiritual affairs as
well have brought with them, that feigned piety is despised as the sign
of an ignoble, weak and wretched character. But there is, with regard to
religious and ecclesiastical questions, a political accommodation which
in its consequences can hardly clear itself of favoring a certain
conventional hypocrisy. In Catholic countries, where cultural
development has taken on not merely an anti-ecclesiastical but even an
anti-religious stamp, there shows itself an endeavor---not expressly to
settle the conflict of principle between the ecclesiastical and the
political by leading it back to the first fundamental questions,
beginning the religious instruction of the people over again, and
guarding with all its might the freedom of spirit that this requires---but
much more to extend, promote and consolidate the purely political
interests, as if civil freedom and true culture could really go arm in
arm with religious unfreedom and medieval mysticism. Such an
accommodation, which brings the enlightened to feign a certain reverence
for what in their quiet thoughts they must reckon as superstition and
obscurantism, causes not merely endless troubles
% --- p. 108 ---
\opage{108}\setcounter{footnote}{0}%
in the outward, but, what is worse, it hollows out the ground of life in
the inward.

That states which have gone through a political revolution without
having previously been prepared by any real reformation of the medieval
church system must feel the smoldering tension as an oppressive fate is
understandable; perhaps the great pressure, when the hour of spirit
strikes, may yet come to furnish conditions for a spiritual rising of
which we can as yet form no representation. But in Protestant countries,
where the religious principle of freedom has long been
acknowledged---at least pro forma---to want to whitewash over the
devices on the ecclesiastical and political coats of arms, and to muddle
through by accommodation, is a procedure little in keeping with the
spirit of the Reformation. The holy requires for its acknowledgment,
fully as much as the profane, a thorough, firm conviction; but a
reflective age cannot possibly arrive at a thorough, firm conviction
without thorough, clear and consistently carried-through thinking.

But this is, if we are not mistaken, precisely one of the essential
hindrances to the spiritual life in the present age, that religious
thinking among the great majority of the cultured has visibly come to a
standstill. People think, and that with a certain perseverance, about
manifold affairs. They think about politics and industry, about art and
poetry, they think about burning questions, foreign as well as domestic:
they think and think about great and small, only not about the
religious. In the time of the Reformation, on the other hand, yes, then
religious problems were thought about; it was not only the learned who
thought, no, the people itself thought along, for then it was in earnest.
When the people's thinking for itself then gradually ceased, because
religious thinking became a learned matter, the theologians became, if
not exactly thinkers for themselves, yet in the
% --- p. 109 ---
\opage{109}\setcounter{footnote}{0}%
official church as good as officially the sole thinkers; for now
orthodoxy was a science. Thereupon the freethinkers began; but the
freethinkers gave offense to the people; the way of thinking of
freethinking was so unpopular that it contributed, even more than the
theological, its share to awakening the people's suspicion of all
thorough religious thinking for oneself. The consequences of this are
noticeable enough in the present time.

We may by no means say that it is either dullness or indifference that
keeps the cultured from taking hold of the tasks of spirit; no, it is
rather despondency, helplessness and lack of overview. People have a
secret fear that nothing will come of all the religious thinking; that
the whole thing will end at last in dissolution, ferment and boundless
confusion. It is in any case so learned, so critical, speculative and
philosophical, that the layman cannot possibly keep up; no, then it is
better to hold to the established order, follow the stream, keep one's
own thoughts to oneself and believe in accommodation.

But is it then in reality so difficult to acquaint oneself with a
religious task? We believe we can convince every cultured person of the
opposite, and we choose an example which shall not be said to lie too
low, in that we choose a determination of the relation between idea and
God, in order thereby to determine the concept of the holy.

The difference between faith and knowledge can be stated briefly thus:
the holy Self is the object of faith; the idea of the holy is the
content of knowledge. The holy Self is God; in the language of faith it
runs: God is holy, righteous, all-good, merciful, etc. Otherwise in
``modern science''; here the view runs thus:
% Rettelser: read "det Absolute selv" for printed "det absolute Selv"; translated per the correction
God, the Absolute itself, cannot be holy; the holy Self cannot be God,
for the holy Self is a contradiction. We find in Strauss's Dogmatics the
following presentation, highly characteristic of negative knowledge:
% --- p. 110 ---
\opage{110}\setcounter{footnote}{0}%
``The definition \ldots{} of the divine holiness (Lev.\ XIX, 2. Ps.\ VI,
5 ff. Mt.\ V, 48. XIX, 17. Eph.\ IV, 24. 1 Pet.\ I, 15 f., 1 John I,
5 ff.), that it is the attribute on account of which God wills what he
ought, and ought what he wills, will hardly be readily accepted on the
ecclesiastical side: and yet it only expresses what is the presupposition
in the usual determinations of the concept. For whether one defines
God's holiness as the highest, flawless purity in God, which also
demands the due purity of creatures; or as the rectitude of the divine
will, on account of which it wills, in accordance with its eternal law,
everything that is right and good; or finally as the attribute by which
God desires and approves only the good: a moral law is always assumed,
according to which the divine will determines itself, but this relation
between a law and a will is otherwise always expressed by an
ought \ldots{} Very rightly, therefore, Weisse says: God himself could not
be called good if evil were not for him too a metaphysical possibility.
But for every being for which evil is possible, the good stands as an
ought, and thus our definition of the divine holiness given above here
comes into its rights. But, just so conceived, this attribute---unless
that possibility of evil and this ought are empty words---abolishes the
concept of the Absolute \ldots{} he who holds fast to the Absolute
dissolves holiness, which is something only in a being placed in
relations; and he who takes holiness seriously encroaches upon the idea
of absoluteness, since it is obscured by the faintest shadow of a
possibility of being otherwise than it is.'' (D. F. Strauss:
Fremstilling
% --- p. 111 ---
\opage{111}\setcounter{footnote}{0}%
af den christelige Troeslære. Trans.\ by H. Brøchner. Vol.~1.
Copenhagen 1842. Pp.\ 513---515).

From the example cited it can be seen how different scientific thinking
is from religious thinking, and what practical significance it has that
this difference be brought to consciousness. For if it were to be
decided by scientific investigations whether a holy, righteous, almighty
God can exist or not, then it is clear that the majority of the cultured
will not be able to follow such investigations, still less the whole
people. The question of the reality of the holy thus becomes an
intricate question, whose answer laymen must leave to the highly
learned, while they, until the final answer to the question arrives, try
to get by with a faith that at most requires a little talkativeness, but
only very insignificant activity of the brain, namely ``childlike
faith.'' So then, by going to an opposite extreme, one has ``childlike
faith'' in ``the childlike heart,'' a faith which is thoughtless insofar
as it is not able to develop concepts, but must keep to emotions, turns
of phrase and images, a faith which gets its edifying quality from a
confession without self-cognition, an inspiration without development of
spirit, a heartfelt assurance without real conviction. The simpler---that
is to say here, the more stupid---the individual is, the easier it will
be for him on these terms to believe; childlike faith is the coziest
affair in the world: of children and childlike faith our Lord, after
all, demands no thoughts at all.

But in a reflective age it is granted only to ``a little flock'' to
preserve faith in complete thoughtlessness. In the greater majority
thoughts must sooner or later be set in motion; and the stronger the
movement becomes, the more insistent also become questions such as
these: how do matters stand with faith and knowledge?
% --- p. 112 ---
\opage{112}\setcounter{footnote}{0}%
Is it really in science, perhaps in theology, that the thoughts of
faith are to seek their last and highest clarification? Might there not
after all be a true religious thinking, independent of knowledge and
science? The significance of such questions every cultured person can
understand, yes, every tolerably awakened person can understand it. But
everyone who is able to see the significance of the questions is also
able to grasp their answer.

When God in science cannot be holy, because he absolutely must be the
gray Absolute: what does science then do with the concept of holiness?
It elevates it to an idea with the predicate divine. Religion, says
Feuerbach, represents God as holy. In the proposition: God is holy,
\textit{God} is subject and \textit{holy} predicate; such a proposition
expresses only a representation, but no concept. In order to dissolve
the representation into its concept, we must turn the proposition
around, make the predicate the subject, and the subject the predicate.
Instead of the proposition: God is holy, we then get the proposition:
holiness is divine; likewise instead of the proposition: God is
righteous, the proposition: righteousness is divine; instead of the
proposition: God is merciful, the proposition: mercy is divine, and so
forth.

How science will otherwise go about holding fast such freely fluttering
divine ideas as holiness, righteousness, etc., each by itself, or
uniting them in one idea, an idea which then cannot possibly be
absolute, since the Absolute itself can be neither holy nor righteous,
we leave to the knowers; here it shall only be emphasized that the
concept of holiness is originally a concept of faith, just as
independent of science as the concept of God. That there is a holy and
righteous God can be decided solely and alone by self-understanding, by
inner experience, by the testimony of the Spirit in
% --- p. 113 ---
\opage{113}\setcounter{footnote}{0}%
conscience, that is, by faith. Certainly it is an essential and
scientific task to make it more and more clear how the relative ideas
relate to the Absolute itself, and above all, how the ethical relates in
principle to the metaphysical; but these investigations do not concern
religion as religion, nor faith as faith. For the believing
consciousness it stands unshakably firm that God is just as surely holy
and righteous as he is almighty. This presupposition is not a
thoughtless assertion, but a conviction deeply grounded in
self-cognition, a determination of concept inherent in the principle of
faith itself.

To hold fast the thought of faith there is required not a new revision
of the history of philosophy, but an ethical concentration in religious
inwardness: the free self worships the Holy One; the faint self is
beatified in metaphysics.

If the holy Self is not, that is to say: if God is not holy, then the
idea of holiness is a possibility that can never become actual; but a
possibility that can never become actuality is an impossibility. Instead
of saying: holiness is divine, we must draw the scientific consequence,
speak out plainly and say: holiness is a self-contradictory concept,
holiness is impossible. As it goes with holiness, so also with
righteousness: if God cannot be righteous, then righteousness cannot be
divine either; it must nicely resign itself to being human. But what
then does this mean? It means that righteousness can only be actualized
as men, by themselves, singly and in society, at last become righteous.
When does that happen? Never! For it belongs, of course, among the
great discoveries of ``modern science'' too, that no man and no human
society can in any actual state whatsoever actualize righteousness.
% --- p. 114 ---
\opage{114}\setcounter{footnote}{0}%
The idea of righteousness, in infinite ought without can, is in endless
discord with its own actualization. To cover over the glaring
contradiction, people have appealed to this, that the righteousness of
individual men was indeed imperfect, but that each could at least
present his side of the righteous being, in other words, contribute his
piece of righteousness. When the many contributions were then gathered
in social life, we should get, by adding the pieces together, an almost
righteous whole. It is, to be sure, only an approximation! Yes, it is in
ethics an approximation to the ideal of the same kind as it would be an
approximation to the ideal in poetry if, for want of a single true,
great poet, one were to let a crowd of moderately gifted people, each of
whom had his bit of poetry, work together in an aesthetic committee in
order to produce a great poem by many small contributions.

He who cannot see that the assurance, the unshakable, unconditionally
decisive assurance of God's holiness and righteousness, is given in and
with the self-concentration corresponding to faith, has, however much he
may otherwise have thought, not yet thought through the principle of
faith, and does well, then, to think it through before he rejects it.
But if anyone, after he has really thought the principle through and
seen that the concept of a holy and righteous God is quite rightly
contained in the thought of faith as an absolute thought of selfhood,
breaks off with the declaration that the absolute thought of selfhood is
for him without reality, since it demands an exertion of the self, an
energy of self-determination, which he cannot endure: then he must seek
the fault in himself and not misuse knowledge as a cover for his own
weakness.

What a difference there is between the holy God of faith and the idea of
God of science can also be seen from this, that while no knower will
think of either praying to
% --- p. 115 ---
\opage{115}\setcounter{footnote}{0}%
or worshiping the idea of his knowledge, the God of faith, on the other
hand, is an object of worship, a God to whom believers pray. Johannes
Climacus (Uvidenskab.\ Efterskr.\ p.~148) has made the following
reflection: ``If one who lives in the midst of Christianity goes up into
God's house, the house of the true God, with the true representation of
God in knowledge, and now prays, but prays in untruth; and when one
lives in an idolatrous land, but prays with the whole passion of
infinity, although his eye rests upon the image of an idol: where then
is there most truth? The one prays in truth to God, although he worships
an idol; the other prays in untruth to the true God, and therefore in
truth worships an idol.'' Where the task is, as with the author cited,
to get subjectivity and inwardness really emphatically stressed, such
reflections are no doubt of great value; but where, as in the present
discussion, it is a matter of bringing out the relation in principle
between worship and the object of worship, the question has only a
fleeting and passing significance. To pray in truth to God and yet
worship an idol is impossible; for if the worship is true, it can only be
true by the idolatrous representation vanishing in the representation of
a true God; but a true representation of God necessarily excludes the
representation of an idol. Likewise it is impossible to worship the true
God in untruth; for when the worship is in untruth, the true God hides
his face from the worshiper and leaves in his place a phantom, an untrue
god.

That worship and its object naturally correspond to one another, and
that even enlightened men of culture often feel a need for something
they can worship, is confirmed by many experiences: if one has nothing
true to worship, well then, one worships something untrue, and the
worship becomes
% --- p. 116 ---
\opage{116}\setcounter{footnote}{0}%
accordingly; if one cannot acknowledge the true God, because the true
God has come into contradiction with metaphysics, then one can, of
course, acknowledge one's own Absolute and set this up as a metaphysical
idol. Idolatry has not ceased because people have ceased to worship sun,
moon and stars. If an individual has not self-importance enough to
worship himself, he can, of course, have his pleasure in worshiping
others. It is in a reflective age by no means unusual that one man
practices idolatry with another. A mother, for example, practices
idolatry with her son when she sees in the son a true wonder of genius
and virtue, a model man, whom the world in its blindness will never know
how to appreciate. That a lover, at least during a certain period,
practices idolatry with his beloved is not anything unheard of either.
That is why she is called, of course, the object, the adored one, and
happiness is in the adoration: happy the adored one, ten times happier
the adorer! But here too it shows itself how impossible it is that
worship can be true when it is only an idol that is worshiped. There is
in all idolatry something mythical; therefore there is also in all
idolatrous worship something mythical. The mythical is ambiguous, and
lasts only for a time. Then one meets the worshiper half a year later,
but where is the worshiped one? It is over; the whole thing was only a
myth. ``That was strange, though,'' says the guileless one, ``for he
really worshiped her; what inconstancy, what faithlessness!'' The best
thing now would be if the accused hit upon shifting the blame onto
science and unburdened himself thus: ``The concept of worship begins
with faith, but must naturally, like other concepts of faith, clarify
itself in knowledge. I really worshiped her, that is true, for I
believed, yes, I believed; I was, after all, at that time only a beginner
in science, and therefore I believed; but after I have
% --- p. 117 ---
\opage{117}\setcounter{footnote}{0}%
become familiar with science, and especially with metaphysics, I have
learned to see that worship is a concept which cannot subsist with the
Absolute, a concept which halves my consciousness, a contradictory
concept, and therefore I have ceased to worship.''

What here sounds like mockery, because it is said of a worship whose
ambiguity everyone sees through so easily, can in full earnest be
applied to the worship of the mythical gods. They could at a given time
be the object of an obscure, provisional, sensuous worship; but no one
has ever worshiped Zeus in spirit and in truth; for a Zeus cannot be
worshiped in spirit and in truth; it goes against his mythical nature.
What I am to worship must be the pure, holy origin of my being,
personality itself in truth, perfection and blessedness: only the God of
spirit, who is spirit himself, can be worshiped in spirit and in truth.

If one raises the question which age, which generation must be called
the unhappiest, then one knows, of course, from history that one
generation has been chastised in one way, another in another; one has
suffered under the scourge of war, another from plague, famine, etc.
There can be great misfortunes in the world; but as long as the holy is
preserved, and the suffering generation has a God to pray to and to
worship, the measure of misfortune is not full. The unhappiest
generation must surely be the one which by a misdirection comes onto a
path of culture that leads it away from faith, away from the holy, so
that there is nothing in which it can believe, nothing to which it can
look up with expectation and hope, nothing, nothing before which it can
bow with inward reverence. Let such an age otherwise make progress in
all directions and enjoy all the temporal goods of civilization; it is
and remains nonetheless the unhappiest. We have a prototype in Rome
under the
% --- p. 118 ---
\opage{118}\setcounter{footnote}{0}%
emperors. If anyone in the great empire could be called a happy man, it
would surely have to be, in the Roman sense, the emperor himself. What
is it, then, that causes him torments? It is not cares for the welfare
of his subjects, for about the welfare of his subjects he did not
trouble himself at all; nor can it be lack of power, for he was in the
human sense almighty; and yet he suffers, he suffers unbelievable
torments, cruel as a Nero, raving as a Caligula, bestial as a Vitellius,
lustful to the point of madness as a Heliogabalus! It is evident that
he suffers. But whence does it come? It comes from this, that such an
almighty one has nothing he can look up to and regard with reverence; it
comes from this, that the emperor is a god, and has, like a god, learned
to despise everything, even himself.

This is the measure of a misfortune of a higher kind: that the holy
vanishes, that all faith vanishes, that there is nothing more to which
the generation can look up with reverence, nothing on which it can build
its hope, nothing that it can worship in spirit and in truth.

\begin{center}\rule{3cm}{0.4pt}\end{center}

\chapter*{VIII.}
\addcontentsline{toc}{chapter}{Lecture VIII}
\markboth{Lecture VIII}{Lecture VIII}
\label{ch:viii}
% --- p. 119 ---
\opage{119}\setcounter{footnote}{0}%
When the talk is of civilization and enlightenment, the present age ought
not to forget that there is a difference between light and light, between
the enlightenment of culture and the enlightenment of the people; the two
kinds of enlightenment are not two sides of the same matter, but things
heterogeneous to one another. In their origin they are related as
knowledge and faith, as science and religion; for the highest of
civilization is in knowledge, and the deepest of popular enlightenment is
in faith. Since the confusion of the concepts of faith and knowledge, as
we have seen, is one of the most essential hindrances to the spiritual
life in the present age, it will be expedient, in order to bring about
their separation, to bring out the difference between the enlightenment
of culture and that of the people somewhat more closely. How different
true popular enlightenment is from what we call the enlightenment of
culture, we shall best be able to see by sharply separating two concepts
which the present age all too easily confuses, namely the folkish and the
popular.

To clear up the difference between the folkish and the popular it will
be necessary to account for what it is that determines the difference
between folkish and popular writings, between folkish and popular men.

There are writings which are suited to general reading by the people
% --- p. 120 ---
\opage{120}\setcounter{footnote}{0}%
and may with reason be called popular, without their therefore ever
being able to become folkish; conversely, there are writings which,
although the majority of the population is not yet rightly able to read
them or find a taste for them, must nevertheless, by their fundamental
character, be called folkish. Literature has a popular astronomy, a
popular geology, agricultural chemistry, and so on, but such books can
never become folkish. And why not? Because they can never be
communicated to readers who are not of the profession in the form
original and proper to the content; for the original form is and must
be scientific. It is not the scientific investigations, but only their
results, that are communicated. But scientific results lose their
original value when they are torn loose from the investigations from
which they proceeded and are communicated in a form in which they cannot
be grounded. The popular lies in the form; but the popular form is a
lower, inadequate form, deviating from real knowledge, and therefore a
more or less shabby one. Instead of leading the reader to independent
insight, such a communication on the contrary accustoms him to rely on
the authority of those better informed and to take what is communicated
on trust. To be educated by reading nothing but popular writings is by
no means to be educated to think for oneself. Were a people to be raised
to independent thinking by the way of the enlightenment of culture, then
it would necessarily have to be raised to scientific insight. But that
is a senseless and unnatural demand. Or would it not be a senseless and
unnatural demand, if one expected that enlightenment should rise so high
that all in the whole people became men of science? Suppose that culture
brought matters so far that all men became astronomers, geognosts,
botanists, zoologists, and so on---what then would become of the women?
They too, presumably, would have to become scientific. Suppose that we
% --- p. 121 ---
\opage{121}\setcounter{footnote}{0}%
one day got nothing but scientific, zoological, botanical, indeed
mathematical women---how then would it go with womanliness? Science, in
the character of science, can never become folkish; for the scientific
is not for all, but the folkish is and must be for all.

What is the folkish? It is the fundamentally human in its native
strength and particularity. What follows from this for the determination
of popular enlightenment and of folkish writings is then easy to see.
Only that thinking, that enlightenment and development of concepts, whose
first presuppositions and conditions stir with life in the personal,
fundamentally human consciousness, can be called folkish. A proposition
such as this, that the straight line is the shortest way between two
points, is indeed in a certain sense both folkish and popular, and yet
there is here an essential difference between the folkish and the
popular apprehension. The apprehension is folkish insofar as the
proposition is evident to sound sense without artificial abstraction, as
when it is said: the straight way is the shortest. The proposition
becomes popular, on the other hand, when a mathematical abstraction,
even if only a provisional one, is expressly demanded in order to hold
fast line as line, point as point. That here, with the popular, we come
away from the folkish, is noticed at once when we ask about the meaning
of the abstraction employed. For the abstraction points out beyond the
immediate and reminds us that the popular understanding is only
provisional and subordinate; for it is science that has laid down the
abstraction, and it is in science alone that its meaning can be
developed. Only one step more is needed on the way of abstraction, and
the popular proposition has ceased to be popular. It is otherwise with a
folkish proposition! Its development, too, must indeed be conditioned by
abstraction; but the abstraction
% --- p. 122 ---
\opage{122}\setcounter{footnote}{0}%
takes place naturally and as of itself: the development is living. That
great heroes have great passions is a proposition which can not only be
made vivid in legend, history and poetry, and is to that extent called
folkish; but it can at the same time be apprehended by thought,
developed and presented in the very finest psychological shadings, and
still be folkish; when the eye falls on the mirror of the saga, the
fundamentally human consciousness itself brings along the conditions
needed for understanding the proposition: Saxo Grammaticus and Snorre
Sturlesen are writers of the people of the first rank, but they cannot be
called popular.

The development of the folkish is constantly exposed to two
misdirections, inasmuch as one either overlooks what is individual in
the nature of a people and holds fast to the merely universally human,
or, with contempt for the universally human, looks exclusively to the
particularity of native strength. The first error is sufficiently
designated by the rationalist enlightenment of the eighteenth century,
its manifold cosmopolitan and philanthropic abstractions; the
``rights of man'' of the Revolution rest on a cosmopolitan foundation.
The cosmopolitan-philanthropic man of understanding, with his otherwise
respectable zeal for education, schooling, the enlightenment of the
common man, and so on, never gets beyond the merely popular. Warned by
the excesses of an empty rationalist enlightenment, one is in our days,
under the influence of the questions of nationality pressing in on every
side, tempted to an opposite extreme, to clinging to immediate
particularity to such a degree that one without further ado lumps the
accidental together with the essential and forgets the universally
human. Instead of awakening and preparing the people for a freer, more
vigorous self-development, one flatters its natural selfishness and
nourishes its boorish inclination toward self-satisfied reserve, toward
% --- p. 123 ---
\opage{123}\setcounter{footnote}{0}%
indifference and foolish contempt for the light that culture spreads.

If the two sides of the folkish nature, the fundamentally human and the
particularity of native strength, are to work united in spirit and in
truth, then this must happen through a thorough transformation, a
rebirth. The regenerating power is conditioned by the revealed word's
being appropriated in faith: the concept of faith on which the emphasis
must here fall is the concept of the image of God.

That man is created in the image of God is the decisive principle from
which all doctrine of people and folkishness must proceed. Herein lies,
in the first place, a reference to the fundamentally human; only as man
so deepens himself in his own being that he therein cognizes the
universally human can he cognize what bears the image of God. But the
essence of the image of God cannot be held fast in empty generality; it
is an image of life, the archetype of man's existence, to which his true
individuality must correspond. When the image of God is realized, it
comes forth womanly in the woman, manly in the man; the unfolding of its
content yields not empty uniformity but rich diversity; in the image of
God every man's particularity is not merely preserved, but furthered and
heightened. What holds of the single human being in the people, namely
that by developing according to the image of God it develops its
original
individuality, the same holds of a whole people; but when all
the individualities in a people are developed, then surely the
individuality of the people is developed.

Yet here a historical misgiving arises: ``If it is really true that the
evangelical archetype by its principle furnishes presuppositions and
conditions for the development of the native-strong particularity of the
peoples, then such a principle of development would surely have to be
% --- p. 124 ---
\opage{124}\setcounter{footnote}{0}%
expressly demonstrable in facts of the history of peoples; but is that
now the case?'' We answer: the testimony of history is insufficient, but
that is not the fault of the principle; the fault lies with those who
have misused and botched the principle.

The Gospel has, from more than one side, been apprehended and applied in
so unevangelical a way that the people and the life of the people have
had to suffer under it. What, then, has been done? Religion and faith
have been used to keep the common man in spiritual tutelage through
centuries. When Gregory the Great wished to defend the holy images, he
appealed to the ignorance of the people:
% the Latin sentence is set in true italic in the Danish
\textit{Legant in parietibus, quod non legunt in codicibus.} The Catholic
Church, as is well known, has not understood this to mean that the
ignorant people should seek edification only by looking at the images
until they could gradually be taught to read in the Scriptures; no, it
has understood it to mean that the Catholic common folk still in our
days have ``the same need'' to look at the images as they had in
Gregory's days.

There is in faith a characteristic mark which the hierarchy has known
excellently how to use to its own advantage, but to the harm of the
people; this characteristic mark is obedience! Just as faith is
inseparable from worship, so faith and worship are inseparable from
obedience. To believe in God is to worship him, to worship him is to
serve him, to serve him is to obey him. The word of revelation is an
authoritative word; the authoritative word demands obedience. The
temptation to lump God's commandment together with the commandments of
% printer's defect in the Danish: "Mennskebud" for "Menneskebud"; translated in the corrected sense
men is in this case great. The authoritative word of revelation lent
itself excellently to procuring obedience for the rulers of the Church,
the blind obedience that makes it so convenient to keep order in
society. The more unconditional the obedience, the less the obedient
themselves need to
% --- p. 125 ---
\opage{125}\setcounter{footnote}{0}%
think. The first demand on a believing Catholic is blind, unconditional
submission to the authority of the Church; blind faith is one with
blind obedience. Why, pray, is heresy a mortal sin? It is not the single
deviations in the single articles on which weight is laid here; no, the
heretic's crime is that he dares to call in question the infallibility
of the Church, that he does not submit to its commandment in
unconditional obedience. Blindness heeds no light: hierarchy and popular
enlightenment are incompatible.

But now Protestantism? Yes, Protestantism! That the Reformation is
folkish from the ground up and began in a folkish way is certain; but
how far has it come? At the Diet of Worms, when the official of Trier
interrupted Luther and demanded a plain and simple answer whether he
would recant or not, there sounded an answer that has now sounded over
a whole world, and the answer runs: ``Since then Your Imperial Majesty
and Your Electoral and Princely Graces desire a simple, plain and
straightforward answer, I will give one that has neither horns nor
teeth, namely thus: Unless I am convinced by the testimony of Holy
Scripture or by open, clear and evident reasons (for I believe neither
the popes nor the councils alone, since it is plain as day that they
have often erred and contradicted one another), so that I am persuaded
by the very passages which I have cited and used, and my conscience is
made captive in the word of God, I neither can nor will recant anything;
for it is neither safe nor advisable to do anything against one's
conscience. Here I stand, I can do no other, God help me. Amen!'' By
this answer the staff is broken over blind obedience; in this answer the
at once religious and folkish principle of the Reformation is fully
expressed---in a great moment, by a great personality! It is in earnest
with the religious, for it concerns
% --- p. 126 ---
\opage{126}\setcounter{footnote}{0}%
the word of God and a good conscience. It is in earnest with the
folkish; for it is precisely the individuality freed by faith which,
by its manly bearing, here presents to us the human, the fundamentally
human in its native strength and particularity. The personal right that
Luther hereby asserts for himself he asserts at the same time for every
believer in the people, and the obligation that he hereby lays upon
himself he lays by the same token upon everyone in the people who shares
his faith. If faith is to lead to personal independence, then it must,
through the spirit and the word, lead to personal self-development, to
the cognition of the truth that makes free, and can then not possibly
be thoughtless. The truth that makes free is indeed a truth of
conscience, not a truth of science; but as truth it must nevertheless be
cognized, and in order to be cognized it must be capable of being
thought. But if it is a matter of open, clear and evident reasons, then
it is in earnest a matter of popular enlightenment and popular culture;
then the people cannot settle down in a common folk's pious simplicity;
then it is not sufficient to believe in the Pope's images of saints, but
neither is it sufficient to break with the Pope and his Latinists only
to remain sitting in the nursery with the figurative talk of childlike
faith and talking baby talk.

When orthodox Lutheranism distorted Luther's work, it was again
obedience, blind obedience, obedience not in conscience toward God's
word but obedience toward the commandments of men, religious obedience
toward ``the high Christian magistracy,'' through which the distortion
took place. How Luther on the whole judged of the peoples' obligations
to obey their lawful princes is well known, especially from his
pronouncements on the Peasants' War. The brutality with which the
oppressed rose against their tyrants brought it about that the
Reformers had to
% --- p. 127 ---
\opage{127}\setcounter{footnote}{0}%
lay a quite extraordinary weight on civil obedience. This was, so far,
quite in order. But before anyone knew it, the duty of obedience in
civil matters was lumped together with the duty of obedience in
religious matters, and now the distortion and the abuses went beyond
all bounds.

It is, however, not so much government and prince as rather the Church
and its theologians who must bear the blame for the Protestant
subjugation of the people's development. In orthodox theology the
dogmas were determined with an objective formal precision which made
them fit indeed for juridical application. Governed by theology, the
prince in his own person governed theologians and Church. We have a
characteristic example in Christian the Fourth. His religious ordinances
and civil laws resound with Bible passages and the fear of God, and
while he not seldom, in a harsh and arbitrary manner, chastises the
individual teachers, he is himself dependent on dogmatic rulings, even
to such a degree that in his domestic quarrels with Christine Munk he
must use the theological faculty as an oracle.

Theology brought with it yet another hindrance, very essential, to
popular enlightenment. As the development of the theological doctrine
was, so naturally were not only the sermons but also the religious
textbooks prescribed by the government. For these textbooks, arranged
for the instruction of the common people, were extracts from dogmatic
systems, and hence unscientific renderings of something that was to
count as science; in this way religious instruction was necessarily
reduced to being popular, without being able to become
folkish. Luther's Small Catechism is a folkish book; but, to name an
example that lies close to us, Balle's textbook is not a folkish book,
it is only popular.
% --- p. 128 ---
\opage{128}\setcounter{footnote}{0}%
Its popular character strikes the eye when we see how the author, in
the doctrine of God's being and attributes, undertakes by way of a
beginning to prove the existence of God. The proof is frail, it bears
the mark of the superficial philosophy of that time; but even if
philosophy went forward and apprehended the matter more thoroughly, the
popular proof would not on that account become better. If the assurance
of God's existence, in a deeper sense, rests on the theological proof's
being fully tested, then the people are not of age; for they can no more
test a theological than a philosophical proof.

So much for the difference between popular and folkish writings.

We shall now cast a glance at the difference between popular and
folkish men, a difference which naturally must be determined according
to the same principle as the one mentioned above. One so easily mixes
the two concepts popular and folkish, and yet they are so unlike that a
man can be popular without ever becoming folkish, and conversely be
folkish without becoming popular. Just as, so far as the single human
being is concerned, there can be a great difference between his
fundamental will, that which by his deeper nature he wills and must
will, and his momentary will, that for which under given circumstances
he has desire and inclination, so too, so far as a whole people is
concerned, a distinction can be made between its finite inclinations,
guided by interests and considerations of prudence, and its fundamental
will. The man who leads a people to carry out what it would most like,
without asking what it must will, becomes popular, but not folkish; he,
on the other hand, who moves the people to carry out not what it would
most like but what it must will, becomes folkish, but not popular: the
popular man
% --- p. 129 ---
\opage{129}\setcounter{footnote}{0}%
makes the people into a crowd, the folkish man makes the crowd into a
people.

The folkish word is an authoritative word, the folkish man an
authoritative man; but the folkish word has in its authority at the same
time something awakening and liberating: the folkish man will awaken the
people and lead it to freedom. In the folkish, weight must be laid both
on the authoritative and on the liberating. Thus Bishop Absalon was for
his time a folkish man, but not popular; only in recent times has he
begun to win some popularity. King Gustav Erikson Vasa in Sweden was a
folkish, but far from a popular, man. How he understood his life, his
reign and his position toward the people, he has set forth in a
gripping and eloquent manner in the last speech he held to the Estates,
where, with death before his eyes, he takes leave of government and the
world and without reserve renders account of his life and of the use of
his royal authority. He says among other things: ``I know that in the
eyes of many I have been a hard king, but the times shall come when
Sweden's children would tear me up out of the mould, if it stood in
their power.'' Those were folkish words, but they were also
authoritative words. The conditions of the time were hard, and hard had
the king to be who with Gustav's means was to reach his goal. But the
goal was,---a goal of liberation, and in the liberation lay the folkish.

Let us go back and cast a glance at antiquity. Already in Homer the
difference between the folkish man and the demagogue is to be found, not
indeed developed in a concept, yet made visible through the contrast
between Odysseus and Thersites: Odysseus is folkish and cares but little
about being popular; but what Thersites aims at is plainly the popular.
When Agamemnon (Iliad. II.), in order to test the mood of the people, has
% --- p. 130 ---
\opage{130}\setcounter{footnote}{0}%
opened to the army the prospect of coming home, the hosts stream down to
the shore: what a scene, what a sight! ``As when a gale blowing from the
west, rushing forth violently, sets the fertile field in waves, so that
the corn bows with its ears, so the people were stirred, with ringing
cries they rushed all together toward the ships.'' The powerful movement
makes vivid the army's enormous longing to return to the fatherland; no
wonder, when one considers what a temptation lay in Agamemnon's words:
``Nine of the high son of Kronos' circling years have passed, the timbers
in the ships are already rotten, and the ropes are decaying: meanwhile
yonder in our home sit our wives and little children, waiting long for
us, who here wholly for nothing waste our strength on the labor for
which we came hither.'' Agamemnon plays a high game; he is silent about
the folkish and tempts with the popular. Odysseus sees with indignation
the clamoring people turn into a crowd; his task is to stop the movement
and turn the crowd back into a people; this can happen only by divine
assistance. Just as Pallas Athene came with a revelation to Achilles
where it was a matter of bringing him personally to his senses, so she
comes here too, where it is a matter of bringing the whole people to
its senses, with a revelation to Odysseus. Now follows a masterly
portrayal of how Odysseus knows how to bend his voice, as with mildness
he admonishes the princes and with severity stops the common men. So he
succeeds, by his authoritative bearing, in bringing the masses to rest
and getting the people to listen. ``All now remained seated and kept
themselves back in their places; only the babbler Thersites let his
chattering voice be heard.'' Thersites is the demagogue in a first crude
sketch. He begins immediately as the single individual in the crowd who
makes a row; but as he gets the floor, he involuntarily speaks in
% --- p. 131 ---
\opage{131}\setcounter{footnote}{0}%
the name of all, and at once presents himself as a leader of the crowd.
The hollowness of his nature makes itself known in quite a
characteristic way in that the speech he makes is an echo of Achilles'
reproaches against Agamemnon: ``While the others fight, toil and suffer,
Agamemnon has plenty of gold and fair maidens.'' The speech rises to
agitating pathos: ``Cowardly, pampered folk, Achaean women! Achaeans no
longer!'' recalls the insult done to Achilles, and ends in parody:
``Achilles is meek, no gall dwells in his body, else, son of Atreus!
this insult of yours would surely have been the last.''

The demagogue's spiritual ugliness is symbolically designated by his
bodily ugliness. ``Thersites was,'' says Homer, ``the ugliest man in the
army that marched against Troy; crooked were both his legs, and he
limped badly on the one, his back was humped, and his shoulders bent
forward toward his chest, pointed was his head at the top, and thin sat
the hair about his crown.'' That the people, with the mythical gods, at
the mythical stage of consciousness, was still uncorrupted is made known
by the fact that the crowd, when its eyes were opened, could see
Thersites' ugliness just as well as the princes. Before the eyes of the
crowd and with the crowd's applause, Odysseus could not only rebuke the
babbler with hard words, but even inflict on him a corporal
chastisement: ``With his staff he struck him over back and shoulder;
Thersites writhed and cringed, and the tears burst from his eyes; under
the golden staff a weal rose on his back, bloody and blue; then he took
his place with trembling and moaning, and with a foolish look he wiped
the tears from his eye. Heartily one had to laugh'' \ldots{} It was a
great victory that the folkish won by this chastisement over the merely
popular.

But let us not forget that it was precisely by the assistance of the gods
that the victory was won. As, little by little,
% --- p. 132 ---
\opage{132}\setcounter{footnote}{0}%
political development displaced the originally religious, the
relations changed. When the popular party in Athens regained its former
influence through the expulsion of the Pisistratids, ostracism was
introduced. By ostracism the crowd was entitled, without proof, even
without an accusation of any crime, to condemn the fatherland's most
deserving men to exile; through ostracism it became impossible for the
most patriotic hero to work in a folkish way without being popular,
whereas it became possible, as an Aristophanes can testify, for a
political ``sausage-seller'' to work popularly without being folkish.

The deeper the thought of life for which one works, the more it
penetrates the fundamentally human, the more strongly, too, the contrast
between the folkish and the popular will be able to come forth. We see
it at the stage of revelation; we see from the story of the golden calf
(1 Moseb. XXXII) what a difference there is in this respect between
% "1 Moseb." as printed; the golden calf is 2 Mos. XXXII --- author's slip, not corrected
Aaron and Moses. Moses delays in coming down from the mountain; the
people gather about Aaron and say to him: ``Up, make us gods, which shall
go before us; for as for this Moses, the man that brought us up out of
the land of Egypt, we know not what is become of him.'' So speak the
people, tempted by flesh and blood; they will fall away from Jehovah,
the God of the spirit; they will fall away from themselves; they are no
longer a people, but a sensual, raw rabble, a crowd. What, now, was
Aaron's duty? As a folkish man to stand against the crowd and see to
getting the people to listen! But instead the priest Aaron bows beneath
the will of the crowd and casts the golden calf. When Moses then comes
down from the mountain and is to hold judgment, there unfolds a highly
characteristic exchange of words between him and Aaron. ``What did this
people unto thee, that thou hast brought so great a sin upon them?''
Aaron: ``Let not
% --- p. 133 ---
\opage{133}\setcounter{footnote}{0}%
the anger of my lord wax hot: thou knowest the people, that they are
set on mischief''; in other words: what was I to do? I could not, after
all, contradict the crowd and make myself unpopular. Aaron's sin against
the people is that he has treated them not as a people but as a crowd.
That he himself feels it we hear from his excusing explanation: ``And
they said unto me, Make us a god, which shall go before us \ldots{} And I
said unto them, Whosoever hath any gold, let them break it off; so they
gave it me: then I cast it into the fire, and there came out this
calf.'' The whole thing happens without Aaron rightly knowing how; the
calf comes out of the fire as if by chance: one would almost think it
had come about by witchcraft. But Moses understands it otherwise. Alone
and at peril of his life he takes his stand at the gate of the camp and
holds judgment: ``Who is on Jehovah's side, let him come unto me!'' so
sounded his authoritative voice; yes, the folkish man is an
authoritative man.

The folkish man is in his way a holder of power just as much as the
prince on his throne and the popular man who rules by popular favor; but
the difference in power must be determined by the difference in
authority. Christ says in the Gospel (Matth. XX) to his disciples: ``Ye
know that the princes of the nations rule over them, and they that are
great exercise authority upon them. But it shall not be so among you:
but whosoever will be great among you, let him be your servant; and
whosoever will be chief among you, let him be subject to you. Even as the
Son of man came not to be ministered unto, but to minister, and to give
his life a ransom for many.'' In these words all power and authority on
earth is essentially characterized. The popular man can at a given time
attain power, but his power is without authority; outwardly seen, he
rules indeed over the crowd, but at bottom it is the crowd
% --- p. 134 ---
\opage{134}\setcounter{footnote}{0}%
that rules over him. As head of the state the prince has indeed both
power and authority; but the power is physical; the authority grounded
on physical power is political, but not in the stricter sense personal.
Personal authority stands in direct proportion to personal superiority,
and personal superiority in turn in direct proportion to self-sacrifice.
But from this it surely follows that the highest folkish authority must
in its innermost ground necessarily be religious. The Savior of the
peoples is the Lord of the peoples, a thorn-crowned, but nonetheless
divine majesty, who alone can teach them obedience. For as the power
and the authority are, so must the obedience be also: toward physical
power obedience is compelled submission, and the highest mitigation is
this, that the submission becomes popular; toward the liberating
spiritual power obedience is free submission; a free and yet
unconditional submission is in the deepest sense folkish, but it is not
popular.

That the difference in principle between the popular and the folkish be
thoroughly cognized is an important condition for the spiritual life of
the present age.

\begin{center}\rule{3cm}{0.4pt}\end{center}

\chapter*{IX.}
\addcontentsline{toc}{chapter}{Lecture IX}
\markboth{Lecture IX}{Lecture IX}
\label{ch:ix}
% --- p. 135 ---
\opage{135}\setcounter{footnote}{0}%
The contrast between Judaism and Christianity, between revelation at the
standpoint of the Old and of the New Testament, has been, and rightly,
determined thus: that Judaism, although a revealed religion, is
nevertheless still only the religion of a single people, and hence
burdened with the limitations which the natural bounds of a people's
life necessarily bring with them, whereas Christianity, by raising itself
above these limitations, has become not one religion among several
others, but religion itself in spirit and in truth, not religion for a
single people but for all peoples, the religion of mankind. With this,
then, agrees also what Paul says, that in Christ there is no difference
between Jew and Greek, that what was hostile is reconciled, that Christ
``hath made both one, and hath broken down the middle wall of partition,
having abolished in his flesh the enmity, even the law of commandments
contained in ordinances; for to make in himself of twain one new man, so
making peace'' (Ephes. II). Faith in Christ thus abolishes the natural
difference both between the nations among themselves and between
individuals: ``There is neither Jew nor Greek, there is neither bond nor
free, there is neither male nor female: for ye are all one in Christ
Jesus'' (Gal. III). That Paul here points to
% --- p. 136 ---
\opage{136}\setcounter{footnote}{0}%
that unity of faith and of life in which all human feelings and thoughts
must meet is clear; but how is the unity to be apprehended? An empty
unity it cannot possibly be, as surely as it really is a unity of life;
for the unity of life is infinitely rich in quickening power, power that
unfolds multiplicity and diversity.

What kind of unity is it, then, of which Paul speaks? It is the unity
which comes forth when the love of the spirit gains mastery over the
selfish. Where egoism reigns, it will in all relations misuse the
differences in gifts and powers, station and circumstances, which the
natural order of things brings with it. Thus the stronger people,
favored by fortune, will not only subjugate the weaker, but, when the
deed of violence succeeds, even flatter itself that it is Providence's
chosen people, a people that knows its mission. If the egoism of a
people now, into the bargain, gets a point of support in the religious,
then there is no limit to hatred and enmity; it is well known how the
Jews in Christ's time, in the conceit that with ``the law of
commandments'' they were, after all, ``the chosen people,'' hated Greeks
and Romans, and were hated by them in turn. This mutual hatred and
enmity of the peoples disappears when the pure self, made mighty by
self-sacrifice, displaces egoism. Where an unbridled egoism reigns in
society, the natural difference between man and woman and the civil
difference between master and servant will naturally bring it about that
the man makes the woman a female slave, and the master makes his servant
a thrall. And the woman and the thrall will both, without nobler
self-respect, apprehend their degradation egoistically, bow in hatred
before the right of the stronger, and, hidden behind the middle wall of
partition, lie in wait for revenge. As a barrier against the enmity and
the lurking revenge ``the law of commandments'' is then set up; but the
law can be broken: liberation is only in faith in him who,
% --- p. 137 ---
\opage{137}\setcounter{footnote}{0}%
by breaking the power of egoism, broke down ``the middle wall of
partition.''

It is the egoistic apprehension and selfish misuse of human
differences whose abolition Paul speaks of in the strong words that in
Christ's kingdom there is neither Jew nor Greek, neither bond nor free,
neither male nor female, but that they are all one. But that
Christianity should relate itself indifferently to the differences
grounded in natural endowments and natural relations, whether between
individuals among themselves or between nations, cannot possibly be
Paul's meaning. If one understands his words to mean that the Gospel
either wholly overlooks or even aims at obliterating national
differences, then one must surely also assume that the Gospel either
wholly overlooks or even aims at obliterating the essential difference,
not only bodily but also of soul, between man and woman. Those who
understand Christianity to mean that, as the religion of mankind, it
looks only to the abstractly human and hovers cosmopolitanly above all
folkish particularities, misunderstand it just as surely as those who
make its effects dependent on the natures of peoples to such a degree
that that middle wall of partition which Christ broke down is walled up
anew.

The spirit of the Gospel will cleanse, rejuvenate and strengthen, in a
word: \textit{regenerate} the fundamentally human. What is the
fundamentally human? It is the personal! In the acknowledgment of
personality the religious in a way coincides with the political. Only as
a person does a human being gain the rights of a citizen in an earthly
society; only as a person does it gain the rights of a citizen in
heaven. The question how the principle of the Reformation in the
articles of faith of 1530 and the principle of the Revolution in the
principles of
% --- p. 138 ---
\opage{138}\setcounter{footnote}{0}%
1789 are to be reconciled coincides with the question of the development
of the principle of personality.

The difference between the civic and the evangelical apprehension of
the essence of personality is now this, that the former acknowledges
personality in its immediate, natural existence and demands merely a
limitation and ennobling of its egoism; whereas the Gospel, on the other
hand, will grasp personality in its innermost spiritual root, and
heighten its independence by purging away its egoism. If there can now
be no doubt which of the two apprehensions is the superordinate and
which the subordinate, then neither can there be any doubt that the
civic must be subordinated to the evangelical, and the original emphasis
fall on the evangelical apprehension.

The evangelical apprehension can be reduced to the following fundamental
determinations: human personality has the divine 1) as its author, 2) as
its exemplar, 3) as its essential bearer and mover. To develop these
fundamental determinations is to develop the apostolic articles of
faith.

The development of thought in the pure, general form must, however, not
be confused with thinking in the actual life of faith: it is one thing
to think faith, another to live in faith. It is not our task in the
present connection to make a confession of faith; but it is expressly
our task to think faith, in order to think personality.

What a difference there is between the way of thinking of faith and that
of knowledge becomes quite vivid to us when it is a matter of the
development of the principle of personality. Already the standpoint,
with the purely human side of the first presupposition, is determined
in quite another way in faith than in science. The demand is that man,
in order to know himself as a person and to assert his personality, must
% --- p. 139 ---
\opage{139}\setcounter{footnote}{0}%
lay weight on himself. How much weight, then, may a human being in the
name of personality lay on his own self? This cannot possibly be settled
by subjective arbitrariness; for in that case the most puffed-up,
presumptuous and conceited individuals would surely have to be the
worthiest representatives of personality. No, the demand of true
personality must be determined by the very principle of selfhood. But
when the question runs: what independence can originally be ascribed to
the human self? then it turns out that knowledge and faith must go
each its own way.

If the matter is to be settled in knowledge, the independence in
principle of the self will prove to be dependent on the organism, on
organic-physiological natural conditions. But if we consult physiology
thoroughly, the self sinks so deep into the circle of conditions that it
can hardly entertain hope of surviving death. Perhaps with the help of
philosophy it can raise itself again in psychology? Yes, indeed it can
raise itself, but only to unity of essence with the idea, an ambiguous
unity of the moment, which now gets the name of a present eternity, now,
more long-drawn-out and abstract, that of a metaphysical immortality.

It is otherwise in faith. The first distinguishing mark characteristic
of the cognition of faith is just this, that the independence in
principle of the self is not made dependent on organic conditions, but
solely and alone on the self's own essence, an essence which comes to
light in the self's concentration into absolute self-determination. How
this self-determination, humanly seen, must look can easily be shown.
The believing self is recognizable by its peculiar demand of life; what
does it demand? An eternal life in spirit and in truth! With this its
unswerving demand, at this highest peak of self-affirmation, the
believing self is easy to tell from a knowing subject.
% --- p. 140 ---
\opage{140}\setcounter{footnote}{0}%
The unswerving demand is no presumption; it is on the contrary an
expression of the soul's innermost longing, its very deepest need. The
determinations that do not concern the demand itself do not concern
faith; those determinations, on the other hand, which are necessary for
the satisfaction of the demand belong originally to the concept of faith
and are decisive for the thinking of faith. Since, then, there can be
talk here of concepts of faith just as well as of concepts of knowledge,
the principle of contradiction will play a role in the development of
thought in faith just as well as in knowledge; only that the
contradictions, in accordance with the heterogeneous principles, must be
heterogeneous. All the contradictions of science are, in their
fundamental character, objective contradictions of concepts; it is
either one concept that contradicts another, or a concept that
contradicts itself. Thus, for example, a round square is an objective,
essential contradiction of concepts. It is otherwise with the
contradictions of faith; they are all recognized, according to their
fundamental character, by this, that through them the self, mediately
or immediately, contradicts itself: the contradictions of faith are in
the deepest sense personal self-contradictions. As a contradiction is,
so must its resolution be also: knowledge can no more resolve the
contradictions of faith than faith can resolve those of knowledge.
Herewith the presupposition necessary for the development of the
principle of personality is determined; we shall now indicate in general
outline the course of the development.

% section head, set in italic (numeral included)
\textit{1) Human personality has the divine as its author.}

How much faith and knowledge must disagree on this point is already
clear from what we have brought out in an earlier connection. When
science cannot acknowledge the personality of God because the concept
of an absolute personality is a self-contradictory concept: how then
could it grant that the
% --- p. 141 ---
\opage{141}\setcounter{footnote}{0}%
personal God can be the author of human personality? Now it is quite
certain, to be sure, that science changes its concept of God every
moment: every original philosophy, the atheistic not excepted, has its
own original concept of God. In Spinoza God is the one infinite
substance, in Leibnitz the monad of monads, in Fichte the moral world
order, in Schelling---at least for a time---the Absolute-Absolute, in
% "Absolut - Absolute" printed with spaced hyphen in the Danish
Hegel ``the concept,'' in Feuerbach zero; he who would rightly cognize
the philosophical God of the times must know the metaphysical rate of
exchange, just as he who would appraise the year's harvest must know
the official grain-price schedule.

To leave philosophy and flee to theology is only to make bad worse; it
is to stir the contradictions of faith and of knowledge together into
one and muddy the whole. The only elements that rise with any
distinctness out of such a stirring together are a multitude of big
words that all end in \textit{ism}: atheism, acosmism, pantheism,
polytheism, monotheism, deism, theism, and the like. But the concept of
personality, however many \textit{isms} one may drive it to, does not
become clearer.

We keep, then, to faith. Every determination in a concept of faith is,
in accordance with the standpoint, to be stated briefly and sharply,
without any long-winded circumlocutions: the corrective lies in the
self-contradiction and its resolution. Whence does the believer know
that God is personal? He does not \textit{know} it, you see, but he
\textit{believes} it! Then his assurance in faith is presumably weaker
than that which he could have obtained in knowledge? By no means! it is
on the contrary far stronger. To deny the personality of God, or even
merely to call it in question, would be for the believer just as
revolting a self-contradiction as it is for the knower a revolting
contradiction of concepts that the round should be square. If the
personal demand of the self has validity, then it cannot,
% --- p. 142 ---
\opage{142}\setcounter{footnote}{0}%
despite its relative dependence on organism and organic conditions,
possibly derive its existence from a selfless nature. With the life of
nature as its foundation it must, in order to live an eternal life,
necessarily have a germ of selfhood from the eternal, absolute self,
i.e.\ from the personal God.

The more living the feeling of self is, the more painfully is felt the
cutting contrast between man's personality and the impersonal being of
nature: a powerless freedom stands here over against an overpowering
necessity. If freedom is in truth to be cognized as the principle of
all existence, then the kingdom of nature, the whole natural world, must
be absolutely dependent on God, a work of his free, almighty will: to
believe in the eternal author of the self is to believe in ``the almighty
Creator of heaven and earth.''

In the question of nature and creation the difference between the
essence of faith and that of knowledge becomes quite conspicuous. For
the principle of knowledge nature as nature is the first, and the
concept of creation must in consequence conform itself to the concept of
nature. Now, if science, by laying its concept of nature at the
foundation, could develop, clarify and make distinct the religious
concept of creation: then it would by the same token be able to develop
all the other fundamental religious concepts; theology would then be
transformed into philosophy, and the scientific reconciliation of faith
and knowledge would be consummated by faith's being dissolved into
knowledge. For the principle of faith, on the other hand, creation as
creation is the first, and the concept of nature must in consequence
conform itself to the concept of creation. Now, if the doctrine of faith,
by laying its concept of creation at the foundation, could develop,
clarify and make distinct the scientific concept of nature: then it
would by the same token be able to develop all the other fundamental
concepts of nature; the scientific reconciliation of faith and knowledge
would then be to be sought and found in ``Christian dogmatics''; and
what then?
% --- p. 143 ---
\opage{143}\setcounter{footnote}{0}%
Why, then we should have---the theological universal monarchy! But the
latter is, according to the conditions of human cognition, just as
impossible as the former (cf. ``Den gode Villie som Magt i
Videnskaben,'' p.~30 ff.).

Since, now, neither is possible, the reconciliation can be brought about
only on the condition that human knowledge cognize its absolute limit,
and that faith, as a principle independent in itself, gain the right to
develop its own concepts.

``But''---it is objected---``for a moment it may well look as if such a
doubling were something; but if we look more closely, it is nevertheless
not anything. Science never grants that outside its own preserves there
can be a coherent, consistent thinking; in particular the clear heads,
the true, sharp-thinking metaphysicians, the holders of pure philosophy,
will soon prove to us that human knowledge cannot possibly have any
absolute limit, and hence cannot grant any riddle of existence either,
and hence cannot acknowledge faith as any principle of thought
independent in itself.'' We answer that if there really arises a
metaphysician who, to the world's great astonishment, proves and makes
it manifest to all who can think that human knowledge has no absolute
limit: then we too shall not fail, as in duty bound, to be astonished.
But until this happens, we use the interval of expectation to impress
upon ourselves the following fundamental demands, which the knowledge
that declares the absolute limit abolished must be able to satisfy
without circumlocution.

Such a knower must show how he goes about: a) thinking, not merely
\textit{successively}, first the one thought and then the other, but
simultaneously, i.e.\ countless distinct thoughts at once, for thus
thinks
% --- p. 144 ---
\opage{144}\setcounter{footnote}{0}%
a godhead; b) seeing through matter and the material so thoroughly that,
by his intellectual power---for to comprehend is, after all, to
master---he becomes able to transform concept into existence, and
existence again into concept, the concept of the horse into the
existence of the horse, the concept of gold into actual gold; for to
overlook the horse and ride on the concept, to keep the gold itself and
pay with the concept, are mere evasions, and such evasions are unworthy
of a metaphysical godhead; c) seeing through the unity of brain activity
and the corresponding thinking, of the moving will and the nerve moved
by the will, so clearly, and making the interaction between soul and
body so evident, so distinct and comprehensible, that even the
physiologists give up their intricate and laborious empiricism in order
to throw themselves upon metaphysics; and so on.

What we demand here the pure metaphysician will perhaps regard as
trifles and petty quibbles; for once one has, with the help of
abstractions, precipitated existence as a sediment of the concept, then
the determinations of power belonging to existence naturally fall away
of themselves, and when the determinations of power fall away, a
powerless knowledge can, of course, easily play godhead. Or might
perhaps---for as regards pure metaphysics one must be prepared for the
incredible---powerless knowledge one day, by ``fine observations,'' be
able to carry things so far that it even deduced, I say and write:
deduced, not the universal concept of power---for the concept of power
in abstracto has been deduced often enough---but power itself, actual
power, I say and write: actual power with the actual determinations of
power? Yes, if that ever happens, then a metaphysical miracle happens.
Let one judge of such a miracle as one will; but every reasonable human
being must surely grant us that the demands set up have their
significance, and that it would nevertheless
% --- p. 145 ---
\opage{145}\setcounter{footnote}{0}%
be rather interesting if, with the help of pure knowledge, one could
some day find out how things go with metaphysical miracles in a
speculative brain.

Yet even if the astonishing thing should happen, namely that pure,
powerless knowledge gets power over actuality itself and deduces, not
the ideal necessity of power---such a necessity has been deduced often
enough---but power's \textit{immediate manifestation} in actual
existence: what then? Then there is still one thing that can be said in
this connection with decisive certainty, and it is this: that such a
metaphysical deduction the people---and it is precisely of the
enlightenment of the people that we are speaking here---will never get
into its head; and that for the good natural reason that a popular head
is not metaphysical. In other words: if the religious concept of
creation is to be dissolved, filtered off, deduced, paralyzed and turned
topsy-turvy into a concept of knowledge by a powerful word of impotence,
so that it is no longer permitted to count as what it originally is and
wants to be---a concept of faith---then the popular light goes out; then
the unscientific must learn to believe, not in the revealed Word, but in
the deductions of metaphysics. They do not, of course, get knowledge
itself, but only a semblance of knowledge in the popular.

To exchange faith in revelation for faith in metaphysics is a mistake
that will hinder the spiritual life not only in the present age but
surely also in the future. If the powers of the spirit are to be
awakened, and the fundamentally human is to stir in all its originality,
then the condition, the indispensable condition, is that the principle
of faith asserts its right, that the fundamental concepts of revelation,
the concept of freedom, the concept of creation and first and last the
\textit{concept of personality}, develop in a form
% --- p. 146 ---
\opage{146}\setcounter{footnote}{0}%
that corresponds at once to religious and to popular thinking. ``Ye
shall all,'' says Christ, ``be theodidacts,'' that is, taught of God.
These words may well be said, in the present connection, to contain a
serious admonition. For just as impossible as it is to get the life of
the spirit in a people developed and clarified by popular after-glimmers
of the torches of the sciences, just so necessary is it to get popular
instruction grounded on personal insight into the essence of
personality. The thought of personality is at once a popular and a
religious fundamental thought. The richer and deeper this fundamental
thought develops in a people's consciousness, the richer and deeper also
will be the influences of the divine Spirit upon the spirit of the
people; the richer and deeper each national spirit, according to its
measure, appropriates the divine instruction, the more will hatred
vanish, the middle wall of partition be broken through, and the
fundamentally human thought of unity triumph through Him who came ``to
make peace.''

\textit{2) Human personality has the divine for its pattern.}

The enlightened of culture will perhaps, out of particular regard for
the great majority of the population, grant us that a religious
development of the principle of personality might well have its
significance when it is a matter of awakening the consciousness of the
fundamentally human; but, they add, with the express restriction that
one keeps to the first article, which expresses the universally
religious, and takes care not to touch the second. ``Instead of
presenting the universal, the second article presents precisely the
altogether particular, indeed something so particular that not even the
Christians among themselves have been able to agree on its
understanding.'' That this article, instead of being a general article
of agreement, has on the contrary become an article of strife, of
discord and of separation, has, in the opinion of these
% --- p. 147 ---
\opage{147}\setcounter{footnote}{0}%
cultured people, its natural ground in this, that a single human being
existing in flesh and blood has been set up as the normal man, indeed
as the God-man has been made an object of faith and worship.

The objection rests on a great misunderstanding. Through the second
article there runs the same, the very same personal thought of unity as
through the first. On closer consideration it will appear that the
first article first acquires its full significance by being seen in the
light of the second.

If in the first article we saw the opposition between man's limited
personality and the power of nature, we must now in the second article
reflect upon a new opposition, the spiritual opposition between egoism
and the holy personality. This opposition points expressly to a
contradiction in man himself. It is the first religious testimony of
conscience that man enters as a debtor into the kingdom of the spirit:
the first consciousness of spirit as spirit goes together with the
consciousness of guilt. The self-contradiction is that the actual
personality stands in conflict with its ideal, that the fundamentally
human is distorted, that the image of God is lost. How is this
self-contradiction to be removed? Philosophy must, insofar as it does
not want to deny every glimmer of spirit, agree with religion that the
actual, existing man is at variance with his own ideal. And yet it is a
demand of the spirit that man shall be judged by his ideal. The demand
is just, but it is severe: there is something horrifying in seeing an
Adam Homo on the defendant's bench and hearing his counsel say that he
``appeals to the ideal.''

The agreement between religion and philosophy ceases, however, the
moment the question is asked about the realization of the ideal. The
dispute on this point is well suited to
% --- p. 148 ---
\opage{148}\setcounter{footnote}{0}%
illuminate the different way in which the principle of contradiction is
applied by faith and by knowledge. Insofar as the opposition between
ideal and actuality is apprehended from the standpoint of science, the
emphasis is laid on the concept of actuality, and the question then
becomes whether the ideal can be realized. If it now turns out that the
existing man cannot realize his ideal, what then? Why, then its
realization, strictly speaking, cannot be demanded of man either. Man in
flesh and blood need not, then, take the demand too much to heart; if he
really cannot fulfill it, he is of course lawfully excused. By coming
into contradiction with the concept of actuality, the concept of the
ideal comes into express contradiction with itself; for an ideal that
demands realization without being able to be realized is a
contradiction. It is otherwise in the religious. When the opposition
between man's personal ideal and the man existing in flesh and blood is
apprehended from the standpoint of religion, the whole decisive emphasis
is laid on the ideal demand: Thou \textit{shalt} realize thy ideal; here
no excuse holds, \textit{Thou shalt!} Instead of placing the
contradiction in the ideal, religion, in order to vindicate the right of
personality, places the contradiction in the
human self and appeals to conscience. If man now wants to withdraw from
the obligation, he feels a painful self-contradiction; for the ideal
shall---it is divine earnest---the ideal shall be realized. If man
acknowledges the obligation, he feels the deepest pain of
self-contradiction; for in this actuality, with ``these members,'' in
this body condemned to death, he cannot possibly fulfill the
obligation. It is a pain as though the personality must burst: what,
then, can man do? In his pain he can understand the Apostle's lament:
% --- p. 149 ---
\opage{149}\setcounter{footnote}{0}%
``O wretched man that I am! who shall deliver me from the body of this
death?''

When the realized ideal, the pattern existing in flesh and blood, is
regarded in the light of knowledge, criticism says consistently: such
an ideal man is impossible, consequently he is unactual. When the
evangelical pattern is regarded in the light of faith, criticism says
consistently: such an ideal man is necessary, consequently he is
actual.

That knowledge bows beneath the power of finite actuality as beneath a
natural power independent of the ideal demand of personality is, after
all, quite in order, since, according to its principle, it has had to
set the concept of nature above the concept of creation: science cannot
acknowledge the second article because it has not been able to
acknowledge the first. With faith the case is the reverse. Since the
% Danish prints a full point after "Fordring" (sense wants none or a
% comma); translated as continuous sense.
thought of faith bows before the holy power of the ideal, it need not
% Printed "hehøver" for "behøver" (need); translated in the corrected sense.
give way before sensuous conditions of nature. By setting the concept of
creation above the concept of nature it has made the sensuous conditions
of nature dependent on the personal will of the Almighty: by
acknowledging the first article it has opened for itself the way to
acknowledge the second.

If we go through the second article member by member, we convince
ourselves that every member corresponds precisely to the ideal demand
of the principle of personality. If the person of the Redeemer is to be
an object of knowledge, then the whole article is a fantastic
composition of mutually contradictory elements: ``Conceived by the
Spirit and born of the Virgin,'' and yet ``suffered under Pontius
Pilate''; ``crucified, dead and buried,'' and yet ``risen from the dead
and ascended into heaven'' \ldots{} what absurdities! If, on the other
hand, the person of the Redeemer is apprehended as an object of
% --- p. 150 ---
\opage{150}\setcounter{footnote}{0}%
faith, then the one member follows from the other with inner
consistency.

As the one who in an actual human life is to fulfill the holy law of
personality, the pattern must have a will perfectly pure and strong,
moved by the Spirit of holiness; but in order that an actual human being
may have a pure and strong and holy will, he must be conceived by the
Spirit and born in purity. ``That is absurd, impossible, in
contradiction with all actuality!'' cries the knower. But faith does not
fear the contradictions of knowledge; it does not heed, and according to
its principle cannot heed, the contradictions that spring from regard
for material actuality; in its concept of creation it has a sure
defense against all such contradictions. On the other hand, faith does
fear one contradiction: that is personal self-contradiction, for in it
is death. But now it would surely be a contradiction in the essence of
personality if its ideal demands could not be realized, and an equally
glaring contradiction that a human being could realize these demands
without a pure and holy will, and then finally also a contradiction
that a human being should be able to have a pure and holy will although
he was, as David says in the pain of his contrition, ``shapen in
iniquity and conceived in sin.''

But now the actual life of the ideal in a sensuous, actual world? Well,
is it not remarkable that the richest course of life can be grasped and
presented in so few and simple words? ``Suffered under Pontius Pilate,
crucified, dead and buried''; in these words the actual life of the
ideal is summed up; for the words make his suffering visible, and the
life of the ideal on earth was suffering and passion from first to
last. The conflict between ideal and actuality is thus removed in
suffering: the self-contradiction vanished when suffering was put in its
place.

% --- p. 151 ---
\opage{151}\setcounter{footnote}{0}%
As regards the descent into Hades (descendit ad inferos), this member is
well suited to show what comes of making Christ an object of knowledge
in the doctrine of religion. The disputes of the theologians about his
supernatural birth, his two natures, his two wills, his union of
impotence with omnipotence on the cross, are no doubt on the whole
offensive and meaningless; but when the dispute turns on the ``descent
into hell,'' the hollowness of the scientific fantasizing becomes
repulsive. If, on the other hand, the pattern is apprehended as the
object of faith, then his presence in the realm of the dead is a
visible representation of his actual death, his victory over death, and
salvation for the captives of death.

How completely the ideal of personality has triumphed over death and
corruptibility is presented, in the visible manner of revelation, by the
personal resurrection and the personal ascension. He who rose is
personally the same as he who was buried; he who was exalted is the
same, personally the same, as he who was in humiliation. The ideal's
course of life is a complete circuit; so we understand his words when he
says of himself: ``I came forth from the Father, and am come into the
world: again, I leave the world, and go to the Father.''

It is the law of the spirit that actual personalities shall be judged by
their true ideal; the realized ideal of personality is therefore the
judge of mankind: ``He shall come again to judge the quick and the
dead.''

\textit{3) Human personality has the divine for its essential bearer
and mover.}

From the different way in which the relation between actuality and the
ideal is apprehended in knowledge and in faith, it follows of necessity
that the spirit in science must be a different one from that in
religion. The concept of spirit is determined by the concept of the
ideal: in order to understand the second
% --- p. 152 ---
\opage{152}\setcounter{footnote}{0}%
article, we must learn to understand the third. Since science,
according to its principle, is unable to cognize the realized ideal of
personality, for the same reason it cannot cognize the divine
personality in the spirit itself either. For the knower, spirit is only
a unity of the idea and consciousness; for the believer
the idea, on the contrary, is a mediating being between his own human
self and the divine self working in him; the knower relates himself
through the idea to himself; the believer relates himself through the
idea (the idea of truth, of holiness, of righteousness, and so on) to
God.

That the divine self has realized the ideal of personality on earth
could not possibly benefit finite personalities if the self from God did
not at the same time will to work toward the ideal's being relatively
realized in every human being. By merely apprehending the pattern as a
historical person we are still so far from getting the
self-contradiction resolved that we rather get it intensified. When the
ideal of personality is realized in this one human being, then this one
human being stands before intuition as the one who upholds the
unrejectable validity of the ideal demands; he fulfills the law of
selfhood, and the law of selfhood shall be fulfilled. ``Think not that
I am come to destroy the law, or the prophets: I am not come to
destroy, but to fulfil. For verily I say unto you, Till heaven and earth
pass, one jot or one tittle shall in no wise pass from the law, till all
be fulfilled.'' But by getting the ideal realized in this way in the
perfect personality, the imperfect personalities have got the actual
ideal outside themselves, and then actually stand outside the ideal.
Were they to be judged by the ideal in such a position, the judgment
would become a condemnation. The contradiction can only be resolved by
this, that the divine
% --- p. 153 ---
\opage{153}\setcounter{footnote}{0}%
self does not confine itself to an outward presentation of the realized
ideal, but expressly posits the outward realization as the
presupposition for the activity by which, as sanctifying Spirit, it
repeats and forms the original pattern in each single human being. As
the finite, imperfect personalities are influenced and moved by the one
and the same Spirit of personality, the congregation is gathered and the
pattern glorified. Without the actual pattern the Spirit cannot develop
the fundamentally human in single individuals; for the development of
the fundamentally human coincides with the appropriation of the pattern.
Therefore Christ says to the disciples (John XVI): ``I have yet many
things to say unto you, but ye cannot bear them now. Howbeit when he,
the Spirit of truth, is come, he will guide you into all truth: for he
shall not speak of himself. He shall glorify me: for he shall receive of
mine, and shall shew it unto you.''

It will be seen from these fundamental features that the humane
carrying-through of the right of personality must essentially rest on a
religious development of the principle of personality. To the
enlightenment of culture in the state corresponds popular enlightenment
in the congregation: a popular state that by its laws of compulsion
protects the civil personality presupposes a people's church that by the
laws of freedom and the means of the spirit educates and forms the
religious personality. A scientific botching of the religious principles
of the people's life is an essential hindrance to the spiritual life;
but the hindrance can be transformed into a condition when the spirit of
personality gets power over the impersonal, half-personal nature of the
enlightenment of culture.

\begin{center}\rule{3cm}{0.4pt}\end{center}

\chapter*{X.}
\addcontentsline{toc}{chapter}{Lecture X}
\markboth{Lecture X}{Lecture X}
\label{ch:x}
% --- p. 154 ---
\opage{154}\setcounter{footnote}{0}%
When popular enlightenment has once advanced so far that it can really
measure itself against a many-sided, richly developed civilization, when
light stands against light, the light of faith against the light of
culture, then the rays from both sides will be thrown back by differing
resistance and refracted in the heterogeneous media of actuality; in the
refraction a many-colored light will be seen, and in the many-colored
light manifold images of life. But will this refraction always be
harmonious? No, not always. Will these images of light and life bear
witness to a single coherent, internally consistent course of life? By
no means! Among the many refracting media there is one in particular
that by its nature must present a difficult matter to settle, and this
refracting medium is sacred history.

Sacred history has this in common with all history, that its content
is to render the actual; but it differs from all other history in this,
that the actual which it renders unites the supernatural with the
natural in an impenetrable way. The actual has thus here taken up
something that seems to stand in conflict with the concept of
actuality, and yet its validity is to be thought. The concept of
actuality is placed in the foreground; in this concept the standard is
given; we shall now see
% --- p. 155 ---
\opage{155}\setcounter{footnote}{0}%
how differently it is applied according as thought enters the service
of civilization or of popular enlightenment.

From the standpoint of civilization, of negative criticism, actuality
will naturally have to be judged according to the universally valid,
unchangeable laws of nature. Civilization moves according to
knowledge; in knowledge the concept of creation is so subordinated to
the concept of nature that the natural cannot take up the supernatural.
From the standpoint of popular enlightenment, on the other hand,
actuality is to be judged according to the law of personality and of
freedom; popular enlightenment moves according to faith: in faith the
concept of nature is so subordinated to the concept of creation that the
natural must have its ground in the supernatural. So once again faith
stands against knowledge as principle against principle. As the
principle is, so must the criticism be; if knowledge has its criticism,
then faith too must have its own. So criticism stands against
criticism, criticism of knowledge against criticism of faith: in the
alternating light of criticism sacred history must change color. The
light that proceeds from the criticism of faith will everywhere clarify
the difference between mythology and revelation, between myth and
sacred history. The light that proceeds from the criticism of knowledge
will everywhere clarify the likeness between mythology and revelation,
between myth and sacred history. We shall begin with the criticism of
knowledge.

The principle is that whatever in a history, be it profane or sacred,
resembles the mythical in all essential features must be mythical. Now,
for example, the likeness between a talking horse and a talking she-ass
is so unmistakable that if the horse is mythical, then the she-ass must
be so too. That the one story is found in Homer, the other in the
Bible, cannot entitle us to regard the latter as a revelation and the
former as a myth; such a procedure would surely be
% --- p. 156 ---
\opage{156}\setcounter{footnote}{0}%
just as uncritical and arbitrary as if a historian took it into his head
to assume something in a document to be untrue merely because it was
written in Greek letters, and something to be true merely because it was
written in Hebrew ones. The task is to determine the fundamental
character both of the biblical and of the Homeric narrative. In Homer's
Iliad (XIX) it is described how Achilles, after Patroclus has fallen and
the reconciliation with Agamemnon has taken place, arms himself for
battle and mounts the war-chariot. He addresses his horses in these
words: ``Xanthus and Balius! ye famous foals of Podarge, forget not now,
as last time, to bring your master home safe from the field to the camp
of the Achaeans when the battle is fought out!'' This conversation is in
earnest, for the speaker receives an earnest answer. Xanthus, ``the
flying horse beneath the yoke,'' bows low and answers with a sigh:
``Home in safety we shall bring thee yet, thou mighty son of Peleus! Yet
it hastens hard toward thy death; the fault is not ours; guilty is only
the mighty god and powerful fate.'' Thus speaks Xanthus, and how has it
come to speak? ``Voice it got from the lily-armed Hera.'' And how does
the speech affect Achilles? It does not surprise him in the least that
his horse starts to talk, indeed even to talk as one initiated into the
secrets of fate; on the contrary, it is as though it went without
saying. Xanthus is deeply moved, Achilles too is deeply moved and
answers sadly: ``Xanthus, why dost thou foretell my death? The
announcement thou mayest spare thyself.'' Now we all say: It is
beautiful, it is touching, but it is only a myth. That a bright child
can take such a story for actuality is very natural; but that a grown
person with sound understanding can become so childish, be so deeply
stuck in ``childhood faith,'' that he too takes the story as actual, is
in the very highest
% --- p. 157 ---
\opage{157}\setcounter{footnote}{0}%
degree unnatural. No one doubts that the horse to which Hera can give
speech, and which the Erinyes can rob of voice, does not belong under
the category \textit{Equus Caballus}, which has its place in zoology,
but simply and solely under mythology.

We have considered Achilles' horse; let us now consider Balaam's
she-ass! The story is found in the fourth book of Moses (XXII). The
children of Israel, it is said, encamped in the plains of Moab, on the
other side of Jordan by Jericho. Balak, the son of Zippor, king of the
Moabites, sees this power, sees in it an invading enemy and is
terrified; in his terror he knows no better counsel than to send his
elders to the prophet Balaam and, by the rewards of divination, prevail
on him to curse the strangers. But in the night Jehovah reveals himself
to Balaam and says expressly: ``Thou shalt not curse the people: for
they are blessed.'' And Balaam shows no disobedience; on the contrary,
he declares that even if Balak would give him his house full of silver
and gold, he could not act disobediently toward the Lord and curse
Israel. They pressed him; another night passed, and Jehovah revealed
himself again to the prophet, but instead of holding him back he now
says: ``If the men come to call thee, rise up, and go with them; but
yet the word which I shall say unto thee, that shalt thou do.'' Balaam
wakes and does as Jehovah has said, and goes along with ``the princes
of Moab.'' But on the way Jehovah's anger is kindled because the
prophet goes along; what has Jehovah become angry about? That is not
known. Here not the least thing is adduced that might point to an
explanation. If one wants an explanation, one must invent it oneself.
The text tells us only that Jehovah has become angry, and in his anger
sends the angel who is to stop Balaam. The anger is very great; the
prophet has fallen into disfavor to such a degree that he is not even
found worthy to see
% --- p. 158 ---
\opage{158}\setcounter{footnote}{0}%
the angel, whereas the she-ass is found worthy. The angel stands in the
hollow way by the vineyard, where there was a stone wall on both sides,
and when the she-ass sees the angel and presses herself against the
wall, the prophet is angered and strikes her three times; now the Lord
gives the she-ass voice, and the speaking she-ass asks quite meekly:
``What have I done unto thee, that thou hast smitten me?''
Instead of being astonished at this address, Balaam listens to it with
the same naturalness with which Achilles listened to Xanthus, and says:
``I would there were a sword in mine hand, for now would I kill thee.''

So striking is the likeness between the story of the talking she-ass
and the story of the talking horse that if both were found in Homer,
then both, of course, would have to be mythical, and if both were found
in the Bible, then both---regarded from the standpoint of faith,
naturally---would have to be revelations. Here one sees, then, what
comes of letting the concept of creation supplant the concept of nature
and of acknowledging an actuality in which the so-called supernatural is
permitted to mix itself into the natural. The believers will no doubt
appeal to the fact that it is only Hera who works the miracle with the
horse, whereas it is Jehovah who works the miracle with the she-ass;
that Hera's miracle belongs only to the imagination, like herself,
whereas Jehovah's miracle is a sign of the power he actually has as the
almighty creator of heaven and earth; in short: one will defend the
actuality of the miracle by appealing to actual omnipotence. But what
comes of letting the concept of omnipotence devour the concept of
nature in unbounded arbitrariness? We learn it from the she-ass's words
to the embittered prophet. ``And the ass said unto Balaam, Am not I
thine ass, upon which thou hast ridden ever since I was thine unto this
day? was I ever
% --- p. 159 ---
\opage{159}\setcounter{footnote}{0}%
wont to do so unto thee?'' Who, now, is really the speaking subject
here? It surely cannot be either the angel or Jehovah who speaks in
such a way that to Balaam it sounds as though the speech came from the
she-ass: such a miraculous ventriloquism would surely be unworthy of
the Godhead. It must, then, be the she-ass that speaks. But what kind of
miracle is this, after all? It is a miracle whereby a dumb animal soul
is in an instant re-created into a rational, self-conscious soul; and
that in such a way that this self-conscious soul in a single moment is
developed to such a degree that it knows the language, the language
Balaam as a child took a longer time to learn; a miracle whereby the
animal soul is not merely re-created from dumb to speaking, and so into
a quite different soul, a quite different being, but on top of that is
re-created in such a way that, in spite of the re-creation, it
nevertheless remains the same being, the same soul. The man on whose
head the hair does not stand on end when he is to acknowledge such an
actuality---his head must have lost either its hair or---its wits; for
the sight is horrible. The speaking she-ass is actual; the dumb one is
also actual: the one crowds the other out, and yet these two actual
animals are in actuality one and the same animal; that is what one may
call actuality!

However, lest it should look as though the criticism of knowledge, if we
may say so, had fastened its teeth in a single, obscure narrative, less
essential to the history of revelation as a whole, we shall, to
illuminate the miraculous actuality, choose a couple of examples that
stand in so intimate an organic connection with the economy of
revelation in its entirety that they cannot be taken away without the
whole collapsing; we mean the legend of Paradise and the Gospel of
Christ's resurrection.

% --- p. 160 ---
\opage{160}\setcounter{footnote}{0}%
The legend of Paradise must, from the standpoint of faith, be
apprehended as a true history and, as a true history, refer to the
actual. There must, then, in the actual Paradise have been an actual
tree of life, an actual tree of knowledge, an actual serpent, and so on,
and so on. Adam and Eve must have lived as full-grown human beings in
holy innocence and perfect blessedness before the fall, and upon the
fall there must have followed an actual curse upon the whole actual
earth. We shall now test what kind of actuality this is. Whatever in all
respects resembles the mythical must be mythical; from this principle we
begin the investigation again. That geology has no place for an actual
Paradise is surely strong testimony against its actuality; in the
present connection, however, the emphasis falls not so much on the facts
of natural science as rather on the very concept of actuality. Just as
little as natural trees, for example apple and pear trees, can be
thought to grow in a supernatural, i.e.\ supraterrestrial, soil, just
as little can such supernatural trees as the tree of life and the tree
of the knowledge of good and evil be thought to grow in a natural soil.
If we now ask what the actual soil in Paradise must have been like, the
answer becomes very doubtful. It is said expressly that the Lord God
Jehovah took the man and put him into the garden of Eden to dress it and
to keep it. The part of the actual garden that Adam was to dress and to
keep may then well be assumed to have been the natural part of
actuality; whereas the section of the garden in which the tree of life
and the tree of knowledge grew hardly needed any cultivation and so
must presumably have made up the supernatural part.

That Adam is actually to dress and to keep an actual
% --- p. 161 ---
\opage{161}\setcounter{footnote}{0}%
garden indicates, as has been said, that Adam is an actual and natural
human being. But how, now, is a natural human being to be thought able
to pluck fruits from a supernatural tree, indeed not only to pluck them
but even to eat them? And yet, alas, first Eve and then Adam, in their
naked naturalness, succeed only too well in plucking and eating the
supernatural; in actuality they succeed so well with the tree of
knowledge that Jehovah becomes downright afraid that they will succeed
with the tree of life as well. ``And the Lord God Jehovah said, Behold,
the man is become as one of us, to know good and evil: and now, lest he
put forth his hand, and take also of the tree of life, and eat, and live
for ever.'' That this may be prevented, Adam with his wife is then
actually driven out of the actual Paradise, and the way to the tree of
life guarded by the cherubim. Before their expulsion from Paradise, our
first parents must thus, in the first actuality, have had a wonderful
mixture of the natural and the supernatural, so that only after being
driven out did they become properly natural.

If we pass over to the animal world and ask the zoologist what he thinks
of the part of sacred history that deals with the animals in Paradise,
he will without doubt answer: Not very much. It is said of the animals
that they were created of all kinds, and that they were all brought to
Adam, who gave them names. Among the paradisal animals thus brought to
Adam the beasts of prey must also have been. In the paradisal state,
however, they could not possibly go hunting for prey; that would surely
in actuality have disturbed the actual blessedness. But, we ask, what is
a beast of prey that does not go after prey? It is a round square. It
lies in the zoological
character of beasts of prey, in their claws, teeth, digestive organs and
so on, that they absolutely must be carnivorous; but
% --- p. 162 ---
\opage{162}\setcounter{footnote}{0}%
how should flesh-eating and blood-drinking animals be able to live when
they dared not go after prey? Might it perhaps have been supernatural
flesh on which they lived? But then they must surely also have been
supernatural beasts of prey. That lions, tigers, leopards and so on did
not yet go after prey means that the lion was not yet a lion, the tiger
not a tiger, and so on. The species of beasts of prey had thus not come
into being; and yet the narrative says that God had already created all
``the beasts of the field.''

It becomes still worse when one comes to the serpent. With the talking
serpent theology has really taken things far too easily, now without
further ado making it the devil himself, ``that old serpent,'' now, more
speculatively, ``the cosmic principle.'' As regards the cosmic
principle, it takes more than human understanding to see how so
metaphysical a being as a principle is to be able either to transform
itself into a serpent or to speak through a serpent. And as regards the
devil, the narrative contains not the slightest hint that the serpent
is possessed by the devil.

When the serpent speaks to the woman, it is the serpent itself in a
quite different sense than when Xanthus speaks to Achilles, or the
she-ass to Balaam. The horse and the she-ass are dumb animals, which
only by a miracle could become speaking. But not so with the serpent;
for ``the serpent was more subtil than any beast of the field which the
Lord God Jehovah had made.'' The serpent belongs to the beasts of the
field, and so to the dumb ones; but---so fruitful is the invention of
subordinating the concept of nature to the concept of creation---though
by its naturalness belonging to the dumb animals, the serpent has
nevertheless from the beginning known how to unite the supernatural with
the actual so cleverly that it has become speaking
% --- p. 163 ---
\opage{163}\setcounter{footnote}{0}%
just as well as man. Should anyone still harbor doubt that, in the
narrative's meaning, it is actually the serpent that speaks, and speaks
of its own, let him but hear Jehovah's curse upon the serpent: ``Because
thou hast done this, thou art cursed above all cattle, and above every
beast of the field; upon thy belly shalt thou go, and dust shalt thou
eat all the days of thy life.'' The serpent has thus by its action drawn
upon itself Jehovah's wrath and punishment, is thus accountable, has
thus acted with freedom and deliberation. If this is not mythical, what
then is mythical? Can anyone well shut his eyes to the likeness between
the serpent's trespass and Adam's and Eve's trespass? Eve is punished,
for she is accountable and has misused her freedom; Adam is punished,
for he is accountable and has misused his freedom. The serpent is
punished; so it too is accountable and has misused its freedom. But the
consciousness that, as regards freedom and accountability, can set a
serpent on an equal footing with a human being is not a consciousness
that grasps the actual, but a mythical consciousness.

What the mythical consciousness has in all naivety presented as a
supernatural fact can, when the content is of deeper ethical
significance, still retain its validity long after the enlightenment of
culture has explained away the supernatural; the phenomena are no
longer, indeed, taken to be actual, but are honored and used as holy
symbols. Instead of a supernatural actuality we get a religious
symbolism, and now it is the higher spiritual task of the teachers of
the people---a task for whose solution they also, of course, prepare
themselves worthily by studying ``sacred theology''---to apply this
symbolism in so ambiguous, clever and ingenious a way that the common
man suspects nothing amiss but takes the symbolic for full actuality.
Thus the legend of
% --- p. 164 ---
\opage{164}\setcounter{footnote}{0}%
Paradise can be used with success in popular instruction as a symbol and
yet look like sacred history; thus the myth of Christ's resurrection can
still, many an Easter morning, have its effect as a historical event and
yet be interpreted as a symbol of the victory of life over death, of
light over darkness, of truth over the lie. But that the story of
Christ's resurrection is just as certainly a myth as the story of
Paradise, the criticism of knowledge will easily be able to prove.

We shall not dwell here on what Strausz has sufficiently shown, that
the evangelical accounts contradict one another in many essential
points; we shall---for the further testing of the actuality that faith
gets out of letting the concept of creation supplant the concept of
nature---keep to a single account, in order to see how contradictory
such an account is in itself.

We choose the account in Luke (XXIV, 36---43). The two disciples who met
the Risen One on the road to Emmaus have already returned to Jerusalem,
where they find ``the eleven gathered together, and them that were with
them.'' Just as those assembled are absorbed in conversation about the
wonderful thing that had happened, Jesus himself suddenly stands in the
midst of them. ``But they were terrified and affrighted, and supposed
that they had seen a spirit''! The first immediate impression is as of
a vision of a spirit. To the disciples' great surprise, however, this
vision of a spirit becomes the Crucified himself in bodily actuality, in
``flesh and bones.'' As yet further proof that in this actuality the
supernatural is united with the natural, the one who reveals himself
puts to them the question: ``Have ye here any meat?'' They give him a
piece of a broiled fish and of a honeycomb. ``And he took it, and did
eat before them.'' That criticism must regard such an appearance as
impossible has simply
% --- p. 165 ---
\opage{165}\setcounter{footnote}{0}%
and solely its ground in this, that it must find what is
self-contradictory impossible. That the Risen One presents himself to
the disciples as a supernatural being, a spirit, is no doubt in formal
agreement with the concept of resurrection, insofar as the resurrection
is a miracle, and he whose life rests on a miracle must, in bodily as
well as in psychical respects, be a miraculous being. But how, now, is
it to be thought that a miraculous being can perform so natural an
action as that of eating ``a piece of a broiled fish, and of an
honeycomb''? It must presumably be thought in faith in the same way as
it must be thought in faith that Jehovah, the almighty creator of heaven
and earth, actually ate actual curds and actual veal with Abraham under
the tree in the grove of Mamre; it must presumably be thought to happen
by a miracle. But to eat is a natural action, and no miracle; what is no
miracle thus happens by a miracle: the supernatural devours the natural;
but precisely by devouring the natural the supernatural shows itself as
the unactual, the mythical.

This is the scientific result of the criticism of knowledge. We have let
the knower express himself so many-sidedly, fully and unreservedly that
no one who is able to understand his utterances will be able to mistake
the fundamental thought; the light with which he gives light is the
light of culture, and he does not put the light under a bushel. So long,
however, as the people---particularly as regards the great mass---can
be kept in a certain peasant-like ignorance, such a criticism from
culture will either not be noticed at all, or else the critic will be
regarded as an ungodly person with whom it is not worth having anything
to do.
% Danish prints no full stop after "med" (printer's defect); supplied.
But as popular enlightenment advances, the matter must change. The
enlightened will see that the attacks of the criticism of knowledge on
religion do not consist in wanton, arbitrary
% --- p. 166 ---
\opage{166}\setcounter{footnote}{0}%
whims, but proceed with inner necessity from a definite fundamental
thought, a distinctive principle. The dispute about the validity of
religion must---it is an unrejectable condition for the spiritual
life---become a question of principle. In proportion as the principle
of culture becomes clear in broad outline to the popular consciousness,
the principle of faith must then also become clear: against the
criticism of knowledge we set the criticism of faith.

The emphasis falls, as has already been indicated, on the difference
between mythology and revelation, between myth and sacred history. This
difference cannot be decisively determined by a finite comparison of
the phenomena; for the phenomenon of revelation has a deceptive likeness
to myth; in the phenomenon as phenomenon a talking horse stands on an
equal footing with a talking she-ass. The essential judgment depends on
the principle.

That the criticism of knowledge rejects every sign of revelation is
quite in order, for the sign of revelation is precisely a phenomenal
interruption of the natural connection of things; if it did not appear
thus, it could not possibly be a sign. But when the scientific critic
imagines that by thinking the sign of revelation through he can
dissolve it into a contradiction and present it as a meaningless breach
of the order of nature, then he oversteps the absolute boundary of
knowledge.

The sign of revelation has nothing whatever to do with the scientific
apprehension of the connection of nature; it is and remains a religious
matter between the God who reveals himself in the sign and the believer
who receives the revelation. From the narrative of the talking she-ass
no natural-scientific conclusion whatsoever can be drawn, either in
attack upon or in defense of the phenomenon. What is the phenomenon of a
talking she-ass? A sign of omnipotence! What is the
% --- p. 167 ---
\opage{167}\setcounter{footnote}{0}%
phenomenon of a burning bush? A sign of omnipotence! What is the
phenomenon of the raising of the dead? A sign of omnipotence! Name what
wonder, what manifestation in the whole of sacred history one will; as
soon as thought tries to penetrate into the essence of the phenomenon,
it strikes everywhere against the same barrier, halts before the same
riddle, gets to all questions the same monotonous answer: it is the
sign of omnipotence!

Were the criticism of knowledge now really able to prove that a sign of
omnipotence is impossible, then it would have to be able to prove that
the almighty will itself is impossible; but proofs of this nature have
validity only for those who, by forgetting personal self-contradiction,
prove that in the deepest sense they have forgotten themselves.

``But,'' it is objected, ``by such signs of omnipotence the most
arbitrary of breaches are surely made in the order of nature.'' This
objection would mean something if the sign, in its capacity as sign of
revelation, actually had anything to do with the natural order of
things; but such a presupposition rests on a scientific
misunderstanding. Each sign is a single, a unique phenomenon, which
cannot possibly be introduced as a member in the series of natural
phenomena; but if it cannot be introduced into the series of natural
phenomena, then neither can it occasion any induction whatsoever that
would disturb natural science: the sign is, each time it is given, a
personal, an exclusively personal matter between the one who reveals
himself by the sign and the one or ones who receive the revelation. The
limit for the imparting of such signs cannot be determined by any law of
nature; but the limit is not arbitrary on that account: it is
determined by divine wisdom, in that it is determined by the divine
will.

% --- p. 168 ---
\opage{168}\setcounter{footnote}{0}%
It follows from the principle of faith that the concept of nature,
namely the religious one, must be subordinated to the concept of
creation; but it by no means follows from the principle that the
religious consciousness should be blind to actual nature. The poetic
contemplation of nature is, after all, still acknowledged as justified
over against the scientific, and why? Because one acknowledges the idea
of beauty as justified over against the idea of knowledge. Why, then,
will one not acknowledge the religious contemplation of nature as
justified over against the scientific? Because the consciousness of
culture has got so well under way with bringing faith under knowledge
that it has overlooked the self-validity of the principle of faith! The
Davidic contemplation of nature, for example, is not scientific, but
should it on that account have no religious validity? The religious
character of the contemplation of nature is determined by this, that
everywhere, even where there is no question of particular miracles, it
refers the things of nature to the creating and sustaining omnipotence.
Instead of dwelling on the natural connection of the phenomena, the
believing eye will everywhere, through the phenomena, discover God
himself. ``He spreadeth out the heavens like a tent; he frameth his
lofty halls in the waters; clouds he maketh his chariot; so he rideth
upon the wings of the storm'' (Psalm.\ IV). Should such a
% "IV" as printed; the verses are Ps. 104 (CIV).
contemplation of nature well let itself be supplanted either by
astronomy or by meteorology? Because the religious view of nature is
here uttered in the form of poetry, should it on that account cease to
be religious? ``The trees of Jehovah stand luxuriant, the cedars which
he planted on Lebanon, where the little birds build their nests''; could
such an utterance perhaps not subsist alongside a scientific botany?
``Thou makest it grow dark,'' says the holy singer of Jehovah, ``then
cometh the night, in it all the beasts of the forest creep forth, the
young lions roar after their prey and seek their meat
% --- p. 169 ---
\opage{169}\setcounter{footnote}{0}%
from God''; it can surely not be the task of zoology to subject this
view of the beasts of the forest to a scientific criticism.

Should anyone be so cultivated as to think that the Psalmist is not in
earnest with the examples of God's activity in nature adduced here, that
they are only figurative utterances of a fleeting mood: then such a one
shows only that out of sheer culture he has forgotten what religion is,
forgotten what emphasis the believer lays on the words: ``O Lord, how
manifold are thy works! in wisdom hast thou made them all!'' ``But,''
one asks, ``can the Psalmist then really himself have believed that
Jehovah now and then rides upon the storm clouds, that he actually and
with his own hand planted the cedars on Lebanon so that the birds might
build nests in them, and that the young lions, when they roar after
their prey, actually pray to God? No, the natural sciences surely cannot
possibly have stood at so low a level even in that time, which was
certainly in other respects by no means enlightened; they must be
images, symbols and images, nothing but symbols and images!'' Yes,
certainly they are images, but yet more than images; certainly they are
symbols, but these symbols have nonetheless actuality, full actuality
for the spirit; in the believing contemplation of nature actuality is
symbol, the symbol actuality. What the symbol is in the contemplation
of nature rich in mood, that again, in a special way, is the sign, the
sign of omnipotence in the miracle. The actuality of the miracle is a
sign, a symbol. ``Only a symbol?'' Yes, but this symbol is at the same
time actuality. The fundamental question is: Can faith count as a
principle or not? If one rejects the principle, then the religious
contemplation of nature can naturally have only subjective, aesthetic
validity, then it gives only images but no actuality; if one
acknowledges the principle, then one is also compelled to grant to the
contemplation of nature corresponding to the
% --- p. 170 ---
\opage{170}\setcounter{footnote}{0}%
principle actuality, fully valid actuality for the spirit.

In order now to investigate more closely how the matter really stands
with the religious contemplation of nature, and especially with the
apprehension of the actuality of the miracle, we must first determine
somewhat more precisely the relation between the experience of faith and
the experience of sense. When believers in sacred history receive a
revelation, the phenomenon often appears with so striking an
immediacy, with so definite a stamp of outward actuality, in so
unmistakable an interaction with natural surroundings, that it looks as
though such a phenomenon could be apprehended by sensuous seeing,
sensuous hearing, sensuous feeling alone, and so without faith. This is
nonetheless a gross error. Just as faith always presupposes revelation,
so revelation in turn presupposes faith: revelation and faith are in
principle one, and cannot possibly be separated. The believer sees,
hears, feels; but in all this sensing faith is present; he sees with the
eye of faith, hears with the ear of faith: all his experience is and
remains an experience of faith. Whether the single individual in the
single case receives a true or a false, an actual or an imagined
revelation---for that the person concerned must himself be answerable.
How a Balaam hears and understands Jehovah's command, for that he stands
personally accountable to Jehovah; nothing about such things can be
decided in pure generality. How a David hears and understands Jehovah's
command is a personal affair between Jehovah and David. In this respect
the two accounts of the census are highly instructive. According to
1 Sam.\ (XXIV) David thinks he hears Jehovah's voice: ``Go, number
% "1 Sam." as printed; the census narrative is 2 Sam. XXIV.
Israel and Judah!'' David thus imagines that in this case he has
received a true revelation; but it is an imagining. It is not Jehovah's
voice
% --- p. 171 ---
\opage{171}\setcounter{footnote}{0}%
that David hears; no, it is (1 Chr.\ XXI.) another voice. The revelation
is false; why? Because David has aroused Jehovah's wrath, and the reason
for this is plain enough. It was, in the religious way of seeing of
David and---as one clearly notices from Joab's words: ``why doth my
lord the king delight in this thing?''---of his contemporaries, a very
questionable matter to undertake a census; it was, for good
psychological reasons, interpreted as a sign of a certain arrogance;
the king wanted by so sensational a step precisely to make it known
that the people of Israel was now so numerous that, without needing
Jehovah's special aid, it could well rely on its own strength. So it
ferments in David's inner being: egoism is in conflict with piety,
ambition with humility, vanity with the thanksgiving owed to Jehovah. In
such a state the believer is unsure of himself and exposed to false
promptings. If, however, the word of revelation could be apprehended by
a sensuous ear without faith, then surely the guilt would have to fall
on Jehovah; but the falsity lies in the hearing, and the guilt falls on
David.

To make sure of the interpretation of an actual phenomenon of
revelation, the question is sometimes raised: must not an unbeliever,
had such a one been present, also have seen the supernatural? From the
answer to such a question one will learn how far the interpreter
regards the manifestation as an actual fact or merely as an inward
vision, a product of the believers' own fantasy; but one does not
notice that the question contains a self-contradiction. If the
supernatural could actually be sensed and seen with the same eyes that
see the natural, then it would surely not be supernatural; then a
revelation could be imparted to a human being without faith, and faith
% --- p. 172 ---
\opage{172}\setcounter{footnote}{0}%
would cease to be a principle. Instead of joining actual seeing and
actual hearing with actual faith, one wants to throw faith away and
seize what is revealed in sensuous, raw immediacy; that is to violate
the holy. One concludes thus: if an unbeliever who was present could not
also have seen the vision, then the vision is not actual; that is a
false conclusion. Instead of placing full assurance in faith, one takes
faith away and places assurance in mere sensation. Let the subject be
Christ's resurrection and the Risen One's various manifestations to the
disciples. The evangelists lay the very greatest weight on the actuality
of the resurrection: ``The Lord is risen indeed, and hath appeared to
Simon,'' say the eleven in Jerusalem to the two who came from Emmaus;
but nowhere will one discover in any of the witnesses the remotest hint
that the Risen One was seen by unbelievers as well. On the contrary,
there are manifold testimonies showing that even the eyes of believers
had to be specially opened before they could know the Transformed One.
Mary Magdalene meets him in the garden; but she does not know him; she
supposed, indeed, that it was the gardener (John XIX., 14---15). Of the
% "Joh. XIX." as printed; the passage is John XX, 14--15.
two who went to Emmaus, when Jesus ``himself drew near, and went with
them,'' it is said (Luke XXIV, 16): ``But their eyes were holden that
they should not know him.'' Even those assembled in Jerusalem (Luke
XXIV, 36---42), who had already been prepared from several sides by the
glad tidings, even these could not at once and directly recognize him
who suddenly stood in the midst of them; ``but they were terrified and
affrighted, and supposed that they had seen a spirit.''

Under the presupposition of faith, on the condition of faith, the true
revelation first acquires fully valid actuality, actuality for the
spirit, spiritual actuality; all phenomena of revelation
% --- p. 173 ---
\opage{173}\setcounter{footnote}{0}%
are in actuality phenomena of the spirit. That actuality here is
actuality for the spirit must not be understood as though the spirit
itself were too powerless to annul the immediately natural barriers. It
is no imperfection in the phenomenon of the spirit that its highest
actuality is exclusively for the spirit; it is on the contrary a
perfection: revelation will confirm not the impotence of the spirit but
its omnipotence, its absolute power over itself. We see it in the Risen
One. As the disciples' eyes are opened, as it becomes clear within them,
so that they see him in faith, hear him in faith, the outward form too
becomes clear and definite, takes on actuality, indeed tangible
actuality. How is it possible that a form which begins by showing itself
as spirit can end by presenting itself in full, tangible actuality?
Yes, how is it possible? With this question the criticism of knowledge
imagines that it can conjure up enormous difficulties and entangle the
believing thinker in enormous contradictions. That is an imagining; from
the standpoint of faith there is in such a question only one
contradiction, one single one, and it lies not in the matter asked
about, but in the question.

As a telling example of what a phenomenon of revelation in full
actuality has to mean, there serves especially the account of the Risen
One's manifestation to Thomas (John XX, 26---29): ``And after eight days
again his disciples were within, and Thomas with them: then came Jesus,
the doors being shut, and stood in the midst, and said, Peace be unto
you. Then saith he to Thomas, Reach hither thy finger, and behold my
hands; and reach hither thy hand, and thrust it into my side: and be not
faithless, but believing. And Thomas answered and said unto him, My Lord
and my God. Jesus saith unto him, Because thou
% --- p. 174 ---
\opage{174}\setcounter{footnote}{0}%
hast seen me, Thomas, thou hast believed: blessed are they that have not
seen, and yet have believed.'' The disciple Thomas appears here in a
quite peculiar light. He is unbelieving, people say. Certainly he
betrays a questionable weakness in faith; we hear him, after all,
before the manifestation takes place, expressly declare: ``Except I
shall see in his hands the print of the nails, and put my finger into
the print of the nails, and thrust my hand into his side, I will not
believe''; a Sadducee could say the same. And yet Thomas is no Sadducee.
% Danish prints "Saducæer" twice; Rettelser: read "Sadducæer". Translated
% in the corrected sense.
He does not deny the possibility of the resurrection; but certainty he
wants, full certainty; he wants to have actuality, full, tangible
actuality, before him. So then he gets actuality too; he gets it as
palpably, as immediately as can well be thought. But where, and how,
does he get it? In the circle of the disciples, and after having heard
the Risen One's greeting: ``Peace be unto you!'' Before this voice doubt
gives way, by this greeting faith is awakened, and now Thomas is
prepared in the spirit to grasp the full actuality. But with the
manifestation there follows a reproach. What is it that Christ
reproaches Thomas for? That he wants full assurance of the actuality of
the resurrection is not held against him; a credulity that relies on
loose rumors Christ never recommended. No, the reproach is that the one
of little faith wants to make his assurance dependent on the testimony
of the senses. Thomas wants, as the saying goes, ``to see it before he
will believe it''; he wants, in other words, ``to relate himself to the
holy facts in a way corresponding to that in which the natural scientist
relates himself to the facts of nature.''

The mistake of which the disciple Thomas here makes himself guilty at
first hand is precisely the one that the orthodox dogmaticians, in
Protestantism just as fully as in
% --- p. 175 ---
\opage{175}\setcounter{footnote}{0}%
Catholicism, have in various ways imitated, and of which they have made
themselves guilty at second hand. Just as Thomas would rather regard the
phenomenon of revelation in a sensuous, outward, empirical way, so the
objective dogmaticians even to this day want first of all to make sure
of the outward historical actuality of the facts; on the outward, firm
actuality they then think, with the aid of the Spirit, to be able to
build their faith as on a walled foundation. But when the outward
actuality of experience is posited as the first thing, the Spirit is
offended and denies the objective believers its testimony. Only to
spiritual facts will the Spirit bear witness; its testimony is,
according to the principle, conditioned by the phenomena of revelation
being apprehended and determined, in their highest actuality, as
phenomena of the spirit. Precisely because the orthodox theologians have
held fast to the actuality of the phenomena of revelation without
spirit, negative criticism has been able to go to the opposite extreme
and offer us, in ingenious myths, hollow phenomena of the spirit without
actuality.

To get the actuality of the phenomena of revelation thoroughly grasped,
and illuminated from every side in the interpretation of sacred
history, is an essential condition for the spiritual life in the present
age.

\begin{center}\rule{3cm}{0.4pt}\end{center}

\chapter*{XI.}
\addcontentsline{toc}{chapter}{Lecture XI}
\markboth{Lecture XI}{Lecture XI}
\label{ch:xi}
% --- p. 176 ---
\opage{176}\setcounter{footnote}{0}%
Sacred history is the history of revelation, the history of the God who
reveals himself in the Word. But the God who reveals himself in the Word
is a speaking God, who differs from the mute gods of the nature religions
and from the silent Godhead of science. The mythical gods are mute gods;
gods of ``silver and gold,'' a work of men's hands. ``They have mouths,
to be sure, but cannot speak; they have eyes, to be sure, but cannot see;
they have ears, but cannot hear'' (Psalm CXV).
% the Danish has a closing quotation mark after "høre" with no opening mark (printer's defect); the opening quote is supplied here
What is said of ``the idols'' from the standpoint of revelation is
literally true, even when we take the idea into account. The poet invokes
the Muse; he begs her for mood, for inspirations, for the assistance of
spirit, for the wondrous gift of language; but how matters stand with this
invocation is easy to see through. For the Muse to be able to give a mood,
the poet must already be in the mood; for the flashes of genius to be able
to begin within him, they must already have begun; it is the movement of
the idea on which everything depends: the Muse as Muse is mute; the mute
goddess makes no one speak. That the poet must nevertheless woo the favor
of the Muse means that under the pressure of the idea he must redouble his
spiritual being, must, in order to be moved, let himself
% --- p. 177 ---
\opage{177}\setcounter{footnote}{0}%
be moved; thus he sees the first mover in the goddess's glance and
receives as a divine gift what originally dwells in himself. So
understood, the invocation of the Muse is a psychological necessity, a
beautiful ambiguity that becomes poetry well. As an image of being that
bears the idea, the Muse is not a goddess; as a goddess, that is: as a
being on its own in independent existence, the goddess is an idol and no
Muse.

If the Muse as goddess cannot speak to the poet, then the Godhead in its
pure concept can still less speak to the thinker. The thinker, to be sure,
in thinking the idea of the Godhead, must also redouble his being, insofar
as he must necessarily think his finite self in distinction from the
absolute idea; but he must in turn be able, with full consciousness, to
cancel the redoubling if he wants to keep himself absolutely within pure,
universal thinking. To think God, says Hegel, is not like thinking an
animal or a stone; for the animal and the stone we represent to ourselves
as having a finite existence outside our concept; our concept of the
absolute being, on the other hand, our true concept, is the absolute being
itself. Although thinking and speaking are so inseparable that all
determinate thinking coincides with an inner speaking, it can never occur
to the thinker either to speak to the absolute being or ever to expect that
the absolute being should speak to him: the Godhead's speaking to man is,
from the standpoint of knowledge, man's speaking within himself.

It is otherwise from the standpoint of faith. The God of faith is a God
who can speak to man, a God to whom man can speak. The God of revelation
is a speaking God: Abraham heard the speaking God and wandered away from
the mute gods; Moses heard the speaking God, and Israel became God's
peculiar people by virtue of all ``these words which God spake.''

% --- p. 178 ---
\opage{178}\setcounter{footnote}{0}%
But for knowledge and science there can scarcely be conceived any concept
that contains a greater contradiction than the concept of a speaking God.
On this occasion we must once again let negative criticism have its say.

Among all the miracles, great and small, of which the Old and New
Testaments tell, the greatest beyond comparison is this, that God in
heaven himself takes the floor and sets about speaking. On the assumption
that he really has a secret, this is of course quite in order; Paul says:
``For what man knoweth the things of a man, save the spirit of man which
is in him? even so the things of God knoweth no man, but the Spirit of
God.'' Just as it depends, then, on a man's good will whether he will come
out with his secret or keep silent about it, so it depends, in a still
stricter sense, on God's good and unsearchable will what and how much of
his hidden counsels he will confide to his elect. But when the miracle of
the speaking God is set up as the fundamental miracle of revelation, it is
seen at the same time how all the contradictions hidden in the essence of
faith involuntarily leap to the eye. If revelation is an arbitrary act,
then the revealed word must necessarily become an expression of
arbitrariness, hence an arbitrary word. That the word of revelation is an
arbitrary word is indeed proved by the Old Testament from beginning to
end.

``And God spake all these words, saying''; thus begins (Exodus XX) the
lawgiving from Sinai. What is characteristic of this lawgiving is thus not
in the first instance its content but its origin; not what is commanded,
but that it is commanded, that it is expressly commanded by God. What this
signifies can be illuminated by the well-known opposition between two
points of view that
% --- p. 179 ---
\opage{179}\setcounter{footnote}{0}%
already found expression in the scholasticism of the Middle Ages:
\textit{it is good because God wills it}, and \textit{God wills it because
it is good}. The latter proposition makes God's will dependent on the idea
of the good and agrees with the demand of knowledge. If God commands only
what is good and forbids what is evil, then I need only use my reason to
know the difference between evil and good, and have no need of arbitrary
revelations. Such commandments as: Thou shalt not steal; Thou shalt not
bear false witness against thy neighbor; thou shalt not kill, and the
like, are universally valid commandments of reason, which can be cognized
without a speaking God's needing to dictate them. But to take the moral
commandments in this way conflicts with the spirit of the Mosaic
lawgiving. Why may an Israelite not steal? Answer: because Jehovah has
forbidden him to. If Jehovah had not forbidden him to? Then he might well
steal. And if Jehovah one day commanded him to? Why, then it would be his
duty to steal. Thus we come to hold fast to the first proposition:
\textit{it is good because God wills it}, as an unconditional proposition
of arbitrariness, an unadulterated proposition of revelation. From this
point of view it cannot surprise us to find moral commandments of reason
placed in a row with wholly special, arbitrary commandments. It says:
``Thou shalt not make unto thee any graven image, or any likeness of any
thing that is in heaven above, or that is in the earth beneath, or that is
in the water under the earth''; such a commandment quite rightly
presupposes a special revelation, for no reasonable man would of himself
hit upon the idea that God should be so sworn an enemy of sculpture and
painting.

When a people is to such a degree made morally a minor that it is not
able by its own reflection to find out the difference between good and
evil, but is to have all
% --- p. 180 ---
\opage{180}\setcounter{footnote}{0}%
its religious, moral and civil precepts expressly dictated, then it goes
without saying that ``the speaking God'' must speak many words and
prescribe everything down to the smallest details. The particular laws
have therefore gone forth just as immediately from Jehovah's mouth as the
fundamental law. It is incredible how far a legislating God can go into
details. ``And Jehovah spake unto Moses (Exodus XXV), saying, Speak unto
the children of Israel, that they bring me an offering of every man \dots\
And this is the offering which ye shall take of them; gold, and silver,
and brass, and blue, and purple, and scarlet, and fine linen, and goats'
hair, and rams' skins dyed red, and sealskins, and acacia wood, oil for
the light, spices for anointing oil,'' etc.\ etc. Further (Exodus XXVI):
``Moreover thou shalt make the tabernacle with ten curtains of fine twined
linen, and blue, and purple, and scarlet.'' If, then, eleven curtains,
for example, were used for the tabernacle, the eleventh curtain would
perhaps be a greater offense to God than false witness against one's
neighbor. ``The length of one curtain shall be eight and twenty cubits,
and the breadth eight cubits''; now if it happened that a curtain came out
half an inch too short, the guilty party would bring down upon himself the
wrath of Jehovah and perhaps be cut off from his people.

Is this not, asks the man of culture, to drag the eternal, the exalted,
down into the finite and the petty? And is it strange that the old
heathens said mockingly: such a speaking God is no God of spirit; he is a
talkative God, who wants to meddle in everything, busy himself with
everything, know all about everything (garrulus, inquietus, impudenter
etiam curiosus)? Where divine arbitrariness reigns, it is impossible to
distinguish between essential and inessential, moral and immoral,
permissible and impermissible. There is a kind of lie which people are
accustomed to call a lie of necessity. If we are asked whether
% --- p. 181 ---
\opage{181}\setcounter{footnote}{0}%
such a lie is permitted, especially on a large scale, when it concerns one
of the serious relations of life, and where its use plainly betrays
cowardice, we answer unconditionally: such a lie of necessity is not
permitted. God cannot possibly approve such a thing; for from the
proposition ``God wills it because it is good'' it of course follows
conversely: ``God does not will it because it is evil.'' Now if one looks
in the twelfth and twentieth chapters of the first book of Moses, what is
one then to judge of Father Abraham? The twelfth chapter begins with the
election of Abraham; Jehovah speaks to him as to his chosen one, and bids
him go out of the land; he will show him the way and give him his
blessing. Imagine, then, a man who is raised so high that the Majesty, the
holy God, deigns to speak to him. Then Abraham comes to Pharaoh in Egypt
with Sarah, and what does he do now? He becomes afraid and says to his
wife: ``Behold now, I know that thou art a fair woman to look upon:
therefore it shall come to pass, when the Egyptians shall see thee, that
they shall say, This is his wife: and they will kill me, but they will
save thee alive. Say, I pray thee, thou art my sister,'' etc. What is it
that Abraham does here? He does that at which a knight would blush; he
conceals that Sarah is his wife, indeed, what is more, he even makes a
profit out of it; for in return for Sarah Pharaoh gave him, as we know,
``sheep, and oxen, and he asses, and menservants, and maidservants, and she
asses, and camels,'' as many as he wanted; Pharaoh in his guilelessness
has no inkling that he is on the point of committing a great crime. So he
must come to know it in another way. To stop the unpleasant consequences
of Abraham's cowardice, Jehovah must step in to help with miracles.

``And the Lord plagued Pharaoh and his house with great plagues because of
Sarai Abram's wife.'' How abashed must Abraham not stand before Pharaoh
% --- p. 182 ---
\opage{182}\setcounter{footnote}{0}%
and listen to his reproachful words: ``What is this that thou hast done
unto me? why didst thou not tell me that she was thy wife? Why saidst
thou, She is my sister? so I might have taken her to me to wife: now
therefore behold thy wife, take her, and go thy way.'' That Jehovah goes
along with the intrigue we can understand at a pinch, since Jehovah is,
after all, Abraham's friend; and what will one not do for a friend? ``There
lies,'' say our believers in everything old, ``a great comfort in seeing
how God can use frail people as pillars of his kingdom. God's guidance of
Abraham, and his employing him in the service of his kingdom, is defense
enough for Abraham.'' But had Abraham not deserved a serious reprimand?
Since he receives no reprimand, he is naturally so far from regretting
what he has done that he rather seems, in all quietness, to congratulate
himself on being the first inventor of the lie of necessity. In the
twentieth chapter he and Jehovah then play the comedy over again with
Abimelech in Gerar. And then people still find it strange that English and
German freethinkers have been able to attack the positive moral teaching
of the Old Testament; and then they cross themselves that a Voltaire,
according to the well-known \textit{taisez-vous, taisez-vous, Habakuk est
capable de tout}, thought he could put the greatest absurdities into the
mouth of a prophet!

With this the critique of knowledge has had its due. ``But,'' someone
might ask, ``why should we sit here and listen to so offensive a critique,
since we know, after all, that it springs from sheer unbelief and can
easily give offense to the simple believers?'' I answer: With the offense
of the simple believers there is no danger; their faith does not sit so
loosely that a little critique by argumentation should be able at once to
shake it. On the other hand there is every possible danger that, under the
pretext of not wishing to give offense to the simple, one keeps silent
about the
% --- p. 183 ---
\opage{183}\setcounter{footnote}{0}%
objections, veils the contradictions, and keeps the whole in so foggy an
indeterminacy that those who think for themselves, who would gladly be
enlightened about how the matter hangs together, can never obtain any
enlightenment. We hear the Savior's words (Matt.\ XVIII, 7): ``Woe unto
the world because of offenses! for it must needs be that offenses come;
but woe to that man by whom the offense cometh!''

Whence comes the offense when the talk is of \emph{the speaking God}? Not
from Luther at Worms with the ``clear reasons'' and the living cognition
of faith, no, not from Luther, but from Lutherolatry, i.e.\ from the
Lutheranism stiffened and crystallized in the theory of the schools; from
bibliolatry, i.e.\ from biblical theology's scientific botching of the
doctrine of inspiration. The Reformers had, after all, carried this
through: the papal Catholic externality had at last been mastered; but in
compensation for the loss of the infallibility of the visible church,
Protestant theology then clung again so convulsively to the biblical
externality that it declared the whole of Holy Scripture to have been
dictated word for word by God himself (Solus Deus, si accurate loqui
velimus, Sacræ Scripturæ autor dicendus est, Prophetæ vero et Apostoli
autores dici non possunt, nisi per quandam catachresin, utpote qui potius
Dei autoris calami fuerint. Quenstedt). By making the authors of the
sacred books into the ``scribes'' of God and Christ (amanuenses, manus,
notarii), the ``pens'' of God and Christ (calami), by letting the Holy
Spirit immediately drive these ``pens'' to write (impulsus ad scribendum),
inspire in them what they were to write (suggestio rerum), and how they
were to write it (suggestio verborum), the human word was made divine and
the divine human to such a degree that not even a glimmer of
% --- p. 184 ---
\opage{184}\setcounter{footnote}{0}%
difference remained. Here is the first ground of the errors, the soil in
which the seed of offense has sprouted so abundantly.

Orthodox theology conceives inspiration mechanically; the fundamental
contradiction is that it makes spirit, in its richest, purest revelation,
spiritless. The old orthodoxy and the modern consciousness meet in the
false presupposition that the divine and the human must coincide in
spirit. The difference is only this, that orthodoxy lets the human perish
in the divine, so that the latter takes possession of the determinate form
of the human word, while the modern consciousness lets the divine be so
absorbed in the human that the latter, in pure thinking, acquires an
abstract divine form. But the unity of the spirit of God and of man is a
unity in essential opposition, a unity of opposition. It is the spiritual
opposition between the divine and the human that is overlooked from both
sides.

We will, with the principle of faith in view, consider the Mosaic
lawgiving. The consciousness of culture can understand excellently that
Moses is the man of his time and of his people, that he is moved by a
great idea, that he has a mission. It can further understand that a man
so mightily moved by the idea must, when he is to give his people laws,
fit the lawgiving to the people's spiritual conditions; that the laws,
sprung from a single great fundamental thought, must branch out widely and
extend to all the relations of life; it can understand that the
fundamental thought, when it is, as in the present case, decisively
religious, must demand that the lawgiver speak, not in his own name, but
in the name of the Godhead, and speak with divine authority. We need only
paraphrase the opening sentence: ``And God spake all these words,'' thus:
``And Moses spake all these words in the name of Jehovah,'' and
% --- p. 185 ---
\opage{185}\setcounter{footnote}{0}%
the consciousness of culture, with its free historical overview, will
admire Moses and see in his lawgiving a work of spirit, one of the
greatest of antiquity. But further than this the consciousness of culture
cannot go. And why not? Because it cannot enter into the relation to God,
does not know how to think the thought of faith. Some features of the
psychology of faith will be able, at least provisionally, to shed light on
the matter.

If we refer the great fundamental idea of the people's liberation and of
the lawgiving, not to the human consciousness in Moses alone, but to
Jehovah's will, so that the whole originally proceeds from him: then the
relation between Jehovah and Moses is not an ordinary relation of idea but
a personal relation of spirit, in which the divine self enters into inner
reciprocal action with the human. On the conception of this reciprocal
action, then, the understanding depends. The strongest, most immediate
expression of such a reciprocal action is a mutual speaking. Speech
% the Danish punctuation after "Talen" is doubtful (full stop or comma); the sense is unaffected
is the immediate form of thought; the two communicate their thoughts to one
another by speaking with one another in one way or another. Jehovah speaks
to Moses, Moses in turn to Jehovah: now what kind of speaking is this? It
is in principle an inward speaking, an original communication of spirit to
spirit. There is an essential difference between inward, spiritual and
outward, sensuous speaking. When two individuals speak to each other, the
speaking is outward; they must understand each other's language in an
outward, finite way. Born compatriots therefore speak together most easily
and most naturally, a Hebrew with a Hebrew, a Greek with a Greek. Each of
those speaking employs the means of language in his own way. One can
resolve the speech into so many inferences, so many judgments, so many
sentences, so many words; one can in an outward, finite way make sure of
what and how much each of the speakers has said. It is otherwise with the
inward speaking, the
% --- p. 186 ---
\opage{186}\setcounter{footnote}{0}%
speaking in which the divine spirit communicates itself to the human. Here
it does not begin immediately with words and sentences; but it begins in
reflection with thoughts and judgments. The divine thoughts and judgments
first give themselves an inner, ideal, mobile form in awakenings,
inspirations, penetrating fundamental representations; the communication
of spirit may in this respect be said to pass through the idea.

But now it makes a difference in principle that the idea as idea is
impersonal, and to that extent selfless, while spirit as spirit is
personally a self. Since Jehovah's spirit is one with the almighty self,
the divine self-determination is expressly at work in the inner
communication. How the awakening inspirations and the strong fundamental
thoughts are more precisely to be determined, apprehended and appropriated
by the human spirit comes to depend, not, as with the mere inspiration of
the idea, on the apprehending consciousness alone, but originally and
decisively on the divine will. The next thing we must emphasize is this,
that he within whom spirit has begun so supernatural an activity
translates the inspirations, the primitive proclamations, into his own
human language; in this the communication of spirit again resembles a
natural, e.g.\ poetic, movement of the idea. And yet there is here,
according to the principle, a decisive difference. He who is moved by the
idea has a content of consciousness that demands and receives its form
from his own consciousness; the language in which the idea is uttered the
poet must expressly acknowledge as his language: the idea does not speak
to him, but only through him, as he speaks out of the idea.

With him who is moved by spirit, God's spirit, the matter stands in
another way. Under the influence of the absolute self, the content
apprehended in the understanding and translated into human language takes
on
% --- p. 187 ---
\opage{187}\setcounter{footnote}{0}%
such a stamp of divine validity that it expressly takes the language with
its words and sentences away from man, makes human language its language,
words and thoughts its judgments. In inner spiritual opposition to what
thus presents itself with the stamp of divine superiority stands the
receiving consciousness with its purely human content, its finite
representations, thoughts and feelings, as a being on its own, a powerless
spirit that is addressed by the Almighty.

If we give up the principle of faith, the immediacy corresponding to such a
redoubling of spirit at once becomes fantastic; in order not to get two
independent spirits in one soul, we then transform the divine self into a
figment of the imagination, and interpret the redoubling mythically. By
virtue of the principle, on the other hand, the redoubling is so far from
being mythical that it rather furnishes the absolute presupposition for
clear personal self-cognition. If
it stands fast that an all-round, complete and sufficiently grounded
cognition of the relative self is to be attained only through a
penetrating cognition of its spiritual relation to the absolute, i.e.\ of
the relation to God: then the redoubling set forth is of course justified
by the same token.

That the relation of inwardness described here must on God's side be a
relation of revelation, on man's side a relation of faith, is clear from
what has gone before. The relation of inwardness reaches its highest
tension in inspiration; only on the condition of inspiration does God
speak personally to man, and man to God. But now the inner must in turn
have its clear, immediate expression in a corresponding outer; the greater
the power of inwardness, the stronger too is the immediacy with which it
breaks out in the outer. Jehovah does not speak
% --- p. 188 ---
\opage{188}\setcounter{footnote}{0}%
immediately to the whole people of Israel, for the people in its
naturalness is not inspired; but he speaks to Moses, for Moses is his
prophet, and the prophet is inspired. From the inner the outer is
understood; the inner reciprocal action in spirit acquires, as revelation,
full actuality, since its essence comes to light in corresponding
phenomena. The inward speaking now becomes an outward speaking; the
outwardness is so strong that to a superficial consideration it looks as
if it were merely a matter of sensuous seeing and hearing. It says
(Exodus XIX) that ``Moses went up unto God,'' and that ``Jehovah called
unto him out of the mountain~.~.~. And Moses came and called for the
elders of the people, and laid before their faces all these words which
Jehovah commanded him. And all the people answered together, and said, All
that Jehovah hath spoken we will do. And Moses returned the words of the
people unto Jehovah. And Jehovah said unto Moses, Lo, I come unto thee in
a thick cloud, that the people may hear when I speak with thee, and
believe thee for ever.'' In a more outward and sensuous way a phenomenon
of revelation can hardly be described; and yet, when we look more closely,
it will turn out that all this outwardness is the outbreak of something
inward, that the phenomenon of revelation, in all its actuality, is a
phenomenon of spirit, a phenomenon that presupposes faith and has
significance only for believers.

If we begin with the most outward, the grandiose spectacle of nature, the
storm with lightning and thunder, then its immediate effect is simply this,
that it strikes terror into the sensuous man. In its natural crudeness the
multitude is sensuous; the people gathered about the mountain are
involuntarily struck with terror and dread. But influenced by the word
that Moses speaks, spirit is awakened in the people; they are attuned in
faith to a religious apprehension of the natural phenomenon; they perceive
that Jehovah is in the
% --- p. 189 ---
\opage{189}\setcounter{footnote}{0}%
darkness of the cloud, and are seized with holy awe. Further than this the
people do not come, and cannot come. The outward arrangement by which they
are (XIX, 12---13) kept away from the top of the mountain corresponds
entirely to their spiritual condition; they are overwhelmed by the
sensuous, divine the holy, and are overcome with horror. ``And all the
people saw the thunderings, and the lightnings, and the noise of the
trumpet, and the mountain smoking (XX, 18): and when the people saw it,
they stood afar off.'' In their dismay the people can now understand that
they need a spokesman, a mediator endowed with spirit between themselves
and God. ``And they said unto Moses (XX, 19), Speak thou with us, and we
will hear: but let not God speak with us, lest we die.''

If in thought we follow Moses, who goes up ``unto God'' on the mountain,
and ask: what would a knower, e.g.\ a bold naturalist, if he had at the
same time made an ascent of Sinai, have
seen and heard? A storm, a quite natural but terrifying storm; nothing
else! That the storm here is a sign of omnipotence, that the Almighty
himself is present in the storm, cannot be seen with the eye of natural
science. ``But Jehovah speaks aloud, does he not; he calls to Moses in the
storm?'' Yes, he calls to Moses; but this call cannot be heard by the ear
of natural science! As Moses ascends the mountain and ``draws near unto the
thick darkness where God is,'' he is not objectively calm like a
naturalist, nor overcome with terror like the crude multitude; but he is
moved in spirit and confident in faith: if his faith wavered and his spirit
grew weak, he would collapse in dismay. It is in the darkness of the storm,
under the immediate impression of the personal presence of the Almighty,
hence in the very highest tension of the crisis of the soul, that the inner
redoubling breaks out in mighty immediacy, that the inward speaking takes
on the stamp of an outward one, that Jehovah
% --- p. 190 ---
\opage{190}\setcounter{footnote}{0}%
calls to him from the cloud. ``And God spake all these words,
saying''~.~.~.

In the Mosaic lawgiving with its many commandments the divine so
coincides with the human that the difference and the opposition are
hidden. This one-sidedness characterizes religion at the Old Testament
stage and shows it to us as a preparatory, not a consummating, revelation.
What we can say about this in all brevity is the following. Had the
redoubling of spirit corresponding to the inner revelation been carried
through completely, then Moses would have had to be either Christ himself
or one of his apostles.

This may perhaps become intelligible if we consider somewhat more closely
the opposition, corresponding to the redoubling, between the absolute and
the relative self in the relation to God. In Christ the relative self is
evangelically determined by this, that ``the Son of Man'' is ``the Son of
God'' and God his ``heavenly Father.'' The relative self is thus, in this
its relation to God, itself divine, is the God revealed on earth himself.
The Son's unity with the Father is expressed in the words (John X, 30):
\textit{I and the Father are one}; whereas the Son's difference from and
dependence on the Father---hence the relativity of the self in the Son of
Man---is expressed in the words: \textit{The Father is greater than I.}

From this now follows a decisive change in the relation between the
outward and the inward in the redoubling of spirit. Christ's revelation is
outward in quite another sense than the revelation through Moses; for
Christ appears as the God revealed in flesh and blood, as the Word itself,
which was from the beginning, whereas Moses on the other hand appears as a
man, as God's spokesman for the people of Israel. But just for this reason
the revelation in Christ is inward in quite another sense than the
revelation through Moses. What
% --- p. 191 ---
\opage{191}\setcounter{footnote}{0}%
Christ thinks and speaks as God, that he thinks and speaks at the same
time as man; in Christ's inner being the finite separation between the
thoughts of God and human thoughts is sublated in personal unity. When
Christ speaks, he speaks in his own name. The Old Testament expression:
``Thus saith Jehovah, the God of Israel,'' is replaced in the New
Testament by the evangelical: ``Verily, verily, I say unto you.'' To a
superficial consideration it now looks as if the Old Testament expression
were far more exalted and divine than the New Testament one; but the
matter stands the other way round. The Old Testament expression is a
divine sanction of the prophetic speech; but the speech itself in its form
and its finite content is a human speech. Moses and the prophets speak
with divine authority in the name of Jehovah; but they speak in a popular,
human way about finite, temporal, human concerns. The Mosaic law with its
many commandments is given by Moses, and Moses has taken every possible
finite consideration of his people's conditions, its destiny, its
upbringing; but this human element in the Mosaic lawgiving is, be it
noted, not so separated from the divine fundamental thoughts and counsels
that Moses himself arbitrarily devised them and then in an outward way
added Jehovah's name. On the contrary! it is in spirit and in truth divine
thoughts that are here humanly thought, divine commandments that are here
expressed in human words. Over against these thoughts of Jehovah, sealed
with divine power and authority, Moses has in his consciousness a circle
of other thoughts, such as he must in a finite sense call his own
thoughts. We must further emphasize, in consequence of the principle, that
the separation a Moses makes between his own thoughts and the thoughts of
Jehovah revealed to him is
% --- p. 192 ---
\opage{192}\setcounter{footnote}{0}%
expressly determined by his faith and by the inspirations corresponding to
faith.

As for the content of the law, finally, it will easily be understood that
the so-called arbitrariness of the commandments is no proof against their
moral worth. Every commandment, even that whose necessity is most clearly
evident to reason, must as a commandment, i.e.\ an order, take on the
stamp of arbitrariness. The idea of morality is in a decisive sense the
idea of the will; the idea has absolute validity only by being posited as
one with the absolute will. But from this it follows in turn that all the
original demands of morality must be expressed in the form of divine
commandments. The moral proposition: ``God wills it because it is good,''
is so far from excluding the religious proposition: ``it is good because
God wills it,'' that the latter is on the contrary necessary to secure the
former. Since in all determinations of morality it depends essentially on
will, it depends in an absolute sense first and last on the divine will.

What has been said here must suffice to characterize the problem of faith
that we can hold fast under the name of the religious problem of
arbitrariness, a problem which for its complete solution will require a
thoroughly worked-out psychology of faith, philosophy of religion, and
ethics.

Of the two propositions set out above, the latter: ``it is good because
God wills it,'' points to discipline; the former: ``God wills it because
it is good,'' to the self-development corresponding to discipline; the one
is the religious proposition of constraint, the other the religious
proposition of freedom; the one corresponds to the spirit of the law, the
other to the spirit of the Gospel. It is then in order that the
proposition: ``it is good because God wills it,'' is one-sidedly put into
force at the Old Testament stage. With this fundamental thought one will
understand why the Mosaic lawgiving with its many divine
% --- p. 193 ---
\opage{193}\setcounter{footnote}{0}%
commandments must go so fully into details. ``Thou shalt'' and ``thou shalt
not'' repeat themselves in such manifold ramifications that all the
concerns of life are, immediately and mediately, drawn in under them.
Everything, not only what belongs to the worship of God with its
sacrifices and ceremonies, but also what concerns domestic and civil
relations, everything is punctiliously determined by Jehovah's will. At
every step the Israelite takes on the road of actuality he meets a divine
commandment; by means of the divine commandment the temporal becomes, link
by link, an account to be settled between him and the Eternal. In
opposition to the priests, who were to watch over the law and the
fulfillment of the commandments, the prophets came forward as spokesmen
for the right of spirit; but the priests outlived the prophets, and
holiness by works, with Pharisaic righteousness, gained the upper hand:
with this the limit of disciplining arbitrariness was set; the time of the
law was over. Through the Gospel ``the law of commandments'' was
transformed into the law of love. With the law of love as its foundation,
self-development could be united with free obedience, and the proposition:
``God wills it because it is good,'' could explain arbitrariness by divine
wisdom.

It seems that people are bent, from every side, on deducing the will from
pure, universal thinking; indeed, it is even demanded that logic should
begin with ``the good will''; that is how far things have come. This is
perhaps the most telling proof yet that the last glimmer of the thought of
faith is about to disappear. That the present age cannot grasp the word of
revelation about ``the speaking God'' and reconcile human thought with the
divine arbitrariness is an essential hindrance to spiritual life, a
hindrance that only the spirit of the will can transform into a condition.

\begin{center}\rule{3cm}{0.4pt}\end{center}

\chapter*{XII.}
\addcontentsline{toc}{chapter}{Lecture XII}
\markboth{Lecture XII}{Lecture XII}
\label{ch:xii}
% --- p. 194 ---
\opage{194}\setcounter{footnote}{0}%
In the second Epistle to Timothy (III, 16---17) a point of view is given
for the conception and interpretation of Holy Scripture and thereby of the
whole of sacred history. ``All scripture,'' it says, ``is given by
inspiration of God, and is profitable for doctrine, for reproof, for
correction, for instruction in righteousness: that the man of God may be
perfect, throughly furnished unto all good works.'' By this
evangelical-apostolic point of view a decisive separation between the
principle of faith and that of knowledge is sufficiently indicated.

Hitherto, however, people have not been willing to heed this separation;
that is to say, the opposition has been paraphrased in such a way that its
terms yield only two sides instead of displaying two mutually opposed
spheres. The doctrine and conviction of which the Epistle speaks is
infinitely different from the doctrine and conviction treated of in
science. It has indeed been pointed out that the expressions ``doctrine
and conviction'' point to the theoretical, while the expressions ``for
correction, for instruction in righteousness'' refer to the practical. But
on the strength of this merely formal difference people have then reasoned
thus: Holy Scripture demands doctrine, thorough doctrine, conviction, true
and living
% the Danish "Overbeviisning" (conviction) renders the Epistle's "reproof"; kept as "conviction" in the discussion that follows
% --- p. 195 ---
\opage{195}\setcounter{footnote}{0}%
conviction; but where should we seek the thorough doctrine, the true and
living conviction, if not in science alone, the science that presupposes
faith, explains and grounds faith, namely theology? Yet we are not to stop
at science alone, we are, after all, also to work for life; to the theory
must correspond a practice; in practice we have the pastoral, the
edifying; it is this that is designated by the words ``for correction, for
instruction in righteousness.'' It has been added that just as the old
orthodoxy one-sidedly kept to doctrine, so pietism has one-sidedly kept to
the practical, ``the edifying.'' ``But,'' it is said, ``we enlightened
theologians keep the balance; we have `doctrine' and `conviction' in the
objective and `correction and instruction' in the subjective; what more
can then be demanded?''

Against this misunderstanding we set up the separation between two
absolutely heterogeneous standpoints, the scientific and the religious, of
which the first conditions the cognition of impersonal truth, the latter
that of personal truth.

When we call scientific truth impersonal, the meaning is of course by no
means that scientific thinking could take place otherwise than in a
thinking consciousness, that there could be a knowledge without a knower, a
conviction without a subjectivity that felt itself convinced; but the
meaning is that the knower, in the development of his science, must give
up every personal regard for himself, must, absorbed in the cognition of
the matter, of things, laws, facts, not let himself be disturbed by any
regard whatsoever for any personal interest whatsoever. A physiologist, for
example, may as a man, as an existing personality, be infinitely
interested in the question of the immortality of the soul; but if he
without further ado lets this interest play a
% --- p. 196 ---
\opage{196}\setcounter{footnote}{0}%
part in purely scientific investigations, he becomes a mediocre
physiologist. An astronomer may as a man, as an existing personality, feel
the need to make sure whether there is in reality a world of a higher order
than this temporal one in which we live; but if, led by this need, he
turns his telescope toward the Milky Way and begins to resolve the nebulae
in favor of the doctrine of a ``kingdom of heaven,'' he will hardly keep
his standing among astronomers for very long. A geologist may as a man, as
an existing personality, feel so strongly attracted, so deeply moved by the
biblical narrative of Paradise that he could truly wish, were it possible,
to find a point in the earth's development of which he could say: here is
the trace of the paradisal state; here, on this spot, the Garden of Eden
lay! But if he lets such longings and endeavors gain influence on his
judgment of the geological facts: then he shows himself a most unreliable
investigator. We hope to be understood, not only by the cultivators of
science, but by all people of culture, when in the sense given here we
repeatedly declare that scientific truth is and must be impersonal.

That we, in opposition to the scientific, call religious truth personal
must not be understood as if the meaning were that the religious cognition
of truth should rest only on individual representations, feelings, moods
and such indeterminable presuppositions; no, the meaning is that selfhood,
absolutely concentrated in the ideal demands of personality, designates the
standpoint. What satisfies the absolute demand of personality is valid;
what conflicts with it is invalid; what neither furthers it nor conflicts
with it is indifferent.

% --- p. 197 ---
\opage{197}\setcounter{footnote}{0}%
To clear up what this involves, we need not look about us for new and
complicated examples; on our way through sacred history in what has gone
before we have already come upon so many knots, so many even hard
stumbling blocks, that it might well be to the purpose to take up the
examples already cited from a new point of view, to see whether a new
light might then perhaps fall on what was earlier enigmatic. Thus we saw in
the doctrine of creation that the astronomical account of the structure of
the universe is in decisive conflict with the account given us in the first
book of Moses; we saw how the knot was tied especially by Christ's words in
Matthew (XXIV, 29), where he foretells the destruction of the world and
expresses himself thus: ``Immediately after the tribulation of those days
shall the sun be darkened, and the moon shall not give her light, and the
stars shall fall from heaven.'' With an expression such as this, that the
stars shall one day fall from heaven, the conflict between science and
religion becomes so hard that it seems as though it could never be
removed, as though a settlement were impossible. If we say: the stars can
fall down from heaven, we contradict astronomy on a decisive point and
thereby break the staff over culture and science. If on the other hand we
say that that word of Christ cannot be fulfilled, then we violate Christ's
majesty, violate religion, and cast contempt on the revealed word. On these
terms every way out toward a reconciliation of faith and knowledge seems
barred.

However, the barrier is more apparent than real. This heaven and these
stars are here seen from wholly different standpoints, with quite different
eyes; therefore they are spoken of in different languages. It is, after
all, known to everyone that what is visible in the world must be seen by
an eye. As the eye is that sees, so is
% --- p. 198 ---
\opage{198}\setcounter{footnote}{0}%
the visual image, so is the visible world; as the ear is that hears, so is
the hearing; as the consciousness is that contemplates, so is the image of
what is contemplated. How different the image of heaven and earth must
become as the standpoint of the beholder changes we can impress upon
ourselves particularly in this connection by distinguishing three
standpoints, namely the scientific, the aesthetic, and the absolutely
religious.

If we place ourselves at the scientific standpoint, what do we then
contemplate as the first thing? The things themselves, their laws and
connection! It is the outward heaven, these determinate stars, these
phenomena of motion with the corresponding laws, that are the object of
contemplation. Does the one contemplating then contemplate himself at the
same time? No, for him who looks scientifically at the starry heaven, the
starry heaven is all in all. During his observations the astronomer does
not think at all of his personality, not of his own wishes and longings.
Such longings he would, if someone were so sentimental as to remind him of
the heart's need, repel with icy coldness. That is why Minerva's breast is
so cold, and her gaze so stern;
science is called objective because the beholder, absorbed in his object,
is to forget his subjectivity. Astronomy is, like every other science,
impersonal, and in its impersonal essence objective, indifferent to the
stirrings of our souls, our deepest longings, our pains and joys, our fear
and our hope, as indifferent as if it were pure fate.

What personal results, then, can be derived from the contemplation of the
structure of the universe? The result of results is that man shrinks to a
nothing. In relation to the earth we become far smaller than the ants in
their hill; the earth itself disappears in the universe. Refreshment and
liberation must therefore be sought in this, that we, instead
% --- p. 199 ---
\opage{199}\setcounter{footnote}{0}%
of holding fast to our own self, lose ourselves in the selfless, that we,
instead of being independent within our personal limitation, free
ourselves in the limitless. It is remarkable how this absorption in the
outward infinity of celestial space, this sublime loss of self, reproduces
an oriental view that belongs already to antiquity; it is Brahmanism we
have in mind.

As is well known, the Orientals are inclined, in profound fantasies, to
draw a line through the finite and to absorb themselves in the vision of
the eternal, the limitless, the infinite. Brahma's priests regard this
world in its finitude as a beautiful deception. Its outward multiplicity,
its many different colors and tones signify that a veil of time has been
spread over the eternal truth; it is the veil of Maya, the glittering veil
of illusion, of deception. It is a task for the wise man to lift this veil.
First he removes the deception of the senses by denying reality to sensuous
things; but this is not yet enough: he must also free himself from the
deception of thought. The deception of thought lasts as long as he
distinguishes determinate thoughts, the one limited by the other. So he
then lets all thoughts flow together into one thought, the thought of
infinity; then he enjoys, in his dissolution of self, the bliss of
contemplation, loses himself in the limitless, and disappears in Brahm.
Here is a liberation that in an inward sense corresponds exactly to the
liberation that astronomers and physicists recommend to us in an outward
sense, when they recommend us to deny what they call our egoism and to
disappear in the infinite universe. This, then, was the first standpoint.

Different from the strictly scientific, objective standpoint is the
aesthetic. At the aesthetic stage personality does not, to be sure, give
itself up, but in the moments of mood it cancels its actual
% --- p. 200 ---
\opage{200}\setcounter{footnote}{0}%
limitation; it expands fantastically and in its self-expansion embraces
heaven and earth. What is characteristic is that the human spirit sees in
the sky ideal figures that once lived and acted on earth; thereby the
constellations known to us acquire a quite peculiar significance. We know,
of course, such constellations in the visible heaven: Cepheus, Cassiopeia,
Andromeda, who is chained to the rock, Perseus, the horse Pegasus, the
kneeling Heracles, the Lernaean serpent, etc. These names remind us of the
mythical heroes and their deeds. What does this signify? It signifies
that as man looks up at the starry heaven, he sees a mirror that reflects
the images of personality, a reflection of personal life through the ages.
``The constellations,'' says J. L. Heiberg (Intelligensblade, vol.~2,
p.~41), ``are true pictures, which antiquity has hung up in the great
gallery where all generations are to contemplate them. Here it becomes
quite evident that the starry heaven, far from being foreign to the earth,
is the reflection of earthly destinies, earthly thoughts, earthly dreams.
One must regard the starry heaven as an illustrated deluxe edition of Greek
mythology in enormous format; one must contemplate it with the poets of
antiquity before one's eyes.'' It is not the single individual who mirrors
himself; it is antiquity and its great figures that he sees in the mirror.
% the Danish lacks the closing quotation mark after "Stjernebillederne" (printer's defect); supplied here

Here, then, contemplation is redoubled and the standpoint transformed; for
when one is to speak of the figures that are seen in the stars, one must
speak not merely from the standpoint of poetry, but in the language of
poetry. We surely all know Heiberg's lovely poem, in which among other
things he has the starry heaven say to the beholder: ``Here the kings and
queens of fable stand, with its monsters, dragons and serpents, whose
names resounded from the gray of antiquity in songs two thousand years
old.'' Here the cold
% --- p. 201 ---
\opage{201}\setcounter{footnote}{0}%
system of phenomena and laws has vanished; when the poet speaks of the
kings of fable, its monsters and serpents, he speaks precisely in his own
language, for the language must correspond to the standpoint. In this
language of mood memory is invoked. Man knows himself to be limited, not
only in relation to outward greatness in space, but also in relation to
time. He lives a brief time; his life is ``as the grass, as the flower of
the field; the sun is risen with a burning heat, the flower withereth and
falleth.'' We live the short life of a day; if personality is to expand
and breathe freely in the kingdom of spirit, it must take refuge in memory.
In memory it turns toward what has vanished, wanders along the paths of the
past, meets with the great ones of antiquity, and understands the heaven
that speaks in the poem: ``With me thou livest a life that is gone, thou
livest with gods and heroes, whose transfigured forms legend bound to my
broad and sparkling belt.''

No one will say that this standpoint of poetry comes into conflict with
astronomy. Here there is no conflict; for the two different standpoints
are wholly heterogeneous. The astronomer would only make himself ridiculous
if in his prosaic language he began to argue against what the poet here
says about the monsters, dragons and serpents of fable, while conversely
the poet would make himself ridiculous if he tried to refute the astronomer
with the kings and queens of fable. There are two standpoints, two
contemplations, two languages, and two different images of the world
contemplated. Each of them has its validity. We will not do without
astronomy, but not without poetry either.

This anyone can understand; but should it not then also become
intelligible when we in a corresponding way take note of the third
standpoint, the religious, not the relatively religious, which is weakened
in human
% --- p. 202 ---
\opage{202}\setcounter{footnote}{0}%
knowledge; no, the absolutely religious, Christ's own standpoint?

If Christ did not stand at a standpoint of his own, he could not, after
all, be Christ; but if he stands at a standpoint which no mortal can
occupy with him, this must surely entail that to the peculiar standpoint of
contemplation there corresponds a distinct image of heaven and earth. What
kind of standpoint this is has already been said earlier: it is that of
personality! The principle of religion is the principle of personality in
the very deepest sense.

It is, of course, characteristic of modern times to have respect for laws,
``the eternal laws''; when anyone wants to contradict the physicist, he
simply points to the eternal laws, and the objections fall silent. Now when
we hear Christ and see what lies in his words, he too points to eternal and
unchangeable laws; but they are, be it noted, not the laws of necessity but
the laws of freedom. Here are two systems of laws, because here, according
to the conditions of human cognition, there must be two different
standpoints. Here are the laws of necessity, corresponding to the kingdom
of nature; but here are also the laws of freedom, of spirit and of
personality, which hold in the kingdom of Christ. It is the laws of freedom
and of personality of which Christ speaks when he says: ``Heaven and earth
shall pass away, but my words shall not pass away.'' What is
characteristic of this standpoint is that it is personal truth itself that
speaks of heaven and earth, and speaks the language of personality,
authoritatively decisive. The one speaking is not impressed by the outward
greatness of the structure of the universe; for the outward world, its sun
and moon and stars, are subject to the laws of necessity, but truth itself
knows with eternal, divine certainty that the laws of necessity are
% --- p. 203 ---
\opage{203}\setcounter{footnote}{0}%
grounded on the laws of freedom. While the natural man must feel himself
shrink to a nothing over against the might and majesty of nature, Christ,
as the absolute representative of personality, stands unshakably firm
on the ground of freedom. Had he let himself be overwhelmed by strong
impressions of nature, by the outward, infinite greatness of the structure
of the universe, he would not have been Christ.

One easily sees, however, that the opposition between the religious and the
scientific view of the structure of the universe is quite otherwise hard
than the opposition between the scientific and the aesthetic. In the
aesthetic view reality is not, after all, taken in full earnest; poetry
does not maintain that ``the kings and queens of fable'' really stand on
the vault of heaven: the aesthetic and the scientific standpoints are, in
other words, only relatively heterogeneous. In the religious contemplation
of nature, on the other hand, reality is taken in just as decisive earnest
as in the scientific. If we ask whether it can be assumed that heaven and
earth will really one day pass away, as Christ has foretold it to us: then
science, reflecting on what is its own, answers with an express No,
religion, reflecting on what is its own, on the other hand with an express
Yes. From the standpoint of knowledge Christ's prediction cannot be true;
and why not? Because its fulfillment is an impossibility! From the
standpoint of faith Christ's prediction must absolutely be true; why?
Because Christ is the truth! So science stands against religion as Yes
against No, as truth against truth: will anyone be able to deny that here
is a dualism?

Those who are bent on undermining positive revealed religion can naturally
take the dualism lightly; they can at their pleasure now cancel it by
sluggish
% --- p. 204 ---
\opage{204}\setcounter{footnote}{0}%
abstractions, now again amuse themselves by making it an object of
caricaturing ridicule and sniggering distortion. But those who want to
hold on to religion without giving up science, the theologians, for
example, what can they be thinking of? Can they really make themselves and
others believe that their ``good will,'' if only it is allowed to come
along into science, will be able to master the dualism?

The dualism is not an arbitrary invention, not a philosophical chicane; it
is a fact, a notorious, incontrovertible fact. For coming to terms with the
dualism, or, as it is commonly put, reconciling faith with knowledge, there
is only one way out, one single one. We shall shed a little more light on
the matter with regard to the example we have set up. Natural science's
conception of the structure of the universe is not religious; Christ's view
of the structure of the universe is absolutely religious: hence the
contradiction. This contradiction can be resolved only by Christ's being
granted unconditional right in the name of personal truth, and science,
with its impersonal truth, scientific right.

In order to grant Christ unconditional right, we must seek to understand
his fundamental thought, the language in which the fundamental thought is
uttered, and the reality corresponding to the utterance.

What Christ utters with absolute personal certainty is this, that the
kingdom of nature with all outward, sensuous things is subject to the
omnipotence of spirit, that the kingdom of spirit, the kingdom whose
founder he is on earth, shall, when the hour comes, be actualized, not only
in the inward, but also in the outward. This certainty he does not have
from any immediate contemplation of nature, any more than from astronomical
observations of sun, moon and stars; no, he has the assurance, the full,
decisive assurance,
% --- p. 205 ---
\opage{205}\setcounter{footnote}{0}%
in himself, through himself; he has it personally. Faith in Christ's words
is on this point too one with faith in Christ himself; in this too he shows
himself as the exemplar, as the actualized ideal of personality, that he
can speak with so infinite a superiority of the power of freedom over
nature bound in necessity.

The frail individualities, the relative personalities entangled in guilt
and self-contradiction, are, in spite of all their striving toward the goal
of spirit, constantly exposed to being overwhelmed by natural hindrances and
cowed by overpowering impressions of sense. It is especially in dark, heavy
moments that the pressure is felt. Faintheartedness comes with doubt; and
whence comes doubt? From a discord within man! And whence comes the
discord? From a conflict between the two absolutely heterogeneous
principles, the principle of necessity, with the thousand proofs from the
reality of the sensuous, material world, and the principle of freedom with
its demands and ideal promptings from the soul's self. Here an
authoritative word is needed, not only for ``doctrine,'' but also for
``edification.''

In Christ's words there is doctrine, and the doctrine is this: If thou
wilt have full assurance of the unconditional power of spirit over
corruptible nature, then lay hold of the principle of freedom and beware of
staring at necessity! in other words: if thou wilt find the way to the
actual kingdom of freedom, of truth and of personality, then do not stray
into the regions of impersonal cognitions, do not seek guidance from
astronomers and physicists and philosophers of nature, but turn to him who
is the way and the truth and the life! In Christ's words there is
``edification.'' What is edifying is that he utters the assurance, not as a
general view of life in a mood awakened and stirred by happy omens and
bright prospects, but in the heavy, terrible moment, in a moment when
% --- p. 206 ---
\opage{206}\setcounter{footnote}{0}%
he had to say to his enemies: this is your hour, and the power of
darkness; it is just such a moment that gives the assurance the inner
weight of personality which counterbalances outward reality.

As the fundamental thought is, so must the language be too. What Christ
sees in the spirit, the mighty visions, the great images in which what is
to come acquires a spiritually immediate presence, he utters in words that
even to this day can make the vision living. His discourses to the
disciples on the destruction of Jerusalem and the end of the world and his
testimony before the high priests in the council are carried on in the
same language. ``And the high priest .... said unto him, I adjure thee by
the living God, that thou tell us whether thou be the Christ, the Son of
the living God. Jesus saith unto him, Thou hast said: nevertheless I say
unto you, Hereafter shall ye see the Son of man sitting on the right hand
of power, and coming in the clouds of heaven'' (Matt.\ XXVI, 63---64). No
faith is needed, but merely sound common sense, to see how absurd it would
be to subject Christ's words to a critique based on meteorological
investigations. Anyone understands that the language Christ speaks is
absolutely heterogeneous with the language natural science speaks.

``But,'' it is objected---and with this we come to the question of
actualization---``Christ's discourse does, after all, point to a real
impending natural event; that is precisely the knot, that if his
prediction is really to be fulfilled, it must come into decisive conflict
with the laws of nature whose unchangeableness science presupposes.'' We
answer that if such a conflict were really to take place, then the reality
of which Christ speaks would have to be a simple outward physical reality;
but the phenomenon of revelation that he describes is a phenomenon of
spirit; its
% --- p. 207 ---
\opage{207}\setcounter{footnote}{0}%
reality is a reality for spirit. ``This reality for spirit, then, has
nothing to do with the physical?'' Yes, it certainly has to do with the
material reality of nature; but the manner in which the actualization in
the physical sense is to take place is not described physically. The
fulfillment of the prophecy presupposes a miracle. How much of this natural
world is to be reshaped and transformed by such a miracle, at what point in
time the transformation, seen temporally, will set in, and what
intermediate stages of reshaping the visible must first run through---of
this the discourse contains nothing.

If on the other hand we place ourselves at the standpoint of natural
science and hold fast to the principle of necessity, then the actual
universe does indeed take on a different appearance, but then the cognition
of truth also becomes of a different kind. What we are to seek along this
road is not a satisfaction of the ideal demands of personality, but the
real connection of things: here truth is impersonal, here no miracles are
held out in prospect. Indeed, just as surely as the truth that expresses
itself in the Pythagorean theorem is impersonal, just so surely is the
whole system of truths that expresses itself in a Laplace's
\textit{Mécanique céléste} a system of impersonal truths. The personal and
the impersonal truths are not merely somewhat different, but absolutely
heterogeneous. The dualism is here; it stands in force. Those who
absolutely want to reject the revealed word and sacrifice their personality
to the gray absolute may well cancel it, but veil it they cannot.

To veil the dualism determined by faith and knowledge, and to beguile the
half-educated with dogmatic phrases and metaphysical ambiguities, is to
hinder and botch spiritual life in the present age.

\begin{center}\rule{3cm}{0.4pt}\end{center}

\chapter*{XIII.}
\addcontentsline{toc}{chapter}{Lecture XIII}
\markboth{Lecture XIII}{Lecture XIII}
\label{ch:xiii}
% --- p. 208 ---
\opage{208}\setcounter{footnote}{0}%
The occurrences in our personal life, what meets us on our way amid
changing surroundings, we are accustomed to designate by different names
according to the different impressions we receive from them. Thus we
speak of happenstance, accident, good fortune, misfortune, the blows of
fate and the trials of Providence; but are these now also different
things in the matter itself? Are there in actuality some occurrences
that must simply be ascribed to fate, and alongside these again others
that are immediately appointed to us by Providence? Have fate and
Providence divided the dominion between them? This question is as
personal as possible.

The personal wants to be understood personally; it wants to be
interpreted and applied in accordance with the peculiar way of viewing
things of the person concerned: only through the understanding and the
application does an event acquire its significance for life. Therefore
one of two things: either everything is under fate, and there is no
Providence, or all things without exception are under Providence, and
there is no fate. Which of the two ways of viewing things is in itself
the right and true one cannot be decided in objective knowledge; it can
be determined only by a personal choice. He who holds to necessity
chooses fate; he who
% --- p. 209 ---
\opage{209}\setcounter{footnote}{0}%
with resolved mind grasps freedom chooses Providence. Here again stands
standpoint against standpoint, interpretation against interpretation; we
shall now see how matters stand with this opposition.

To place oneself at the standpoint of natural necessity and choose fate
is, in other words, to build one's view of life on objective knowledge
and to apply the impersonal truth as though it were personal. What this
means we see in many of the wise men of antiquity, especially in the
Stoics. For the noble Stoic, natural necessity is one with the necessity
of reason; relying on the necessity of reason, he has with firm will and
full consciousness placed himself under fate. His motto, ``Lead me, O
fate,'' thus means: Guide me, thou eternal necessity of reason, in thee
is wisdom, in thee is truth! and so: Lead me, O eternal truth! How, then,
is he led by fate, and how does he understand its leading? First he
makes a distinction between the essential and the indifferent. To this
distinction among the concerns of existence there corresponds in his own
consciousness a sharp and clear personal distinction between ``that
which depends on us'' and ``that which does not depend on us.'' What,
then, is it that does not depend on us? Whether we are to be sick or
healthy does not depend on us; it stands in the power of chance and
belongs to the indifferent. Whether we are to occupy a low or a high
rank in society, whether we are to belong to the weak or the strong, the
lowly or the mighty, does not depend on us; it hangs on chance and
belongs to the indifferent. Thus he continues to go through the finite
conditions of life, and levels the differences in the one phrase: ``It
is indifferent.''

That all this indifferent matter nevertheless acquires a kind of
validity lies solely and entirely in the significance it has for
character. Therefore Epictetus says: ``The world is like a
% --- p. 210 ---
\opage{210}\setcounter{footnote}{0}%
great stage; men are the various actors, and the Deity assigns each his
particular role.'' By this he by no means wants life to be conceived
either as a blind play of forces or as an arbitrary play of intrigue.
When Epictetus speaks of playing one's role well, he takes the word
role in another and deeper sense than an Emperor Augustus, who on his
deathbed says with an ironic smile to his intimates: ``Have I not played
my role well?'' Epictetus wants to impress the proposition that every
human being attains his destiny when with his powers and gifts he fills
the place on which he was set according to the order of things.
Therefore the wise man bows beneath necessity as beneath a holy power.
When you sit in the theater and watch a tragedy, your heart must not
beat uneasily, and you must not weep; for you shall not wish that things
went otherwise for the hero than they go. It goes as it must go, and
with that you shall be content. However dearly you love your wife and
son, impress upon yourself constantly that a precious earthen vessel can
be broken; if it is then broken, your misfortune will not surprise and
pain you. If you have in good time seriously put it to yourself that
your wife and son are mortal, then you will be able to see them die with
the calm that is worthy of fate and of reason.

There is in the Stoic resignation something grand and exalted. A Stoic
resigns on an altogether different scale than one otherwise resigns in
everyday life. Every reasonable person whom a misfortune strikes must,
to be sure, try with a certain composure to accommodate himself to the
unavoidable; but in finite resignation one renounces only some goods, in
quiet expectation of getting them replaced by others. The Stoic's
resignation, on the other hand, is infinite; he resigns once and for
all
% --- p. 211 ---
\opage{211}\setcounter{footnote}{0}%
all finite goods and gives up all temporal wishes, so that he shall not
become a ball for the whims of fortune. One would think that this was
the most exalted ethics there was in the world; and assuredly it is also
exalted characters that we meet among the Stoics of antiquity. Yet on
closer consideration it is easy to see that this moral greatness has
something strained, something inhuman about it, that in its innermost
nature a contradiction has hidden itself.

The Stoic is cold and strong; for he is the man of fate, and fate is
cold and strong. He needs no one, loves no one, hates no one, admires no
one, despises no one; he hopes in no one and fears no one. He cannot be
bent, he can only be crushed; to him apply the strong words:
\textit{Si fractus illabatur orbis, impavidum ferient ruinæ}. The Stoic
sets store exclusively by character.
% printer's defect: "au" printed for "an" (lægger an paa); translated in the corrected sense
But what is character other than pronounced personality? And what is
pronounced personality without a personal inwardness? But the principle
of the Stoic personality is impersonal: here is the contradiction. The
innermost of the Stoic's nature is not a self, but the cold, selfless
universal.

Stoicism affords a remarkable example of what comes of wanting to
assert the self on a selfless foundation. In the human sense, the Stoic
makes the greatest exertions to hold fast to himself, to be in agreement
with himself, to be sufficient to himself; but if we look more closely,
this self-preservation rests indissolubly on an essential absorption in
the selfless. The Stoic self is without a kernel of selfhood, is only a
point of transit for the world-reason, is against the power of fate
only semblance and shadow. Its highest is to will, not what a personal
self must reasonably will: to assert its original validity---no, but to
will what fate
% --- p. 212 ---
\opage{212}\setcounter{footnote}{0}%
wills, to make itself into a nothing. This becomes especially
recognizable where conflicts arise that bring the Stoic to seek
deliverance in death. Suicide is otherwise usually regarded as an
immoral act, indeed as a great crime. But from the Stoic point of view a
suicide motivated by the collision of circumstances and carried out with
calm deliberation can be regarded as the consummation of virtue. Cato at
Utica put himself to death---``after having read Plato's Phaedo and
slept calmly until midnight.'' Plutarch portrays him as a pattern of
virtue. It is the cold consequence of the doctrine of fate that the
self, after having exercised all its powers and played its role on the
motley stage of life, ends the role by blowing out the lamp and
vanishing into the selfless.

Let ``modern science'' exert itself as much as it will to paint up and
embellish the old doctrine of fate by extolling ``the Idea'' and ``the
Absolute'' and man's ``absolute relation to the Absolute'': it cannot
possibly get away from the Stoic self-contradiction. That a personal
self should have its ground, its essence, in the selfless is a desperate
point that does not bear the light of ethics, a metaphysical ambiguity
that makes the innermost of consciousness mythical.

The dialectic of self-contradiction that corresponds to this has been
strikingly demonstrated by S. Kierkegaard in his work
\emph{Sygdommen til Døden}, in which he among other things treats
(pp.~65 ff.) ``the despair of willing in despair to be oneself,'' and
(p.~67) characterizes ``the despairing self'' in the following manner:
``It acknowledges no power over itself; therefore it lacks in the last
resort earnestness, and can only conjure up a semblance of earnestness
when it itself bestows upon its experiments its very highest attention.
This is, however, a counterfeit earnestness; like the fire Prometheus
stole from the gods---so this is
% --- p. 213 ---
\opage{213}\setcounter{footnote}{0}%
to steal from God the thought, which is earnestness, that God is looking
at one, instead of which the despairing self contents itself with
looking at itself, which is now to bestow on his undertakings infinite
interest and significance, while it is precisely this that makes them
experiments.\,\ldots{} There is nothing firm in the whole dialectic
within which it acts; what the self is does not stand firm for a single
moment, that is, eternally firm.\,\ldots{} The self is its own master,
absolutely, as it is said, its own master, and precisely this is the
despair, but also what it regards as its pleasure, its enjoyment. Yet
on closer inspection one easily makes sure that this absolute ruler is
a king without a country; he actually rules over nothing; his
condition, his dominion, is subject to the dialectic that at every
moment rebellion is legitimacy. For this rests in the final analysis
arbitrarily on the self itself.''

From the standpoint of religion, all occurrences, all actual events, all
chances that meet us, insofar as they can acquire essential, decisive
significance for man's personal life, are without exception to be
conceived and determined as dispensations of the governing Providence.
This cognition does not rest on insight into the naturally necessary
connection of events; on the contrary! it rests on insight into the
essence of religion; it proceeds with inner consequence from the
principle of faith. Only by thinking this principle through will it be
possible to arrive at a penetrating understanding of the signs of
revelation in the Old and the New Testament. For what, after all, is a
sign of revelation according to its religious concept? It is a fact, an
event, which, by having stepped out of the immediate connection of
nature, expressly signifies a personal transaction between the God who
reveals himself and the one chosen in faith who receives the revelation,
and so a dispensation of
% --- p. 214 ---
\opage{214}\setcounter{footnote}{0}%
Providence. Thus the phenomenon of the burning bush is a personal
transaction between Jehovah and Moses, a dispensation for Moses. That
Moses had in general believed in the God of the fathers and his
governing Providence before the revelation occurred must be granted,
just as it must also be presupposed that in the solitude of the
wilderness he had pondered the deliverance of his people. But there is,
in the life of faith too, a great step from the universal to the
particular, from wish to fulfillment, from thought to action. At this
stage of transition the miracle enters as an express sign. Now the hour
has come when the great work of deliverance is to begin; now Moses is to
recognize that he and no other is the man whom Providence has chosen as
its instrument. ``Come now therefore,'' so sounds Jehovah's voice
(Exod.~3:10), ``and I will send thee unto Pharaoh, and bring forth my
people, the children of Israel, out of Egypt.''

But however definite the calling sounds, however unambiguous the
dispensation seems to be, it must nevertheless be conceived in faith and
appropriated in faith. One would suppose, according to an external
consideration, that Moses must immediately have seen the possibility
that, despite his human frailty, he would be able to accomplish what was
laid upon him, since Jehovah after all expressly promises to stand by
him. But in actuality it does not go so easily with grasping and
understanding the divine dispensations. It is indeed quite
characteristic that Moses, although he stands before the almighty
Providence and hears the promise of divine assistance in repeated
confirmations, nevertheless, with the dangers before his eyes, shows
himself disheartened and irresolute. Can there be any more telling proof
that the miracle, however outwardly it appears, is nevertheless
essentially a phenomenon of the spirit, an actuality for the spirit?
That the dispensations of Providence, even where they enter in the most
unexpected
% --- p. 215 ---
\opage{215}\setcounter{footnote}{0}%
manner as a help in the hour of need, even then, when they show
themselves as signs of omnipotence, must be conceived and interpreted in
faith---of this sacred history affords many quite striking examples.

In the second book of Moses (XVII) the following is told: ``And the
people thirsted there for water; and the people murmured against Moses,
and said, Wherefore is this that thou hast brought us up out of Egypt,
to kill me and my children and my cattle with thirst? And Moses cried
unto Jehovah, saying, What shall I do unto this people? they be almost
ready to stone me. And Jehovah said unto Moses, Go on before the people,
and take with thee of the elders of Israel; and thy rod, wherewith thou
smotest the river, take in thine hand, and go. Behold, I will stand
before thee there upon the rock in Horeb; and thou shalt smite the rock,
and there shall come water out of it, that the people may drink. And
Moses did so in the sight of the elders of Israel.''

If anyone would see in this a proof that the actuality of a miracle can
become knowable also for unbelievers, then he is in a great error. The
narrative, after all, proves precisely the opposite. The great crowd,
which despaired in its need, accepted the help, drank from the spring
(cf.\ 1~Cor.~X), quenched its thirst, was perhaps startled for a moment
at so unexpected a piece of good fortune, but here nothing stands about
its having received any impression of the miracle as miracle.

Another, still more telling testimony we have in the New Testament
(John~VI). After Jesus has fed 5000 men in the neighborhood of Tiberias,
he goes to Capernaum; ``(howbeit there came other boats from Tiberias
nigh unto the place where they did eat bread, after that the Lord had
given thanks:) when the people therefore saw that Jesus was not there,
neither his disciples, they also took shipping, and came to Capernaum,
seeking for Jesus. And
% --- p. 216 ---
\opage{216}\setcounter{footnote}{0}%
when they had found him on the other side of the sea, they said unto
him, Rabbi, when camest thou hither? Jesus answered them and said,
Verily, verily, I say unto you, Ye seek me, not because ye saw the
signs, but because ye did eat of the loaves, and were filled.'' To these
people, then, it is quite indifferent where the bread comes from;
whether it happens by a stroke of luck, a fabulous occurrence, by
sorcery or by a special arrangement of Providence, if only they can get
bread.

We learn from these examples not only that the sign as sign can be
cognized only by believers, but we learn as well that the believers who
see in the sign a dispensation of Providence must will to conceive the
dispensation in such a way that it can acquire its significance for
life: the appropriation and the understanding bring with them a personal
responsibility. This relation between the miracle and the dispensation
can in turn cast light on Abraham's conduct before Pharaoh and before
Abimelech.

The history of Abraham must, like most of the narratives in the Old and
the New Testament, be \emph{interpreted}; that is to say: there must be
added intermediate determinations, propositions and judgments, which
either immediately or mediately belong to determining the point of view
from which the narrative wants to be conceived. The history of Abraham
is a popular history; it is communicated in broad strokes; the
intermediate determinations recede into the background: the people's
gaze is fixed on the beauty of the ancestress of the tribe, the faith of
the patriarch and the grace of Jehovah; in the popular presentation the
religious is contained. As Jehovah's chosen one, Abraham is in an
altogether special sense a man of Providence, his way through the world
a way of Providence. But precisely as the man of Providence, Abraham is
responsible for how he interprets and uses the dispensations. So he
comes as a defenseless stranger to Egypt, to Pharaoh's court, and sees
the danger for himself and his wife. Had he now,
% --- p. 217 ---
\opage{217}\setcounter{footnote}{0}%
amid the threatening surroundings, been without further ado entitled to
rely on supernatural assistance; had he dared to be quite certain of
understanding the thought of Providence when he said to himself: ``I
have received the election; it is impossible that the Egyptians can
kill me; sooner shall the Almighty by a miracle kill them all,'' then he
would indeed have had an easy time playing the ``knight.'' But to
understand the dispensations is a personal task. The matter presented
itself to him in another light; the thought on which he comes to agree
with Sarah can be expressed thus: ``We must gain time, avert the
momentary danger and let events develop; who knows what Jehovah will
do?''
% printer's defect: the Danish quotation opened at "Vi maae give Tid" is never closed; closed here in the translation
If this is now correct with respect to the doctrine of Providence, then,
as regards the criticized ``white lie,'' we need only take note of
Abraham's words to Abimelech (Gen.~20:12): ``And yet indeed she is my
sister; she is the daughter of my father, but not the daughter of my
mother; and she became my wife''; for in other words this means: ``I
have kept something from you, but what I have said is true.''

What holds of the wondrous signs and dispensations in sacred history
holds in a corresponding manner of the events in every human being's
personal life. Faith in Providence requires a doubling of consciousness.
Here is presupposed an actual world, an order of things, in which the
occurrences are related as links in a chain; this is the one side. But
here it is presupposed as well that the whole actual order of things is
borne, moved and pervaded by God's almighty will; each single occurrence
shows us the opposition between faith and knowledge in the clearest
possible way. Let the occurrence be, for example, a sad death; a hero
dies at the moment he is on the point of delivering his people; a father
of a family dies and leaves wife and children
% --- p. 218 ---
\opage{218}\setcounter{footnote}{0}%
helpless, and so on. If these cases of death are regarded from the
point of view of knowledge, then there is surely no one who doubts that
each of them must have had its natural causes. But if we stop at this
conception, then there can be no talk here of any Providence; each
event is, to recall the Stoics, in its way woven into the great
necessity: everything falls under fate. If a Providence is to show
itself here, the occurrence must be conceived for itself as an event
into which the divine will has expressly laid a personal task for those
concerned.

It is now easy to see as well that the conception corresponding to
faith by no means excludes the natural one. The natural causes of the
hero's death may be known to everyone in the land, and still the people
can in earnest be called upon to take note of the will of Providence in
the sad event. It must be expressly emphasized that when faith
presupposes a unity of the natural causes and God's will, it presupposes
as well that the all-governing will can at every moment alter the
co-operating causes and let events occur precisely as is most
serviceable to man. But science cannot possibly engage in such a way of
viewing things. When a physician conceives a death as a physician, and
a clergyman the same death as a pastor of souls, then the two
conceptions are absolutely heterogeneous. If now the two absolutely
heterogeneous ways of viewing things were incompatible in one
consciousness, then it would surely be as impossible for the physician
to understand the clergyman as it would be impossible for the clergyman
to understand the physician. Is this now the case? By no means! The
dualistic way of regarding things has grown together with our
religious-moral representations to such a degree that it would arouse
the greatest offense in a house of mourning if the physician declared
that the death in question
% --- p. 219 ---
\opage{219}\setcounter{footnote}{0}%
did not concern Providence, since the patient had demonstrably died a
natural death, while the clergyman, in pious zeal to show the finger of
Providence, expressly asserted that God himself had literally killed the
man.

The religious mourner unites the consciousness of naturally operating
causes with the consciousness of God's unsearchable will, thus unites
two absolutely heterogeneous ways of viewing things in one
consciousness; and then people still cry Dualism, Dualism! and make a
noise and a fuss as though it were something unheard of, terrible and
self-contradictory---what, nevertheless, when it is seen in actual life,
they will reckon among the most ordinary things of all.

The miracle and the everyday event resemble each other in that both can
be conceived as dispensations; but just as impossible as it is without
faith to cognize the everyday event as a dispensation, just so
impossible is it without faith to cognize the most wondrous thing of all
as a miracle. A multitude sunk in coarseness could, according to the
account of sacred history, be brought for a moment to be startled and
terrified when the wondrous happened; but it soon forgot the surprise
and felt no impulse to thank Providence. In a corresponding manner it
goes with dispensations even to this day. Something may meet a human
being so joyful, so sorrowful, that in the moment of feeling it is as
though he saw the eye of Providence; but if there is no power of
inwardness, no earnestness of faith in the personality, what happens
then? Then what is not unusual happens: the feeling dwindles, the mind
loses its elasticity, what began by showing itself as a higher
dispensation ends by transforming itself into a quite natural event.
The difference between the knowing ones of the present age, who could
not see the dispensation in a natural event, and those ignorant ones,
who could not see it in the signs of omnipotence,
% --- p. 220 ---
\opage{220}\setcounter{footnote}{0}%
is only this, that the knowing are blinded by thoughts, the ignorant by
thoughtlessness.

To what extreme the more recent theology, infatuated and bewitched by
``scientific'' thinking, has gone in undermining the doctrine of
Providence can be seen from Schleiermacher's dogmatics. Schleiermacher
says of preservation (Der christliche Glaube. 1ster Band. Berlin 1835.
S.~222): ``The pious self-consciousness, according to which we place
everything that stirs us and acts upon us in the absolute dependence on
God, coincides entirely with the insight that precisely all this is
conditioned and determined by the connection of nature.'' So completely
does the divine will, according to this explanation, become identified
with the operating natural causes that fate itself becomes Providence,
and Providence fate. It is the old Stoic doctrine of fate, freshened up
with a little evangelical paint. Schleiermacher's ``pious
self-consciousness'' and Epictetus' pious resignation are here brought
into complete congruence. This congruence has long been in preparation:
Holberg already saw it.

``Nothing,'' says Holberg, ``can move more than the resignation to God's
will that is made by Epictetus. His words are these: I wish with my
whole heart that death at my last may find me prepared, so that without
confusion, without hindrance and constraint, I may end my life, and that
I may say to God: Lord! have I transgressed thy commandments? have I
misused the gifts with which thou hast endowed me? have I not subjected
my will to thy will? have I ever complained of my evil fate? have I ever
accused thy divine Providence? I have been sick, because it has been thy
will, just as it has also been my will; I have been poor, because it has
so pleased thee, and I have been content in my
% --- p. 221 ---
\opage{221}\setcounter{footnote}{0}%
poverty; I have lived in contempt, because thou hast willed to have it
so, and I have not striven to raise myself up. Thou hast never seen me
distressed in my poor condition. Thou hast never found me impatient and
grumbling. I am still ready to submit to all that thou findest good to
lay upon me. The least sign that thou givest is for me a holy and
inviolable order. Thou willest that I shall leave this glittering stage
of the world: behold! I leave it, and render thee a thousand thanks that
thou hast deigned to set me upon it, to show me thy works, and to place
before my eyes the excellent order with which thou governest all created
things.'' (Hundrede og tyve af Holbergs Epistler ved Fabricius.
Kjøbenhavn 1858. S.~114).

At a first fleeting glance it must surely seem as though there were
contained in Epictetus' utterance a deeply satisfying doctrine of
Providence; but precisely the opposite is the case. It is the selfless
being, the universal necessity of reason, that here in the pious
contemplation is personified; there is ascribed to fate a wisdom, a
self, a will, which by the nature of the case it neither has nor can
have. Precisely by the impersonal element in the government of the world
taking on in this manner the semblance of personality, it comes about
that the pious resignation, far from nourishing and strengthening the
personal self in man, consumes the kernel of personality. The pious one
wills what the personified fate wills; by making his will one with the
will of fate, he sacrifices the noblest powers of personality to the
impersonal.

The pious doctrine of fate and the evangelical doctrine of Providence
proceed from absolutely heterogeneous principles. Just as revelation as
a whole is characterized by subordinating the concept of nature to the
concept of
% --- p. 222 ---
\opage{222}\setcounter{footnote}{0}%
creation, so the revealed doctrine of Providence is characterized in
particular by subordinating all operating natural causes to God's
governing will. Instead of identifying God's will with the natural
causes, faith on the contrary presupposes the difference.
(Observandum quod Deus non solum vim agendi dat causis secundis\,\ldots,
sed quod immediate influit in actionem et effectum creaturæ. Qvst.):
the cause that is freedom itself determines the events which, to judge
by the phenomenon, proceed immediately from natural causes.

He who cannot conceive the miracle as a sign of omnipotence cannot
conceive any event whatsoever as a dispensation of Providence either.
The unity of the will of Providence and natural causes that is
characteristic of the everyday event receives in faith its light from
the opposition between the will of Providence and the natural causes
that is made visible in the miracle; the unity must be illuminated by
the opposition, and the opposition by the unity: the miracle and the
dispensation presuppose and explain each other mutually.

The true doctrine of Providence is an unconditional doctrine of
personality; the pattern, the true personality, is himself the revealed
Providence. What, then, has the revealed Providence taught us about
dispensations and trials, about courage and patience, about fear and
danger? We hear His word, which says (Matt.~10:28--30): \emph{Are not two
sparrows sold for a farthing? and one of them shall not fall on the
ground without your Father. But the very hairs of your head are all
numbered. Fear ye not therefore, ye are of more value than many
sparrows}.---``If these words were only human words, then we should
surely have to regard them as an exaggeration in an outburst of
enthusiasm; for in the world of dangers that lies between man and God, a
human being has, after all, much to fear.
% --- p. 223 ---
\opage{223}\setcounter{footnote}{0}%
But when he who speaks the word is the Word itself, then we must learn
from him where the danger is, and what we have to fear. There is in the
world of dangers only one unconditional danger, only one thing that is
fearful, only One who is to be feared. Therefore the Word says: `I will
forewarn you whom ye shall fear.' And Providence is human. He who
speaks of the danger and shows us whom we shall fear; he who surveys all
the dangers of finitude and says: `Fear not!'---he knows what danger
is; his life was familiar with danger. Danger sought him in his cradle;
it followed him to his grave. Everywhere, in the desert as in the
cities, on the plains of Galilee, among the mountains of Judea,
everywhere he went, danger went after him and lay in wait for him.
Watchful as Providence, confident as he who is under Providence, the Son
of Man goes along the way of dangers and fears nothing in the world. He
works, and no one stops him; he goes to his goal, and no one lays a hand
on him before his hour comes. But when his hour had come, he went
willingly to meet the danger; it did not surprise him as an accident, it
did not come over him as a fate; it came because he let it come. Nor
did it come, then, in outward tumult, not in a storm of human passions,
not in a clamor of hostile weapons; it came in all stillness, alone to
the lonely one; only thus did he let it come. Alone that night in the
still garden he let the danger come to him and strove with it, not as
the strong strive, but as the frail, so that in flesh and blood there
should be no difference of fate between the strong and the frail, but a
victory of devotion in frailty. Alone that night in the still garden he
let the danger pass as a mortal anguish through his soul; he did not
defy it, he did not despise it; he did not proudly lift his head like
one who sacrifices himself to fate; he fell on his
% --- p. 224 ---
\opage{224}\setcounter{footnote}{0}%
face and bowed himself humbly beneath the Father's will. Alone that
night in the still garden he let all the dangers of finitude transform
themselves in inwardness into one danger. So he strove in inwardness and
conquered in inwardness; and when he then rose from the strife, behold,
then it was over with the danger, then there was no danger at all.''
(Om Tilværelsens ideale Magter: Skjæbne og Forsyn. Andet Oplag.
Kbhvn.\ 1867. S.~30 flg.).

The doctrine of Providence is dualistic. If the dualism is abolished,
then Providence vanishes. But when Providence vanishes for the leaders
of the people and for the people itself: how will it then go with the
spiritual life in the present and the future?

\begin{center}\rule{3cm}{0.4pt}\end{center}

\chapter*{XIV.}
\addcontentsline{toc}{chapter}{Lecture XIV}
\markboth{Lecture XIV}{Lecture XIV}
\label{ch:xiv}
% --- p. 225 ---
\opage{225}\setcounter{footnote}{0}%
It was about hindrances and conditions for the spiritual life of the
present age that these lectures were essentially to treat. But for an
assessment of the spiritual life, its phenomena and powers, it was
first and foremost necessary to choose a standpoint from which the whole
could be tolerably surveyed, and to secure a standard under which the
mutually different, intersecting tendencies could be brought. It is the
separation between faith and knowledge, between religion and science,
that marks our chosen standpoint. From various sides this standpoint has
gradually been illuminated in the foregoing; but we must admit it, it is
only a partial illumination! Here an all-round settlement is required, a
last concluding account.

Such an account must surely at first glance seem bound to become very
intricate, since it is to be rendered before two tribunals, each with
its own procedure, that of religion and that of science; the values
must in this settlement be reckoned in two kinds of coin: it becomes, as
has been said mockingly, a double bookkeeping. Nonetheless the matter
can from the beginning be laid out very simply. For since there are
here two circles of cognition, each with its principle, it suggests
itself that each of the parties must punctually have its own; but when
% --- p. 226 ---
\opage{226}\setcounter{footnote}{0}%
each gets its own, then surely all strife and complaint cease. If
religion, with its ideals of revelation, its sacred facts and personal
thoughts of will, is permitted to constitute a kingdom of the spirit
for itself, a region of representations and concepts that are developed
with inner consequence in their own connection from their own
presuppositions: what more can faith then demand? And now to the other
side! what greater concession can one well make to science than when
one ascribes to it a standpoint of its own, a method of its own, laws
and consequences of its own, and besides freedom to develop in all
directions, to draw everything before its tribunal, to investigate and
test everything that can in any way fall within a human horizon? One
would think that science for its part would have to find itself
satisfied, just as well as religion. And yet, we dare not conceal it,
the account has its difficulties. The point is that the two circles of
cognition, although they are mutually absolutely heterogeneous, are to
be united in one consciousness, and this in such a way that the
personality developed to self-understanding becomes able to move with
certainty in knowledge and yet find rest in faith. This can only happen
by its learning to see through the fundamental error, the first false
presupposition from which all the objections of knowledge against faith
proceed, so clearly that it is able everywhere to remove the
misunderstanding without violating science. But in this endeavor the
personality is exposed to temptation. ``The temptation is that faith
and knowledge are so easily made somewhat homogeneous, although they
are absolutely heterogeneous. If the absolute in the heterogeneity is
stressed sharply enough, then a screaming contradiction is heard:
\emph{what is truth in faith is untruth in science; what is truth in
science is untruth in faith}. But in the contradiction lies the
possibility
% --- p. 227 ---
\opage{227}\setcounter{footnote}{0}%
of its solution. It is the absolute limit of our knowledge, it is the
riddle of existence, on which the attention of the knower as well as of
the believer must here fix itself. If all human knowledge has an
absolute limit, if there is here really a mystery of the spirit, a
riddle of existence, then there is here also something that must be
\emph{believed}, because it cannot be \emph{known}. Now when that which
is hidden in the riddle, that which is solely and entirely destined to be
an object of faith, is nonetheless drawn under knowledge and science,
what then? Then knowledge oversteps its absolute limit. And when
knowledge has overstepped its absolute limit, what then? Then science
sets about, blithely, to fathom, to ground and to comprehend that which
by its nature neither can nor shall be scientifically comprehended. But
when science takes it upon itself to comprehend the incomprehensible,
what then? Then it necessarily arrives at contradictory concepts. What
can only be expressed in contradictory concepts must be untrue; the
mysteries of faith can only be expressed in contradictory concepts:
therefore they are untrue. That something can be true in faith and
untrue in science now follows directly from this, that \emph{science,
by wanting to comprehend what lies outside its limit, itself becomes
untrue}.''

``But how can something that is really true in science be untrue in
faith? Answer: Because it is not only knowledge, but faith too, that
has an absolute limit: that knowledge has faith as its limit
necessarily presupposes that faith has knowledge as its limit. But that
faith has knowledge as its limit means that a concept of faith cannot
possibly be developed into a concept of knowledge. In relying on the
miracle, faith presupposes an actual nature with actual laws. But how is
nature related to the miracle? Why, that is precisely what
% --- p. 228 ---
\opage{228}\setcounter{footnote}{0}%
is the riddle of existence. If now faith, forgetting its absolute
presupposition, namely the riddle, oversteps its limit and begins to
want to overhaul the scientific concept of nature for the benefit of the
miracle, then the concept of nature of natural science becomes not only
an untruth in faith, \emph{but faith, which has engaged in something it
is not equal to, itself becomes untrue}.'' (Om den gode Villie
S.~31--33.)

The next thing that presents difficulties is the form of the account.
This difficulty does not concern in the first instance the religious,
but the scientific. The fundamental thoughts of the spiritual life can
really be presented in a humane form of communication. It would here be
out of place to make excuse for what is deficient in such a way as to
shift the blame onto the difficulty of cognition. That would be a
self-contradiction according to the chosen standpoint; for it lies, as
already shown, in the nature of personality that its decisive
conditions of life must be capable of being developed clearly,
comprehensibly, popularly, and intelligibly for all.

Otherwise with the scientific side of the account; here the humane form
of communication is insufficient. A scientific account can only be
rendered in strict scientific form, and every attempt to make it fully
recognizable and comprehensible for outsiders, for the humane
consciousness in cultured circles, is an attempt to convert the
strictly scientific into something popular; but how much the popular
form departs from the scientific requirements has already been
sufficiently emphasized in the
% printer's error in the Danish ("videnskalige" for "videnskabelige"); translated in the corrected sense
foregoing. Among the stupefying elements in the spiritual life of the
present age one may without doubt reckon the disproportionate trust in
scientific communications at second hand, in conception in popular
forms.

I must therefore forgo presenting the
% --- p. 229 ---
\opage{229}\setcounter{footnote}{0}%
scientific account; but so that it shall not look as though under this
pretext I were evading a subsequent account, I will go as far as is
tolerably possible, in that I will briefly point out the points of view
and orienting points of departure on which, in the rendering of the
scientific account, it must essentially depend. I indicate these points
of view by referring to written statements, insofar as these are
available. The matter is, after all, in its first coming-to-be; I have
not yet had occasion---what is nevertheless necessary for the complete
elucidation of the whole---to give a detailed lecture course on the
philosophy of religion.

When the talk is of separation between faith and knowledge, it suggests
itself to separate religion from science by making religious cognition
into a paradox. This is what Kierkegaard with wonderful power and genius
has carried through. In resolutely breaking the staff over every attempt
to want to justify the truths of religion by the scientific way, he made
the break so decisive that every cognition from the side of religion
turned against science with a pricking sting, faith being posited as
paradox. I shall remark on this that my chosen standpoint does not stand
in conflict with, but neither does it coincide with, the Kierkegaardian
view. This marks a stage, a standpoint by itself; this stage had to be
reached, and the knot of the paradox drawn properly tight, before we
could get further in religious cognition. How different the chosen
standpoint is from the Kierkegaardian can be made plain in a few
strokes. The opposition between faith and knowledge is in Kierkegaard
transformed into an opposition between passion and science. Science and
passion are the terms of the fundamental opposition; he complains
precisely that the theologians in our theocentric nineteenth
% --- p. 230 ---
\opage{230}\setcounter{footnote}{0}%
century are about to forget passion over science. He wants, then, to
make passion count again, especially the holy passion, faith. When faith
is posited as holy passion, one sees at once how impossible it is to
present it in a doctrine, in an objective dogmatic form of teaching.
Therefore Kierkegaard impresses again and again that the truths of
religion cannot be taught; Christianity, he says, is no doctrine, it is
a communication of existence.

For the sake of an overview I have published a work: ``Paa
Kierkegaardske Stadier, et Livsbillede.'' I have not published this work
to spare the reader from going to the sources; but the Kierkegaardian
literature is so extensive, many-sided and ramified that one easily goes
astray in the various forms of productivity as in a primeval forest.
Therefore there are many who read large portions of Kierkegaard without
having a sure overview of the whole productivity. Some lose themselves
in ``Enten Eller'' and appropriate him as a witty, brilliant aesthete;
others skip over the aesthetic, keep to ``Frygt og Bæven,'' to
``Kjærlighedens Gjerninger,'' to ``Indøvelse i Christendom,'' and so on,
and take him without further ado for the stern preacher of repentance.
The little work therefore presents the stages, the aesthetic, the
ethical and the religious, in their mutual relation---for an overview.

``Love,'' says Kierkegaard, ``has after all its priests in the poets,
and at times one hears a voice that knows how to keep it in honor; but
of faith not a word is heard; who speaks to the honor of this passion?
Philosophy goes further. Theology sits painted at the window and courts
its favor, offers its loveliness for sale to philosophy.'' The
presentation of the ideals of faith ought by no means to be doctrinaire;
Kierkegaard does everything here
% --- p. 231 ---
\opage{231}\setcounter{footnote}{0}%
to break the scientific form. Thus in his ``Smule Philosophie'' he
keeps the positive content of Christianity in the form of an
experiment, and in the scientific Efterskrift presents the whole so
humorously and with such light play that one would think it was only a
jest, although with this jest it is in dead earnest.
% the Danish reads "den videnskabelige Efterskrift" (not "uvidenskabelige"); rendered literally
To present the ideals of faith is for the defender of the paradox rather
a poet's task than a scientific concern. ``If it is so that you can
present faith, that only proves that you are a poet, and if you do it
well, that you are a good poet, but nothing less than that you are a
believer. Perhaps you can also weep when you present faith; that would
then prove that you were a good actor.'' To develop the paradox
theologically would in Kierkegaard's eyes be one of the most
wrong-headed things one could undertake, for the paradox does not allow
itself to be determined from a traditional or objective doctrinal
standpoint. The whole manner of presentation must, according to his way
of viewing things, essentially be a presentation of moods. The
Kierkegaardian dialectic is a dialectic of moods, his concepts are
concepts of mood. As far as I know, no one in the present age has been
able to measure himself with Kierkegaard in developing the religious
concepts to the highest pathos. I choose an example from ``Frygt og
Bæven.'' The pseudonymous author speaks of Abraham's faith, and speaks
in mood. To characterize Abraham's spiritual greatness he compares him
with other heroes and thereby comes to depict for us characters of
different spirit, different kinds of greatness. ``No one,'' he says,
``shall be forgotten who was great in the world; but each was great in
his own way, and each in proportion to the greatness of that which
\emph{he loved}. For he who loved himself became great by himself, and
he who loved other men became great by his devotion, but he who loved
God became greater
% --- p. 232 ---
\opage{232}\setcounter{footnote}{0}%
than all. Each shall be remembered, but each became great in proportion
to his \emph{expectation}. One became great by expecting the possible,
another by expecting the eternal; but he who expected the impossible
became greater than all. Each shall be remembered, but each was great
wholly in proportion to the magnitude of that with which he strove. For
he who strove with the world became great by overcoming the world, and
he who \emph{strove} with himself became great by overcoming himself;
but he who strove with God became greater than all. Thus there was
striving in the world, man against man, one against a thousand, but he
who strove with God was greater than all. Thus there was striving on
the earth: there was he who overcame all by his power, and there was he
who overcame God by his powerlessness. There was he who relied on
himself and gained all, there was he who, secure in his strength,
sacrificed all, but he who believed God was greater than all. There was
he who was great by his power, and he who was great by his wisdom, and
he who was great by his hope, and he who was great by his love; but
Abraham was greater than all, great by the power whose strength is
powerlessness, great by the wisdom whose secret is foolishness, great by
the hope whose form is madness, great by the love that is hatred of
oneself.''

Here, then, we have the paradox, in the highest pathos, the richest
mood. One easily sees that this speech rich in mood does not at all
permit any scientific reasoning. How would not every formal,
metaphysical objection fall dead and powerless to the ground in the
face of an utterance such as this, that \emph{he who expected the
impossible became greater than all}! To expect the impossible, what
does that mean? When one is to reflect, one can well grant that in
faith man expects the improbable; but who can expect what is in itself
% --- p. 233 ---
\opage{233}\setcounter{footnote}{0}%
impossible? And yet the impossible is posited so sharply in the place
of the improbable. Much may seem impossible for men, but precisely the
fact that it is impossible for men lays a tacit stress on its being
possible for God. Every accord between what is impossible for man and
what is possible for God is, however, cut off by the paradox. The
direct understanding is rejected, all finite reasoning must fall
silent; the force and the momentum lie in the ideal pathos: the
communication of existence does not take time for theoretical
negotiations.

Kierkegaard has strikingly seen that the question of the truth of
Christianity cannot possibly be decided in science, since subjectivity
in science must always be contemplative, and the contemplative
subjectivity objective. ``The modest, objective subjectivity keeps
itself outside with applauded heroism; it is at your service in being
willing to accept the truth as soon as it has been procured. Yet it is
a distant goal that is striven for (undeniably, for an approximation
can go on as long as it pleases)---and while the grass grows, the
observer dies, calmly, for he was objective. O thou objectivity, not
praised for nothing, thou art capable of everything; not the most
believing one has been so certain of his eternal happiness, and above
all so sure of not losing it, as the objective one! Unless it should be
that this objectivity and modesty were out of place, that it was
unchristian. Then it would indeed be precarious to enter in that manner
into the truth of Christianity. Christianity is spirit, spirit is
inwardness, inwardness is subjectivity, subjectivity is in its essence
passion, at its maximum an infinitely, personally interested passion for
its eternal happiness.'' (Afsluttende uvidenskabelig Efterskrift til de
philosophiske Smuler af Joh.\ Climacus. Kbhvn.\ 1846. S.~19--20).

% --- p. 234 ---
\opage{234}\setcounter{footnote}{0}%
When Kierkegaard develops his concepts, we must not take them in the
merely general form; for when one separates these concepts from the
corresponding mood and strong pathos and draws them under contemplation,
one has distorted the whole, and it is precisely this that Kierkegaard
inveighs against so strongly. Metaphysical objections, critical
nit-picking remarks have nothing to signify here.

That this standpoint does not coincide with the one chosen by us can
easily be seen; for we have, after all, precisely in direct opposition
to the Kierkegaardian rejection, striven again and again to draw the
whole system of religious cognitions under the level of contemplation.
We grant that Kierkegaard is unconditionally right in the assertion that
faith must be a paradox for the non-believer. We grant that he who,
although he accepts the presuppositions of faith, nevertheless desires
to clarify these in knowledge and science, and thinks he has need of a
science of faith, is repelled by the paradox. We grant that the
existing individual, the one hard pressed in the crises of existence,
must be stimulated by the sting of contradiction, must again and again
be wounded by the paradoxical. But now a question remains, and it is the
one that has moved through our investigations, namely this: Is faith,
then, also a paradox for itself? Is the believer so blind to his own
moments of thought that they cannot possibly develop for him into clear
concepts coherent in themselves? This must be denied. Faith is in
itself a principle of cognition; the word of faith is the word of light
and of life, and from this proceeds all clarity about the content of
religion. Faith is indeed ``a holy passion,'' but when the holy passion
is reflected through, it shows itself as a pathetic expression of the
will striving after holiness: the principle of faith is in the deepest
sense a principle of will, faith itself an act of will, an
% --- p. 235 ---
\opage{235}\setcounter{footnote}{0}%
innermost action (nemo credit, nisi volens): the thoughts of faith are
thoughts of will, all of them together decisive thoughts of
personality. When the whole, internally coherent sum of religious
thoughts of will is developed and clearly presented in a system of
concepts of faith, we get, not a paradoxical communication of existence,
but a doctrine of religion valid from the standpoint of faith, a
systematically coherent doctrine of religion. So much with regard to the
Kierkegaardian stage. We come now to the next question.

For if it is granted that there is a system of religious cognitions,
that we must once again get a doctrine of religion, then the chosen
standpoint seems, although we have declared ourselves decisively
against all theology, nevertheless to fall back under the theological.
I will begin with a stupid remark, a remark which otherwise from so
many sides presents itself as the offspring of rare cleverness; the
remark is this: ``to give a systematic presentation of the religious
cognitions is to make the doctrine of religion objective; to make the
doctrine of religion objective is to place oneself at the standpoint of
the observer and to make the doctrine scientific, and so, when one is
nevertheless to hold to the unconditional validity of subjectivity and
reject science, an obvious contradiction.'' Why I call this a stupid
remark I can easily explain. If it were not possible to develop and
ground the doctrine of religion because this can only happen on the
condition that one places oneself at the standpoint of the observer,
then it would, after all, not be possible either---what has nevertheless
notoriously happened in ``Uvidenskabelig Efterskrift''---to give a
humorous, experimenting presentation of the Christian problems.
Johannes Climacus himself stands, after all, from first to last at the
standpoint of the observer; and yet there is no contradiction whatever
in his, standing at this standpoint, smiling at the ``applauded
heroism'' of the
% --- p. 236 ---
\opage{236}\setcounter{footnote}{0}%
``modest objective'' subjectivity, and sharpening the subjective. For
here there is an enormous difference between the contemplative
consciousness that argues and the consciousness of which, from which and
to which the argument is made. In the doctrine of religion as well as in
the humorous experimenting, the argument must throughout the whole
contemplation be made from the selfhood concentrated in itself, the
unconditionally valid subjectivity. Now follows another, more reasonable
objection. The systematic development of a doctrinal concept
corresponding to the principle of faith is, and must by its nature
necessarily be, scientific; but what is a scientific development of the
content of positive religion other than theology? ``You can,'' people
say, ``choose your standpoint where you will: if you really solve the
task of presenting the content of the revealed religion in a coherent
development flowing from one principle, then you are and remain a
theologian, and it is no use your protesting.''

To see how matters stand with this assertion, we need only cast a
glance at the theological standpoints.

The first that presents itself is then the more recent
speculative-dogmatic one. This standpoint has, after all, wanted to take
up all the presuppositions of faith, has wanted to build the cognition
of faith on the life of faith, has wanted to develop the concepts not
from pure knowledge but from ``the good will,'' and in this way has
wanted to make the principle of will into a principle of science. But to
make the principle of will into a principle of science is impossible.
Should theology by virtue of its good will really be science, then it
would have to be the highest of all sciences; it would have to be able
in the name of the will to regulate the other sciences, which stand
lower in relation to the holy and are therefore
% --- p. 237 ---
\opage{237}\setcounter{footnote}{0}%
subordinate; it would, in short, have to be able to re-establish its
medieval universal monarchy.

``But a new theological universal monarchy, what an enormous thought!
Can a thought of this height, without stooping so low that it becomes
quite crooked, enter by the door of actuality? I too have met the lofty
thought; it has many times come to meet me when I regarded dogmatics,
not from the shadow side, but from the light side; for although,
according to my task, I have most often had to wander through its
shadowy parts, I have not therefore forgotten where I can still find the
sunlit places. But when, after having beheld the ideal in all its
glory, I then seriously asked myself: can such an ideal be realized? I
have always received a negative answer. It is surely so that the idea of
the good, conceived as the idea of love, of the holy will, is in divine
knowledge one with the eternal idea of truth; but from this it by no
means follows that the idea of truth in a human knowledge is one with
the idea of the good in human existential willing. It is surely so that
in wisdom itself no gulf can be thought between faith and science: but
why is God no man of science? Because God is no believer. To think of
the divine knowledge as a scientific knowledge, be it philosophical or
theological knowledge, is as blasphemous as to speak of a believing God.
Why must a human being apprehend the absolute in faith? Because the
purest and highest development of human knowledge is, not divine, but
scientific.
Why must a human being apprehend the absolute in knowledge? Because
faith is a human, not a divine, cognition. Why does the absolute that is
apprehended in human knowledge show itself so heterogeneous with the
absolute that is apprehended
% --- p. 238 ---
\opage{238}\setcounter{footnote}{0}%
in faith? Because the relative forms of apprehension grasp the
absolute, not in its central depth, as it must be thought by us
\textit{in abstracto} to be, with infinite concretion, in the divine
self, but in two polar extremes, of which the one is marked by universal
thinking grounded in rational necessity, the other by the punctual life
of will germinating in conscience. Here there is in the absolute
something I must absolutely believe, because there is in the absolute
something I must absolutely know.'' (Den gode Villie som Magt i
Videnskaben S.~6--7).

That there is a relation between will and science cannot be denied; he
who is to think must will to think; he who is to cultivate a science
must will to cultivate it. Indeed, the relation is even personal,
ethical, insofar as everyone who spends his time cultivating a science
must be responsible for how he employs his powers, whether cases do not
arise when for higher considerations he must break off; and so on. But
these and similar considerations have nothing whatever to do with our
present question. The personality, which is answerable for everything
it thinks and does, is answerable also for the manner in which it
conceives the relation between the will and science. The good will is
not a power in science, but it is a power in religion. The religious
problems of thought are at bottom all problems of will. In every problem
of will the argument runs from the self, by the self, to the self;
whereas in every problem of knowledge the argument runs from the thing,
by the thing, to the thing: to argue from nature is to argue from the
thing, to argue from creation is to argue from the self; to argue from
the physical-organic individuality is to argue from the thing, to argue
from the pure heart and sinless will of the ideal man is to argue from
the self; to argue
% --- p. 239 ---
\opage{239}\setcounter{footnote}{0}%
from an inner necessity of the idea is to argue from the thing, to
argue from the freedom of the spirit is to argue from the self, and so
forth.

The confusion in the Christian-speculative dogmatics betrays itself in
each of the chief dogmas, creation, the incarnation, inspiration, in
that two half thoughts, one half from the self and one half from the
thing, one half from the will and one half from science, one half from
the subjective and one half from the objective, rise up against each
other, clash with each other and end by neutralizing each other. This
\textit{concordia discors} of the heterogeneous thoughts has already
been described in several ways and can be described in still more.

``One knows how it has gone with `the Christian dogmatics,' what
mishaps it has had with the antirational in the great simple chief
dogmas. Take now the dogma of creation; why, that would be a very
profound and exceedingly scientific dogma, if only there were not the
snag, the antirational, about it, that as surely as dogmatics is quite
honestly to make it come out that the world is created from nothing,
then there is no nature in the world, and if it is to acknowledge that
here is an actual nature, then it cannot possibly make it come out that
the world is created from nothing. (Cf.\ `Den gode Villie' S.~96--104.)
To lump the whole together in speculative common stock, to let nature
be creation and creation nature, so that faith may be knowledge and
knowledge faith, is so humble an inconsistency that it will surely not
fail to humble theology. But the incarnation, the dogma of the God-man?
Yes, the dogma of the God-man, the great dogma which through the
centuries, with the rational, the antirational and the irrational, has
set so manifold thoughts as well as passions in motion, would
undoubtedly be the favorite dogma of the faculty, the palladium of
theology, the central dogma of the heart, if `the Christian
% --- p. 240 ---
\opage{240}\setcounter{footnote}{0}%
dogmatics' had not had that sorrow at the last, the great sorrow with
the Son. What sorrow? Alas, his consciousness, his dogmatic
Logos-consciousness, is, \textit{horribile dictu}, cracked, dualistically
cracked; the one Son has---through the crack---become two sons: one
whose consciousness is from eternity but who has not become incarnate,
and another, a younger, who has become incarnate but whose
consciousness is not from eternity. (Cf.\ `Den gode Villie'
S.~91--92.) To be sure, theology holds as much as it can to `the unity
of consciousness,' in the incarnation too; but the crack, the crack!
The dogma of the incarnation is a very difficult dogma. But the dogma
of inspiration, that is surely not so difficult? If there is any dogma
that has the testimony of the Spirit, then it must surely be the dogma
of the outpouring of the Spirit. Since it is now, moreover, to be quite
certain that dogmatics too, by virtue of the dogma of the outpouring of
the Spirit, has `the testimony of the Spirit,' we may well expect that
the dogma of the outpouring of the Spirit must be as though made for
dogmatics. Yes, it really would have been as though made for dogmatic
speculation, if `the Christian dogmatics,' in stumbling over the
antirational, had not in the fall had the mishap that its Pentecost
miracle came to stand on its head. But to turn the Pentecost miracle,
the great `church-founding miracle,' upside down to such a degree that
the hearers hear many \emph{different} tongues although the apostles
together speak only one \emph{single} tongue; that the uppermost thus
becomes lowermost, the lowermost uppermost---that is more than
reformatory, it is revolutionary. That dogmatics with its miracle
revolution has meant well, heartily well, both with the apostles and
with the Church and the feast of Pentecost, we shall gladly grant; but
revolutionary it is, that is certain (Cf.\ `Den gode Villie' S.~84).
The dogma of inspiration is therefore a very difficult
% --- p. 241 ---
\opage{241}\setcounter{footnote}{0}%
dogma. Yes, and the dogma of baptism, with its conscious `mystical-magical'
element, is a very difficult dogma; and the dogma of the speculative presence
of Christ's body in the Supper is a very difficult dogma; and the dogma of
eternal damnation is a difficult, a very difficult, dogma. There is in eternal
damnation something so antirational that `Christian dogmatics' has
expressly had to demand a non-damnation (apokatastasis), so that damnation and
non-damnation, in a terrible collision, might end the whole with a solemn
antinomy, and the antinomy (the conflict of reason) might stand as a
testimony, a testimony concluding the system, that there is unity in the
dogmatic consciousness, that theology truly masters the antirational. ---
So long as we frivolously slip away from every particular discussion of any
particular dogma whatsoever by saying, when things get tight: yes, here of
course there are difficult points, but for the rest, etc., it is an easy
matter to play at reconciliation with faith and knowledge. But if a stand is
to be made on definite points, if the particular, definite dogmas are to be
accounted for in earnest, then one will soon come to acknowledge that all
dogmas without exception are equally difficult, equally antirational, equally
at odds with scientific comprehension. The thoughts of faith all bear the mark
of faith, are all furnished with the antirational sting, so that whoever would
crumple them scientifically will surely feel the prick in the end.'' (Om
Holbergs Kirkehistorie og Theologie. Kjøbenhavn 1866. S.~35--37).

The same thing is said in another way in my work Den gode Villie som Magt i
Videnskaben, S.~92--93: ``What is peculiar to dogmatic thinking is that one
thinks scientifically within faith: \textit{credo, ut intelligam.} That one
thinks scientifically within faith means that one does not think, but
represents, intuits, fantasizes in a
% --- p. 242 ---
\opage{242}\setcounter{footnote}{0}%
metaphysical manner for the elucidation of the mysteries, for the benefit of
revelation, for the joy of faith. Dogmatics, which objectively builds on the
word of Scripture, sees in the passages of Scripture something quite other and
deeper than either Christ himself or his apostles ever saw: it sees the
metaphysical. `He would not,' it says (Dogm.\ \S~133) of Christ's
self-abasement, `merely live in the pure glory of the Godhead, in metaphysical
majesty; but he emptied out his fullness in the lowly form of a servant in
order to become the head in the kingdom of love, the reconciler and the
redeemer. But this self-abasement must at the same time' --- here comes the
proper positive explanation of the miracle --- `be seen as his
self-perfection,' the self-perfection, namely, of the Logos that is in itself
eternally perfect. Inasmuch as he, the one and same Logos, `lives in human
fashion, and as the Son of Man possesses his Godhead only within the
limitation of human individuality, in the bounded forms of human
consciousness, he must indeed be said to live in abasement and poverty, since
he has renounced the majestic glory in which, as the omnipresent Logos, he
shines through the whole creation.' (Dogm.\ \S~133). If we now take note of
the two expressions `he has renounced' and `he shines through,' and remember
that it is the same, the very same Logos, the one and very same Logos-I, that
has here renounced the majestic glory, which he nevertheless has not for a
single moment renounced, how then could we doubt that it is precisely the
higher Phlyaric, the deeper Phlyaretic, that by this explanation is posited in
its positive moment.''

The result of the investigations can be held fast in the following
little-edifying reflection: ``The mediation of faith and knowledge can of
course only be carried through on the
% --- p. 243 ---
\opage{243}\setcounter{footnote}{0}%
condition that both parties approach each other with all possible diplomatic
courtesy, that is, in all sincerity, disinterestedness, animated only by the
warmest wishes for peace and good understanding. In the course of the mutual
approach a happy change takes place in both of them: faith becomes genial,
knowledge good-natured. That faith becomes genial means that it gives up its
old-fashioned orthodox fastidiousness, is ashamed of what is to the Jews an
offense and to the Greeks foolishness, smiles at philosophy and falls in love
with speculation. It can of course be no ordinary faith, not the popular,
simple faith of refuge, that behaves so charmingly; no, it is a cultured,
well-bred faith, a faith that knows its mystics and has attended lectures at
the Institute; only such a faith, a subjective faith with an objective heart,
a faith that would not love practically if it did not love theoretically, can
open its arms, its deep, dogmatic arms, to the higher knowledge. That the
knowledge which is to be able, without nausea, to give itself up to so
cultured, so Institute-cultured a love must absolutely be \emph{good-natured}
is surely self-evident. For here there is a difference, a very considerable
difference, between the good-natured and the not good-natured knowledge. The
not good-natured knowledge we may with reason call the teasing knowledge,
namely the dogmatically teasing, because it always takes its pleasure in
teasing dogmatics. The teasing consists quite simply in this, that it pursues
with merciless consistency every contradiction into which dogmatics in its
innocence has entered. Out of a couple of innocent ``Must-not-be-thought''s it
conjures up, with the help of the logical-physical, a demonic swarm of
``Must-not-be-thought''s; by critical-dialectical witchcraft it gets the
dogmatic contradictions, if we may say so, to breed.
% --- p. 244 ---
\opage{244}\setcounter{footnote}{0}%
Yes, what is most outrageous: the very deepest positive, that of which
dogmatics expressly says that \emph{it must be thought}, it transforms into a
restless negative, a thoroughly unthinkable, a ``Cannot-be-thought'': that is
malicious. The malicious knowledge, the knowledge that will not let go of the
logical-physical, is a sinful knowledge, a knowledge that cannot say
\textit{credo, ut intelligam}, an \emph{unbelieving} knowledge. It is
otherwise with the good-natured knowledge: blessed are they who will not
criticize the good-natured knowledge. The good-natured knowledge does
sometimes look somewhat thoughtless; but, by dogmatics, it is not
thoughtless, for it goes about lost in thought. Why does it go about lost in
thought? It is in love, subjectively in love with the objective heart of
faith: it is, after all, a believing knowledge. The believing knowledge is a
knowledge that longs --- out beyond the logical into the supralogical: the
metaphysics of miracle is the goal of its longing. If it does not reach the
goal of its longing, it dies; but if it reaches the goal of its longing, it
lives, lives on the positive of faith, snaps its fingers at the negative of
criticism and lulls all contradictions to sleep in absolute joy over the
relative mysteries. Should one or another contradiction be so unruly as to
cry out from the nest of the mysteries, what does that matter? `The French
don't take it so precisely.' The believing knowledge is good-natured:
good-naturedly it sinks to the breast of dogmatics, good-naturedly it rests at
the heart of faith. Practically, good nature is love; theoretically, love is
good nature. If it must be said in the practical sphere that love covers the
multitude of sins, then it must
be said here, where the deepest practical is at the same time the highest
theoretical, that good nature, the incredible good nature of the believing
knowledge, covers the multitude of contradictions. On these terms the union of
faith and
% --- p. 245 ---
\opage{245}\setcounter{footnote}{0}%
knowledge is accomplished; the Phlyaric is posited in speculative
perfection.''

Here, then, there can be no doubt that there is a conflict in principle
between my standpoint and the speculative-dogmatic one.

Then a third question arises: How does the chosen standpoint relate to the old
orthodox doctrine? For when one departs from the paradox by forming a
doctrinal concept, when one departs from the modern foundation by protesting
against mediation, then it seems we must surely come back to the old orthodox
presentation. Is this now the case? I answer: the chosen standpoint can by no
means merge with the old orthodox one; for the orthodox standpoint is
characterized by pointing up the concepts of faith quite rightly but
motivating them very badly. What is characteristic of the concepts of faith
the old dogmatics knows how to find out with rare acuteness, but it is
beguiled by science and takes purely human knowledge up into its development.
Its fundamental contradiction is that it too makes the doctrine of faith into
science, and the consequence is that this doctrine then in turn falls prey to
scientific criticism and is dissolved by science.

In order to become science, theology must necessarily use reason; but in order
to be a doctrine of faith, an orthodox doctrine of faith, it must at the same
time take reason captive under the obedience of faith. To use reason is to
acknowledge the consequences that follow with logical necessity from the
premises set up in the name of faith; but to take reason captive, in other
words, not to use reason, is to deny a consequence even though it manifestly,
and with logically irrefutable necessity, follows from its premise. As a
science of faith,
% --- p. 246 ---
\opage{246}\setcounter{footnote}{0}%
theology is thus both to acknowledge and not to acknowledge that the
consequence follows from the premise, is both to use and not to use reason;
this is the self-contradiction, a self-contradiction as old as theology.
There is not a single dogmatic dispute in the whole history of the Church that
does not bear witness to this. A more telling proof of the scientific
self-contradiction of Lutheran orthodoxy than its protest against Calvin's
doctrine of election can hardly be conceived. It is openly admitted that
Calvin is scientifically consistent, and just as openly his consistency is
condemned; and yet theology is a science!

``The beneficently inconsistent way of representing things in the Formula of
Concord, in relation to the blasphemously consistent Calvinist one, is,'' says
a Lutheran orthodox church historian, ``in brief this: sanctification and
salvation are according to both only a work of divine grace, damnation a
consequence of one's own guilt. But whereas the Formula of Concord humbly
stops at this last proposition, and therefore expressly ascribes the guilt to
man alone, and not to the lack of grace \ldots{}, Calvin argues further: But
that guilt occurs only because God does not here too powerfully impart his
grace; but God does not here too impart his grace because from eternity he has
decreed damnation, and even the fall of man has had to serve to realize this
decree, just as in general everything is decreed from eternity --- God has
thus himself willed the sin of men. This last necessary consequence, however
much it may be embellished, the Formula of Concord especially abhors as
contrary to the holiness of God and to the divine word in Scripture, although
it also does not fail to recognize the difficulties in the
% --- p. 247 ---
\opage{247}\setcounter{footnote}{0}%
whole connection of this doctrine, which no human understanding is able to
resolve.''

If, after what has been made clear here, there should be any sense in applying
the laws of reason in some cases and abolishing them in others, then there
would surely have to be in theological knowledge an infallible divine rule
according to which the application and the abolition could be carried out.
But where should we turn to find such a rule? It is no use applying to reason;
for if reason itself could teach us in which cases we ought, for faith's
sake, to prefer ``the humble inconsistency'' and suspend the law of reason,
then the application of reason would precisely be confirmed just as it was to
be abolished: reason would in this way be judge in its own cause and would get
an absolute dominion over faith. If, on the other hand, it is faith separated
from reason, and so a faith that by its separation from reason is unreasonable,
that is both to annul and to confirm the law of reason, well then faith
certainly keeps an unconditional dominion over reason; but what then becomes
of the scientific character of theology with this faith that is in itself
unreasonable? To speak of a certain reasonable suprarational, a mystical
something, which by standing both above reason and above faith should regulate
the relation between them, is to talk off into the fog; for the ``humble
inconsistency'' of the Formula of Concord proves precisely that here there is
at bottom neither regulative nor consistency. In the dispute over election the
Lutheran theologian solemnly declares Calvin's ``inference'' to be ``contrary
to the holiness of God and to the divine word in Scripture.'' On which side,
now, is consistency? Naturally on both sides. Calvin proceeds with the
Lutheran from the premise that sanctification and salvation are only a work
of divine grace, and then concludes in his fashion
% --- p. 248 ---
\opage{248}\setcounter{footnote}{0}%
consistently that the failure of salvation must absolutely be a consequence
of the failure of grace, and this failure in turn a consequence of the divine
decree. The defender of the Formula of Concord proceeds with Calvin from the
premise that damnation is a consequence of man's own guilt, and then concludes
from this quite consistently that it cannot possibly be a consequence of God's
eternal decree. On which side is inconsistency? Naturally on both sides.
Calvin is inconsistent when he posits damnation as a consequence of sin and
would make God the eternal author of damnation without, however, making him
the author of sin. The defender of the Formula of Concord is inconsistent in
that, in order to make God alone the author of salvation, he declares that man
does indeed have the ability to resist grace when it is a matter of
perdition, but on the other hand no ability to cooperate with [grace],
% printed "samvirke med / altsaa ingen Evne": the object of "med" is wanting;
% sense supplied in brackets.
and so no ability to resist grace when it is a matter of deliverance and
salvation. The Lutheran's inconsistency consists in this, that he now adjusts
the principle to the outcome, now again the outcome to the principle: if the
outcome is damnation, then man has been able to resist grace, for man must
himself be to blame for his damnation; if, on the other hand, the outcome is
deliverance and salvation, then man has not been able to resist grace. Had
man's will been able to resist saving grace, then this man's will would,
precisely by not counteracting grace, have cooperated with grace. But the
thought that man's will should in any way be able to cooperate with divine
grace,
is a synergism which the Formula of Concord can only pronounce upon with
condemnation.

If now both, the strict Calvinist and the strict Lutheran, are equally
consistent and equally inconsistent, then it is surely impossible for science
to award the one the
% --- p. 249 ---
\opage{249}\setcounter{footnote}{0}%
advantage over the other. On what, then, does Lutheran theology ground its
advantage over the Reformed? Not on its consistent character (for it itself
admits its inconsistency); nor directly on its inconsistency, but on what is
humbling in the inconsistency, and so on humility. When the Calvinist is
inconsistent, the Lutheran is permitted to use his reason, appeal to logic and
hold to consistency, for what is the Calvinist inconsistency but an
unscientific self-contradiction, a dogmatic prostitution? But when the
Lutheran is not merely caught, but manifestly catches himself, in notorious
inconsistency: what is it then? Why, then it is not self-contradiction, not
scientific prostitution, but humility, pure, sheer, theological humility. And
what, then, is the Calvinist's consistency? Naturally pride, sheer pride, the
ungodly conceit of uncircumcised reason, of the unregenerate understanding.

If at least we could in the end descry a sign, a mark, an unfailing
distinguishing mark by which one were able in a scientific manner to
distinguish between the unscientific, and so prostituting, and the humble,
and so not prostituting, inconsistencies: then shipwrecked theology might
perhaps still hope for a plank of rescue; but where do we find such a
distinguishing mark? ``In Holy Scripture,'' the Lutheran answers; ``for Holy
Scripture is God's word, and it stands expressly in Holy Scripture: God will
have all men to be saved and to come unto the knowledge of the truth.'' The
Calvinist's pride is thus seen from this, that his consistency is in conflict
with God's word. But it also stands expressly in Holy Scripture that ``God
hath mercy on whom he will have mercy, and whom he will he hardeneth.'' If we
now proceed from the latter passage, then the Calvinist's consistency becomes
humble, while the Lutheran's inconsistency,
% --- p. 250 ---
\opage{250}\setcounter{footnote}{0}%
by coming into open conflict with the word of Scripture: ``Nay but, O man, who
art thou that repliest against God? Shall the thing formed say to him that
formed it, Why hast thou made me thus? Hath not the potter power over the
clay, of the same lump to make one vessel unto honor, and another unto
dishonor?'' --- is transformed and becomes a sign, not of humility, but of
pride, \emph{the} most outrageous pride, the pride of wishing to reply against
% Danish "dét" (acute doubtful in the print, perhaps a speck); rendered with
% stress on "the".
God and deny him the right from eternity to elect --- ``one vessel unto honor,
another unto dishonor.'' This must suffice to illuminate the original
fundamental contradiction of orthodox theology.

This fundamental contradiction must be kept sharply in view if one would
understand the relation between orthodox theology and the rationalism that
followed.

Rationalism, however much it conflicts with orthodox doctrine, has developed
as a scientific consequence of orthodoxy's own method. Had orthodoxy not
called to its aid what it called reason, captive reason, so that it fell into
the contradiction of making use of reason and yet holding it captive,
rationalism would never have penetrated into the Church. This, as already
indicated, I have sought to clarify in a little work, ``Om Holbergs
Kirkehistorie og Theologie, et Bidrag fra Fortiden til at belyse Nutiden,''
from which work I will further bring forward the following: All knowledge and
science is in its essence rational; here no confessional charters are
tolerated, no high-priestly guardianship \ldots{} without reason the Church
could not possibly erect its edifice of doctrine; reason was thus the real
master builder, and the master builder the only one who knew the technique
and could assess the durability of the building. Certainly reason had not
from the beginning erected the great building either on its own initiative or
on its own account; it was, on the contrary, forced
% --- p. 251 ---
\opage{251}\setcounter{footnote}{0}%
labor and at faith's expense; the faith of the Church was, if we may use this
image, the real owner of the building. Had the owner now wished to secure
himself forever against the demolition of the building, then faith, like that
barbarous despot who, by having the master murdered, secured himself against
another still more glorious masterpiece being erected, would have had to
secure itself against reason's tearing down the old half-timbering, clearing
the site and erecting a new, solidly walled building in a new style, by
getting reason strangled. But for such a misdeed the faith of the Church,
however despotically it otherwise behaved, had neither courage nor strength.
So it came about that reason, as soon as it had finished the edifice of
doctrine on faith's account, immediately began to tear it down again on its
own account.

Since now, as is well known, work goes on inside a building just as much when
it is being torn down as when it is being built up, it easily looked to those
standing outside as if the work of demolition which reason carried out on the
old theological building at the end of the 18th and the beginning of the 19th
century was the work not merely of the theologians but of theology, as if the
process of dissolution that destroyed theology were a reformation that
theology, with inner consistency and from its own scientific impulse,
accomplished upon itself: a great literary illusion of the senses! The time of
theology lies behind Holberg; the busy thoughts we meet in him are a
demolition corps that he has just assembled and set to work. When, then, the
reason that had begun its work on its own outside the theological faculty
gradually, though not without considerable resistance and all manner of
trouble, obtained an appointment in the faculty, what happened then? Did it
perhaps end with the theologians, after having let reason continue the work
of dissolution for somewhat more than a generation in the
% --- p. 252 ---
\opage{252}\setcounter{footnote}{0}%
faculty itself, arriving at more positive results than Holberg? No, it ended
with a purging of the old leaven of orthodoxy down to the last remnant; it
ended with a pure, clear, altogether water-clear rationalism.

In the spiritless sensibleness of rationalism the life of faith perished.
Therefore on the side of faith what was needed was a break, not merely with a
worldliness in which a doctrine of utility smeared out into trivialities had
effaced every trace of the ideal, but first and last with a churchliness whose
edifying discourses might well distinguish themselves even by recommending
the use of church towers as windmills and, quite con amore, varying Master
Mannegrim's favorite themes: superstition and crop rotation and stall-feeding
with potatoes. On these terms the protests had to become sharp, and even the
judgment on Holberg fall stern and severe; for, it cannot be denied: with
Holberg rationalism makes its appearance in Denmark.

At the beginning of our century what was needed was something more than
becoming a little witty in a spiritless age; what was needed was to seize the
banner of faith, break with the nature of unbelief and make the break so deep
and incurable that it could never again be healed. So a break was made with
all the powers of the age, made for faith's sake in the name of the
congregation; it came to battle, inspired battle for the right of life, for
the light of spirit, ``for the riddle in the word divine that creates what it
names, that fills the valleys round about on
earth, and levels the rocks.'' What was needed was not to come to terms with
the formal enlightenment, adorn oneself with the charms of reason and get
oneself a halo by recommending a new middle way; no, the time of middle ways
was past; what was needed was a breakthrough. Had it been the task, in an
ingenious, supple manner agreeable to the cultured, to smuggle in a quantum
satis of the
% --- p. 253 ---
\opage{253}\setcounter{footnote}{0}%
old faith without offending the new reason: then a tasteful oratorical fusion
of much that is subjective with something objective, a little mysticism and a
little philosophy, a little John and a little Jacobi, would have been fully
sufficient. But here something more was at stake; what was needed was
precisely not to glide over knotty places, but to press on the plow and plow
so deep that the sprouting seed-corn could take root, and the congregation
see its fruitful ground with the holy springs. The enlightened age, to be
sure, held that a reawakening of the simple faith of the fathers was as much
a popular impossibility as a scientific absurdity; but the enlightened age was
mistaken. The old faith rose up ``like ears of corn in the field, like May in
the beech woods''; it really rose up again, to everyone's astonishment. An
Easter morning, yes, how incomprehensible, an Easter morning gave the dead
life: ``an Easter morning came and took --- the black cross from the grave.''

In the last words connoisseurs will hear an echo of the Grundtvigian tones and
see the mighty figure in the background. Thus the break has taken place,
essentially and from the ground up. With the break, dualism is posited; but
the personal unity of consciousness is by no means broken on that account. Far
from it! It is precisely the highest intellectual task of personality to
recognize the worth that must be ascribed to science and yet hold fast the
ideal demands that proceed from faith. The investigations belonging here are
to be carried out not in any theology --- for theology must go --- but rather
in a philosophy of religion.

``But what, then, is philosophy of religion? To what sphere must it be
reckoned? Insofar as philosophy of religion is a science, it must naturally
belong to the domain of knowledge; but insofar as it has religion for its
% --- p. 254 ---
\opage{254}\setcounter{footnote}{0}%
content, it must orient us in the sphere of faith. By this two-sided task its
two-sided character as a \emph{boundary science} is determined. The
philosophizing consciousness is, as goes without saying, the subjective
medium in which the heterogeneous concepts meet; from this it follows that the
unity of consciousness must presuppose a formal unity of concept corresponding
to the determination of the boundary, that, in other words, the absolutely
heterogeneous concepts must in the boundary itself become formally
homogeneous. What such a formal homogeneity has to signify can easily be
shown. We speak of absolutely heterogeneous principles, but a principle is,
after all, \textit{in abstracto} always a principle. We speak of heterogeneous
cognitions, heterogeneous concepts, etc.; but when we abstract from what is
peculiar, cognition $=$ cognition, concept $=$ concept.

As a boundary science, philosophy of religion gets its concepts of faith from
the life of faith and its concepts of knowledge from science. In the latter
respect it becomes apparent what a logical theism has to signify; in the
former respect it depends on religious psychology. It is of no use for the
philosopher of religion to try to help science along with assurances about his
own faith, his own individual experiences. The question here is not where the
philosopher has his experiences from; the question here is the psychological
value of the experiences in their relation to the principle of faith. The
apparent contradiction, that philosophy of religion on the conditions set out
must become a science that treats both of what falls \emph{outside} and of
what falls \emph{inside} science, is easily removed. Should anyone maintain
that there cannot possibly be a knowledge of a cognition of faith when the
cognition of faith by its nature is to fall outside knowledge, let him then
begin at the beginning and maintain the impossibility of there being a
% --- p. 255 ---
\opage{255}\setcounter{footnote}{0}%
consciousness of the earth when the real earth by its nature is to fall
outside consciousness. An objection such as that the concepts of faith, by
being reflected in knowledge, must absolutely become concepts of knowledge,
signifies nothing; the task of science is precisely to know the heterogeneous
as heterogeneous.'' (Den gode Villie. pp. 136---37).

\begin{center}\rule{3cm}{0.4pt}\end{center}

\chapter*{XV.}
\addcontentsline{toc}{chapter}{Lecture XV}
\markboth{Lecture XV}{Lecture XV}
\label{ch:xv}
% --- p. 256 ---
\opage{256}\setcounter{footnote}{0}%
As often as there is talk of the spiritual life, its conditions, requirements
and obstacles, in particular of the conflict between faith and knowledge, so
often must we think of Luther's words: Faith is a restless thing. And which of
the Church's great teachers could say this with stronger emphasis than Luther
himself? It was faith that gave him peace! Yes, but faith also prepared a
struggle for him, a struggle that lasted his whole life through, a struggle
first in the inward and then in the outward as well. In the darkness of night,
in dread of the world and of himself, the young man flees to the Augustinian
monastery in Erfurt, knocks and is let in. One would now think that he who,
with a thunderstorm behind him, went in at the monastery door escaped the
storms of life, the unrest of the world, and entered into peace; but it was
not so. He entered not into peace but into strife, the most terrible strife,
strife with himself and with God, strife in conscience, strife with doubt in
order to win faith. And when he had won faith, when he recognized that it was
not by papal good works but by faith alone that a man is justified, how does
he then characterize faith? Does he perhaps say that in faith there is nothing
but peace, and that it costs strife and unrest only before one reaches faith?
No, of the faith that
% --- p. 257 ---
\opage{257}\setcounter{footnote}{0}%
gives peace, of just that faith, he says that it is a restless thing. And it
went on being so for him. For it was not enough for him to have won peace in
his own inner being. What had given him inward peace wanted to come forth on
his lips, wanted to be proclaimed with his mighty voice; then the strife began
in the outward. When the tempest then came in the outward, when for faith's
sake he found himself in the midst of the stream of events, in the midst of the
wild struggle, roared about by storms, then he surely had to say with
strengthened emphasis: Faith is a restless thing.

Luther's words about faith as a restless thing are confirmed in the history of
the Church from first to last. It is known to those versed in the matter how
the ancient Church under bloody persecutions fought its way forward to
victory. In the time of dangers, when the persecutions truly raged against the
Christians, then the strength of faith grew, then more and more came forward
with ever-rising enthusiasm, not only vigorous men but youths, women and old
men, all with frank confession, all prepared to sacrifice their lives in the
name of the Crucified. If, on the other hand, the Church had had peace for a
time, then slackness was felt. The longer the interval had been in which
faith was, as it were, on the point of becoming a restful thing, so that one
could arrange oneself with it too for comfort and well-being, the less
prepared one was when the storm suddenly broke loose. Under the Decian and
later under the Diocletian persecutions so many fell away precisely because
the Church had had peace too long.

That faith is and must always be a restless thing follows directly from its
being the highest personal concern and, as the highest personal concern,
grounded in a dualistic consciousness. Or could perhaps the opposition between
the worldly and the holy, between the flesh and the spirit, between the old
and the new man, ever in the world cease to be dualistic? It is
% --- p. 258 ---
\opage{258}\setcounter{footnote}{0}%
another matter that the dualism can show itself differently at different
times, and the personal strife in consequence take on a different appearance.

It makes, as regards the restlessness, a very considerable difference whether
the strife is outward-going or merely inward-going, whether it is waged on
barbarous, half-barbarous or humane terms.

The time of the persecutions was for the first confessors a barbarous time;
under the first immediate opposition between Christianity and paganism the
decisive trial of faith had necessarily to be: to be able to sacrifice one's
life for the confession of Christ's name; the ideal was to become a
blood-witness. So glowing was the spirit, so flaming the holy passion, that
the craving for the martyr's crown even became a kind of vain point of honor.
The Redeemer's words: ``Fear not them which kill the body, but are not able to
kill the soul,'' found ample fulfillment: the young ``soldiers'' of faith went
bravely into the fire. But however magnificent the testimonies to the power of
spirit and the victory of freedom that came forth in this way, the fulfillment
of the first elementary fundamental demand was nevertheless not sufficient for
an all-round development of the spiritual life as a whole. An unceasing
continuation of persecutions and blood-witnesses would surely have led, not to
higher personal self-understanding, but to wild enthusiasm and dark
fanaticism. The excesses of the Montanists and the Donatists bear witness to
this.

It was, then, for the sake of the development of spirit, to the benefit of
religion as much as of culture, of the new man as well as of the old, that the
Church at last got outward peace.

But merely outward peace could not bring rest, for faith is a restless thing.
The outward resistance in inflamed paganism now transformed itself into an
inward resistance in heresies and sectarianism. If
% --- p. 259 ---
\opage{259}\setcounter{footnote}{0}%
during the persecutions it had been the most essential condition for eternal
salvation to be able to sacrifice one's life for the faith, then now, after the
persecutions had ceased, it became an inescapable condition to secure the
separation between the true and the false faith. Hereby not only the passions
but also the thoughts were set in motion. The violent dogmatic disputes in the
Greek Church from the fourth to the eighth century were all waged on
half-barbarous terms. The jealousy of the bishops, the fanaticism of the
monks, the passions of the people and the emperor's words of power give the
course of development so worldly a direction that the Church's doctrinal
concept seems bound to become a simple product of human arbitrariness; and yet
a development of ideas takes place. As the many voices are heard, and all
means are employed, arbitrariness falls away: through the din of strife the
spirit of the Church moves forward with inner assurance toward a consistent
shaping of the content of faith.

It is a first step on the way of thinking for oneself that the Church gets its
propositions of faith set up and brought into a system. But precisely with
this first step it made a false step that brought it to apostasy from the
apostolic way of thinking: in the further formation of the doctrinal concept
it came onto the way of objective knowledge, it got theology founded.

Through theology the doctrine became outwardly objective in the same sense as a
positive legislation can be said to be outwardly objective; through the favor
that Christianity on the whole enjoyed from the emperors from the time of
Constantine the Great, the Church's position in the State likewise became
outwardly objective: the emphasis falls in all directions on the objective,
the outwardly objective.

The main thing in the outwardly objective is that the Church
gets a doctrinal concept of dogmas fixed once and for all; for by such a
doctrinal concept its ideal
% --- p. 260 ---
\opage{260}\setcounter{footnote}{0}%
subsistence ``through changing ages'' is far better secured than by the
congregations' ``uncertain and fluctuating, poorly guaranteed'' appropriation
of the spirit of the Gospel according to the merely apostolic confession of
faith. The main thing is --- in the outwardly objective --- ``that the
Church's revenues are in grain value according to the capitular price
assessment, and that the Church possesses landed estates''; for by the
Church's possessing landed estates ``its material subsistence through changing
ages is far better secured than by uncertain and fluctuating, poorly
guaranteed money matters.'' The main thing is that the whole ecclesiastical
system, spiritual as well as temporal, comes to rest in the objective. But
faith is a restless thing.

Why, now, was it precisely necessary that to the apostolic confession of faith
there should be added the Nicene, concerning the Son's unity of essence with,
and eternal generation from, the Father; to this in turn the
Constantinopolitan, concerning the three hypostases in one essence; to this in
turn the Chalcedonian, concerning Christ's two natures, which must neither
merge together nor keep apart; to this in turn the confession of Christ's two
wills, etc.? Why, we say, was such a heaping up of subtle dogmatic
determinations, whose interpretation and grounding had to require superhuman
acuteness, necessary at all, unless to give faith a truly intricate,
entangled and learned foundation? Why should faith have so intricate,
entangled and learned a foundation? Naturally for the sake of objectivity, so
that the teachers, standing on this objectively firm foundation, could through
changing ages with force and vigor interpret to the simple believers Christ's
words: I thank thee, O Father, because thou hast hid these things from the
wise and prudent, and hast revealed them unto babes! And why must the Church
now absolutely have its revenues in grain value according to the capitular
price assessment? Why must it
% --- p. 261 ---
\opage{261}\setcounter{footnote}{0}%
for its material subsistence necessarily possess landed estates (Patrimonium
Petri)? Naturally for the sake of objectivity, so that Christ's vicars
through changing ages might know themselves secured! How should the preachers
of the Word be able to give the believers comfort against cares for temporal
things if they did not themselves feel secured against temporal cares; but how
should they feel secured, how should they with full confidence interpret the
words of the Gospel: ``Take no thought for the morrow, for the morrow shall
take thought for the things of itself,'' if they did not know within
themselves that through changing ages they received their ecclesiastical
revenues in grain value according to the capitular price assessment?
``Consider the lilies of the field!'' Yes, but when the Church has no landed
estates, then it has no fields, and when it has no fields, then it has no
lilies either. Were a preacher of the Gospel through changing ages, without
fields and without lilies, to be secured only by uncertain and fluctuating,
poorly guaranteed money matters, how then could he with a calm mind go out and
consider the lilies of the field!

There are different ways of considering the lilies of the field. It is one
thing to consider them from the standpoint of the Gospel, another to consider
them --- according to the capitular price assessment; one is the
consideration of the new man, another that of the old; the two considerations
are absolutely heterogeneous; the absolutely heterogeneous considerations make
consciousness dualistic! the dualistic consciousness makes faith a restless
thing.

But is it not, then, a law of civilization that human society is to work its
way forward to order and rest? We answer: it is not only a law of
civilization but also a law of religion; only one must
% --- p. 262 ---
\opage{262}\setcounter{footnote}{0}%
fundamentally distinguish between order and rest in civil
% printed "af gjøre" (wrong sort) for "at gjøre"; translated in the intended sense.
society and order and rest in religious society. Just because in religious
society a different spirit must reign from the one that reigns on human terms
in civil society, order and rest must also be brought about by other means and
on other conditions in religious than in civil affairs. In civil society laws
of compulsion hold, in religious society laws of freedom. It is a crude
confusion of laws of freedom with laws of compulsion that makes the medieval
spirit of the Church half-barbarous.

On the hierarchy and the papal power the great order of things, civil as well
as religious, was to be grounded in the West. All, great and small, laymen and
learned, were religiously and civilly interested in there being able to rise
above the confusion an ordering, governing power, an established order capable
of holding unrest down and, in human fashion, in the name of the holy, of
stifling the revolt of thoughts; but faith is a restless thing. In the name of
faith and for faith's sake the papal Church got the great tracts of land, the
many temporal goods with their many glories, so that Christ's kingdom might
truly be glorified in the visible; but the enjoyment of these glories had such
questionable consequences for the faith of the clergy that an Arnold of
Brescia took it deeply to heart and the Fraticelli despaired over it. Unrest
came, moreover, from another side. In the name of faith and for faith's sake
the Church at last got such supremacy over the State that Innocent III could
compare the papacy with the sun and the empire with the moon, which gets its
light from the sun. But the same tendency of spirit that gave the Church
political power also tempted the hierarchy to misuse the power. When the Pope,
in passionate struggle against Germany's medieval Emperor,
% --- p. 263 ---
\opage{263}\setcounter{footnote}{0}%
had struck the Hohenstaufens to the ground, he came, as is well known, into
dependence on France, which was striving politically forward toward a new
order of things, and as a consequence had to bring himself to leave his
eternal Rome and move to Avignon. Thus began what church history calls the
Babylonian captivity of the Pope, or the Exile. The consequence of the Exile
was the great papal Schism; the consequence of the Schism was an enormous
offense. Pope stood up against Pope, Pope condemned Pope; to whom was
Christendom to hold? Yes, faith is a restless thing. The unrest went not only
outward in disputes about the worldly, it went also inward in disputes about
the spiritual, in the mutual strife of thoughts about the eternal truths. The
thoughts of knowledge from a classical, pagan culture were to be reconciled
with the thoughts of faith in the Church's dogmas; Greek philosophy was, in
the name of theology, to spread the light of reason over the mysteries of
revelation; Aristotle, a Christian before Christ (Christianus ante Christum),
was supposed already to have comprehended everything in the Church's doctrine
that could be comprehended. So there was reasoning and disputing; in the
universities of scholasticism and dialectic men contended and fenced in
knightly fashion as at a tournament. During this continued feud many were
restless, some on behalf of knowledge, others on behalf of faith.

The unrest increased when at last, after the conquest of Constantinople, amid
the steadily growing interest in science that characterizes the period of
Italy's history called the age of the humanists, men more and more immersed
themselves in the classics, and so also in Plato and Aristotle. Then it became
clear that here there must be two systems of cognition: one in philosophy, the
Greek-classical, grounded on reason; another in theology, the
orthodox-ecclesiastical, resting on the authority of the Pope. These two could
not be reconciled; therefore certain thinkers, who
% --- p. 264 ---
\opage{264}\setcounter{footnote}{0}%
did not dare break with the Church while on the other hand they would not give
up philosophy either, sought to escape the impending danger by proclaiming a
doctrine of two truths, a philosophical and a theological. If only they might
have their philosophical truths in peace, they were willing in every way, as
regards authority and faith, to bow beneath the doctrine of the Church. So
Pomponatius in the 16th century and Vanini in the 17th. But it did not work.
Pomponatius discovered that the classical Aristotle had quite another doctrine
of the nature and immortality of the soul than the scholastic one. The
scholastic Aristotle was supposed to agree with the doctrine of the Church and
furnish proof of the immortality of the soul; the classical one, on the
contrary, furnished precisely proof of the mortality of the natural soul. The
Inquisition was set in motion; Pomponatius was saved through the intervention
of powerful friends (cardinals), but Vanini was burned. They burned the
confessor of the two truths, and at that point in time it was quite natural.

It is understandable that the Catholic Church could not be content with that
division of truth; for our part, we cannot either. If the division were to
mean anything, it would have to be understood thus, that only philosophy was
science and theology a self-contradictory monstrosity; but for such a break
with the dogmatic science of miracle the time had not yet come. That conflict,
then, was also soon succeeded by one far more deep-going. It came so far that
the truth of faith became divided against itself at the Reformation. What
happens? Out of the old Church's bosom a new Church proceeds, and yet the old
Church does not die. So Church stands against Church, each with its doctrine
of faith. What is true in Protestantism is untrue in Catholicism, and
conversely.
% --- p. 265 ---
\opage{265}\setcounter{footnote}{0}%
If the Catholic's faith was true, then the Protestant's was untrue. So one
could then see that faith is a restless thing.

The unrest went further still. As the evangelical Protestant Church broke with
the papal Catholic one, it was split again dualistically into two great
oppositions: the Lutheran and the Reformed type, with the traditions belonging
to them. Indeed, the splitting went so far that the Catholic Church even
became divided against itself, as the Jesuit Pelagianizing tendency set about
stifling the Augustinian. The forces were set in motion: assertion stood
against assertion, passion against passion, word of power against word of
power; but the principle of faith itself was not clarified, the thought of
faith not developed.

The evangelical Protestant fundamental thought, as it was apprehended and held
fast by Luther himself, would surely have been carried through far more
consistently, and the conflict with the papists as well as with the Reformed
would have led to a different result, had not the theological idée fixe
intervened that the dogmas ought after all, at least to a certain degree, to
be explained and grounded scientifically. We see it from the disputes over the
Supper. Against the Catholic doctrine of transubstantiation Luther set the
words of institution. But then Carlstadt comes forward and twists the words of
institution, because he cannot comprehend how it is possible that Christ's
body and blood can be in bread and wine, whereas he can quite well comprehend
that Christ, pointing to his body, may have said: this is my body, and
thereupon, pointing to the bread, added: take this --- the bread, which is not
my body --- etc. What Luther would answer to this had already been stated in
his work on the Babylonian captivity of the Church, in which he says: ``Truly,
if I cannot comprehend how the bread is Christ's body, yet I will take my
understanding captive and hold simply and solely to his word.'' What he adds
by way of
% --- p. 266 ---
\opage{266}\setcounter{footnote}{0}%
elucidation: ``If iron and fire, two substances, are so united in the one
glowing iron that every part is iron and fire, why then should not the body of
the almighty God-man, the glorified Lord, much more be able to be in the
substance of the consecrated bread?'' is certainly not put forward as any
dogmatic grounding. ``No reason,'' says Luther in his ``Manifesto against the
heavenly prophets,'' ``is so feeble that it would not rather believe that
plain bread and wine are present here than that Christ's flesh and blood are
hidden here. For that one needs no spirit; it is easy for anyone to believe,
and nothing more is needed here than that someone who has a little standing is
bold enough to preach it, and he already has disciples enough. But if one will
deal with our faith in such a way that we first put our own preconceived
opinion into Scripture and then twist it according to our mind and look only
to what is straightforward for the common darkened understanding, then no
article of faith will remain standing; for there is not one in Scripture that
God has not set above reason.'' --- So here, then, according to Luther's
express testimony, there is dualism of faith and reason: faith is a restless
thing.

Luther was right, he felt it deeply; it was this feeling that made him so
strong in confession, so bitter in strife. Yes, certainly Luther was right;
modern Christendom is hardly conscious of what the Church owes to the man who
on this point stood so hard upon the words of the text and defied a world with
the Gospel of faith. Yes, Luther was right, yet not as theologian but as
priest. His conception of the Supper is explained precisely by his having
originally seen the sacrament not with the theologian's but with the priest's
eye. What Luther as theologian put forward about the doctrine of the Supper is
significant only for his time; his theologumena
% --- p. 267 ---
\opage{267}\setcounter{footnote}{0}%
are vanishing, but his priestly testimonies will remain as long as faith is
--- faith.

Luther himself was not clear about the relation in principle between religion
and theology. The new doctrine needed the assistance of science; Melanchthon
was a theologian, and the theologian then thought that Calvin was right, for
Calvin was, after all, a theologian too. But the Lutheran Church, which was
unconditionally right in faith, now wanted at the same time to be
unconditionally right in theology as well, and so it went from bad to worse.
Lutheran theology worked itself weary in unspeculative speculations, and has
long since had to see its favorite theories of \textit{communication
idiomatum and ubiquitas corporis Christi} given
% "communication" as printed (for "communicatio").
over to dissolving criticism. (Evangelietr.\ og Theolog.\ pp. 117---118.)

Already at the end of the 17th century there had formed \textit{two
altogether opposed views; according to the one of these one should give up
outright the connection of faith and reason, religion and science; according
to the other one should defend and assert it by all possible scientific
means}. The first-named view in turn took two directions, in that on the one
side one confined oneself to operating scientifically by pointing out
contradictions and awakening doubt, on the other side conceived religion
practically as life and relied on feeling: the representative of the first
direction is Bayle, of the latter Spener.

In Bayle's historical and critical dictionary (Dictionaire historique et
critique) the conflict between faith and knowledge, religion and science, is
pointed out with great acuteness. The dogma wants to be believed but not
comprehended; \emph{it is and must be a mystery to human reason}. ``The
secrets of revelation belong to a supernatural realm; they rest on the highest
authority,
% --- p. 268 ---
\opage{268}\setcounter{footnote}{0}%
the authority of God, and he has not set them before us so that we should
comprehend them, but so that we should believe them with all the humble
submission that we owe to the infallible highest Being. The incompatibility of
revelation and philosophy makes it necessary to choose between the two and to
decide for one of them. He who will believe nothing but what is clear and in
agreement with common notions, let him seize philosophy and let Christianity
go; for it is impossible to possess clarity and incomprehensibility at once;
the union of these two things is no less impracticable than finding the
quadrature of the circle. He for whom a round table is not convenient, let him
have a square one made; but he must not demand that one and the same table
shall afford him the conveniences of both a round and a square table.'' Bayle
pleads with equal facility now the cause of faith, now that of reason. In
favor of faith he adduces: ``In order not to grasp blindly in the choice, one
must make oneself closely acquainted with the nature of both sides, of faith
and of science. At first glance, when one hears that faith is a certainty
whose object remains obscure, but that science, on the other hand, besides
giving certainty, also makes what is known clear and distinct, it would seem
as if the decided advantage were on the side of science. But science comes
about through reason, and what if reason should not be equal to such a task?
But this is in reality the case; reason is only a principle of destruction,
not of edification; it knows better what things are not than what they are; it
is fit only for awakening doubt and for twisting to right and left so as to
draw out the dispute into infinity. It is a vagabond who cannot come to rest,
a Penelope who continually unravels her own
% --- p. 269 ---
\opage{269}\setcounter{footnote}{0}%
web. It is good only for bringing to men's consciousness the darkness in which
they grope, their incapacity to work their way out of it, and the need for a
revelation.'' For reason he adduces: ``There are principles against which
doctrines of Holy Scripture, however expressly and definitely stated, would be
unable to accomplish anything, as for example the principles that the whole
is greater than a part of it, that two mutually opposed propositions cannot
possibly both be true, and the like. This the orthodox theologians of all
confessions also concede, in that they do not admit that the Trinity and the
Incarnation are dogmas that contain a self-contradiction; whereby they in
fact acknowledge reason as the highest court, since they do not believe that a
dogma is sufficiently secure unless it is, so to speak, registered and
formally ratified by the supreme court of reason. God willed to bestow on the
soul an infallible means of distinguishing the true from the false, and this
means is the natural light, the metaphysical principles, in which we possess
an original rule by which we can distinguish whether the various doctrines we
find in books or hear expounded by our teachers are correct or incorrect. The
only criterion of truth we therefore have in the agreement of objects with
this universal and original light which God has poured out into the souls of
all men.''

One would now suppose that such deeply penetrating discussions must especially
have attracted the attention of the theologians, since it is manifestly the
scientific standing of theology that is at stake here. This was indeed also
partly the case in other countries, notably in Holland, where a Clericus came
forward against Bayle; but it was by no means the case in Denmark. In Denmark
the theologians, after all, had enough on their hands day to day; what should
they want with investigations of principle in general science?
% --- p. 270 ---
\opage{270}\setcounter{footnote}{0}%
Bayle's scribblings were, after all, an ungodly business, some philosophical
trumpery that did not concern God-fearing men. Certainly they had an uneasy
feeling that the time was not favorable to them, that something here was
amiss; but where the trouble lay, they did not comprehend. It must be the
wickedness of the world, the unbelief of the generation and the age, the pride
of reason; so they easily agreed to regard the many uneasy phenomena, however
diverse these otherwise were, as manifestations of one and the same ungodly
nature. Whatever did not agree with their theological views was, under the
most diverse designations, lumped together in one single great heap of
heretics. Naturalists, indifferentists, syncretists, atheists, pietists ---
Pontoppidan, out of sheer godliness, does not even dare write the blasphemous
names out in full, but abbreviates them to A\,.\,.\,.ists and P\,.\,.ists ---
fanatics, enthusiasts, deists and rationalists fly about one's ears when one
starts stirring in the heap. So the theologians went round in the old
treadmill, and the old treadmill went as best it could. (Om Holbergs
Kirkehistorie og Theol.\ pp. 41---43.)

If the difficulties could be blown away by the theologians' shutting their
eyes so as not to see them, then they would long since have been blown away;
so the partridge no doubt thinks that by sticking its head in a hole in the
stubble it becomes invisible to the hunter; but fact is always fact. If
theology does not have its opposition outside itself, then it has it within
itself. The old orthodoxy held on to faith by fighting against reason; that
was dualism. Theological rationalism held on to reason by fighting against
faith; that too was dualism. Dualism recurs under the most diverse shapes;
even the separation between faith and faith, between
% --- p. 271 ---
\opage{271}\setcounter{footnote}{0}%
faith in revelation and faith of reason, is, after all, dualistic. We see it
in Strausz.

% The quotation opened here closes on p. 273; the inner quotation of Strauss
% closes on this page.
``Although Strausz, by assuming a faith of reason, is not yet quite clear
about the relation between rationalism and philosophy, between the religion of
reason and science, it is nevertheless clear that his faith of reason suits
only those he calls the knowing, not the believing. `Here,' he says, `is a
gulf fixed between two classes of human society, the knowing and the people,
i.e.\ those who do not philosophize, in the higher as well as in the lower
estates, and it will perhaps never be filled\,... So let the believer let the
knower, as the latter the former, go his way in peace; we let them keep their
faith, so they must also let us keep our philosophy; and if the overweeningly
pious should succeed in excluding us from their church, we shall count it a
gain; of false attempts at union enough has now been done; only a sharpening of
the oppositions can lead further.' With this the dualism, the dualism between
religion and science pointed out by Bayle, and with it the invalidity of
theology, is pronounced as the result of our century's scientific testing.

And with that, our learned theologians will perhaps say, it is supposed to be
proved that the death sentence on theology has been irrevocably passed. Who
does not now see that the juxtaposition of Bayle and Strausz proves precisely
the opposite. Bayle, and with him numerous freethinkers, were in their time
just as sure of the so-called dualism, i.e.\ of the insurmountable gulf between
faith and knowledge, as Strausz, with numerous other freethinkers both hidden
and open, may be in our time. Nevertheless it came about that faith and
knowledge, under a constantly advancing theological science, went on working
together a good century after Bayle: what now, we ask, should
% --- p. 272 ---
\opage{272}\setcounter{footnote}{0}%
prevent faith and knowledge, once the passing day's paradoxical squalls with
``stray thoughts, aphorisms, notions and flashes'' have gone by and it clears
up again in the heaven of the sciences, from continuing to work together on
new terms in a new century and giving theology a new shape? Experience surely
teaches that one should not be in a hurry to declare a scientific standpoint,
still less a science that has the prescription of centuries, something
absolutely overcome: \textit{multa renascuntur, quæ jam cecidere, caduntque,
quæ nunc sunt} etc. --- Reflections of this nature strongly recall the
well-known: \textit{qui nimium probat, nihil probat}. By adducing limited
truths in limitless generality one ends either in the meaningless or in the
trivial. Insofar as all things in the world go up and down, theology must
naturally go up and down. --- Theological consolations of this sort strongly
recall the people's wise fool, who always laughed when the road went uphill,
because after an adversity going up there must follow a prosperity going down;
or, perhaps better still, the drunken poet who, when everything was going
round, stood still with the key and waited for the door to come to him in
due course, without his needing, in his infallibility, to change standpoint.
Theology has long gone up and down, down and up, as in a swing; but precisely
because theology has sat so long in the swing, it has now at last, despite the
many alternations in movements and oscillations, become possible to trace its
path and determine the equation of the curve.

If we compare the dualism of religion and science as it appears in Bayle with
the dualism as it has come to a breakthrough in more recent times, then it is
clear that the latter is more than a simple repetition of the former. When
Bayle set religion and philosophy against each other, he took, without further
ado,
% --- p. 273 ---
\opage{273}\setcounter{footnote}{0}%
religion as it then existed in the Church's confessional writings, in
theologically positive, determinate form. In his time it had not yet become
clear to the thinker what would happen when one broke the crystallized forms,
melted down the religious and recast the content in another form; it had, in
other words, not yet dawned on scientific consciousness how the matter would
look when the content of faith took on the form of the concept. To carry out
the labors of this great experiment in melting down and recasting required no
less than a whole century, the period between Bayle and Hegel. Hegel's
philosophy of religion is the deepest consequence of rationalism, its
transformation into science proper. But scarcely was the transformation
carried through before it stood clear to the knower as to the believer that
faith is not knowledge, religion not science. The last glimmer of faith of
reason, an afterglow of the old rationalism, which Strausz still let remain in
his doctrine of faith, has vanished in the name of knowledge in
L. Feuerbach, in the name of faith in S. Kierkegaard. Kierkegaard and
Feuerbach, however diametrically opposed they are, are both agreed in
declaring the characters of religion and science incompatible.'' --- Even if
knowledge can come to rest in the idea, faith nevertheless cannot come to rest
in knowledge: faith is a restless thing.

But --- it is asked --- will not the whole ecclesiastical edifice of doctrine
collapse; will not the religious representations fall into ferment and much
confusion; will not the foundation necessary for the clergy's more thorough
culture, higher development of spirit and freer overview of the people's
religious and civic conditions be undermined, when the scientific preparatory
school is closed, when, in other words, theology is thrown on the scrap heap?
No, let us think of a reform of theological study,
% --- p. 274 ---
\opage{274}\setcounter{footnote}{0}%
if such a thing is really needed at present; there is sense in that; but to
want to abolish theology as theology is a crazy notion, a project that may
perhaps find acceptance with those who want disturbance and unrest, but not
with those who want peace and order.

To meet such objections, we must first and last be agreed with ourselves
about what theology is, or rather, gives itself out to be: that theology
expressly gives itself out to be a science, an outright science of faith, a
science of the supernatural, the suprascientific, the miracle, that it is thus
a self-contradictory science of miracle; that what is self-contradictory in
theological science is only veiled and concealed by its having an apparatus
of historical, philological, dogmatic, metaphysical, semi-metaphysical
components for use and critical treatment, insofar as there can be talk of
objective use by subjective judgment and critical treatment without principle.
This is the contradiction, an age-old, ingrown fundamental contradiction,
which must be completely laid bare and brought forth if we are to understand
properly what theology is.

``Although philology, criticism and history, which in and for themselves
belong to science, constitute components of theology, they by no means
constitute what is theological in theology. Insofar as the disciplines named
are treated scientifically, they are treated untheologically; insofar, on the
other hand, as they are treated theologically, they are treated
unscientifically. What is theological in theology is known, namely, by this:
that the cause from which all effects in the realm of nature as in that of
spirit are derived and explained, partly immediately, partly mediately, is
sought not in the nature of things or the essence of reason, but in the
absolute will, in the Almighty, who creates out of nothing. But even the
purest, most acute and most consistent application of such a
% --- p. 275 ---
\opage{275}\setcounter{footnote}{0}%
principle (occasionalism; Malebranche) makes science unscientific. --- When a
principle is at issue in the theory of science, the matter cannot be settled
by a more or less, as if someone, by following ``the middle way,'' would
think that one could still get through scientifically by assuming a little of
the supernatural, provided only one did not make too much of it. The question
of a principle is a question of an unconditioned presupposition that
conditions everything; such a question must therefore always be answered with
yes or no. Either, then, theology must once and for all avow supernatural
causality, and has thereby at once repudiated the objectively rational
foundation, or else once and for all avow this foundation, and has thereby at
once renounced every footing in the supernatural. In the first case theology
will thus cease to be science, in the latter case to be theology.'' (Cf.
Philosoph.\ Propædeutik. Copenhagen 1857. pp. 70---80).

A glance at the spiritual conditions of the present age teaches that the
theologians' many-colored dogmatics are no longer able to master the
contradictions breaking in in ever-growing numbers and with renewed strength,
or to give the religious instruction of the people any firm, universally valid
foundation. A revision of the doctrine of faith in its entirety, a development
and presentation of the essence of Christianity that could show how the
personal truths in their mutual connection proceed from one single principle,
valid in itself, must surely be said to be needed as far as religious thinking
is concerned. But if such a revision is to be carried through, then theology
must go; for theology botches the principle. If a younger generation of
popular teachers of religion is to find the right tone, hit the right mark and,
with the spiritual superiority that only assurance in thinking can give the
preachers of the Word, stand up against
% --- p. 276 ---
\opage{276}\setcounter{footnote}{0}%
a reflecting cultural consciousness: then a stricter preparatory school,
expressly designed for the solution of life's tasks, is altogether necessary.
On theological terms such a preparatory school is impossible: if the priests
are to come forward, the theologians must go; for theology botches the
principle.

With the dualism of faith and knowledge there is no danger when those whose
calling it is to proclaim the word of faith have first learned to think the
thoughts of faith. But to think the thoughts of faith is not possible without
essential insight into the corresponding thoughts of knowledge. Even if the
knowing have hitherto neglected to think through the principle of faith, the
spokesmen of faith must above all see to it that they do not in return
neglect to think through the principle of knowledge. Only when it comes so
far that the believing thinker can truthfully say to the knower: I understand
you well enough, but you do not yet understand me --- only then will the
superiority be reached that can be reached by the way of reflection for the
benefit of faith. Will faith then at once be secured? No, for faith is more
than a thinking; its unrest is more than an unrest of thoughts, it is also an
unrest of hearts: faith is a restless thing.

The unrest of faith, however, does not stem from its principle but from the
natural frailty of believing individualities. When even a Paul (Phil.\ III,
13---14) must say to the Philippians: ``Brethren, I count not myself to have
apprehended. But this one thing I do: forgetting those things which are
behind, and reaching forth unto those things which are before, I press toward
the mark,'' he speaks in ``fear and trembling'' of the unrest of faith and of
the struggle of personality for the holy ideal; but when he then, as in the
Epistle to the Romans (VIII, 37---39), adds: ``Nay, in all these things we are
more than conquerors through him that loved us. For I am persuaded, that
neither death, nor life, nor angels, nor principalities, nor powers, nor
things present, nor
% --- p. 277 ---
\opage{277}\setcounter{footnote}{0}%
things to come, nor height, nor depth, nor any other creature, shall be able
to separate us from the love of God, which is in Christ Jesus our Lord'' ---
then he speaks of the assurance in which the unrest is annulled, the ground of
inwardness in which the believing mind finds rest. At no time in the temporal
can the doubleness of rest and unrest disappear without faith itself
disappearing;
% no punctuation printed after "forsvinder"; the sense wants a semicolon.
but the incessantly continuing strife can be waged on ever more humane terms.

From many sides and in many ways the powers of faith have been tried; from the
first persecution of the Christians right up to the latest Jesuit-reactionary
attempts to root out the last shoots of evangelical freedom the trial has been
withstood, often heroically and under the most barbarous ill-treatment; but
only in the present age, the age with the humane conditions, has the
development of spirit come so far that not only the powers of faith but, what
is decisive in relation to our cultural life, the thoughts of faith as well
can be subjected from the ground up to a trial.

Should anyone perhaps, on the pretext that Christianity is \emph{a life}, not
\emph{a theory}, dismiss the intellectual trial as an impractical,
hair-splitting scholastic dispute, lay about him with phrases and assure us
that nothing comes of disputing about faith, that one should at once declare
oneself a believer and fight for faith, etc.: then I know no more fitting way
to conclude than by recalling a statement with which seventeen years ago I
introduced my work \emph{Evangelietroen} og \emph{Theologien}: ``If those days
when the question of the blessed life was the highest question of the race
were to return immediately; if the Church's great jubilee year were now so
near that it rang in our ears, as though we already heard the sound of the
bells that are to ring it in: why, then our awakened ones might well dismiss
`a merely theoretical
% --- p. 278 ---
\opage{278}\setcounter{footnote}{0}%
scholastic dispute,' and direct to life's `trial by fire' those who ---
`instead of acting accordingly!' --- only write about it, insofar as they
write about faith. `But' --- adds a Christian psychologist --- `shall faith be
found on the earth?' .\,.\,. Should it at last come so far that one did not
even find it worth the trouble to \emph{dispute about} faith, would that too be
a good sign? Would it not much rather be a sign that one despised the human
strife because at bottom one did not respect the divine either; a sign that
one found it still less worth the trouble to \emph{fight for} faith: a sign
that faith was no longer found on the earth?''

So, then, when the talk is of the spiritual conditions of the present age, we
must abide by Luther's words: Faith is a restless thing!

\begin{center}\rule{3cm}{0.4pt}\end{center}

\chapter*{XVI.}
\addcontentsline{toc}{chapter}{Lecture XVI}
\markboth{Lecture XVI}{Lecture XVI}
\label{ch:xvi}
% --- p. 279 ---
\opage{279}\setcounter{footnote}{0}%
It is ``a closer scientific cooperation between the Nordic universities''
that I have had the honor of inaugurating with the lectures now delivered.
If, however, one lays particular weight on its being a cooperation between
the universities that is in question, and on this cooperation's being
scientific, a natural objection arises. Why, it will be asked, precisely
where a contribution to a scientific cooperation is at stake, choose a
form of presentation that is not scientific? It is, as it seems, a
contradiction to want to inaugurate something scientific in an unscientific
manner. In conclusion I owe the highly esteemed assembly an account of my
choice.

If a cooperation between the different universities is to mean anything,
if the teachers from these universities are to be able to entertain the
hope of contributing to such a cooperation, then, so far as I can see, two
things must be kept in view. It is, in the first place, natural and
necessary that every university teacher choose his subject from his own
field and in accordance with his scientific position. It is natural that
the jurist choose a juridical subject, the physicist a physical one, the
historian and the philologist a historical and philological one, and so
on. That is my first point. The next is that no visiting teacher will
choose his subject in such a
% --- p. 280 ---
\opage{280}\setcounter{footnote}{0}%
way that he communicates something that must be presupposed to be known to
everyone at the other university. For it surely cannot in truth be called
inaugurating any higher cooperation, that a lecturer from another
university comes and communicates to the students what they already know.
I would thus have to regard it as less fitting to lecture to the Norwegian
students on a survey of some general part of philosophy or a fragment of
propaedeutics, in short, something which, given the thoroughness, purity,
and rigor with which philosophical science is cultivated here, the
students could hear just as well from their appointed teachers. No, to
offer something of that sort to the awakened young men whose lively
interest in a freer development of spirit has in particular set the matter
here under discussion in motion, I would regard not merely as a sign of
illusion on my own part, but as an insult.

Had I chosen a philosophical subject, I would have had to choose it from
another standpoint and in such a way that what is peculiar to me in it
found expression; and such a subject I could easily have chosen; for so
far as I know, my endeavors diverge, in a fairly decisive way I believe,
from those of the philosophers in the various literatures, when the talk
is of method and in particular of the interaction between philosophy and
the investigative sciences.

For a long series of years I have made it my particular task to study the
encyclopedia of the sciences, since I have considered it altogether
necessary that philosophical science, if it is to be carried a step
further, must first give itself an account of its relation to the other
sciences, the so-called special sciences. Had I, then, chosen a
philosophical subject peculiar to me, I would have had to limit it to
such a degree that the subject itself would undoubtedly have
% --- p. 281 ---
\opage{281}\setcounter{footnote}{0}%
interested only a very small circle of the students. Just imagine what it
means to have to carry through, in a comparatively small number of hours,
a strictly scientific analysis from which it could be seen how philosophy,
for example, relates to mathematics. How many listeners would have
followed these analyses with undivided interest; in how many could I have
presumed an equal knowledge of philosophy on the one side and of
mathematics on the other? Had I chosen another, related subject, for
example to set forth the relation of philosophy to chemistry, particularly
with regard to the so important doctrine of atoms, and then also with
regard to the problem of change and the relation between the two concepts
of change and transformation, then again I would have been confined to a
very limited circle of listeners.

If, now, for the reasons stated, I had to give up choosing a strictly
scientific subject, I was on the other hand directed to the great circle
of humane ideas of which the science of spirit has so rich a store at its
disposal; and as I looked about in the circle of these ideas, I truly did
not lack subjects that might come into consideration. But among all the
subjects there was one that came to the fore again and again, the question
of faith and knowledge, a question that is on the agenda, indeed, I might
say, at present a burning question in my native land. This question has
precisely a peculiar main aspect by which its treatment falls under the
humane point of view, and I was in this respect particularly tempted to
let the choice fall on just this subject.

But here again a natural objection arises. Why, it will be asked, not
simply announce lectures on the relation between faith and knowledge and
deliver them, since they were after all meant to promote a scientific
cooperation between the universities, in the
% --- p. 282 ---
\opage{282}\setcounter{footnote}{0}%
academic, strictly scientific form? I answer: that was impossible for me!
As far as faith is concerned, I could indeed, as I have tried to show,
develop the great fundamental concepts even for those who were not quite
clear with themselves about the principles of knowledge; but when, on the
other hand, it was a matter of assigning to knowledge and science their
strict limit, it was altogether impossible for me here to give any survey
that could lay claim to scientific stringency; I would then have had
--- it lies in the nature of the case --- to give a detailed account of my
philosophical concept of knowledge. The peculiar limitation that,
according to my way of seeing, I must give to science is motivated
precisely by the encyclopedic studies mentioned before and by the
investigation of the relation of philosophy to the exact and positive
sciences. Since I could not, as I said, expect to present anything on this
last point that would be penetratingly intelligible, I could not get the
concept of knowledge treated with sufficient thoroughness, and so, as far
as this subject is concerned, had to give up the scientific mode of
presentation in the narrower sense.

To this was added yet another consideration: what is most important to me
at the moment --- for with knowledge and science I believe I shall
certainly manage --- is not the concern of knowledge but that of faith. It
is not the purely scientific but the popular development of the opposed
concepts on which I have to lay particular emphasis. Now, in order not to
incur the inconsistency, when I expressly put the two concepts, knowledge
and faith, side by side and thereby obliged myself to elucidate knowledge
just as clearly as faith, of leaving science in the shade, I recast the
title of my lectures and announced them, not under the heading Faith and
Knowledge, but under the more general and more humane point of view:
Hindrances and Conditions for the Spiritual Life in the Present
% --- p. 283 ---
\opage{283}\setcounter{footnote}{0}%
Age. Under this point of view I could in a freer way bring forward much
that belongs partly to civilization, partly to the life of the people, and
particularly to religion. Here I could move freely in the domain of free
cognitions. So much for the choice.

If, now, the choice was --- as I believe I have shown --- in a way
prescribed for me, there was added to this the by no means inessential
consideration that on these terms I could entertain the hope of gathering
a circle of cultured listeners; indeed, not merely of listeners, but,
and I must lay particular weight on this circumstance, also of women
listeners. These are not flattering words; for flattering words there is
no place here, no mood; it is my full conviction that fundamental
thoughts that move in the life of faith, fundamental thoughts with a
personal and popular stamp, must on the highest as on the lowest rung of
culture be heard equally by man and woman. In treating such tasks it is
important to have the feminine individuality present, especially for the
sake of the presentation, where it is precisely a matter of distinguishing
sharply between the scientific, the popular, and the folk-like. Only
before a feminine tribunal are we quite sure not to mix the heterogeneous
forms with one another. As to the scientific treatment, in the first
place, surely no one who has resolved to speak about the general concerns
of culture to the cultured will be in danger of straying into the learned
preserves; but should it happen to him, he would become aware of the
mistake with a certain embarrassment on seeing himself surrounded by
attentive women listeners. It is probably only rarely that a gifted woman
feels herself drawn to learned investigations. Of course: no rule without
exception, and where talent is, be it in woman or in man, we must hold it
in honor. We know, in individual cases, the peculiar
% --- p. 284 ---
\opage{284}\setcounter{footnote}{0}%
gift in noble and beautiful women. Madame Chatelet, for example, had so
eminent a talent for mathematics that, besides playing the role of the
amiable great lady in the salons, she passed for a scientific authority
among the great mathematicians of her time. But how rare, indeed, to the
good fortune of society, how exceedingly rare is such a phenomenon! If we
keep to the rule, then a lecture that is to be able to count on feminine
attention had best be freed from learned baggage and artificial formulas;
but lectures of this nature are, after all, best suited of all for the
communication of humane, aesthetic, ethical, and religious content.

The next thing on which I must lay weight is the distinction between the
popular and the folk-like. For the popular, woman in a certain sense has
just as much sense as man; she has, and must have, her share in culture
and civilization: without her participation the so-called higher culture
would look meager. But the merely popular will not easily beguile her; if
she has not been twisted into a seminarist-governess by all too
encyclopedic cramming at the institute, she will always with fine tact
note the limit where the ensouled leaves off and the unintelligible
begins; the fleeting, light, seemingly superficial way in which she
intimates a view is often a telling proof that with the certainty of
instinct she has at bottom hit upon the right thing. Even if she cannot go
quite so far into popular astronomy as the cultured man, she can in return
as a rule go all the more deeply into poetry and psychology. The
folk-like, the religious, in short: all that is personal a woman
apprehends deeply and essentially, just as deeply and essentially as a
man. But when we say that woman is capable of understanding the highest
truths of personality just as originally as
% --- p. 285 ---
\opage{285}\setcounter{footnote}{0}%
man, we do not therefore say that she understands them in the same way as
he. No! The difference between man and woman is essentially a difference
of soul. It has often been asked, naturally not without learned
self-importance, whether woman, this weaker sex, could now really also in
actuality think thoroughly, reflect thoroughly, argue thoroughly, and so
on. We answer Yes! but she does not reflect, she does not think, she does
not argue like a man; she is, after all, a woman.

This takes us a step further; for precisely by touching on the difference
between man's and woman's appropriation of the highest truths we come at
the same time to reflect on the double conception of popular instruction,
of what was characterized in an earlier connection as the enlightenment
of the people. Should it really be the case, as ``the modern
consciousness'' assures us, that knowledge must in the end dissolve faith,
then woman will sink back down to a very subordinate rung in society; for
this is certain: on the terms of knowledge woman can never become man's
equal. Should it on closer reflection prove impossible for a many-sidedly
cultured man to unite knowledge and faith in one consciousness, then a
chasm of soul will open between man and woman; for this is certain: a
woman who for the sake of knowledge will give up having her life in the
mystery of faith must forfeit her womanliness and wither like a plant cut
off from its root. A man, even if he is not properly a man of science, can
at a pinch scrape along with the help of the popular; but a woman whose
whole culture is built on reasonings and principles taken from popular
writings is a prosaic magnitude, an individuality bedraggled in soul. A
man can at a pinch become a cosmopolitan and perhaps come out of it
tolerably well;
% --- p. 286 ---
\opage{286}\setcounter{footnote}{0}%
but when a woman becomes cosmopolitan, what does she then become? A barren
abstraction, a schema painted in many rubrics, an outline of the
universally feminine without light and life. Even on the highest rung of
culture the feminine individuality stands guard for the religious as well
as for the folk-like. This Grundtvig has strikingly seen, and what
Grundtvig asserts in robust immediacy S. Kierkegaard's subtle reflection
has fully confirmed through aesthetically, ethically, and religiously
nuanced transitions.

As the characteristic mark common to the feminine and the folk-like we
name the deeper simplicity. From the simple all true personal culture, all
development of spirit, must begin; to the simple the whole development,
however rich and manifold it may otherwise be, must return again. Popular
instruction begins simply; for it begins with the children, and it begins
with the common people. Herein lies that it must begin precisely in such a
way that in its rudiments it has possibilities enough, has the whole
wealth of the endowment in its first shoots. For just as little as a child
is destined to live out its time in a perpetual childhood, just so little
are the common people destined to remain common people forever. The common
people are the people in a state of spiritual childhood; in the common
people the life of the people has its root, its endowments, its instincts;
in the common people the spiritual powers stir in an obscure but original
and naturally peculiar way. The common people are the people itself, not
in complete being but in becoming and growth. If we forget the growth and
put a static being in place of a continual becoming, then we must
necessarily run aground on one of two extremes: either simply make the
people into the common people or regard the common people merely as a
mass. Those who in a crude way identify the people with the common people
% --- p. 287 ---
\opage{287}\setcounter{footnote}{0}%
distinguish themselves especially by a peculiar hatred of everything that
suggests higher culture, richer development of spirit, intellectual
superiority: science, art, poetry are regarded askance by the leaders of
the masses. Those who, in the blindness of an aristocracy of spirit, look
down on the common people as on that part of the population which fate
has once and for all destined to toil in tutelage, and relish the
well-known words: ``Zuschlagen kann die Masse, da ist sie respectabel,
urtheilen aber gelingt ihr miserabel'' --- forget the dearly bought
experiences of history, how terribly such contempt has nevertheless
avenged itself, and will go on avenging itself. The instruction of the
people must not stop at what belongs to the common people; it must just as
little restrict itself to the religious alone as to the civic. There is in
the people, we repeat it, a double need: need of religion and need of
enlightenment, need of the soul's salvation and need of culture. Step by
step the common people must be raised to human insight into what is
human, to ever higher self-cognition, to a clearer understanding of
actuality, to deeper personal conviction.

Without personal conviction there is no inner independence, without inner
independence no real freedom. There is, and can be, a free people, but
there is not, and cannot be, a free common people; for the rung of the
common people is, spiritually seen, the rung of tutelage, hence of
unfreedom. In popular instruction it is a matter of bringing the ordinary
man so far that he will be able: to think his own thoughts clearly, to
give himself a sufficient account of his own actions; to see and
understand that no one in society can have rights without at the same time
having duties; to see and understand that the interest of the individual
man is furthered in a reasonable way only by the interests of society
being taken care of on the large scale; to see and
% --- p. 288 ---
\opage{288}\setcounter{footnote}{0}%
understand that the interests of society on the large scale cannot
possibly be furthered unless the different working classes, not merely
those who work in the bodily sphere, but also those who work in the
spiritual, enjoy full recognition, each according to the terms grounded in
the nature of the circumstances. Only he who is able to give himself
reasons for what he wills and what he does will be able to understand the
reasons that are given him by others. The difficulty with the
emancipation of the ordinary man is not that he lacks either noble
self-respect or sound understanding or honest will; but the difficulty is
that his thinking is unclear, his horizon restricted, his mind suspicious.
The difficulty is that he does not rightly know how either to give reasons
or to accept reasons, because he is essentially unfree. The ordinary man
is unfree as long as his inner being is obscure; but that his inner being
is obscure shows itself especially in his being dark, shut up, in the most
personal concerns. Inner obscurity and reserve give the character, under
the pressure of life, a perceptible tinge of melancholy. On the surface the
life of the people may indeed in certain sunshine moments be festivity
itself, an image of natural carefreeness, immediate joy, and inexhaustible
zest for life; but innermost there is something dark, and in the dark
dwells melancholy. A thorough connoisseur of the folk melodies from the
various regions of the world has said that on the whole they bear the
stamp of its not being joy but sorrow that taught the peoples to sing.
This heavy, oppressed, shut-up nature will disappear only as the warmth
and light of the sun of freedom penetrate further into the masses.

And who, then, is to instruct the common people, provide for their growth,
educate and form them into a people? However many changes the forms of
society may undergo, there can
% --- p. 289 ---
\opage{289}\setcounter{footnote}{0}%
be no doubt that pastors and schoolteachers must continue to be the real
teachers of the people. How these teachers conceive and tend their calling
is then of the very greatest significance, not merely for the state, that
it may maintain peace and order, but for the people itself, that it may
secure spiritual freedom for itself. One sometimes hears remarks such as
these: to learn a science requires a good head, to become an artist
requires talent; even to become a riding master requires natural
aptitude. But to become a village pastor? Well, one can surely become a
pastor in the country with a small head, without talent. And a
schoolteacher? For that, natural aptitude is hardly needed! Such
judgments are in a high degree skewed and confusing. When there is
disputing about how much endowment is required for preaching to the
common people, it is from the outset a thoroughly false beginning to seek
the standard in the good head. The standard is theological, and precisely
for that reason false. One proceeds theologically from the presupposition
that a pastor at bottom ought to be a learned man; or at any rate, so to
speak, a kind of popular man of science, if he is to be a pastor of
standing. Now this is the ideal. But since ideals are very seldom realized
in the world, one must in many cases turn a blind eye to the scientific
deficiencies and console oneself that a pastor, although he is neither
properly learned nor a man of science, can nevertheless, if he is
otherwise a believing man who knows his Christianity, very well be a good
pastor of souls. Of course, as a rule he cannot get the great living, the
one with many and rich revenues, when he does not have the great
theological grade; but for all that he can surely get a considerable,
burdensome living, namely a living with many poor souls, if, as was said,
he is good enough to be a pastor of souls. In the manifold
% --- p. 290 ---
\opage{290}\setcounter{footnote}{0}%
cases of practical life there can hardly be shown any greater
contradiction between the standard and what is to be measured than the one
that strikes the eye when the question is of judging the qualifications of
the pastors of souls paid by the state. No one, of this I am convinced,
can feel this more deeply than the pastors themselves.

For, let us speak quite plainly, how many pastors actually become men of
science? Not very many --- to the good fortune of the clergy, indeed to
the good fortune of the congregation! It is surely in truth not the
pastor's calling to immerse himself in learned scientific studies; he has
another work to do, a work of the spirit, for which there is required, not
only in the divine but also in the human sense, a special endowment, a
particular preparation, upbringing, and culture. Let us ask the pastors
round about the country whether most of them will not admit that the
extracts of learned theology that were impressed upon them in a so-called
scientific manner belong precisely to what they would most like to manage
to forget, the sooner the better, in order to be able to begin afresh?
Would not most of them have noticed a difference, an essential difference,
between the theological and the religious way of thinking, between the
learned and the edifying discourse? Instead of making do with learned
fragments and scientific surveys, they are directed precisely to
observations of actual life, to personal apprehension and appropriation of
the apostolic word, to drawing from the well of inner experiences.

Certainly there may be need of an objective rule by which the
qualifications of the people's teachers can be judged, as surely as the
state is not merely to pay its officials but also, as far as possible, to
fit the pay to the service; for a subjective, pietistic opinion about the
individual's ``sealing'' is quite certainly unreliable. But precisely in
this
% --- p. 291 ---
\opage{291}\setcounter{footnote}{0}%
respect the absurdity shows itself when theory goes in one direction and
practice in another. Were the present school for pastors really fitted to
give its students the scientific culture, the religious cognition, the
spiritual consecration that in a special degree make them fit to educate
the people and to proclaim the Gospel in the congregation: then a
testimony as to how the individual had passed his final examination could
quite certainly be orienting; but is this now the case? There can be no
doubt that the pastors themselves, when it becomes clear to them that the
way that is to lead them to the goal is quite another than the
objectively dogmatic one, will be the most zealous in working against the
theological nuisance. For it weighs on them, this nuisance; it dulls the
sense and confuses the eye; it paralyzes the powers of the spirit by
cramming the memory with bits and pieces of learned traditions, literary
relics, into which not even a demon shall be able to breathe either life
or spirit, unless it happen by a speculative miracle. To form pastors in a
school that effaces the characters of revelation, overlooks the principle
of faith, and denies the original validity of religious thinking --- what
does that lead to? It leads to a quiet, systematic, steadily advancing
undermining of the position of the clergy in civilized society. What was
it, then, that really grounded the power and authority of the clergy
in the Middle Ages? It was its clearer eye for spirit, its deeper insight,
its ideal advantage. In the present age, whose spiritual direction is on
the whole predominantly intellectual, reflecting, criticizing, the
clergyman must not merely manage to keep up, but, as far as the religious
is concerned, expressly assert his intellectual superiority. His
contemporaries must get the impression that he is precisely the man, the
right man, who knows what is what when the question is
% --- p. 292 ---
\opage{292}\setcounter{footnote}{0}%
of the decisive fundamental thoughts of personal life, the man who in
religious reflections can master the freethinker, settle accounts with
those who know, and orient the believers. If this impression can become
the predominant one in society, then the clergy can assuredly be
untroubled about all the rest. Among warring powers, the one that has the
greatest forces, the strictest discipline, and the best weapons will
always be surest of victory. It will surely be confirmed in the strife of
spirits that on the path of history there is no spirit so mighty as the
spirit of the Gospel, no thought so penetrating as the thought of the
will, no word so powerful as the word of faith, and that then, too, no
combatant who knows his weapons is so well equipped as he who fights for
the right of personality, its true life in the eternal kingdom. But the
priestly coat of arms is not theology; to clothe the clergyman in a
theological harness is in our day just as preposterous as it is
preposterous to puff him up into a ``witness to the truth'' and expose
him to satire.
% printer's defect: the Danish prints "Váaben" (a wrong sort) for "Vaaben", weapons

One more point shall be touched on here. In the spiritual life of the
present age it is a matter of interaction between science and the
enlightenment of the people, a higher interaction which will undoubtedly
bring with it that what has already in my native land been the object of
very comprehensive and serious deliberations will soon attract attention
in Norway and Sweden too, namely, whether the spiritual conditions of the
present time should not require that, over against the university, the
learned high school, there should arise another high school on a
popular foundation on a large scale. It is not enough with the small
peasants' high schools; what is needed here is an institution of learning
resting on so broad a foundation that it can comprise the various
elements of culture and of science, mark well, in such a way that the
fundamental thought of the enlightenment of the people
% --- p. 293 ---
\opage{293}\setcounter{footnote}{0}%
penetrates each single element and governs the whole. It is my hope that
it will come to that, as I at the same time presuppose that such an
opposition of folk high school and university, far from oppressing the
university and putting it in the shade as an antiquated institution,
will on the contrary mightily contribute to raising its scientific
standing, giving it new impulses, freeing it from extraneous burdens,
setting ideas in freer circulation, and rejuvenating the blood in the old
limbs. There are those who think that the universities are antiquated
institutions that have essentially played out their role. This opinion I
cannot join; whereas I can well understand that the highest point of
union of the sciences cannot suffice as the sole focal point of the
enlightenment of the people, once it stands firm that this enlightenment
must have a double point of departure, one in faith and one in knowledge.

There can surely be no doubt that the particular special sciences do best
to group themselves into one great whole, just as it cannot be doubted
either that the enlightenment of the people, with its peculiar tasks,
had best have its own center. As for philosophy, as the science of spirit
it will deal equally with both; on the one side it will immerse itself in
the strict sciences and take scientific elements up into the doctrine of
ideas, on the other side enter into the fundamental thought of popular
instruction and take religious problems up into the theory of cognition,
thus forming a connecting link between the two schools. If, now --- so far
as we can judge --- the two schools are so far from one-sidedly
restricting, stunting, and harming each other that on the contrary they
will strengthen and help each other forward by their opposition, then it
appears from this point of view too what the redoubling that we have here
sought to assert essentially has to mean. So much with regard
% --- p. 294 ---
\opage{294}\setcounter{footnote}{0}%
to the main matter. As for the more detailed carrying-through of the
fundamental thought, we must, in the practical as in the theoretical,
refer to the future itself. For it goes with ideal tasks in such a way
that he who solves one always gets two back. This, that the more one
works at a task of the spirit the more manifold it becomes, is precisely
what frightens so many, both those who choose theory and those who tread
the practical way of experience. And yet this is precisely the law of the
development of spirit. The one comprehensive task, to further and promote
the spiritual life, contains innumerable tasks; the fewest will probably
be solved in the present age; most belong to the future. Of the tasks of
the future one had best speak briefly, since otherwise speech easily
becomes tiring, in becoming hovering, groping, indefinite: to use the
present in the best way is to prepare what is to come.

\begin{center}\rule{3cm}{0.4pt}\end{center}

It was a great and, as we hope, fruitful thought of the future that
gathered this honorable circle of listeners, men and women. How far an
utterance has now been heard in the lectures delivered, a thought of life
developed, that the future itself will approve, is a question that must
no doubt wait some time yet for its answer. Only this I can say, and I
say it openly, that I trust in the weight of the reasons. I have spoken
without reserve, from inner prompting, decidedly; but I have never
thought or meant that any of my propositions should count as articles of
faith. I have spoken in order to submit the matter to the consideration
of those competent to judge; I have spoken against the encroachments of
civilization, but not against civilization, against the missteps of
reason, but not against reason, against scientific delusion, but not
against
% --- p. 295 ---
\opage{295}\setcounter{footnote}{0}%
science. I have spoken of the people's need and the people's right to a
peculiar enlightenment independent of science and learned research; I have
spoken freely and openly, for I was conscious that I was speaking to an
assembly in whose midst men of high culture, thorough learning, practical
and theoretical expertise, some among the first men of the realm, to
whom the people must look up with esteem and confidence, had for the
sake of the cause decided to choose a seat. I was convinced that one
could speak to men of science, not merely in praise of science, but also
of scientific encroachments, without being misunderstood; I was
convinced that one could speak to men of the people about the people's
need and want and right without occasioning ill feeling. I have spoken
freely; but I take with me the conviction that however much may have
been defective in the particulars, my free speech can hardly have obscured
or botched any good cause.

Now we stand at the goal; I confess with deep gratitude that these
lectures will be characteristic of my activity as a teacher in more than
one respect, that they mark an essential turning point for it. It is a
leaf of the book of my life, a leaf of remembrance, that I have here got
written; among many leaves it is in truth one of the richest, I confess
it, one of the most beautiful. So I will then imprint this deeply in my
mind, dedicate it to a grateful remembrance; so then let the cause be
inaugurated, so then let the inaugural lectures herewith be concluded: may
the honored assembly of listeners, men and women, receive my warmest, my
most respectful thanks!

\begin{center}\rule{3cm}{0.4pt}\end{center}

% --- p. 296 ---
% [p. [296] is BLANK: the verso after the end of Lecture XVI. No folio,
%  no text, so no \opage.]

\chapter*{Postscript.}
\addcontentsline{toc}{chapter}{Postscript}
\markboth{Postscript}{Postscript}
\label{ch:efterskrift}
% --- p. 297 ---
\opage{297}\setcounter{footnote}{0}%
However much I have striven in these lectures to state my view of the
relation between faith and knowledge quite clearly and definitely, it
has nevertheless not been possible to avoid several, even very essential,
points being exposed to misunderstanding. I perceive this partly from
individual inquiries from listeners, partly from articles in Norwegian
papers. In particular there is one point for whose further clarification
I find myself called upon first of all on this occasion; it concerns the
\emph{actuality} that must be ascribed to the sacred facts, particularly
the resurrection of Christ.

As regards this matter, I can refer to an article to be found in
\emph{Aftenbladet} for the 24th of December 1867 under the heading:
Professor C. P. Caspari on ``Faith and Knowledge.'' It is in ``the form of
confession and testimony'' that the Professor states what he ``with a
little-fitting word'' must call his ``standpoint.'' ``I do so,'' he says,
``because on account of views and thoughts which in recent times have
gradually begun to stir and assert themselves among us, \ldots\ I must
consider it to be not so entirely out of place.'' Since the ``views and
thoughts'' that Professor Caspari here means to meet by making his
confession of faith are not directly connected with my lectures, it
naturally cannot occur to me to want to burden
% the Danish prints "Vidnetsbyrdets" for "Vidnesbyrdets" (testimony); translated in the sense
% --- p. 298 ---
\opage{298}\setcounter{footnote}{0}%
him with any scientific exchange with me; I merely make use of his
statement because it gives me a convenient occasion to set down some
further determinations of what I understand by \emph{phenomenon of spirit,
actuality of spirit, actuality for the spirit}.

Professor Caspari confesses the faith ``that the revelation of salvation
of which the Holy Scriptures tell belongs to external historical
actuality, that it is not merely true but also actual, that it possesses
not merely an inner actuality for faith --- which would be the same as
that it exists only in the idea, or that it is at bottom only an
imagination --- but also an external actuality in itself, and that not
merely in its main features and fundamental facts but also in all its
particulars. And on the external actuality of this revelation of
salvation our salvation rests; without it we would be and remain to all
eternity in our sins, be and remain eternally lost. Mere ideas of
salvation cannot give the human heart, restless in sin, peace and
blessedness. \ldots\ But'' --- the Professor continues --- ``the objections
of science to the fundamental facts and particulars of the revelation of
salvation? --- I am myself, by the grace of God, a man of science, but I
must openly confess that they have never caused me a single sleepless
night. As regards in particular the objections to the miracles of the
revelation of salvation, I must on the contrary set against them the
historical actuality of these miracles and above all of the resurrection
of Christ. If there is any fact that is sufficiently, indeed
superabundantly, historically attested, it is the resurrection of Christ;
if it is not to be called historical, then we may as well strike a line
through everything that is called history. \ldots\ And then there is,
after all, a miracle that happens to this very day, that is a fruit of
those old miracles of the revelation of salvation and therefore also
bears witness to them, and that everyone can experience in himself, I
mean the miracle of rebirth, the great inner miracle, a miracle which ---
if it is permitted in the domain of
% --- p. 299 ---
\opage{299}\setcounter{footnote}{0}%
miracles to speak of greater or lesser --- is just as great as the
resurrection of Christ. Every true Christian knows himself as a living
miracle; the actuality of the miracles is therefore also just as certain
to him as his own existence as a new man, as a Christian.''

In this statement there is something about which I must at once declare
myself unconditionally in agreement with Professor Caspari, and it is:
that the facts in sacred history are not ``imaginations,'' that ``mere
ideas of salvation cannot give the human heart, restless in sin, peace
and blessedness''; with how much reason I emphasize this here, anyone who
with an unbiased eye has read through not merely the present lectures but
any one whatever of my other writings will be able to tell himself. It is
precisely the actuality of the sacred facts, their religiously marked
character, their full, undistorted reality, their unweakened validity in
themselves, whose acknowledgment I seek from every side and in every
possible way to enjoin. But when it is in earnest a matter of asserting
actuality, full, self-valid actuality, then surely the question presses
itself upon us: what is actuality, the true actuality, fully valid in
itself? It is in the answer to this important question that I do not
know whether I dare declare myself in agreement with Mr.\ Caspari.
% AUDIT doubtful: "fuldgyldige" (fully valid) could be a blotted "fuldgyldjge"; the sense is unaffected

In the first place it seems to me as if the Professor, where the talk is
of the phenomena of revelation, lays rather much weight on the sensuous
exterior, when he makes our salvation and blessedness dependent on it and
enjoins that this exterior of actuality is to hold, not merely of certain
facts, but of all, not merely of ``the fundamental facts,'' but of ``their
particulars.'' Let me choose some examples. Isaiah describes his calling
as a prophet (Isa.\ VI): ``In the year that King Uzziah died I saw the
Lord sitting upon a throne, high and lifted up, and his train filled the
temple. Above it stood the seraphim, each one had six wings''
% --- p. 300 ---
\opage{300}\setcounter{footnote}{0}%
and so on. The prophet Amos describes his visions (Amos VII): ``This vision
Jehovah let me see: he was about to form grasshoppers. \ldots\ This
vision Jehovah let me see: Behold, the Lord Jehovah called upon the flame
as his advocate; it licked up the mighty deep and devoured the portion.
\ldots\ This vision Jehovah let me see: the Lord standing upon a wall
made by a plumb line, with a plumb line in his hand. And Jehovah said to
me: What seest thou, Amos? I said, A plumb line. Then said the Lord: now
I will cast the plumb line into the midst of my people Israel'' and so on.
To Ezekiel the divine voice sounds (Ezek.\ III): ``Son of man! eat what
thou here receivest, eat this roll of a book; then go and speak to the
house of Israel. And I opened my mouth, and he gave me this roll to eat.
And he said to me: Son of man; let thy life be consumed, and fill thy
inward parts with this roll that I give thee; and I ate, and it was in my
mouth as honey for sweetness.'' \ldots
% printer's defect: the Danish has no opening quotation mark before "Menneskesøn!" though the quotation is closed after "Honning"; supplied here

That these prophetic visions are not empty imaginations but true
revelations, that these revelations have actuality, full actuality for
the spirit, must be held fast; but, we ask, is it now permissible, with
such phenomena, to set up a separation between the prophet's vision and
what the prophet sees? May we really assume that the throne on which
Jehovah lets himself be seen by Isaiah in the temple has the same
outward actuality as the temple of Jerusalem? that the grasshoppers that
Jehovah forms in the vision for Amos have the same outward actuality as
natural-historical grasshoppers? that the plumb line Jehovah casts into
the midst of Israel has an outward actuality like the plumb weight that
hangs on a Bornholm wall clock? that the roll of a book that Ezekiel eats
--- in the vision --- has the same outward actuality as the Bible that
lies on the table?

The prophetic visions are seen in the spirit. He who says that the
visions the prophets see have the same outward actuality as the things
that are seen with bodily eyes, what does he
% --- p. 301 ---
\opage{301}\setcounter{footnote}{0}%
do? He makes revelation spiritless. Can a spiritless revelation be called
a ``revelation of salvation''? Can anyone, without mocking what is holy,
really maintain that all salvation and blessedness rests on the external
actuality of the revelation of salvation, that he who, as regards the
visions named, does not lay the truth into the exterior of the external
actuality of ``the throne,'' of ``the grasshoppers,'' of ``the plumb
line,'' of ``the roll of a book,'' in short: of ``all the particulars,''
must ``remain to all eternity in his sins, be and remain eternally
lost''?

It is by no means my intention to want to force upon Professor Caspari
the consequence of his own statement; I only want to draw the attention
of thinking readers to the dangerous misdirection that lies in so
one-sided a ``testimony.'' It is one thing to acknowledge the exterior
insofar as in the phenomenon of revelation it belongs to determining the
interior, another to be so dazzled by the exterior, the sensuous exterior,
that one makes all salvation and blessedness dependent on the external
actuality of the particulars. That the Professor by no means wants to deny
the interior he makes clear enough by the emphasis he lays on the miracle
of rebirth, ``the great inner miracle, which --- if it is permitted in the
domain of miracles to speak of greater or lesser --- is just as great as
the resurrection of Christ.'' By this ``testimony'' decisive weight is
now quite certainly laid on the interior; but how does the latter now
relate to the former? Of the miracle of rebirth it holds, after all, that
it has inner actuality, actuality for faith. Take faith away, and what
then becomes of the ``great miracle'' of rebirth? The assurance of this
miracle's validity is in faith alone; or should the believer perhaps be
able to behold his own rebirth in an external actuality, an actuality
with which it stands, to use the Professor's expression, ``as with the
sun, which indeed first comes to be for me by my seeing and feeling it,
but does not thereby first come to be at all, but exists before all my
seeing
% --- p. 302 ---
\opage{302}\setcounter{footnote}{0}%
and feeling, has an existence independent of it, and would exist though
no human eye had ever seen it or any human body felt it''?

If every thinking human being --- consequently Professor Caspari too ---
must concede that the miracle of rebirth has its true validity in itself
by means of the spirit through faith, indeed for faith and the spirit
\emph{alone}: then it is surely herewith conceded that there is an inner,
and yet in itself valid, actuality, an actuality of spirit, an actuality
for the spirit. But how are we now to understand the Professor's
``confession and testimony,'' for he confesses and testifies that the
revelation of salvation that possesses merely an inner actuality in faith
``exists only in the idea, or that it is at bottom only an imagination.''
In truth, a terrible testimony --- against ``the great miracle of
rebirth''! Since, now, Professor Caspari has here borne not one but two
testimonies to the actuality of the miracle of rebirth, namely one
\emph{for} and one \emph{against}, both equally forceful, the question is:
how are we, who would gladly accept his ``testimony,'' to make sense of a
``confession'' so at odds with itself?

We shall surely be compelled to concede: 1) that the actuality of
revelation is first and last actuality of spirit, actuality for the
spirit; 2) that the actuality of spirit always goes from the interior
through the exterior back to the interior, so that its decisive mark must
always be sought in the interior; 3) that the exterior which the spirit,
the self-revealing Spirit of God, masters, pervades, and takes up into
the peculiarly determined actuality of the sacred facts, is of very
different kinds, now bodily, now incorporeal, now natural, now
supernatural, now immediately sensuous, so that the phenomena of
revelation, precisely for that reason, if we are really to be able to
``give ourselves up to them and immerse ourselves in them,'' must become
clear to us through a criticism that answers, not to ``our own arbitrary
% --- p. 303 ---
\opage{303}\setcounter{footnote}{0}%
representations,'' but to the self-valid principle of revelation and of
faith.

I shall briefly illuminate each of the propositions set out. Thus:

\emph{The actuality of revelation is actuality of spirit, actuality for
the spirit}. Herein lies that nothing can have validity for faith except
that which the believer, by virtue of the testimony of the Spirit, must
ascribe validity in itself. According to Professor Caspari's ``testimony''
one would surely have to assume that something could have inner actuality
for faith and yet perhaps be a pure and simple imagination; this must be
denied. That something has inner actuality for faith means that for the
believer it has validity in itself; as soon as the believer becomes
conscious that that in which he believes has no validity in itself, he
also ceases to believe in it. He who believes in God ascribes to God
validity in himself. Were anyone to say: God, to be sure, does not exist
in and for himself; but precisely for that reason I believe in him, that
I may give him an inner actuality in my faith: what would one then have to
judge of such a person's mental state? Were anyone to say: I believe in
God, that is: I give him a provisional inner actuality in my faith; but
whether the God in whom I believe also has actuality in himself, that I
must leave undecided until I can in one way or another get to see him in
an external actuality, for example sitting on a throne in the temple or
standing on a wall with a grasshopper or with a plumb line in his hand:
what would one then have to judge of such a faith? No, to believe in God
is to ascribe to God validity, absolute validity in himself; he for whom
God's personal being does not have validity in itself does not believe
in God. When Professor Caspari, instead of ``mere ideas of salvation,''
wants to have ``persons, divine persons, the personal triune God'' to hold
to, we can be wholly and entirely in agreement with him in this;
% the first thesis carries no "1)" in the Danish (cf. "2)" p. 304, "3)" p. 306)
% --- p. 304 ---
\opage{304}\setcounter{footnote}{0}%
but the actuality of the triune God is, after all, actuality of spirit,
actuality for faith, inner actuality for the spirit. It is not with the
assurance of God's existence as with the assurance of the sun's existence.
A human being can, even if he becomes bodily blind and is thereby put out
of the condition to see the sun, still continue to feel assured of the
sun's existence; but a human being cannot in spiritual blindness cease to
believe in God and yet continue to feel assured of God's existence: the
original assurance of the existence of ``the divine persons'' is not given
by external actuality for the senses, but by inner actuality for faith.

2) \emph{The actuality of spirit goes from the interior through the
exterior back to the interior, so that its decisive mark must always be
sought in the interior}. When the question is of the external actuality in
which God has personally revealed himself, we must think particularly of
Christ's manifestation in the flesh. With regard to this fact, Professor
Caspari has then also quite rightly referred to 1 John 1:1 ff. With this
external actuality it is so much in earnest that the evangelist declares
every spirit that testifies against it to be a false spirit. ``Beloved!
believe not every spirit, but try the spirits whether they are of God;
for many false prophets are gone out into the world. Hereby know ye the
Spirit of God: every spirit that confesseth Jesus Christ to be come in the
flesh is of God. And every spirit that confesseth not Jesus Christ to be
come in the flesh is not of God, and this is the spirit of Antichrist,
whereof ye heard that he cometh, and even now already is he in the
world'' (1 John 4:1--3). But let us now see how the exterior in this
manifestation stands to the interior. Is the relation this: that the
certainty of the manifestation of the Son of God begins with the exterior
and from the exterior passes over into the interior? Is the relation
% --- p. 305 ---
\opage{305}\setcounter{footnote}{0}%
not rather this, that the certainty begins from the interior, dwells for
some time with the exterior, and then at last, after this exterior has
vanished, returns with reinforced emphasis to the interior? If the
cognition that Jesus of Nazareth is the Christ originally rested on the
exterior, then
all who saw the man Jesus with bodily eyes must surely also have been able
to see from his exterior that he really was the Son of God. But that this
was not the case the Gospel history teaches. The sensuously present Jesus
of Nazareth was in actuality the Son of God only for those who saw him
with the eye of faith, heard his word with the ear of faith. What it means
that the actuality of Christ revealed in flesh and blood is an actuality
for the spirit is seen, among other places, from Matt.\ 16:13--17. ``When
Jesus came into the coasts of Caesarea Philippi, he asked his disciples,
saying: Whom do men say that I, the Son of man, am? And they said: Some
say that thou art John the Baptist; some, Elias; and others, Jeremias, or
one of the prophets. He saith unto them: But ye, whom say ye that I am?
Then Simon Peter answered and said: Thou art the Christ, the Son of the
living God. And Jesus answered and said unto him: Blessed art thou, Simon
son of Jonas, for flesh and blood hath not revealed it unto thee, but my
Father which is in heaven.'' Had Christ been directly recognizable by his
exterior, there would surely not have been such divergent opinions about
him, and there would have been no reason to call the disciple blessed
when he saw only what anyone, even a Herod and a Pilate, must have been
able to see. How the exterior in the actuality of Christ's manifestation
relates to the interior we learn from Jesus' own words at the parting
from the disciples (John 16:7 and 13--15). ``Nevertheless I tell you the
truth: it is expedient for you that I go away; for if I go not away, the
Comforter will not come unto you'' \ldots\ This we must surely understand
thus: it is expedient for you that my external actuality
% --- p. 306 ---
\opage{306}\setcounter{footnote}{0}%
vanish from your sight, for only on this condition can the Spirit enrich
you in inner actuality with the full truth. The conversion of the exterior
into the interior is expressly designated in such a way that in Jesus' own
words we have a sure guide for apprehending and understanding sacred
history. Here it is said that the Spirit ``shall not speak of himself; but
whatsoever he shall hear'' (from the mutual speaking of the Father and the
Son, the interior) ``that shall he speak.'' The Spirit's presentation of
Christ's life on earth is not to be restricted to a finite copying and
registering of external occurrences; the Spirit of truth is to explain,
illuminate, and present the history of the Savior in such a way that the
exterior receives its significance from the interior; in this interior is
his glory, a glory that only the eye of faith can see in the light of the
Spirit. ``He shall glorify me: for he shall receive of mine, and shall
show it unto you.''
% the Danish reads "Talen det Indre" as printed (sense perhaps "Talen i det Indre", speaking in the interior); not emended

3) \emph{The exterior which the Spirit of revelation takes up, pervades,
and transfigures is of very different kinds, so that the facts in sacred
history on this account assume a very different stamp of actuality}. ``If
there is,'' says Professor Caspari, ``any fact that is sufficiently, indeed
superabundantly, historically attested, it is the resurrection of Christ;
if it is not to be historical, then we may as well strike a line through
everything that is called history.'' We say the same, but on the condition
that this fact is not drawn down into the circle of the events that belong
to profane history and are subject to the criticism of knowledge, but is
acknowledged for what it in truth is, a supernatural fact, a fact whose
actuality, with all its exterior, is actuality of spirit, actuality for
the spirit.

That sacred history takes up and unites absolutely heterogeneous elements
of actuality becomes evident at once when we reflect that, in connection
with natural events, it contains miracles. The crucifixion, for example,
is a natural
% --- p. 307 ---
\opage{307}\setcounter{footnote}{0}%
event, the resurrection on the other hand a supernatural one; this makes
an essential difference when the talk is of an evaluation of the
historical testimonies. That Jesus of Nazareth really was crucified under
Pontius Pilate is a fact that no historian, whether Jewish or heathen, has
called into doubt; on the contrary, it was with scorn and mockery cast up
against the first Christians, by Jews and heathens alike, that they
believed in a crucified malefactor. It is otherwise with the resurrection;
this fact has been assumed to be historically actual by neither Jews nor
heathens; it is, says the critic of profane history, by no means
sufficiently attested. How does this stand?

The resurrection is a miracle, one of the most strongly attested miracles;
but by whom is a miracle attested? Only by believers, never by
non-believers! This is no accident; it lies in the nature of the case; it
follows with irrefutable consequence from the character and significance
of the miracle.

From the side of the act of power, nothing can, on the presupposition of
faith, be objected against a miracle being performed in such a way that
even the most obdurate would have to concede that they had really seen a
sign of omnipotence. Had Christ, for example, when he meant to raise
Lazarus, had it solemnly announced some days in advance, so that those who
wished to make observations, Pharisees, Sadducees, an Annas, a Caiaphas, a
Herod, a Pilate, could be sufficiently prepared; had he further seen to it
that each of the observers could go into the tomb and assure himself that
the corpse had really gone into putrefaction; had he next placed these
observers in a circle around the tomb and in their hearing uttered the
words: Lazarus, come forth! with such an effect that the dead man was
really seen to rise and come out, bound hand and foot with graveclothes,
and so on; had he then finally, after this had happened, prevailed on the
observers, among whom there would have to be not merely adherents but also
opponents, cultured
% --- p. 308 ---
\opage{308}\setcounter{footnote}{0}%
men whose testimony had weight and significance in all circles, solemnly
to attest, one for all and all for one, what they had seen and heard: then
the criticism of profane history would absolutely have had to declare
itself satisfied, then every inquiring historian, be he otherwise
believing or unbelieving, would have had to concede that if any fact ``is
sufficiently, indeed superabundantly, historically attested,'' it is the
raising of Lazarus. But the story of the raising of Lazarus is not
attested in that way, and the story of the resurrection of Christ is not
attested in that way either.

Why, now, does there rest, for profane eyes, a mystical veil over the
miraculous in every one of the miracles of sacred history? Because it
conflicts with the purpose of the miracle to come forth so outwardly
visibly that an observer without faith should be able, by the mere
exterior, to assure himself that it is a will with supernatural power that
makes itself known in such a work. If true submission to God's will could
be effected by an indisputable sign of omnipotence in the exterior, then
Christ would easily, with a couple of miracles, have brought both an Annas
and a Caiaphas, both a Herod and a Pilate, to throw themselves in the dust
before the divine majesty. Pilate, who is so cowardly, so afraid of
falling into disfavor with the emperor, because, by keeping to the
exterior, proceeding from the exterior, and staking the truth on the
exterior, he is assured that the emperor has the power, the power that
makes poor and makes rich, the power that exalts and that casts down ---
that same Pilate would certainly have crawled like a dog at the feet of
the man whose almighty word could strike a legion of Roman soldiers to the
ground. By coming forth in such an external sphere of action, a miracle
of Christ would quite certainly have become an indisputable fact of world
history; but if a single such miracle had happened, the revelation of
omnipotence would at one stroke have crushed all inner independence, all
human
% --- p. 309 ---
\opage{309}\setcounter{footnote}{0}%
personality, and let a Messiah after the Jews' mind, infinitely mightier
than the emperor but like the emperor, set up an earthly kingdom and rule
as divine despot over slaves. Quite certainly the actual miracle has its
exterior, indeed, according to its different purpose, a very different
exterior; but the exterior does not here have its significance in and
through itself; it has, and must have, its true significance through the
interior, solely and alone through the interior: it presupposes faith, is
apprehended in faith, strengthens faith; it presupposes the spirit, is
attested by the spirit, strengthens the spirit.

To demand a miracle in order to ground one's faith on the exterior solely
as exterior is a crime against the spirit; it is \emph{to tempt God}. To
venture the uttermost for faith's sake and then, in the midst of dangers,
trust
in God's wondrous help and assistance is not --- as soft-hearted preachers
sometimes, in commending ``Christian wisdom,'' let the congregation
understand on the quiet --- to tempt God; but to demand miracles with so
prominent a historical exterior that all salvation and blessedness could
be built on immediate assurance of all the historical particulars in this
historical exterior, because one imagines that without all this external
actuality we would be and remain to all eternity in our sins, be and
remain eternally lost: that is in truth to tempt God. There is one
miracle, but only one single one, of which the believer must say that its
exterior, when it really comes to pass, will be raised above all doubt
even for non-believers; it is Christ's manifestation to judgment; but
this manifestation cannot then become an object of historical criticism
either, for with it history will end.

Only on one condition can it be said, but then with full right, that the
resurrection of Christ is sufficiently, if not ``superabundantly,''
historically attested, namely this, that its actuality is not one-sidedly
placed in the exterior, but, with all the determinations of an exterior,
is referred to and united with its interior; that it is, in other words,
not degraded to stand on an equal
% --- p. 310 ---
\opage{310}\setcounter{footnote}{0}%
footing with other historical actuality, but raised to its proper height
by being apprehended and determined as actuality of spirit, actuality for
the spirit. The true, complete actuality of spirit is neither an exterior
by itself nor an interior by itself, but an exterior answering to its
interior, an interior determined by its exterior; of such an actuality of
spirit, expressed in fact, testimony can be borne only in faith: all the
witnesses of the resurrection are witnesses of faith.

After these explanations it can surely be seen that, if it is not
Professor Caspari himself, neither is it I who rates the actuality of
Christ's resurrection too low; I am conscious that the view set forth
here agrees in each and every respect with the one stated in the
lectures, and can therefore conclude with the assurance that the
representations of ``mere ideas of salvation'' which the Professor wants to
combat can be ascribed to my ideas only by misunderstanding. Of course,
ideas can well make a person sleepless. It would, however, be
regrettable if ideas, even in the form of ``ideas of salvation,'' should
now all at once come up and cause unrest and disquiet a man who has
hitherto had it so quiet that as a ``man of science by the grace of God''
he must openly confess that ``the objections of science to the
fundamental facts and particulars of the revelation of salvation'' have
never yet ``caused him a single sleepless night.''

Even before that statement of Professor Caspari's was inserted in
\emph{Aftenbladet}, \emph{Morgenbladet} had furnished a series of readable
articles under the heading: \emph{Epilog} til R. N.s \emph{Forelæsninger}
[Epilogue to R. N.'s Lectures]. These articles, which all bear the traces
of a practiced pen, are composed with a free scientific overview, in a
benevolent mood, a well-sounding, humane key, and an easy language; but,
we ask, for what purpose? Is it to \emph{guide} the public or to
\emph{misguide} it? Obviously the latter! In assuming this without further
ado, a knowledgeable person will at once understand that it is not
censure but
% printer's defect: the Danish prints "ndrykket" for "indrykket" (inserted); translated in the sense
% --- p. 311 ---
\opage{311}\setcounter{footnote}{0}%
on the contrary praise that I herewith express; for as \emph{guiding}
the articles would be rather mediocre, as \emph{misguiding} they are
excellent. ``But misguiding? That is surely an insult to the author!'' By
no means; it is on the contrary a compliment. The honored author has by
his appearance paid me so many compliments that I would have to lack all
culture and good manners if I omitted most dutifully to pay him my
counter-compliment: before a display bout one salutes, after all, with
the rapier.

For \emph{misguiding} and \emph{misguiding} are in philosophical
discussions two highly different things. There is a misguiding that
harms the truth, the clumsy, disagreeable kind that arises when an
immature person wants to orient the public without being oriented
himself; under this category the honored author may by no means be
reckoned. But there is also a misguiding that by dialectical side-lights
lures the public along into the investigation, leads the reader out onto
slippery ice, and gives him the notion that he who now really wants to
get into the matter must see with his own eyes and think for himself. The
author's misguiding is, as I hope to be able to prove, a higher
intellectual guiding; it is the old, well-known art of midwifery that he
employs: the method --- anyone who knows a little of Greek philosophy sees
it with half an eye --- the method is \emph{Socratic}.

Right at the beginning, where the author (Epil.\ II) expresses himself
on the introductory lecture, one suspects that something must lie behind
this. On the occasion of a figurative expression I am supposed to have
used in the course of the lecture, namely ``\emph{the breath of the
spirit},'' he expresses himself in the following manner: ``When the
spirit thus by its mere blowing transforms hindrances into conditions,
it may on the one hand seem impertinent to ask again about the condition
for the spirit's doing this, or about the reason why it does not always
blow. For it might be said that the spirit bloweth where it listeth, and
we hear its voice, but no one knows
% --- p. 312 ---
\opage{312}\setcounter{footnote}{0}%
whither it goeth. And that question can besides, as one easily sees, go
on endlessly. But when, on the other hand, the spirit is such a sorcerer
that, unconditionally, when and where it will, it transforms hindrances
into conditions or the reverse, what then is the point of the whole
question of hindrances and conditions for the spiritual life? Both seem to
have anything to do only where there is a reasonably grounded connection
and a law-bound activity; over against the absolutely magical, both
hindrances and conditions are empty words. So on the one side there
looms an endless heaping up of condition upon condition, on the other a
dismissal of the whole question! This certainly seems to be a hindrance
to the investigation begun. Will the speaker's spirit perhaps blow on it
and transform it into a condition, or will he blow it off and jump over
it?''
% the opening quotation mark of this quotation from the Epilogue is on p. 311
--- See, that was now the yield of the first lecture, hence the first
guidance.

Those readers who, without having acquainted themselves with my
introductory lecture, entrusted themselves to such guidance must, by
taking it straightforwardly and in dead earnest, indeed get a strange
notion of my appearance at the foreign university. Just imagine a
visiting lecturer coming blowing along with some phrases about ``the
blowing of the spirit,'' and that without even having so much going on in
his head that in the midst of the stream of words it might occur to him
to ask ``why the spirit does not always blow''! Just imagine a
philosophical professor, who by his position must after all be obliged
to know a little about the relation between the conditions and the
unconditioned, lacking reflection to such a degree that one must remind
him afterward that ``over against the absolutely magical, both hindrances
and conditions are empty words''! Just imagine a circle of cultured
listeners, among them even highly esteemed men of science, who for all of
sixteen hours have been able to hold out listening to something so
ungrounded,
% --- p. 313 ---
\opage{313}\setcounter{footnote}{0}%
so little coherent in itself, as such a lecture on ``the blowing of the
spirit'' leads one to expect! Just imagine the ill will that must now be
awakened afterward in the public when, guided by the Epilogue, it has to
admit to itself that it let itself be taken by surprise, that the phrases
by which in an unguarded moment it let itself be beguiled contained only
windy witticisms, castles in the clouds driven by the blowing of the
spirit. ``Lamentable blowing, disagreeable --- but nevertheless
salutary --- guidance!''

That the guidance is salutary I shall willingly concede, once we,
without puff and blowing, first get the guiding transformed into a
misguiding. This happens quite simply by paying heed to what our guide in
the Epilogue calls ``a reasonably grounded connection.'' There must, after
all, be a reasonably grounded connection between what has really been
said in a lecture and what is said about it afterward. That my
introductory lecture, as it now stands, is in all essentials the same as
the one that was delivered, my listeners can attest; but once this stands
firm, anyone, even on a cursory reading-through, will be able to convince
himself that what stands in the Epilogue as guiding is precisely
misguiding. But this misguiding has again, as by a smile from blessed
gods, become for us a true guidance. The spirit spoken of in the
introduction is ``the absolutely divine Spirit,'' hence the almighty
Spirit. Now the author of the Epilogue knows very well that when the
omnipotence of the spirit is to be spoken of, a listener, man or woman,
can easily mishear, and afterward think that the talk was of a spirit that
worked magically and therefore needed to heed neither hindrances nor
conditions, and that this is precisely a disturbing misunderstanding when
the lecture is expressly to treat of hindrances and conditions. To
forestall such intellectual after-pains, the honored author has wanted,
in a reversed style, to enjoin what is needful. Remember what is said
here:
% --- p. 314 ---
\opage{314}\setcounter{footnote}{0}%
\emph{The spirit in its working is by no means indifferent to conditions,
not unconditioned, nor does it create the conditions directly out of
nothing; no, it forms them out of natural presuppositions, it forms them
by a transformation} (p.~8). This transformation is not magical, for
remember what is said here! The old, worn-out question: \emph{if the
spirit must itself inspire its servants and itself bring the conditions
along, is it not then incomprehensible that it will not always and
everywhere work with equally forceful presence}? is not cursorily
dismissed, but answered in detail by an illumination of freedom and of
the irrefutable conditions of freedom (pp.~9--16); and that this
illumination, necessary for the fundamental thought, constitutes precisely
the main content of the introductory lecture --- see, this the clever
author of the Epilogue has obviously wanted to enjoin. But how, then, has
he gone about it? In order to impress on the listener how dangerous it is
to begin with a complete misunderstanding, he has himself begun with a
complete misunderstanding; in order to show the comical in a listener's
forgetting the objections that in a lecture are not merely raised but
also met, he has himself --- comically enough --- begun by forgetting
them; in order to forestall anyone with a short memory making himself
important by having forgotten the most important thing, the author begins
by pretending that it was he who became important by having forgotten
the most important thing. In truth, a guidance that demands resignation
and self-denial: one warns against misunderstandings by making oneself
guilty of misunderstanding; it is as if a father wanted to warn his sons
against drunkenness by getting himself tipsy.

In Epilogue IV the following rather misguiding hints for guidance occur.
``No one shall, because he says A in science, have the right to say B in
faith. \ldots{} If only, the moment this is conceded, logic at any rate
does not poke its nose into this system of concepts of faith, and so the
magic circle by
% --- p. 315 ---
\opage{315}\setcounter{footnote}{0}%
which all science was to be kept away is broken through.'' Hereby the
author gives a Socratic flick to all those who, although it is shown again
and again that thinking, consequently logic too, has validity just as
well in the doctrine of faith as in science, are caught in the old
ingrained prejudice that faith as faith is thoughtless and must borrow all
its thoughts from science. To drive this misunderstanding to the point of
caricature, he now depicts quite ingeniously what comes of giving oneself
up to a kind of scientific fantasizing about the heterogeneous circles of
thought of faith and knowledge. ``I do not know,'' he says, ``to what
degree I may demand mutual agreement in the domain of faith. How does
faith relate to Holy Scripture? Are these perhaps also heterogeneous, so
that what holds in the one must not be carried over to the other? --- And
Holy Scripture itself? Mutual contradictions have, after all, been shown,
for example between 2 Samuel and 1 Chronicles. Can these then be removed
by the magic formula imparted, that A and B must not belong to
heterogeneous domains? Is the Bible, then, itself so heterogeneous within
itself that the one thing may well be true in the books of Samuel while
the opposite is true in the books of Chronicles?''

From the conspicuous superiority with which the author here gives free
play to the misguiding, one notices how easy it would have been for him to
keep to what is stated in the lectures and give us a straightforward
guidance. That it must have cost him self-conquest not to do so I may well
assume; for the temptation has truly been great. How far a conclusion can
be drawn from given premises does not, after all, depend on your pleasure
and mine; it depends on the laws of thought. Were the guidance not, in
accordance with the method, absolutely bound to be misguiding, then the
author of the Epilogue would naturally have strongly emphasized that in
the lectures no magic formula at all has been set up, but that it has
been shown in definite examples what it leads to when one draws
consequences for knowledge
% --- p. 316 ---
\opage{316}\setcounter{footnote}{0}%
from the presuppositions of faith and consequences for faith from those
of knowledge. The question is: can it be done or not? If it really can be
done, then let the lecturer be shown how it can be done and what his
mistake consists in; if it cannot be done, then the lecturer is surely
right, and what, then, is the use of the many hovering, over-hovering,
extravagantly hovering deliberations concerning magic formulas where
there is no magic formula?

The tactic that the author has employed with so much success to misguide
the readers concerning the inner agreement of the concepts of faith within
their circle, the same tactic he employs with equally brilliant effect to
confuse the representations concerning the inner coherence of the concepts
of knowledge. ``There has,'' he says, ``as far as I remember, been nothing
said about mutual contradictions in the domain of science. Perhaps it is
presupposed that they will not occur. But what if they occurred
nevertheless? I at any rate cannot feel quite secure. For I
have after all heard talk of disputes between different sciences which,
when they are set forth sharply, might easily seem just as radical and
inaccessible to mediation as the dispute between Samuel and Chronicles,
indeed as several of the contradictions that seem to obtain between faith
and science in general.\ldots{} I am in general not at all without fear
that the absolute separation between faith and science set up by the new
doctrine will consistently --- if consistency is not in the end a quite
empty word --- lead to dissolving the inner unity of both faith and
science and splitting both into atoms, so that there will be neither any
faith nor any science any more, but only an indeterminable multitude of
detached representations of faith and scientific concepts or propositions.
Indeed, there already seems to appear a prevailing inclination at least to
conceive the sciences exclusively one by one, while the thought of a
general organism of science, of which the individual fields are
integrating members,
% --- p. 317 ---
\opage{317}\setcounter{footnote}{0}%
comparatively recedes into the background.'' Here we have the Socratic
roguery again.

The ironizing author wants to make guileless readers believe that because
religion is heterogeneous, absolutely heterogeneous, with science, it
should follow that all the sciences together, when they noticed that
faith was no longer able to keep them in check, might take it into their
heads to become refractory and declare themselves mutually absolutely
heterogeneous. He paints a dark picture. It might then happen that even
each single science, once it had tasted the absolutely heterogeneous, got
an ache in its limbs and began to dismember itself; that arithmetic and
geometry, for example, when they got wind that something could be true in
2 Samuel (XXIV) which was untrue in 1 Chronicles (XXI), and conversely,
took the resolution that something should in future be able to be true in
arithmetic which was untrue in geometry, and conversely. These are
threatening prospects! Some ``lover of knowledge'' or other, not yet
properly oriented in the republic of letters, might at such ill-boding
omens be seized with anxiety, get a fright in his blood, imagine that we
already had anarchy, that the sciences were in full revolt, that there was
fire in the philosophical quarter, and start crying Fire!

For the calming of tempers I must point out that the matter is not so
dangerous. The sciences are, after all, all of a sober, calm, and at
bottom peaceable character; certainly they have various mutual disputes;
these disputes may well also now and then become somewhat vehement and
passionate; but the vehement and passionate proceeds from the human
``lovers of knowledge,'' not from the sciences. It is precisely through
the quiet impersonal power of scientific truths over the passionate
personalities that the lost equilibrium is restored and harmony preserved.
Should anything bring the sciences into revolt, it would have to be that
they were brought back under the guardianship of faith, under
% --- p. 318 ---
\opage{318}\setcounter{footnote}{0}%
the ``good will'' of dogmatics, into the theological universal monarchy. I
have really investigated the matter, I have been round and made inquiries
of the individual sciences, and can from each and every one (cf.\ ``Den
gode Villie som Magt i Videnskaben,'' pp.~9--108) procure a certificate
that they are all unconditionally agreed in not wanting to take up
miraculous elements. On the other hand, I was to give their regards and
say that if they might only be free of revelations, miracles, and
prophecies, they would never say an ill word either against faith or
against the believers, but in all personal concerns conduct themselves
morally and reasonably, and well know how to observe the dignity they owe
themselves over against religion. Our Socratic's fear that ``science
without Christianity, in its extreme consequence, would have to abolish
itself and perish in skepticism'' is thus a misguiding guidance.

So that it shall not look, however, as if on the occasion of the fright I
had only now hit upon the thought that the sciences constitute an
organism, I shall here permit myself to refer to ``Philosophisk
Propædeutik'' i Grundtræk, Copenhagen 1857, pp.~68--117, and
particularly to emphasize the following (p.~71): ``From the relation
between the theory of cognition and the theory of science it follows that
the true comprehensive theory of science can come about only by the
fundamental interests of the heterogeneous sciences being represented in
the discussions from which it issues. In `the comprehensive discussion and
debate of all against all' (Sibbern), each science thus has its vote and
makes its contribution. That philosophy works together, clarifies, and
organizes these contributions into an intellectual whole proves that in
this respect it does not decide the matter on its own authority, but
works as the common organ of the sciences, as their own universal
consciousness, the being that is their kinship, the light in which they
know and recognize one another mutually.''

The fixed idea at which our Socratic ironizes most finely of all, and
rightly, since it still haunts
% --- p. 319 ---
\opage{319}\setcounter{footnote}{0}%
so many brains, is precisely the representation of the necessity of
Christianity for grounding, maintaining, and furthering the sciences,
from which in turn follows the necessity of the sciences for grounding,
maintaining, and furthering Christianity. ``Supposing,'' he says (Epil.\
V), ``that there were some truth in the thought from which I have not yet
been able to free myself, namely that there is after all an inner and
essential connection between science and faith, then there would thus lie
in it a reasonable answer to the \ldots{} question why --- by and large
--- science and faith have not been able to leave each other alone, but
have either sought to be united or fought each other. There may certainly
have been much sophistry both in the strife and in the attempts at union,
and in particular many an intended encroachment on the part of science;
but the sophistry and the encroachment would be without motive if faith
could be indifferent to science as a whole. If, on the other hand, it is
so, as history too seems to show, that certain forms of faith will exclude
science, and that its flourishing is in general conditioned by a certain
prevailing religious standpoint --- so that, for example, science and the
whole culture bound up with it have never been able to attain the depth
and extent on Mohammedan as on Christian soil --- then science surely has
good reason not to let the content and form of faith be a matter of
indifference to it. Only the faith that is essentially pure and true and
free can therefore also let science develop freely; but precisely the
true faith has no need either --- for example by a doubtful exceptio fori
or other means --- to withdraw from the criticism of mature science. For
the full truth bears the full light from every side.''
% printer's defect: the Danish prints "ndvikle" for "udvikle" (develop); translated in the sense

The historical and the unhistorical, the true and the half-true, the loose
and the firm, the author has here driven in among one another in a most
ingenious way, as under all manner of remarks he spins the reader round
once, twice, three times and helps him most adroitly to mistake what
essentially matters --- naturally in order to put his
% --- p. 320 ---
\opage{320}\setcounter{footnote}{0}%
independent thinking to the test. What history teaches about the many
disputes between faith and knowledge and the corresponding attempts at
reconciliation I have, after all, sought to show in such definite
features that no one who will take the trouble to go through my writings
will easily mistake the results. He who can still think that, after so
many shipwrecked attempts at mediation, there is the least prospect here
either of making science believing or faith scientific must surely be
without all knowledge of the theological and philosophical literature of
the last century.
% printer's defect: the Danish prints "Toen" for "Troen" (faith); translated in the sense

It is surely this too that the author wants to remind us of in his
Socratic-backward way. ``There is,'' he says, ``an inner essential
connection between science and faith; faith cannot be indifferent to
science.'' He who takes these propositions literally is fooled; for ---
that we should repeat it once more --- there is in our day no true
science, be it empirical or a priori, that in any way either directs
itself or will direct itself by faith, namely the positive faith of
revelation, and it is of that we speak; on the other hand, meaning is
pressed into the phrases when one understands them thus, that \emph{it
can by no means be indifferent to personality how the religious and the
scientific truths relate to each other}.

``The true faith has no need to withdraw from the criticism of mature
science'' --- this proposition is misguiding when it is taken as
straightforwardly guiding. What becomes of the Gospel of Christ's birth,
of Christ's resurrection, of Christ's ascension, of the outpouring of the
Spirit, and so on, when sacred history is really to be subjected to ``the
criticism of mature science''? Give me a sample of such a piece of
criticism, a criticism that equally satisfies both faith and science, and
I shall say with Socrates: ἀγαιμην ἀν θαυμαστως,
ω Ζηνων. If, on the other hand, we understand the proposition as
misguiding,
% Greek (Plato, Parm. 129c) copied as printed, unaccented
% --- p. 321 ---
\opage{321}\setcounter{footnote}{0}%
then there is meaning in it; the meaning is then this: the true faith has
no need to fear the criticism of mature science, for if the believer is
at the same time a knower, then he must surely know that the criticism of
knowledge precisely dissolves the sacred facts, because these facts have
their truth not in knowledge but in faith.

In confirmation of my view that the author of the Epilogue guides
backward by misguiding, I could cite a great many more, even brilliant,
examples; but time does not permit it. In order to repel the stupid
tittle-tattle with which the theologians so often regale us, that
certainty about the sensuous rests on faith, that astronomy rests on
faith, that induction and with it all empirical knowledge and science
rests on faith, and that such a faith is then again akin to religious
faith, I have (particularly in Lecture IV) strongly emphasized the
\emph{certainty} of empirical knowledge. To make this certainty evident to
outsiders, I appealed to a natural \emph{instinct for truth}; this our
Socratic wittily twists into the backward assertion that I was about to
want to ground science on ``unmediated instinct''; he reminds us of
Descartes's proposition that all human certainty rests on confidence in
God's truthfulness. Hereby he has in an indirect and very fine way
referred those readers who misunderstand my presentation, as if I built
science on raw instinct, to my theory of science, particularly to
Grundideernes Logik II (pp.~176--275).

What ironic roguery runs through the author's maieutic distortions I
illustrate with yet one more example, in which humor and wit rise almost
to exuberance. I have (Lecture V, p.~81) expressed myself thus: If it is
asked, for example, what granite and flint are, then these concepts, with
corresponding facts, can be fixed so definitely that anyone who has
% --- p. 322 ---
\opage{322}\setcounter{footnote}{0}%
sound understanding, preparatory education, and a sense for scientific
investigations, be he otherwise Jew, Greek, Mohammedan, regardless of
national differences and faith, is compelled to concede the validity of
the concepts: ``This must be granite, this must be flint.'' \ldots{} What
has our Socratic not managed to get out of this! ``There is held up to
me,'' he says, ``yet another mark by which I am to tell the concept of
faith from the scientific one, namely that the latter can be made
comprehensible to every human being, while the former presupposes a
certain religious standpoint. The concept of oxygen and the propositions
of geology must be able to be grasped and appropriated by Christians,
Mohammedans, and Brahmins without distinction, while one will seek in
vain to convince others than Christians of the resurrection of Christ.''
The adroit distortion now consists in this, that whereas I make ``sound
understanding, preparatory education, and a sense for scientific
investigations'' in every human being express conditions for
understanding a scientific doctrine, our Socratic briskly skips over the
conditions and has me speak quite generally of ``every human being.'' The
exaggeration that \emph{every human being} should without further ado be
able to be convinced of a scientific truth is then (Epil.\ VI) illustrated
by examples such as these: ``that it might perhaps be just as difficult to
get a Chinese to grasp European law and European history as to grasp
Christianity; that it is difficult to get the Germans to comprehend that
the greater part of Schleswig is Danish, and that a Negro who believes he
can raise or still a storm with the help of a fetish will have difficulty
enough in grasping the thought of unchangeable laws of nature.''

Herewith I close these discussions, and thank the honored author of the
Epilogue for the essential service he has rendered me on this occasion.
The Epilogue has a broad stream; the stream runs through a whole nine
articles, until in a
% --- p. 323 ---
\opage{323}\setcounter{footnote}{0}%
tenth it at last flows out into an ``Epilogue to the Epilogue.'' Precisely
this is what is Socratically excellent. Were I to clarify in detail and
point by point every one of the subtly feigned misunderstandings, I would
have to write a thick book. And what would the result then be? That the
public, having had enough of the articles, would decline with thanks to
read the book. This the author of the Epilogue has very cleverly
calculated, and precisely for that reason wanted to give me occasion to
get it established.

There can, as is well known, always be said various things for and
against anything, and especially for and against a new philosophical
doctrine. That contradictions arise against my philosophy is so far from
surprising me that the opposite would on the contrary have prepared me a
by no means pleasant surprise. ``\emph{What is philosophy}? A science full
of contradictions! This definition, however much, regarded from one side,
it seems designed to break the staff over all philosophy, is nonetheless
grounded both in the history of science and in the nature of the case.
For any other science such an admission would surely be annihilating; for
philosophy it indicates an essential point of departure, a reliable point
of support, a firm, unshakable foundation. That philosophy must contain
contradictions is, after all, a matter of course, since it is precisely
the science that everywhere, both in existence and in sound common sense's
apprehension of existence (Propæd.\ pp.~40--44), discovers and points out
contradictions; with regard to this, philosophy is precisely to be
conceived as a \emph{general doctrine of contradiction}. As doctrine of
contradiction, however, philosophy must not merely discover and point out
the contradictions, it must also test the contradictions; for there is a
difference among contradictions. Some contradictions, namely, are
superficial, unnecessary, easy to avoid, others on the contrary pressing
and irrefutable; some contradictions creep in through thoughtlessness,
others are artificially produced by hair-splitting, others are exposed
% --- p. 324 ---
\opage{324}\setcounter{footnote}{0}%
by acuteness; some contradictions play as in a game of thought, others
mark essential knots, difficult problems, hard to resolve.'' (Forelæsninger
over ``Philosophisk Propædeutik'' from the academic year 1860--61,
Copenhagen 1862, pp.~1--2.)
% the Danish prints "Universitetsaeret" for "Universitetsaaret" (academic year); translated in the sense

The circumstance that philosophy moves through contradictions of different
kinds and very different depth makes it so exceedingly easy to conduct a
polemic of opinion against any philosophical doctrine whatever; for the
public is easily dazzled when, without deeper insight into the connection
of the matter, it sees contradictions set out in a few propositions and
supported by detached quotations. If sound common sense gets a
contradiction under one illumination, it can see through its nullity at
once and is not beguiled; if on the other hand it gets the same
contradiction under another illumination and in another connection, it is
impressed and believes that it is something. If I say: ``God and the world
cannot stand in opposition to each other without limiting each other; they
cannot limit each other without restricting each other. God must thus be
restricted by the world just as well as the world by God; but a
restricted God is a contradiction, therefore God does not exist'': then
sound common sense smells a rat. One need only recall the difference
between limitation and restriction, such as that every head must be
limited, but that it by no means follows that every head is a restricted
head; that one art can have its limit in another without the one
therefore restricting the other, or \emph{doing the other damage}, and so
on --- and anyone who has natural understanding grasps at once the
difference between \emph{limiting}, i.e.\ \emph{determining}, and
\emph{restricting}, i.e.\ \emph{doing damage}. If, on the other hand, the
talk is of faith and knowledge, hence of an opposition that is still
rather unclear to the general consciousness, then it is easy to impress
by the contradiction: ``Faith and knowledge
% --- p. 325 ---
\opage{325}\setcounter{footnote}{0}%
limit each other, for faith is not knowledge, knowledge not faith, and yet
they are not to restrict each other; but to limit is surely to restrict;
therefore: faith and knowledge are to restrict each other by not
restricting each other. What a contradiction, indeed what a monstrous
contradiction!'' The reader's understanding stands still, the polemic of
opinion flourishes, and the public believes that it is something.

To expose all the hollownesses, distortions, and crude misunderstandings
that can be contained in a pamphlet of a few sheets, I would, as I said,
each time have to write a voluminous book; it is a fruitless and
unmanageable labor. Here comes one in the juggling, gallant manner of wit
with ``dualism'' and ``genius'' and ``collier's faith'' and presents the
public with one and a half nimbly fabricated contradictions, and then I am
to answer that! Another, in the dogmatically trumping high-pressure
manner, with a theological universal monarchy on the tip of ``the good
will,'' steps up and shows the public a blind man who is not riding on a
lame one, a figurative, emblematic contradiction, and then I am to answer
that! A third, in the speculative-turkey gobbling manner, announces
himself with twelve blasts on a mystic trumpet; each blast is a
contradiction, a turkey-cock contradiction, and then I am to answer that!
A fourth, in the well-known nice-man manner of self-important narrowness,
comes stalking along with something on a critical lime-rod that he calls
a conviction; the conviction is a contradiction, and the lime-rod is a
contradiction too; he strikes at my edifice of doctrine with the
conviction and the contradictions and the lime-rod and says Ruin! then I
am to answer that! Here comes a fifth, a sixth \ldots. --- it exceeds my
powers. Perhaps the gentlemen pamphleteers would very well live on
answers; but it is certain that were I to answer them all and on
everything, I should have to die of answers.

From these prolixities, troubles, and difficulties
% --- p. 326 ---
\opage{326}\setcounter{footnote}{0}%
my Socratic benefactor has now, by his feigned misconstructions,
misconstructions that in a light and pleasant way caricature what the
public otherwise so readily takes for serious objections, freed me for the
time being. The author of the Epilogue has held out a helping hand to me;
I thank him, moved, and press his hand with gratitude.

\begin{center}\rule{3cm}{0.4pt}\end{center}

\backmatter
% --- Rettelser (PDF 338), unnumbered leaf after p. 326: no page marker.
\chapter*{Errata.}
\addcontentsline{toc}{chapter}{Errata}
% The Danish readings are kept as printed; the English text above
% follows the corrections, with a % note at each site.

\noindent
\begin{tabular}{@{}l r l l@{}}
P. & 50  & line 1 from top & vente at finde den \emph{read}: vente det\\
-  & 70  & --- 1 - -      & stridt mod \emph{read}: stridt med\\
-  & 109 & --- 3 from bottom   & det absolute Selv \emph{read}: det Absolute selv\\
-  & 174 & --- 9 and 10 from top & Saducæer \emph{read}: Sadducæer\\
\end{tabular}

\begin{center}\rule{3cm}{0.4pt}\end{center}

\end{document}
